% Lutte contre l'érosion — SVT 1ère C — Cours signé OnBuch+
% Programme officiel MINESEC. Compiler avec : tectonic NCT08-lutte-erosion.tex
\newcommand{\DOCMATIERE}{SVT}
\newcommand{\DOCNIVEAU}{1ère C}
\newcommand{\DOCMODULE}{Module 3 : L'Éducation à l'Environnement et au Développement Durable}
\newcommand{\DOCLECON}{NCT08}
\newcommand{\DOCTITRE}{Lutte contre l'érosion}
% OnBuch+ — préambule « cours signés » ENRICHI (compilé avec tectonic / XeLaTeX)
% Système visuel « Pop » d'OnBuch+ : fond crème (jamais blanc), bordures encre
% épaisses, ombres « dures » décalées, titres Archivo Black en capitales.
% Version MATHÉMATIQUES : graphes (pgfplots/tikz), boîtes pédagogiques
% (définition, propriété, méthode, attention, le savais-tu, exercice résolu,
% exercice, TP, bilan), page de garde et signature OnBuch+ ancrée en bas.
%
% Macros attendues dans main.tex AVANT \input{preamble} :
%   \DOCMATIERE  \DOCNIVEAU  \DOCMODULE  \DOCLECON (numéro)  \DOCTITRE
\documentclass[11pt]{article}

\usepackage{fontspec}
\usepackage[french]{babel}
\usepackage[a4paper,margin=2cm,top=3.4cm,bottom=2.8cm,headheight=1.4cm,footskip=1.3cm]{geometry}
\usepackage{amsmath,amssymb,amsfonts}
\usepackage{mathtools} % \xmapsto, \coloneqq... souvent utilisés par les modèles
\usepackage{cancel} % simplifications barrées (bilans de forces)
% \abs{z} (module, valeur absolue) et \norm{u} (norme), utilisés par les modèles
\providecommand{\abs}{}\let\abs\relax\DeclarePairedDelimiter{\abs}{\lvert}{\rvert}
\providecommand{\norm}{}\let\norm\relax\DeclarePairedDelimiter{\norm}{\lVert}{\rVert}
\usepackage[table]{xcolor} % [table] : fournit aussi \rowcolors (alternance de lignes)
\usepackage{pagecolor}
\usepackage{tikz}
\usetikzlibrary{arrows.meta,positioning,calc,shapes.geometric,shapes.misc,shapes.arrows,decorations.pathmorphing,decorations.pathreplacing,decorations.markings,patterns,shadows,mindmap,trees,fit,backgrounds,matrix,intersections,angles,quotes}
\usepackage{pgfplots}
\usepackage[version=4]{mhchem} % équations nucléaires \ce{^{235}_{92}U}
\usepackage[european,siunitx]{circuitikz} % schémas électriques
\pgfplotsset{compat=1.18}
\usepgfplotslibrary{fillbetween,groupplots} % aires entre courbes (calcul intégral)
\usepackage{enumitem}
\usepackage{fancyhdr}
% [most] charge toutes les bibliothèques tcolorbox (listings, minted, raster,
% documentation...) et gonfle inutilement la mémoire TeX sur un document long
% (~300 boîtes) : on ne charge que ce qui est réellement utilisé.
\usepackage[skins,breakable]{tcolorbox}
\usepackage{graphicx}
\usepackage{colortbl}
\usepackage{booktabs}
\usepackage{tabularx}
\usepackage{diagbox} % case d'en-tête diagonale des tableaux à double entrée
% Colonnes extensibles centrée / gauche / droite (C, L, R), souvent utilisées par les modèles
\newcolumntype{C}{>{\centering\arraybackslash}X}
\newcolumntype{L}{>{\raggedright\arraybackslash}X}
\newcolumntype{R}{>{\raggedleft\arraybackslash}X}
\usepackage{array}
\usepackage{multicol}
\usepackage{siunitx}
% parse-numbers=false : accepte aussi \SI{2,5 \times 10^{-3}}{...}
\sisetup{locale=FR,output-decimal-marker={,},group-minimum-digits=4,parse-numbers=false}
\DeclareSIUnit\ohm{\ensuremath{\symup{\Omega}}} % U+2126 absent de la police
\DeclareSIUnit\division{div} % réglages d'oscilloscope (ms/div, V/div)
\DeclareSIUnit\tr{tr} % tours (vitesse de rotation, ex. \SI{1500}{\tr\per\minute})
\DeclareSIUnit\ch{ch} % cheval-vapeur (puissance moteur, unité usuelle hors SI)
\DeclareSIUnit\franccfa{FCFA} % franc CFA (coûts, contexte camerounais)
\DeclareSIUnit\dioptre{\ensuremath{\delta}} % vergence d'une lentille (dioptrie)
\usepackage{qrcode}
% \degree : commande fréquemment utilisée par les modèles pour un angle ou une
% température en dehors de \SI{}{} (non fournie par les paquets chargés).
\providecommand{\degree}{^{\circ}}
% \qed (fin de démonstration), utilisé par les modèles sans amsthm
\providecommand{\qed}{\hfill$\square$}
% \barre : double barre des valeurs interdites dans un tableau de variations (tkz-tab)
\providecommand{\barre}{\ensuremath{\|}}
% environnement proof (amsthm non chargé) utilisé par les modèles
\ifdefined\proof\else\newenvironment{proof}[1][Démonstration]{\par\noindent\textit{#1.}\ }{\qed\par}\fi
\usepackage{lastpage}
\usepackage{hyperref}
\hypersetup{hidelinks}

% ============================================================
% Palette « Pop » OnBuch+ — mêmes hex que la classe P (plus_ui.dart)
% ============================================================
\definecolor{poporange}{HTML}{F59321}
\definecolor{popdark}{HTML}{E07A0C}
\definecolor{popcream}{HTML}{FFF6E8}
\definecolor{popink}{HTML}{1C1714}
\definecolor{poppurple}{HTML}{7A5AE0}
\definecolor{popgreen}{HTML}{2E9E6A}
\definecolor{poppink}{HTML}{E0457B}
\definecolor{popblue}{HTML}{2F7DE1}
\definecolor{popgold}{HTML}{C98A12}
\definecolor{popmuted}{HTML}{8A7F76}
\definecolor{popline}{HTML}{EADFD2}
\definecolor{popgray}{HTML}{8A7F76}
\colorlet{popgrey}{popgray}
% Alias tolérants (le modèle nomme parfois les couleurs différemment)
\colorlet{popwhite}{white}
\colorlet{popblack}{popink}
\colorlet{popred}{poppink}
\colorlet{popyellow}{popgold}
% Teintes claires pour fonds de tableaux / zones
\colorlet{popblueL}{popblue!10!white}
\colorlet{poporangeL}{poporange!14!white}
\colorlet{popgreenL}{popgreen!12!white}
\colorlet{poppurpleL}{poppurple!12!white}
\colorlet{poppinkL}{poppink!12!white}
\colorlet{popgoldL}{popgold!14!white}

\pagecolor{popcream}

% ============================================================
% Typographie
% ============================================================
\defaultfontfeatures{Ligatures=TeX}
\setmainfont{PlusJakartaSans-Medium.ttf}[Path=../../../fonts/,BoldFont=PlusJakartaSans-Bold.ttf]
% Maths en sans-serif (Fira Math) avec chiffres et lettres droites de Plus
% Jakarta, pour que les formules se fondent dans le texte.
\usepackage[math-style=ISO,bold-style=ISO]{unicode-math}
\setmathfont{FiraMath-Regular.otf}[Path=../../../fonts/]
\setmathfont{PlusJakartaSans-Medium.ttf}[Path=../../../fonts/,range={up/num,up/{Latin,latin},\mathup}]
\usepackage{icomma} % virgule décimale sans espace en mode math
% Caractères mathématiques Unicode absents des polices de texte (Plus Jakarta,
% Archivo Black) : rendus par leur équivalent mathématique, en texte comme en
% formule — sinon ils s'affichent en carré vide (ex. « Arithmétique dans ℤ »).
\usepackage{newunicodechar}
\newunicodechar{ℝ}{\ensuremath{\mathbb{R}}}
\newunicodechar{ℤ}{\ensuremath{\mathbb{Z}}}
\newunicodechar{ℂ}{\ensuremath{\mathbb{C}}}
\newunicodechar{ℕ}{\ensuremath{\mathbb{N}}}
\newunicodechar{ℚ}{\ensuremath{\mathbb{Q}}}
\newunicodechar{𝔻}{\ensuremath{\mathbb{D}}}
\newunicodechar{α}{\ensuremath{\alpha}}
\newunicodechar{β}{\ensuremath{\beta}}
\newunicodechar{γ}{\ensuremath{\gamma}}
\newunicodechar{δ}{\ensuremath{\delta}}
\newunicodechar{ε}{\ensuremath{\varepsilon}}
\newunicodechar{θ}{\ensuremath{\theta}}
\newunicodechar{λ}{\ensuremath{\lambda}}
\newunicodechar{μ}{\ensuremath{\mu}}
\newunicodechar{σ}{\ensuremath{\sigma}}
\newunicodechar{τ}{\ensuremath{\tau}}
\newunicodechar{φ}{\ensuremath{\varphi}}
\newunicodechar{ψ}{\ensuremath{\psi}}
\newunicodechar{ω}{\ensuremath{\omega}}
\newunicodechar{Σ}{\ensuremath{\Sigma}}
% majuscules produites par \MakeUppercase dans les titres (α-aminés -> Α)
\newunicodechar{Α}{\ensuremath{\alpha}}
\newunicodechar{Β}{\ensuremath{\beta}}
\newunicodechar{Γ}{\ensuremath{\Gamma}}
\newunicodechar{Δ}{\ensuremath{\Delta}}
\newunicodechar{Ε}{\ensuremath{\varepsilon}}
\newunicodechar{Θ}{\ensuremath{\Theta}}
\newunicodechar{Λ}{\ensuremath{\Lambda}}
\newunicodechar{Μ}{\ensuremath{\mu}}
\newunicodechar{Τ}{\ensuremath{\tau}}
\newunicodechar{Φ}{\ensuremath{\Phi}}
\newunicodechar{Ψ}{\ensuremath{\Psi}}
\newunicodechar{Ω}{\ensuremath{\Omega}}
\newunicodechar{↦}{\ensuremath{\mapsto}}
\newunicodechar{∈}{\ensuremath{\in}}
\newunicodechar{∉}{\ensuremath{\notin}}
\newunicodechar{∧}{\ensuremath{\wedge}}
\newunicodechar{⇔}{\ensuremath{\Leftrightarrow}}
\newunicodechar{⟺}{\ensuremath{\Longleftrightarrow}}
\newunicodechar{⇒}{\ensuremath{\Rightarrow}}
\newunicodechar{⟹}{\ensuremath{\Longrightarrow}}
\newunicodechar{∘}{\ensuremath{\circ}}
\newunicodechar{∪}{\ensuremath{\cup}}
\newunicodechar{∩}{\ensuremath{\cap}}
\newunicodechar{≡}{\ensuremath{\equiv}}
\newunicodechar{⊂}{\ensuremath{\subset}}
\newunicodechar{⊥}{\ensuremath{\perp}}
\newunicodechar{∃}{\ensuremath{\exists}}
\newunicodechar{∀}{\ensuremath{\forall}}
\newunicodechar{∛}{\ensuremath{\sqrt[3]{\,}}}
\newunicodechar{∜}{\ensuremath{\sqrt[4]{\,}}}
\newunicodechar{⃗}{\ensuremath{{}^{\to}}}
\newunicodechar{ⁿ}{\ifmmode^{n}\else\textsuperscript{\ensuremath{n}}\fi}
\newunicodechar{ˣ}{\ifmmode^{x}\else\textsuperscript{\ensuremath{x}}\fi}
\newunicodechar{ᵇ}{\ifmmode^{b}\else\textsuperscript{\ensuremath{b}}\fi}
\newunicodechar{ʳ}{\ifmmode^{r}\else\textsuperscript{\ensuremath{r}}\fi}
\newunicodechar{ᵘ}{\ifmmode^{u}\else\textsuperscript{\ensuremath{u}}\fi}
\newunicodechar{ʸ}{\ifmmode^{y}\else\textsuperscript{\ensuremath{y}}\fi}
\newunicodechar{ᵃ}{\ifmmode^{a}\else\textsuperscript{\ensuremath{a}}\fi}
\newunicodechar{ᵉ}{\ifmmode^{e}\else\textsuperscript{\ensuremath{e}}\fi}
\newunicodechar{ᵏ}{\ifmmode^{k}\else\textsuperscript{\ensuremath{k}}\fi}
\newunicodechar{ᵖ}{\ifmmode^{p}\else\textsuperscript{\ensuremath{p}}\fi}
\newunicodechar{ᴺ}{\ifmmode^{N}\else\textsuperscript{\ensuremath{N}}\fi}
\newunicodechar{ᶜ}{\ifmmode^{c}\else\textsuperscript{\ensuremath{c}}\fi}
\newunicodechar{ʰ}{\ifmmode^{h}\else\textsuperscript{\ensuremath{h}}\fi}
\newunicodechar{ᵠ}{\ifmmode^{q}\else\textsuperscript{\ensuremath{q}}\fi}
\newunicodechar{ₙ}{\ifmmode_{n}\else\textsubscript{\ensuremath{n}}\fi}
\newunicodechar{ₐ}{\ifmmode_{a}\else\textsubscript{\ensuremath{a}}\fi}
\newunicodechar{ᵢ}{\ifmmode_{i}\else\textsubscript{\ensuremath{i}}\fi}
\newunicodechar{ᵣ}{\ifmmode_{r}\else\textsubscript{\ensuremath{r}}\fi}
\newunicodechar{ₖ}{\ifmmode_{k}\else\textsubscript{\ensuremath{k}}\fi}
\newunicodechar{ᵦ}{\ifmmode_{\beta}\else\textsubscript{\ensuremath{\beta}}\fi}
\newunicodechar{ₓ}{\ifmmode_{x}\else\textsubscript{\ensuremath{x}}\fi}
\newfontfamily\popdisplay{ArchivoBlack-Regular.ttf}[Path=../../../fonts/]
\newfontfamily\poplogo{SpaceGrotesk-Bold.ttf}[Path=../../../fonts/]
\newfontfamily\popsemi{PlusJakartaSans-SemiBold.ttf}[Path=../../../fonts/]
\newfontfamily\popxbold{PlusJakartaSans-ExtraBold.ttf}[Path=../../../fonts/]
\newfontfamily\popxboldls{PlusJakartaSans-ExtraBold.ttf}[Path=../../../fonts/,LetterSpace=3.0]

\setlist[enumerate]{leftmargin=1.7em,itemsep=5pt,topsep=4pt}
\setlist[itemize]{leftmargin=1.7em,itemsep=4pt,topsep=4pt,label={\color{poporange}\textbullet}}
\setlength{\parindent}{0pt}
\setlength{\parskip}{5pt}
\renewcommand{\arraystretch}{1.25}

% pgfplots : style Pop par défaut
\pgfplotsset{
  every axis/.append style={axis line style={popink,line width=1pt},tick style={popink},
    grid style={popline,line width=0.5pt},label style={font=\small},tick label style={font=\footnotesize}},
  popcurve/.style={poporange,line width=1.8pt,no markers,smooth},
}
% Même style disponible dans un \draw TikZ simple (hors pgfplots)
\tikzset{popcurve/.style={poporange,line width=1.8pt}}
\tikzset{popgrid/.style={popline,very thin}}
\colorlet{popcurve}{poporange} % le nom du style est parfois utilisé comme couleur

% ============================================================
% Pill / squiggle
% ============================================================
\newcommand{\poppill}[3]{%
  \tikz[baseline=-0.55ex]\node[draw=popink,line width=1.2pt,rounded corners=2.6pt,%
    fill=#2,inner xsep=6.5pt,inner ysep=3pt]{%
    \popxboldls\fontsize{7}{7}\selectfont\color{#3}\MakeUppercase{#1}};%
}
\newcommand{\popsquiggle}[1]{%
  \tikz[baseline=-0.4ex]{
    \draw[#1,line width=2.6pt,line cap=round]
      plot [smooth] coordinates {(0,0.09) (0.34,0.01) (0.68,0.21) (1.02,0.01) (1.36,0.10)};
  }%
}
% Mot-clé surligné (marqueur orange sous le texte)
\newcommand{\cle}[1]{{\popxbold\color{popdark}#1}}

% ============================================================
% En-tête / pied
% ============================================================
\pagestyle{fancy}
\fancyhf{}
\renewcommand{\headrulewidth}{0pt}
\renewcommand{\footrulewidth}{0pt}
\lhead{\raisebox{-0.15ex}{\poplogo\fontsize{13}{13}\selectfont\color{popink}OnBuch\color{poporange}+}}
\rhead{\poppill{\DOCMATIERE\ \textbullet\ \DOCNIVEAU\ \textbullet\ Leçon \DOCLECON}{white}{popink}}
\lfoot{\popxbold\fontsize{7.6}{7.6}\selectfont\color{popmuted}\MakeUppercase{Cours signé OnBuch+}\ \textbullet\ {\popsemi\DOCTITRE}}
\rfoot{\tikz[baseline=-0.6ex]\node[rounded corners=8.5pt,fill=popink,minimum height=17pt,minimum width=17pt,inner xsep=5pt,inner sep=0]{\popxbold\fontsize{8}{8}\selectfont\color{popcream}\thepage\,/\,\pageref*{LastPage}};}

% ============================================================
% Titres
% ============================================================
\newcommand{\chaptitre}[2]{%
  \begin{tcolorbox}[enhanced,colback=poporange,colframe=popink,boxrule=2.6pt,arc=11pt,
      drop shadow=popink,left=15pt,right=15pt,top=12pt,bottom=11pt,before skip=3pt,after skip=18pt]
    \poppill{#2}{popink}{popcream}\par\vspace{9pt}
    {\raggedright\hyphenpenalty=10000\exhyphenpenalty=10000\popdisplay\fontsize{22}{24}\selectfont\color{popcream}\MakeUppercase{#1}\par}
  \end{tcolorbox}
}
% Section numérotée : pastille orange avec le numéro + titre Archivo
\newcounter{popsec}
\newcommand{\coursec}[1]{%
  \stepcounter{popsec}\setcounter{popsub}{0}%
  \par\vspace{16pt}\noindent
  \begin{minipage}{\linewidth}
  \tikz[baseline=-0.7ex]\node[draw=popink,line width=1.6pt,fill=poporange,rounded corners=4pt,minimum size=22pt,inner sep=0]{\popdisplay\fontsize{12}{12}\selectfont\color{popcream}\thepopsec};%
  \hspace{8pt}{\popdisplay\fontsize{15}{17}\selectfont\color{popink}\MakeUppercase{#1}}\par
  \vspace{4pt}\noindent{\color{poporange}\rule{36pt}{3.6pt}}
  \end{minipage}\par\nopagebreak\vspace{8pt}
}
\newcounter{popsub}
\newcommand{\courssub}[1]{%
  \stepcounter{popsub}%
  \par\vspace{9pt}\noindent{\popxbold\fontsize{11.5}{13}\selectfont\color{popink}{\color{poporange}\thepopsec.\thepopsub}\hspace{6pt}#1}\par\nopagebreak\vspace{4pt}
}

% ============================================================
% Boîtes pédagogiques — même recette : fond blanc, bordure épaisse colorée,
% ombre dure encre, badge plein. Titre optionnel [#1].
% ============================================================
\tcbset{popbase/.style={breakable,enhanced,colback=white,boxrule=2.3pt,arc=9pt,
  shadow={3pt}{-3pt}{0pt}{popink},left=14pt,right=14pt,top=11pt,bottom=11pt,before skip=11pt,after skip=15pt,
  fontupper=\popsemi\fontsize{10.6}{14.5}\selectfont\color{popink},
  fontlower=\popsemi\fontsize{10.6}{14.5}\selectfont\color{popink}}}
% Coche dessinée (le glyphe ✓ n'existe pas dans les polices du document)
\newcommand{\popcheck}[1]{\tikz[baseline=-0.5ex]\draw[#1,line width=1.6pt,line cap=round,line join=round](-0.13,0.0)--(-0.04,-0.09)--(0.13,0.1);}
\providecommand{\coche}{\popcheck{popgreen}}
\providecommand{\resultat}[1]{\par\smallskip\noindent\fcolorbox{popgreen}{popgreenL}{\parbox{\dimexpr\linewidth-2\fboxsep-2\fboxrule\relax}{\textbf{Résultat.} #1}}\par\smallskip}
\newcommand{\popboxhead}[3]{\poppill{#1}{#2}{white}\ifx\relax#3\relax\else\hspace{6pt}{\popxbold\fontsize{10.6}{12}\selectfont\color{popink}#3}\fi\par\vspace{7pt}}

\newtcolorbox{aretenir}[1][]{popbase,colframe=popgreen,before upper={\popboxhead{À retenir}{popgreen}{#1}}}
\newtcolorbox{exemplebox}[1][]{popbase,colframe=poporange,before upper={\popboxhead{Exemple}{poporange}{#1}}}
\newtcolorbox{definition}[1][]{popbase,colframe=popblue,before upper={\popboxhead{Définition}{popblue}{#1}}}
\newtcolorbox{propriete}[1][]{popbase,colframe=poppurple,before upper={\popboxhead{Propriété}{poppurple}{#1}}}
\newtcolorbox{methode}[1][]{popbase,colframe=popgold,before upper={\popboxhead{Méthode}{popgold}{#1}}}
\newtcolorbox{attention}[1][]{popbase,colframe=poppink,before upper={\popboxhead{Attention}{poppink}{#1}}}
\newtcolorbox{experience}[1][]{popbase,colframe=popblue,colback=popblueL,before upper={\popboxhead{Expérience}{popblue}{#1}}}
% « Le savais-tu ? » : carte violette pleine (contexte, culture, Cameroun)
\newtcolorbox{savaistu}[1][]{popbase,colback=poppurple,colframe=popink,
  fontupper=\popsemi\fontsize{10.4}{14.2}\selectfont\color{white},
  before upper={\poppill{Le savais-tu\,?}{white}{poppurple}\ifx\relax#1\relax\else\hspace{6pt}{\popxbold\color{white}#1}\fi\par\vspace{7pt}}}
% Exercice résolu : énoncé en haut, \tcblower puis la solution sur fond crème
\newcounter{popexo}\newcounter{popexr}
\newtcolorbox{exoresolu}[1][]{popbase,colframe=poporange,colbacklower=popcream,
  segmentation style={popink,line width=1.2pt,dashed},
  before upper={\stepcounter{popexr}\popboxhead{Exercice résolu \thepopexr}{poporange}{#1}},
  before lower={\poppill{Solution}{popink}{popcream}\par\vspace{6pt}}}
% Exercice (non résolu en place) — la correction est dans la partie Corrigés
\newtcolorbox{exercice}[1][]{popbase,colframe=popink,
  before upper={\stepcounter{popexo}\popboxhead{Exercice \thepopexo}{popink}{#1}}}
% Corrigé
\newtcolorbox{corrige}[1][]{popbase,colframe=popgreen,colback=popgreenL,
  before upper={\popboxhead{Corrigé}{popgreen}{#1}}}
% Objectifs / prérequis / bilan
\newtcolorbox{objectifs}{popbase,colframe=popink,colback=poporangeL,
  before upper={\popboxhead{Objectifs de la leçon}{popink}{}}}
\newtcolorbox{prerequis}{popbase,colframe=popink,colback=popblueL,
  before upper={\popboxhead{Prérequis}{popblue}{}}}
\newtcolorbox{bilan}{popbase,colframe=popink,colback=white,boxrule=2.8pt,
  before upper={\popboxhead{Fiche bilan}{poporange}{L'essentiel à connaître pour l'examen}}}
% Figure encadrée avec légende
\newtcolorbox{popfigure}[1][]{enhanced,breakable=false,colback=white,colframe=popline,boxrule=1.4pt,arc=8pt,
  left=10pt,right=10pt,top=10pt,bottom=8pt,before skip=10pt,after skip=12pt,halign=center,
  lower separated=false,halign lower=center,
  fontlower=\popsemi\fontsize{9}{11.5}\selectfont\color{popmuted},
  #1}
\newcommand{\legende}[1]{\par\vspace{6pt}{\popsemi\fontsize{9}{11.5}\selectfont\color{popmuted}#1}}

% ============================================================
% Signature OnBuch+
% ============================================================
\newcommand{\poplogoinline}[1]{{\poplogo\fontsize{#1}{#1}\selectfont\color{popink}OnBuch\color{poporange}+}}
\newcommand{\signatureonbuch}{%
  \begin{tcolorbox}[enhanced,colback=popink,colframe=popink,boxrule=2.2pt,arc=10pt,
      drop shadow=poporange,left=14pt,right=14pt,top=10pt,bottom=10pt,before skip=0pt,after skip=0pt]
    \tikz[baseline=-0.7ex]\node[circle,fill=popgreen,draw=popcream,line width=1.3pt,minimum size=22pt,inner sep=0]{\popcheck{white}};%
    \hspace{9pt}\begin{minipage}[c]{0.62\linewidth}
      {\popxbold\fontsize{12}{13}\selectfont\color{popcream}Signé\ \textbullet\ {\poplogo OnBuch\color{poporange}+}}\par\vspace{2pt}
      {\raggedright\popsemi\fontsize{8}{10.5}\selectfont\color{popline}\MakeUppercase{Équipe pédagogique OnBuch+}\\\MakeUppercase{Conforme au programme officiel MINESEC}\par}
    \end{minipage}\hfill
    {\poplogo\fontsize{22}{22}\selectfont\color{popcream}OnBuch\color{poporange}+}
  \end{tcolorbox}%
}

% ============================================================
% Page de garde
% #1 titre · #2 matière · #3 niveau · #4 module · #5 n° leçon · #6 durée ·
% #7 description / objectifs (texte ou itemize)
% ============================================================
\newcommand{\pagegarde}[7]{%
  \thispagestyle{empty}
  \newgeometry{a4paper,margin=1.7cm,top=1.6cm,bottom=1.4cm}%
  \begin{tikzpicture}[remember picture,overlay]
    \fill[poporange!18] ([xshift=-3.2cm,yshift=-3.6cm]current page.north east) circle (4.6cm);
    \fill[poppurple!14] ([xshift=2.4cm,yshift=7.6cm]current page.south west) circle (3.8cm);
  \end{tikzpicture}%
  {\poplogo\fontsize{30}{30}\selectfont\color{popink}OnBuch\color{poporange}+}\hfill
  \raisebox{0.6ex}{\poppill{Cours signé \textbullet\ Premium}{popink}{popcream}}\par
  \vspace{1pt}\popsquiggle{poppurple}\par
  \vspace{8pt}
  \begin{tcolorbox}[enhanced,colback=poporange,colframe=popink,boxrule=2.8pt,arc=13pt,
      drop shadow=popink,left=18pt,right=18pt,top=12pt,bottom=13pt,after skip=10pt]
    \poppill{#2 \textbullet\ #3}{popink}{popcream}\hspace{5pt}\poppill{#4}{white}{popink}\par\vspace{8pt}
    {\popdisplay\fontsize{14}{14}\selectfont\color{popink}LEÇON #5}\par\vspace{4pt}
    {\raggedright\hyphenpenalty=10000\exhyphenpenalty=10000\popdisplay\fontsize{25}{27}\selectfont\color{popcream}\MakeUppercase{#1}\par}
    \vspace{8pt}
    {\popxbold\fontsize{9.5}{9.5}\selectfont\color{popink}\MakeUppercase{Durée conseillée : #6}}
  \end{tcolorbox}
  \begin{tcolorbox}[enhanced,colback=white,colframe=popink,boxrule=2.2pt,arc=10pt,
      drop shadow=popink,left=16pt,right=16pt,top=10pt,bottom=10pt,after skip=8pt]
    \poppill{Dans cette leçon}{poppurple}{white}\par\vspace{5pt}
    \popsemi\fontsize{9.3}{12.6}\selectfont\color{popink}#7
  \end{tcolorbox}
  \noindent\begin{minipage}[t]{0.487\linewidth}
    \begin{tcolorbox}[enhanced,colback=popgreen,colframe=popink,boxrule=2.2pt,arc=10pt,
        drop shadow=popink,left=11pt,right=11pt,top=10pt,bottom=10pt,height=3.1cm,valign=center]
      \tcbox[colback=white,colframe=popink,boxrule=1.3pt,arc=5pt,boxsep=0pt,nobeforeafter,
        left=3pt,right=3pt,top=3pt,bottom=3pt,valign=center]{\qrcode[height=1.9cm]{https://play.google.com/store/apps/details?id=com.luvvix.onbuch&hl=fr}}%
      \hspace{8pt}\begin{minipage}[c]{0.5\linewidth}
        {\popdisplay\fontsize{11}{12.5}\selectfont\color{white}\MakeUppercase{Télécharge OnBuch}}\par\vspace{4pt}
        {\raggedright\popsemi\fontsize{8.4}{11}\selectfont\color{white}Gratuit \textbullet\ Google Play\\Tuteur IA, annales, concours, résultats\par}
      \end{minipage}
    \end{tcolorbox}
  \end{minipage}\hfill
  \begin{minipage}[t]{0.487\linewidth}
    \begin{tcolorbox}[enhanced,colback=poppurple,colframe=popink,boxrule=2.2pt,arc=10pt,
        drop shadow=popink,left=11pt,right=11pt,top=10pt,bottom=10pt,height=3.1cm,valign=center]
      \tikz[baseline=-0.7ex]\node[circle,fill=white,draw=popink,line width=1.3pt,minimum size=1.6cm,inner sep=0]{\popdisplay\fontsize{24}{24}\selectfont\color{poppurple}+};%
      \hspace{8pt}\begin{minipage}[c]{0.6\linewidth}
        {\popdisplay\fontsize{11}{12.5}\selectfont\color{white}\MakeUppercase{Passe à OnBuch+}}\par\vspace{4pt}
        {\raggedright\popsemi\fontsize{8.4}{11}\selectfont\color{white}Tous les cours signés, Tuteur IA illimité, fiches PDF — onglet OnBuch+ de l'app\par}
      \end{minipage}
    \end{tcolorbox}
  \end{minipage}
  \vspace{8pt}
  \popcontenu
  \vfill
  \signatureonbuch
  \restoregeometry
  \newpage
}

% Rangée « Ce cours contient » (page de garde)
\newcommand{\poptile}[3]{% #1 couleur #2 gros texte #3 légende
  \begin{tcolorbox}[enhanced,colback=#1,colframe=popink,boxrule=2pt,arc=9pt,shadow={2.5pt}{-2.5pt}{0pt}{popink},
      left=6pt,right=6pt,top=9pt,bottom=9pt,halign=center,valign=center,height=1.75cm,width=\linewidth,nobeforeafter]
    {\popdisplay\fontsize{12.5}{13}\selectfont\color{white}\MakeUppercase{#2}}\par\vspace{3pt}
    {\centering\popsemi\fontsize{7.8}{9.5}\selectfont\color{white}#3\par}
  \end{tcolorbox}}
\newcommand{\popcontenu}{%
  \par\noindent{\popxboldls\fontsize{7.5}{7.5}\selectfont\color{popmuted}CE COURS SIGNÉ CONTIENT}\par\vspace{7pt}
  \noindent\begin{minipage}[t]{0.235\linewidth}\poptile{popblue}{Cours}{complet, illustré, pas à pas}\end{minipage}\hfill
  \begin{minipage}[t]{0.235\linewidth}\poptile{poporange}{Méthodes}{et exercices résolus}\end{minipage}\hfill
  \begin{minipage}[t]{0.235\linewidth}\poptile{poppink}{Exercices}{type examen + corrigés}\end{minipage}\hfill
  \begin{minipage}[t]{0.235\linewidth}\poptile{popgreen}{Fiche bilan}{l'essentiel à retenir}\end{minipage}\par}

% Sommaire
\newtcolorbox{sommaire}{enhanced,colback=white,colframe=popink,boxrule=2.2pt,arc=10pt,
  drop shadow=popink,left=18pt,right=18pt,top=13pt,bottom=13pt,before skip=6pt,after skip=20pt,
  before upper={\poppill{Sommaire}{poppurple}{white}\par\vspace{9pt}\popsemi\fontsize{10.6}{14.5}\selectfont\color{popink}}}

% Signature de fin de cours — poussée tout en bas de la dernière page
\newcommand{\findecours}{%
  \par\vfill\nopagebreak
  \begin{center}\popsquiggle{poporange}\end{center}\vspace{4pt}
  \signatureonbuch
}
\AtBeginDocument{\def\llbracket{\mathopen{[\mkern-3.5mu[}}\def\rrbracket{\mathclose{]\mkern-3.5mu]}}} % intervalles d'entiers (glyphe ⟦ absent de la police)
% Glyphes absents de Fira Math : redéfinis à partir de glyphes disponibles
\AtBeginDocument{%
  \def\setminus{\mathbin{\backslash}}\let\smallsetminus\setminus
  \def\vdots{\mathord{\vbox{\baselineskip3.2pt\lineskiplimit0pt\kern4pt\hbox{.}\hbox{.}\hbox{.}}}}%
  \def\ddots{\mathinner{\mkern1mu\raise6.5pt\hbox{.}\mkern2mu\raise3.6pt\hbox{.}\mkern2mu\raise0.7pt\hbox{.}\mkern1mu}}%
  \def\triangle{\mathord{\scalebox{0.9}{$\symup{\Delta}$}}}%
  \def\nabla{\mathord{\raisebox{\depth}{\scalebox{1}[-1]{$\symup{\Delta}$}}}}%
  \def\checkmark{\popcheck{popdark}}%
  \def\star{\mathbin{\ast}}\def\bigstar{\mathord{\ast}}%
  \def\diamond{\mathbin{\scalebox{0.6}{$\Diamond$}}}%
  \def\bigcup{\mathop{\mathchoice{\raisebox{-0.3em}{\scalebox{1.7}{$\cup$}}}{\scalebox{1.25}{$\cup$}}{\cup}{\cup}}\displaylimits}%
  \def\bigcap{\mathop{\mathchoice{\raisebox{-0.3em}{\scalebox{1.7}{$\cap$}}}{\scalebox{1.25}{$\cap$}}{\cap}{\cap}}\displaylimits}%
  \def\lesseqqgtr{\mathrel{\lessgtr}}\def\bigoplus{\mathop{\oplus}\displaylimits}%
}
\providecommand{\ssi}{si et seulement si} % abréviation fréquente
\AtBeginDocument{\providecommand{\wideparen}{\overparen}} % arc de cercle

%% FIGFIX-BEGIN v3 — correctifs globaux des figures (docs/onbuch-generation/tools/figfix)
\usepackage{adjustbox}\usepackage{etoolbox}
% toute figure TikZ/pgfplots plus large que la ligne est réduite à la largeur disponible
\BeforeBeginEnvironment{tikzpicture}{\begin{adjustbox}{max width=\linewidth}}
\AfterEndEnvironment{tikzpicture}{\end{adjustbox}}
% légendes pgfplots lisibles par-dessus les courbes
\pgfplotsset{every axis/.append style={legend style={fill=white,fill opacity=0.92,text opacity=1,draw=black!15,inner sep=3pt}}}
% étiquettes d'arêtes (syntaxe edge["..."]) avec fond blanc
\tikzset{every edge quotes/.append style={fill=white,inner sep=1.2pt}}
% bulles de cartes mentales : texte à la ligne dans la bulle, bulle un peu plus grande
\tikzset{every concept/.append style={text width=2.0cm,align=center,minimum size=2.3cm,inner sep=2pt}}
%% FIGFIX-END
\begin{document}

\pagegarde{Lutte contre l'érosion}{SVT}{1ère C}{Module 3 : L'Éducation à l'Environnement et au Développement Durable}{NCT08}{4 h}{Cette leçon étudie les types d'érosion (hydrique et éolienne), leurs facteurs favorisants et leurs conséquences sur les sols et l'environnement. Elle présente ensuite les techniques de lutte contre l'érosion (reboisement, terrasses, cultures en courbes de niveau) applicables aux réalités agricoles du Cameroun.
\vspace{4pt}
\begin{multicols}{2}\small
\begin{itemize}
\item À la fin de cette leçon, je sais définir l'érosion et distinguer l'érosion hydrique de l'érosion éolienne.
\item À la fin de cette leçon, je sais identifier les facteurs naturels (pente, climat) et anthropiques (déforestation, pratiques agricoles) qui favorisent l'érosion dans une situation donnée.
\item À la fin de cette leçon, je sais décrire les conséquences de l'érosion sur les sols, les rendements agricoles et l'environnement.
\item À la fin de cette leçon, je sais interpréter un schéma ou une photographie montrant un phénomène érosif (ravin, sol dénudé, dune).
\item À la fin de cette leçon, je sais expliquer le principe des techniques de lutte : reboisement, terrasses, cultures en courbes de niveau, bandes enherbées.
\item À la fin de cette leçon, je sais proposer des techniques de lutte adaptées à un milieu donné (relief, climat, type de culture).
\end{itemize}\end{multicols}}

\begin{sommaire}
\begin{multicols}{2}
\begin{enumerate}
\item Généralités et définition de l'érosion
\item Les types d'érosion
\item Les facteurs favorisants l'érosion
\item Les conséquences de l'érosion
\item Les techniques de lutte contre l'érosion
\item Interprétation de documents et étude de cas
\item Activité d'intégration
\item Méthodes et exercices résolus
\item Exercices
\item Corrigés des exercices
\item Fiche bilan
\end{enumerate}
\end{multicols}
\end{sommaire}

\begin{objectifs}
\begin{itemize}
\item À la fin de cette leçon, je sais définir l'érosion et distinguer l'érosion hydrique de l'érosion éolienne.
\item À la fin de cette leçon, je sais identifier les facteurs naturels (pente, climat) et anthropiques (déforestation, pratiques agricoles) qui favorisent l'érosion dans une situation donnée.
\item À la fin de cette leçon, je sais décrire les conséquences de l'érosion sur les sols, les rendements agricoles et l'environnement.
\item À la fin de cette leçon, je sais interpréter un schéma ou une photographie montrant un phénomène érosif (ravin, sol dénudé, dune).
\item À la fin de cette leçon, je sais expliquer le principe des techniques de lutte : reboisement, terrasses, cultures en courbes de niveau, bandes enherbées.
\item À la fin de cette leçon, je sais proposer des techniques de lutte adaptées à un milieu donné (relief, climat, type de culture).
\item À la fin de cette leçon, je sais sensibiliser ma communauté sur la protection des sols contre l'érosion.
\end{itemize}
\end{objectifs}

\begin{prerequis}
\begin{itemize}
\item Notion de sol : composition, horizons et fertilité (classes antérieures)
\item Le cycle de l'eau et les précipitations
\item Notions de relief : pente, versant, courbe de niveau (lecture simple de carte)
\item Les activités humaines et l'environnement : déforestation, agriculture (Module 3, leçons précédentes)
\item Notion de développement durable vue en début de module
\end{itemize}
\end{prerequis}

\newpage

\coursec{Généralités et définition de l'érosion}

\textbf{Situation d'accroche.} Après les fortes pluies de septembre, les habitants du quartier Nkolbisson à Yaoundé constatent que l'eau de ruissellement emporte la terre rouge des collines, creuse des ravins et dépose du sable dans les maisons en contrebas. Sur les plateaux de l'Ouest, des agriculteurs voient chaque année leur couche de terre fertile diminuer et leurs rendements de maïs chuter. Au Nord, le vent de l'harmattan soulève la poussière et dénude les champs.

\textbf{Problématique.} Comment expliquer ces phénomènes ? D'où vient cette terre qui s'en va, et où va-t-elle ? Quelles techniques les populations peuvent-elles utiliser pour protéger leurs sols ? Cette leçon nous donne les outils pour comprendre l'\cle{érosion} et proposer des solutions adaptées au contexte camerounais. Mais avant de lutter contre un phénomène, il faut d'abord le définir et comprendre son mécanisme : c'est l'objet de cette première partie.

\courssub{Qu'observe-t-on ? Une première approche du phénomène}

\textbf{Observation 1.} Sur un chemin de terre battue non revêtu, après une pluie, on voit de petits sillons parallèles creusés par l'eau. Quelques jours plus tard, à la pluie suivante, ces sillons sont plus profonds et plus larges. En bas du chemin, on trouve un amas de sable et de graviers.

\textbf{Observation 2.} En saison sèche à Maroua, lorsque l'harmattan souffle fort, l'air devient jaunâtre, chargé de poussière. Les champs nus perdent leur fine couche de terre superficielle, qui s'accumule contre les clôtures et les murs.

\textbf{Question.} Ces deux observations ont-elles quelque chose en commun ? Oui : dans les deux cas, des particules du sol sont \textbf{arrachées} à leur place d'origine, \textbf{déplacées} par un agent (l'eau ou le vent), puis \textbf{abandonnées} ailleurs. Ce phénomène porte un nom : l'érosion.

\begin{definition}[Érosion]
L'\cle{érosion} est l'ensemble des processus d'\textbf{usure}, d'\textbf{enlèvement} et de \textbf{transport} des particules du sol et des roches sous l'action d'agents externes, principalement \textbf{l'eau} (pluie, ruissellement, cours d'eau) et \textbf{le vent}.
\end{definition}

\begin{definition}[Ruissellement]
Le \cle{ruissellement} est l'écoulement de l'eau de pluie à la surface du sol, lorsque celle-ci ne peut pas s'infiltrer (sol saturé, imperméable, dénudé ou en pente). C'est le principal agent de l'érosion hydrique sous les climats pluvieux comme ceux du Centre et de l'Ouest camerounais.
\end{definition}

\begin{attention}[Érosion ou lessivage : ne pas confondre]
L'\textbf{érosion} est un enlèvement \textbf{horizontal} de particules solides : la terre quitte littéralement le champ, emportée par l'eau ou le vent. Le \textbf{lessivage} est un entraînement \textbf{vertical}, en profondeur, des éléments nutritifs dissous (sels minéraux) par l'eau qui s'infiltre : la terre reste en place, mais elle s'appauvrit. Au Probatoire, confondre ces deux notions est une erreur classique : retiens que l'érosion déplace le \emph{sol}, le lessivage déplace les \emph{substances dissoutes} dans le sol.
\end{attention}

\courssub{Les agents érosifs : l'eau et le vent}

L'érosion est provoquée par des \textbf{agents externes} au sol. Les deux agents majeurs sont :

\begin{itemize}
\item \textbf{L'eau}, sous deux formes complémentaires :
\begin{itemize}
\item l'\textbf{impact des gouttes de pluie} : une goutte de pluie tropicale frappe le sol nu avec une force suffisante pour éclater les agrégats de terre et projeter de fines particules à plusieurs dizaines de centimètres de distance ; c'est l'\emph{érosion par splash} (ou érosion par impact) ;
\item le \textbf{ruissellement} : l'eau qui s'écoule en nappe ou en filets arrache les particules détachées et les entraîne vers le bas des pentes.
\end{itemize}
\item \textbf{Le vent} : dans les zones sèches et dénudées (Grand Nord camerounais), le vent soulève les particules fines (sables, limons, poussières) et les transporte parfois sur de longues distances. Les grains de sable, plus lourds, roulent ou sautent à la surface (\emph{saltation}) et usent le sol par abrasion.
\end{itemize}

\begin{exemplebox}[Deux visages de l'érosion au Cameroun]
À Foumban (Ouest), ce sont les pluies intenses de la saison humide qui emportent la terre rouge des champs de maïs plantés sur les pentes. À Kousséri (Extrême-Nord), c'est l'harmattan de décembre à février qui balaie les champs nus après la récolte et emporte la couche fertile en poussière. Même résultat — un sol appauvri — mais des agents différents, donc des techniques de lutte différentes.
\end{exemplebox}

\courssub{Les trois étapes du phénomène érosif}

Quel que soit l'agent, l'érosion se déroule toujours selon trois étapes successives qu'il faut savoir nommer et repérer sur un schéma ou une photo.

\begin{enumerate}
\item \textbf{Le détachement (ou arrachement)} : les particules du sol sont désagrégées et arrachées à la masse du sol, sous l'impact des gouttes de pluie, de la force du courant ou du vent. Cette étape est d'autant plus facile que le sol est nu, sec ou mal structuré.
\item \textbf{Le transport} : les particules détachées sont entraînées par l'eau de ruissellement (en suspension ou en roulant au fond des filets d'eau) ou par le vent (en suspension pour les poussières, par saltation pour les sables).
\item \textbf{Le dépôt (ou sédimentation)} : lorsque la vitesse de l'agent diminue (bas de pente, zone plate, obstacle), l'eau ou le vent ne peut plus porter sa charge : les particules se déposent. On parle de \cle{sédimentation} lorsque le dépôt se fait dans un cours d'eau, un lac ou un bas-fond.
\end{enumerate}

\begin{popfigure}
\begin{center}
\begin{tikzpicture}[scale=1.0, font=\small]
% Versant en coupe : remplissage exact sous la courbe de pente
\fill[popgoldL] (0,3.4) -- (4,2.6) -- (8,1.4) -- (10,0.9) -- (10,0) -- (0,0) -- cycle;
% Courbe de pente (surface du sol)
\draw[popdark, thick] (0,3.4) -- (4,2.6) -- (8,1.4) -- (10,0.9);
% Pluie
\foreach \x in {0.6,1.3,2.0,2.7}{
  \draw[popblue, thick, ->] (\x,5.3) -- (\x-0.25,4.6);}
\node[popblue] at (1.6,5.6) {Pluie (impact des gouttes)};
% Filets d'eau (ruissellement)
\draw[popblue, thick, ->] (1.2,3.35) .. controls (2.5,2.9) .. (3.8,2.55);
\draw[popblue, thick, ->] (4.2,2.5) .. controls (6,1.95) .. (7.8,1.35);
\node[popblue] at (5.9,2.6) {Ruissellement};
% Particules transportées (points noirs)
\foreach \p in {(2.2,3.05),(3.1,2.8),(5.0,2.1),(6.3,1.75),(7.2,1.5)}{
  \fill[popdark] \p circle (0.05);}
% Dépôt en bas de pente (zone orange)
\fill[poporange!60] (8.2,0.9) .. controls (8.8,1.05) .. (9.6,0.9) -- (9.6,0.55) .. controls (8.8,0.7) .. (8.2,0.55) -- cycle;
\draw[popdark] (8.2,0.9) .. controls (8.8,1.05) .. (9.6,0.9);
% Annotations des trois étapes (flèches fines vers légendes)
\draw[popline, thin] (1.6,3.3) -- (1.6,6.1);
\node[align=center, popink] at (1.6,6.5) {\textbf{1. Détachement}\\ particules arrachées};
\draw[popline, thin] (5.5,2.15) -- (5.5,4.9);
\node[align=center, popink] at (5.5,5.3) {\textbf{2. Transport}\\ par l'eau de ruissellement};
\draw[popline, thin] (8.9,0.75) -- (8.9,-0.6);
\node[align=center, popink] at (8.9,-1.0) {\textbf{3. Dépôt}\\ (sédimentation en bas de pente)};
% Légende sol
\node[popdark] at (2.5,1.0) {Sol du versant};
\end{tikzpicture}
\end{center}
\legende{Coupe schématique d'un versant montrant les trois étapes de l'érosion hydrique : (1) détachement des particules sous l'impact des gouttes de pluie, (2) transport par le ruissellement le long de la pente, (3) dépôt des sédiments en contrebas lorsque la vitesse de l'eau diminue.}
\end{popfigure}

\begin{experience}[Observer les trois étapes dans la cour du lycée]
\textbf{Matériel} : un arrosoir, une planche inclinée recouverte de terre fine, une bâche ou une cuvette en bas de pente.\\
\textbf{Protocole} : verser lentement le contenu de l'arrosoir en haut de la planche ; observer ce qui se passe sur la surface, le long de la pente et dans la cuvette. Recommencer en versant plus vite.\\
\textbf{Observations} : l'eau éclate la surface de la terre (détachement), des filets boueux descendent la planche (transport), et un dépôt de sable et de limon s'accumule dans la cuvette (dépôt). Plus le débit est fort, plus la quantité de terre transportée est grande.\\
\textbf{Interprétation} : la force du ruissellement détermine la quantité de particules érodées ; c'est pourquoi les pluies violentes tropicales sont si érosives, et pourquoi ralentir l'eau est la clé de la lutte contre l'érosion.
\end{experience}

\courssub{Érosion naturelle et érosion accélérée}

\textbf{Observation.} Le Mont Cameroun, comme toutes les montagnes, s'use lentement sous l'action de la pluie et du vent depuis des millions d'années, même là où aucun homme ne vit. En revanche, une colline déboisée pour un champ peut perdre en quelques années l'équivalent de plusieurs siècles d'érosion naturelle.

\begin{definition}[Érosion naturelle et érosion accélérée]
L'\cle{érosion naturelle} (ou géologique) est l'usure lente des sols et des reliefs par les agents climatiques, en l'absence d'intervention humaine. Elle se mesure en fractions de millimètre par an et est compensée par la formation naturelle du sol à partir de la roche.\\
L'\cle{érosion accélérée} est l'érosion fortement amplifiée par les activités humaines (déforestation, cultures sur pente, surpâturage, feux de brousse). Elle peut atteindre plusieurs tonnes de terre perdues par hectare et par an, bien plus vite que le sol ne se régénère.
\end{definition}

\begin{aretenir}[L'essentiel à retenir]
L'érosion naturelle existe sans l'homme : elle fait partie du fonctionnement normal de la planète et sculpte les paysages à l'échelle des temps géologiques. Mais les activités humaines l'\textbf{accélèrent fortement} : en supprimant la couverture végétale qui protège le sol, l'homme transforme un phénomène lent et équilibré en un phénomène rapide et destructeur. C'est contre cette érosion accélérée que portent les techniques de lutte.
\end{aretenir}

\begin{popfigure}
\begin{center}
\begin{tikzpicture}[scale=0.95, font=\small]
% Panneau gauche : forêt (érosion naturelle)
\begin{scope}
% Sol
\fill[popgreenL] (0,2.6) -- (5,2.2) -- (5,0) -- (0,0) -- cycle;
\draw[popdark, thick] (0,2.6) -- (5,2.2);
% Arbres (troncs + houppiers)
\foreach \x in {0.7,1.8,3.0,4.1}{
  \draw[popdark, thick] (\x,2.5) -- (\x,3.3);
  \fill[popgreen] (\x,3.55) circle (0.32);}
% Pluie
\foreach \x in {1.0,2.3,3.6}{
  \draw[popblue, ->] (\x,4.6) -- (\x-0.2,4.0);}
\node[popblue] at (2.5,4.85) {Pluie};
% Petit filet d'eau discret
\draw[popblue, ->] (1.5,2.55) .. controls (3,2.3) .. (4.6,2.15);
\node[align=center] at (2.5,-0.7) {\textbf{Érosion naturelle}\\ sol protégé par la forêt\\ perte très faible et lente};
\end{scope}
% Panneau droit : sol nu (érosion accélérée)
\begin{scope}[xshift=7cm]
% Sol
\fill[popgoldL] (0,2.6) -- (5,2.2) -- (5,0) -- (0,0) -- cycle;
\draw[popdark, thick] (0,2.6) -- (5,2.2);
% Souche (trace de déforestation)
\draw[popdark, thick] (1.0,2.52) -- (1.0,2.9);
\draw[popdark, thick] (0.85,2.9) -- (1.15,2.9);
% Pluie
\foreach \x in {1.0,2.3,3.6}{
  \draw[popblue, ->] (\x,4.6) -- (\x-0.2,4.0);}
\node[popblue] at (2.5,4.85) {Pluie};
% Gros ruissellement (très épais)
\draw[popblue, very thick, ->] (0.8,2.55) .. controls (2.5,2.25) .. (4.7,2.0);
\draw[popblue, very thick, ->] (1.2,2.5) .. controls (3,2.2) .. (4.9,1.9);
% Particules transportées
\foreach \p in {(2.0,2.35),(3.0,2.15),(4.0,2.0)}{
  \fill[popdark] \p circle (0.06);}
% Ravin (creux en V)
\draw[popdark] (2.6,2.42) .. controls (2.75,2.0) .. (2.7,1.6);
\draw[popdark] (3.1,2.38) .. controls (2.95,2.0) .. (3.0,1.6);
\node[popink] at (3.9,1.5) {Ravin};
\node[align=center] at (2.5,-0.7) {\textbf{Érosion accélérée}\\ sol dénudé (déforestation)\\ perte rapide et massive};
\end{scope}
\end{tikzpicture}
\end{center}
\legende{Comparaison de deux versants soumis à la même pluie. À gauche, la couverture végétale intercepte les gouttes et freine le ruissellement : l'érosion naturelle est très faible. À droite, le sol dénudé subit l'impact direct des gouttes et un ruissellement violent qui creuse des ravins : c'est l'érosion accélérée par les activités humaines.}
\end{popfigure}

\begin{savaistu}[Des ravins qui grandissent à vue d'œil]
Sur les plateaux de l'Adamaoua, certains ravins peuvent s'approfondir de \textbf{plusieurs mètres en une seule saison des pluies} lorsque le sol est dénudé par les feux de brousse et le surpâturage. Ces ravins, appelés localement «~canyons~», coupent parfois les pistes rurales et isolent des villages entiers. À l'échelle du pays, on estime que plusieurs tonnes de terre fertile par hectare sont perdues chaque année sur les terres cultivées en pente non protégées : c'est l'avenir agricole du Cameroun qui part littéralement à la rivière.
\end{savaistu}

\courssub{Le vocabulaire de l'érosion}

Pour décrire correctement un document (photo, schéma) au Probatoire, il faut maîtriser un vocabulaire précis.

\begin{definition}[Mots-clés du phénomène érosif]
\begin{itemize}
\item \textbf{Ruissellement} : écoulement de l'eau de pluie à la surface du sol, sans infiltration.
\item \textbf{Ravinement} : creusement de sillons profonds (ravins) par un ruissellement concentré et violent ; c'est la forme la plus spectaculaire de l'érosion hydrique.
\item \textbf{Sol dénudé} : sol privé de sa couverture végétale (forêt coupée, champ après récolte, terrain brûlé) ; il est directement exposé aux agents érosifs.
\item \textbf{Sédimentation} : dépôt des particules transportées lorsque l'agent érosif perd de sa vitesse (bas de pente, lit d'une rivière, lac, retenue de barrage).
\end{itemize}
\end{definition}

\begin{exoresolu}[Identifier les étapes de l'érosion sur une description]
\textbf{Énoncé.} Après une pluie à Bafoussam, on observe sur une colline cultivée : (a) des gouttelettes de boue projetées autour des plants de maïs ; (b) des filets d'eau rouge descendant la pente ; (c) une couche de sable rouge recouvrant la route en bas de la colline. Pour chaque observation, indique l'étape de l'érosion concernée et justifie.
\tcblower
\textbf{Solution.} (a) Il s'agit du \textbf{détachement} : l'impact des gouttes de pluie éclate les agrégats du sol nu et projette les particules. (b) Il s'agit du \textbf{transport} : le ruissellement entraîne les particules détachées vers le bas de la pente, d'où la couleur rouge de l'eau, chargée de terre. (c) Il s'agit du \textbf{dépôt} (sédimentation) : arrivée sur la route plane, l'eau ralentit et abandonne les particules qu'elle transportait. La colline perd donc sa terre fertile au profit de la route, ce qui illustre à la fois l'appauvrissement du sol en amont et les nuisances en aval.
\end{exoresolu}

\begin{exercice}[À toi de jouer]
Dans le quartier Nkolbisson, une photo montre une maison en contrebas envahie de sable rouge après une pluie, et en amont une colline déboisée pour la construction. 1) Nomme le phénomène responsable. 2) Identifie l'agent érosif. 3) Explique en trois étapes comment le sable est arrivé dans la maison. 4) Pourquoi la déforestation de la colline a-t-elle aggravé le phénomène ?
\end{exercice}

\begin{corrige}[Exercice : Nkolbisson]
1) Il s'agit de l'érosion hydrique (avec ravinement probable sur la pente). 2) L'agent érosif est l'eau : impact des gouttes de pluie puis ruissellement. 3) Détachement : les gouttes frappent le sol nu de la colline et arrachent les particules de terre rouge ; transport : le ruissellement les entraîne vers le bas ; dépôt : en contrebas, l'eau ralentit et dépose le sable dans et autour de la maison. 4) La végétation interceptait les gouttes, retenait le sol par ses racines et favorisait l'infiltration ; en la supprimant, l'homme a transformé une érosion naturelle lente en érosion accélérée destructrice.
\end{corrige}

\begin{aretenir}[Bilan de la section]
L'érosion est l'usure, l'enlèvement et le transport des particules du sol par l'eau et le vent. Elle comporte trois étapes : \textbf{détachement}, \textbf{transport}, \textbf{dépôt}. Elle existe naturellement, mais les activités humaines (déforestation, cultures sur pente) l'accélèrent fortement. Les mots-clés à maîtriser : ruissellement, ravinement, sol dénudé, sédimentation. La section suivante précisera les grands types d'érosion : hydrique et éolienne.
\end{aretenir}

\coursec{Généralités et définition de l'érosion}

L'\textbf{érosion des sols} est un processus naturel de modelé de la surface terrestre, mais sa \textbf{vitesse actuelle}, largement amplifiée par les activités humaines, en fait une menace majeure pour la sécurité alimentaire et la durabilité des écosystèmes au Cameroun.

\begin{definition}[Érosion des sols]
L'érosion est l'ensemble des phénomènes de \textbf{détachement}, de \textbf{transport} et de \textbf{déposition} des particules du sol (minéraux, matière organique) sous l'action d'agents naturels externes : l'\textbf{eau} (pluie, ruissellement, cours d'eau) et le \textbf{vent}. Elle se distingue de l'altération (météorisation) qui dégrade la roche \emph{in situ} sans transport significatif.
\end{definition}

\begin{aretenir}[Le cycle érosif : trois étapes indissociables]
\begin{enumerate}
    \item \textbf{Détachement (éclatement) :} Rupture des agrégats structuraux par l'énergie cinétique des gouttes de pluie (\emph{splash}) ou par cisaillement éolien. Les particules fines (argile, limon, MO) sont les premières mobilisées.
    \item \textbf{Transport :} Déplacement des particules détachées. L'eau les transporte en suspension, saltation ou traction (roulement) selon la vitesse du ruissellement. Le vent les transporte en suspension (poussières), saltation (sable) ou reptation (graviers).
    \item \textbf{Déposition (sédimentation) :} Arrêt du transport quand l'énergie de l'agent érosif diminue (rupture de pente, obstacle végétal, élargissement du lit, baisse de vitesse du vent). Cela forme des \textbf{cônes de déjection}, des \textbf{alluvions} ou des \textbf{dunes}.
\end{enumerate}
\textbf{Érosion géologique vs Érosion accélérée :} L'érosion géologique (naturelle, lente, $\approx$ 0,1--1 t/ha/an) façonne les paysages sur le temps long. L'érosion accélérée (anthropique, rapide, $>$ 10--100 t/ha/an) détruit le capital-sol en quelques décennies.
\end{aretenir}

\coursec{Les types d'érosion}

Deux grands types menacent les sols camerounais, selon l'agent moteur dominant : l'érosion \textbf{hydrique} (Sud, Ouest, Littoral, Adamaoua) et l'érosion \textbf{éolienne} (Nord, Extrême-Nord). Elles partagent un facteur aggravant commun : la disparition du \cle{couvert végétal}.

\courssub{L'érosion hydrique : la force de l'eau}

\begin{experience}[Observation de l'impact d'une goutte de pluie et du ruissellement sur sol nu]
\textbf{Matériel :} Bac de sol meuble non végétalisé (terre de surface tamisée), pulvérisateur d'eau (simulant la pluie), règle graduée, chronomètre, seau collecteur, balance de précision.\\
\textbf{Protocole :}
\begin{enumerate}
    \item Placer le bac horizontalement. Pulvériser de l'eau à faible intensité (simulant une pluie douce, $\approx$ 10 mm/h) pendant 2 minutes. Observer la surface du sol et mesurer la hauteur des éclaboussures.
    \item Incliner le bac à 15\textdegree (pente $\approx$ 27\%). Pulvériser à intensité modérée (pluie orageuse, $\approx$ 50 mm/h) pendant 5 minutes. Observer l'écoulement de l'eau (infiltration vs ruissellement) et l'aspect du sol (apparition de sillons).
    \item Recueillir l'eau de ruissellement dans le seau. Laisser décanter 30 min, puis mesurer le volume de sédiments ou la turbidité.
\end{enumerate}
\textbf{Observations :}
\begin{itemize}
    \item \textbf{Étape 1 (horizontal) :} Formation de micro-cratères d'impact. Éclaboussures de particules (sable, limon, argile) jusqu'à 30--50 cm de hauteur. La surface devient lisse, brillante et imperméable : formation d'une \textbf{croûte de battance} (scellement structural).
    \item \textbf{Étape 2 (incliné) :} L'eau ne s'infiltre plus (porosité bouchée par la battance) ; elle ruisselle en nappe fine, entraînant une couche uniforme de terre fine (\textbf{érosion en nappe}). Rapidement, le ruissellement se concentre dans des sillons parallèles qui s'approfondissent (\textbf{rigoles}). L'eau collectée est très trouble (charge solide élevée, $>$ 50 g/L).
\end{itemize}
\textbf{Interprétation :} L'énergie cinétique des gouttes de pluie détruit la structure superficielle du sol (\emph{splash}). Sur pente, l'eau de ruissellement devient l'agent de transport principal, passant d'un transport diffus (nappe) à un transport concentré (rigoles, puis ravins).
\textbf{Conclusion :} L'érosion hydrique est un processus en cascade : \emph{splash} $\rightarrow$ érosion en nappe $\rightarrow$ érosion en rigoles $\rightarrow$ érosion en ravins.
\end{experience}

\begin{popfigure}
\begin{tikzpicture}[scale=0.9, every node/.style={font=\small}]
    % Sol nu
    \fill[poporange!30] (0,0) rectangle (12,-0.5);
    \draw[thick] (0,0) -- (12,0);
    % Gouttes de pluie
    \foreach \x in {1,2.5,4,5.5,7,8.5,10,11.5} {
        \draw[popblue, thick, -{Stealth[length=3mm]}] (\x,4) -- (\x,1);
        \node[popblue, above] at (\x,4) {Goutte};
    }
    % Splash (éclaboussures)
    \foreach \x/\y/\r in {1/0.2/0.3, 1.2/-0.1/0.2, 0.8/-0.1/0.2, 2.5/0.2/0.3, 2.7/-0.1/0.2, 2.3/-0.1/0.2} {
        \draw[poporange, thick, -{Stealth[length=2mm]}] (\x,0) -- (\x+\r,\y);
        \draw[poporange, thick, -{Stealth[length=2mm]}] (\x,0) -- (\x-\r,\y);
    }
    \node[poporange, align=center] at (2,-1.2) {Splash (éclaboussures)\\(détachement)};
    
    % Ruissellement en nappe (flèches horizontales)
    \draw[popblue!70!black, thick, -{Stealth[length=3mm]}] (0.5,-0.3) -- (5.5,-0.3);
    \draw[popblue!70!black, thick, -{Stealth[length=3mm]}] (0.5,-0.4) -- (5.5,-0.4);
    \node[popblue!70!black, align=center] at (3,-1.2) {Ruissellement diffus\\(érosion en nappe)};
    
    % Rigoles (courbes creusées)
    \draw[popdark, thick] (6.5,0) .. controls (7,-0.3) and (8,-0.4) .. (8.5,-0.5);
    \draw[popdark, thick] (7.5,0) .. controls (8,-0.3) and (9,-0.4) .. (9.5,-0.5);
    \draw[popdark, thick, -{Stealth[length=3mm]}] (8.5,-0.5) -- (10,-0.5);
    \node[popdark, align=center] at (8,-1.2) {Concentration en rigoles\\(érosion linéaire)};
    
    % Ravins (plus large)
    \draw[poppurple, very thick] (10.5,0) .. controls (11,-0.6) and (11.5,-0.8) .. (12,-1.0);
    \draw[poppurple, very thick] (11.5,0) .. controls (12,-0.6) and (12.5,-0.8) .. (13,-1.0);
    \node[poppurple, align=center] at (11.5,-1.5) {Ravin\\(érosion massive)};
    
    % Légendes techniques
    \node[draw, fill=popcream, rounded corners, align=left, anchor=north west] at (0,4.5) {
        \textbf{Mécanismes :}\\
        1. Impact goutte $\rightarrow$ \emph{Splash}\\
        2. Ruissellement diffus $\rightarrow$ Nappe\\
        3. Concentration $\rightarrow$ Rigoles\\
        4. Élargissement $\rightarrow$ Ravins
    };
\end{tikzpicture}
\legende{Processus de l'érosion hydrique sur sol nu : de l'impact des gouttes (splash) à la formation de ravins en passant par l'érosion en nappe et en rigoles.}
\end{popfigure}

\begin{definition}[Érosion hydrique]
L'\textbf{érosion hydrique} est l'ensemble des processus de détachement, de transport et de déposition des particules du sol causés par l'action mécanique de l'eau (impact des gouttes de pluie et ruissellement de surface). Elle se manifeste sous trois formes successives selon l'énergie disponible :
\begin{itemize}
    \item \textbf{L'érosion en nappe (ou érosion diffuse) :} Enlèvement plus ou moins uniforme d'une fine couche de terre sur toute la surface du sol par un ruissellement laminaire. C'est la forme la plus \textbf{insidieuse} car peu visible, mais elle appauvrit durablement l'horizon arable (perte de 1 mm = 10--15 t/ha).
    \item \textbf{L'érosion en rigoles :} Le ruissellement se concentre dans des sillons peu profonds (quelques cm à dm) que le labour peut encore effacer morphologiquement. C'est le stade intermédiaire d'alerte.
    \item \textbf{L'érosion en ravins :} Les rigoles s'élargissent et s'approfondissent (plus de 30--50 cm, parois verticales), formant des entailles permanentes que le labour ne peut plus combler. Le réseau hydrographique se structure, entraînant une perte massive de sol et une dissection du paysage.
\end{itemize}
\end{definition}

\begin{aretenir}[Le mécanisme du \cle{Splash} (éclaboussure)]
L'énergie cinétique d'une goutte de pluie (diamètre 2--5 mm, vitesse de chute 5--9 m/s) projetée sur un sol nu pulvérise les agrégats du sol. Les particules fines (argile, limon, matière organique) sont projetées en l'air (jusqu'à 60 cm de haut, 1,5 m en horizontal) et bouchent les pores superficiels à leur retombée, formant une \textbf{croûte de battance} imperméable. Cela réduit drastiquement l'infiltration (conductivité hydraulique $\downarrow$) et déclenche le ruissellement. \textbf{Le couvert végétal (ou un paillis) amortit cet impact, préserve la porosité et la structure.}
\end{aretenir}

\courssub{L'érosion éolienne : la force du vent}

\begin{experience}[Simulation de l'érosion éolienne (soufflerie simple)]
\textbf{Matériel :} Deux bacs identiques : Bac A = sable sec fin (sol nu, simulant un sol sableux ferrugineux nu) ; Bac B = même sable recouvert de brindilles/herbes sèches (simulant un couvert végétal/résidus), sèche-cheveux (air froid, débit réglable) ou ventilateur, balance de précision (0,01 g), anémomètre.\\
\textbf{Protocole :}
\begin{enumerate}
    \item Peser les deux bacs ($m_0$). Mesurer la vitesse du vent à 30 cm du sol (cible : 15 km/h $\approx$ 4,2 m/s).
    \item Souffler à vitesse constante pendant 5 min sur chaque bac à distance fixe (30 cm).
    \item Peser à nouveau ($m_1$). Calculer $\Delta m = m_0 - m_1$. Observer la surface (défaut, dépôts) et le dépôt au vent (au-delà du bac).
\end{enumerate}
\textbf{Observations :}
\begin{itemize}
    \item \textbf{Bac A (Sol nu) :} Perte de masse significative ($\Delta m > 5$ g). Surface "pelée", grains grossiers résiduels apparents (\emph{lag gravel}). Nuage de poussière visible pendant le soufflage. Dépôt de fines particules (poussières) au vent.
    \item \textbf{Bac B (Sol végétalisé) :} Perte de masse négligeable ($\Delta m \approx 0$). Surface intacte, résidus en place.
\end{itemize}
\textbf{Interprétation :} Le vent exerce une contrainte de cisaillement ($\tau$) au sol. Si la vitesse de frottement $u_*$ dépasse le \textbf{seuil de fluidification} (environ 0,2--0,3 m/s, soit vent $\approx$ 4--5 m/s à 30 cm pour du sable fin), les particules sont arrachées. Le transport se fait par \emph{saltation} (sauts de grains 0,1--0,5 mm, énergie principale de l'érosion), \emph{reptation} (roulement grains $>$ 0,5 mm) et \emph{suspension} (poussières $<$ 0,1 mm portées sur km). La végétation brise la vitesse du vent au ras du sol (augmentation de la rugosité $z_0$, abaissement du profil logarithmique) et piège les particules.
\textbf{Conclusion :} L'érosion éolienne sélective appauvrit le sol en éléments fins (argile, limon, MO), ne laissant qu'un désert de graviers (reg). La végétation est le rempart principal.
\end{experience}

\begin{popfigure}
\begin{tikzpicture}[scale=0.9, every node/.style={font=\small}]
    % Profil de vent
    \draw[popblue, thick, ->] (0,0) -- (0,5) node[left] {Hauteur (m)};
    \draw[popblue, thick, ->] (0,0) -- (6,0) node[below] {Vitesse du vent (m/s)};
    % Courbe log sans veg
    \draw[poporange, very thick, domain=0.01:4.5, samples=100] plot (\x, {2.5*ln(\x/0.01)});
    \node[poporange, right] at (5,3.8) {Sol nu : $z_0 \approx 0,01$ m};
    % Courbe log avec veg
    \draw[popgreen, very thick, domain=0.5:4.5, samples=100] plot (\x, {1.2*ln(\x/0.5)});
    \node[popgreen, right] at (5,2.2) {Sol végétalisé : $z_0 \approx 0,5$ m};
    % Seuil érosion
    \draw[popred, dashed, thick] (0,4.5) -- (6,4.5) node[right] {Seuil déflation ($\sim$5 m/s)};
    
    % Particules transport
    \begin{scope}[xshift=7cm, yshift=0.5cm]
        \draw[popdark, thick, ->] (0,0) -- (4,0) node[below] {Distance};
        \draw[popdark, thick, ->] (0,0) -- (0,2.5) node[left] {Hauteur};
        % Saltation
        \foreach \i in {0,1,2,3} {
            \draw[poporange, thick] (\i*1.1,0) parabola[bend at end] (\i*1.1+0.5,0.8) parabola (\i*1.1+1,0);
        }
        \node[poporange] at (2,1.5) {Saltation (sable 0,1-0,5 mm)};
        % Reptation
        \draw[poppurple, thick, decorate, decoration={snake, amplitude=0.5mm, segment length=3mm}] (0.2,0.1) -- (3.8,0.1);
        \node[poppurple] at (2,0.5) {Reptation (graviers > 0,5 mm)};
        % Suspension
        \draw[popblue, very thick, dashed] (0,1.5) -- (4,2.2);
        \node[popblue] at (2,2.5) {Suspension (poussières < 0,1 mm)};
    \end{scope}
\end{tikzpicture}
\legende{À gauche : Profil de vitesse du vent (loi logarithmique) selon la rugosité du sol ($z_0$). La végétation augmente $z_0$, réduit la vitesse au ras du sol sous le seuil de déflation. À droite : Modes de transport éolien des particules (saltation, reptation, suspension).}
\end{popfigure}

\begin{definition}[Érosion éolienne (ou déflation)]
L'\textbf{érosion éolienne} est le processus d'arrachement, de transport et de dépôt des particules du sol sec et non protégé par l'action du vent. Elle opère une \textbf{tri granulométrique} sélective :
\begin{itemize}
    \item \textbf{Suspension :} Particules $< 0,1$ mm (argiles, limons fins, matière organique, nutriments) $\rightarrow$ transportées sur des centaines de km (poussières du \cle{Harmattan} visibles jusqu'au Sud-Cameroun).
    \item \textbf{Saltation :} Particules $0,1$--$0,5$ mm (sables fins) $\rightarrow$ sauts balistiques (hauteur 10--50 cm), érodent le sol par abrasion (effet "papier de verre" sur les plants, polissage des rochers/troncs).
    \item \textbf{Reptation (ou roulement) :} Particules $> 0,5$ mm $\rightarrow$ roulent au fond, ne quittent pas le sol, forment un \textbf{pavé éolien (lag gravel / desert pavement)}.
\end{itemize}
Le résidu grossier (pavé) peut temporairement protéger le sol sous-jacent (auto-stabilisation), mais le capital fertilité (fines + MO) est perdu à jamais.
\end{definition}

\courssub{Répartition géographique au Cameroun : une dualité climatique}

Le Cameroun, par sa position intertropicale et sa diversité de reliefs, présente une nette zonation des risques érosifs.

\begin{popfigure}
\begin{tikzpicture}[scale=1.1]
    % Contour Cameroun simplifié
    \draw[popdark, very thick] 
        (0,0) -- (1,0.5) -- (1,2) -- (0.5,3) -- (1,4) -- (2,4.5) -- (3,4) -- (4,4.5) -- (5,4) -- (5.5,3) -- (6,2) -- (6,1) -- (5.5,0) -- (4.5,-0.5) -- (3,0) -- (2,0.5) -- cycle;
    
    % Zones hydriques (Sud) - hachures bleues
    \foreach \y in {0,0.3,...,4.5} {
        \draw[popblue!40, thick] (0.5,\y) -- (5.5,\y);
    }
    \node[popblue, font=\bfseries\large, align=center] at (3,2.5) {ÉROSION\\ HYDRIQUE};
    \node[popblue!70!black, font=\small, align=center] at (3,1.8) {Régions : Sud, Centre, Est,\\ Littoral, Sud-Ouest, Ouest,\\ Adamaoua (sud)};
    \node[popblue!70!black, font=\small, align=center] at (3,1.0) {Facteurs : Pluies > 1500 mm/an,\\ Reliefs accidentés (Monts),\\  Sols latéritiques lessivés};
    
    % Zones éoliennes (Nord) - points orangés
    \foreach \x in {0.5,0.8,...,5.5} {
        \foreach \y in {-0.5,0,...,1.5} {
            \fill[poporange!50] (\x,\y) circle (0.05);
        }
    }
    \node[poporange, font=\bfseries\large, align=center] at (3,-0.8) {ÉROSION\\ ÉOLIENNE};
    \node[poporange!70!black, font=\small, align=center] at (3,-1.5) {Régions : Extrême-Nord,\\ Nord (bénoué), Adamaoua (nord)};
    \node[poporange!70!black, font=\small, align=center] at (3,-2.2) {Facteurs : Saison sèche 7-8 mois,\\ Vent d'Harmattan (NE),\\  Sols sableux/ferrugineux nus};
    
    % Villes repères
    \node[popdark, circle, fill, inner sep=1.5pt, label=above:\tiny Yaoundé] at (2.5,2.8) {};
    \node[popdark, circle, fill, inner sep=1.5pt, label=above:\tiny Douala] at (2.2,1.2) {};
    \node[popdark, circle, fill, inner sep=1.5pt, label=below:\tiny Maroua] at (4.5,0.2) {};
    \node[popdark, circle, fill, inner sep=1.5pt, label=below:\tiny Garoua] at (4.0,1.0) {};
    \node[popdark, circle, fill, inner sep=1.5pt, label=above:\tiny Ngaoundéré] at (3.5,3.2) {};
    \node[popdark, circle, fill, inner sep=1.5pt, label=above:\tiny Buea] at (1.5,1.5) {};
    
    % Légende
    \node[draw, fill=popcream, rounded corners, anchor=south east, font=\small, align=center] at (6.2,4.7){
        \textbf{Légende :}\\
        \textcolor{popblue}{\rule{0.5cm}{2pt}} Zone à risque hydrique dominant\\
        \textcolor{poporange}{\rule{0.5cm}{2pt}} Zone à risque éolien dominant\\
        $\bullet$ Villes de référence
    };
\end{tikzpicture}
\legende{Carte simplifiée du Cameroun montrant la répartition des types d'érosion dominants. L'érosion hydrique domine au Sud (fortes précipitations, reliefs), l'érosion éolienne au Nord (sécheresse, Harmattan, sols nus). La zone de transition (Adamaoua) subit les deux séquentiellement.}
\end{popfigure}

\begin{savaistu}[Les ravins de la côte du Wouri et les sols dénudés de l'Extrême-Nord]
\textbf{Cas 1 : Les ravins de la côte du Wouri (Douala, Région du Littoral).} L'urbanisation galopante et l'exploitation anarchique des carrières de sable latéritique ont mis à nu des pentes fragiles sur latérite. Les pluies torrentielles (3000--4000 mm/an) creusent des ravins spectaculaires (jusqu'à 20 m de profondeur, 50 m de large) qui menacent les habitations, la route nationale N3 et le pipeline Tchad-Cameroun. C'est un exemple type d'\textbf{érosion en ravins} anthropique, coûtant des milliards de FCFA en protection (busages, parpaings).
\textbf{Cas 2 : La plaine de Maroua (Extrême-Nord) en saison sèche (novembre--mars).} Le vent d'Harmattan, sec et chargé de poussières (visibilité $<$ 100 m), balaie les champs après récolte (mil, sorgho, coton) et les parcours surpâturés. Les cultures de contre-saison (tomates, oignons) sont décapitées par l'abrasion (saltation). Les toits et cours d'eau (Mayo Tsanaga) se colmatent. Les sols ferrugineux tropicaux (ferralsols), privés de leur argile et MO, deviennent stériles. C'est l'\textbf{érosion éolienne} dans sa forme sahélienne.
\end{savaistu}

\coursec{Les facteurs favorisants l'érosion}

L'érosion résulte de l'interaction entre des \textbf{facteurs naturels} (intrinsèques au milieu) et des \textbf{facteurs anthropiques} (liés à l'usage du sol). Au Cameroun, l'action humaine est le déclencheur principal de l'emballement érosif.

\begin{propriete}[Facteurs naturels de risque érosif]
\begin{itemize}
    \item \textbf{Climat (Agressivité) :} \textbf{Érosivité des pluies (indice $R$ USLE/RUSLE) :} Intensité, quantité, énergie cinétique. Au Cameroun : $R$ très élevé au Littoral/Sud-Ouest ($>$ 8000 MJ mm/ha/h/an) et Mont Cameroun ; modéré à l'Adamaoua ; faible au Nord. \textbf{Vent :} Vitesse, fréquence, durée (Harmattan : sec, NE, 5--10 m/s, 7 mois/an au Nord).
    \item \textbf{Topographie (Facteur $LS$) :} \textbf{Pente} (longueur $L$ et déclivité $S$). Vitesse du ruissellement $\propto \sqrt{S}$. Pentes $>$ 10\% = risque fort (Monts Mandara, Mont Cameroun, reliefs de l'Ouest). \textbf{Forme du versant :} Versants convexes (concentration) vs concaves (dépôt).
    \item \textbf{Sol (Érodibilité $K$) :} Texture (limon/sable fin $>$ argile $>$ gravier), structure (agrégats stables), MO (stabilise), perméabilité, profondeur. Sols ferrallitiques lessivés (Sud) = $K$ modéré mais profondeur forte = pertes massives. Sols ferrugineux tropicaux (Nord) = $K$ élevé (texture sableuse, faible MO, structure instable).
    \item \textbf{Couvert végétal naturel (Facteur $C$) :} Forêt dense ($C \approx 0,001$) $\ll$ Savane arborée ($C \approx 0,01$) $\ll$ Culture annuelle sans résidus ($C \approx 0,3-0,5$) $\ll$ Sol nu ($C = 1$).
\end{itemize}
\end{propriete}

\begin{propriete}[Facteurs anthropiques aggravants (Pratiques non durables)]
\begin{itemize}
    \item \textbf{Déforestation / Déboisement :} Conversion forêt $\rightarrow$ culture/patûre. Perte de l'interception, des racines, de la MO. Cause n°1 au Sud (cacao, huile de palme, bois d'œuvre) et au Nord (feu de brousse, bois de chauffe).
    \item \textbf{Agriculture sur pente sans aménagement :} Labour \textbf{perpendiculaire aux courbes de niveau} (sens de la pente) = autoroute pour le ruissellement. Absence de \cle{courbes de niveau}, de \cle{terrasses}, de \cle{bandes enherbées}.
    \item \textbf{Brûlis et labour à nu :} Destruction de la MO, destruction de la structure, exposition directe au \emph{splash} et au vent. Sol nu pendant la saison des pluies (avril-mai) ou sèche (nov-mars).
    \item \textbf{Surpâturage / Pâturage itinérant :} Piétinement (compaction, battance), arrachage des plantes, mise à nu du sol. Problème majeur dans l'Adamaoua et le Nord.
    \item \textbf{Urbanisation non maîtrisée / Carrières :} Imperméabilisation (ruissellement $\uparrow$ 10x), déblais/remblais instables, extraction de sable/latérite (Wouri, Mfoundi).
    \item \textbf{Mauvaise gestion de l'eau :} Absence de drainage des pistes/rues, busages sous-dimensionnés, concentration des eaux de toiture sans dissipation.
\end{itemize}
\end{propriete}

\begin{methode}[Analyser les causes de l'érosion dans une situation donnée (Terrain / Photo / Étude de cas)]
\begin{enumerate}
    \item \textbf{Identifier le type d'érosion} (hydrique / éolienne / mixte) via les formes (ravins, dunes, battance, pavé).
    \item \textbf{Lister les facteurs naturels présents} : Climat (pluie/vent), Pente (\%, longueur), Type de sol (texture, structure, profondeur).
    \item \textbf{Lister les facteurs humains} : Occupation du sol (culture, forêt, ville), Pratiques culturales (labour, résidus, brûlis), Aménagements (existence/absence de terrasses, fossés, brise-vent), Densité de population/cheptel.
    \item \textbf{Hiérarchiser} : Quel est le facteur \textbf{déclencheur} (souvent : mise à nu du sol) ? Quels sont les facteurs \textbf{aggravants} (pente, intensité pluies) ?
    \item \textbf{Conclure} sur le couplage Nature/Homme : "L'érosion observée résulte de la conjugaison d'une forte érosivité climatique (pluies/Harmattan) et d'une grande sensibilité du sol (pente/texture), \textbf{déclenchée et amplifiée par la suppression du couvert végétal par l'homme} (déforestation, labour à nu, surpâturage)."
\end{enumerate}
\end{methode}

\coursec{Les conséquences de l'érosion}

Les impacts se font sentir \textbf{sur site} (là où le sol part) et \textbf{hors site} (là où le sédiment arrive), avec des coûts économiques, écologiques et sociaux majeurs.

\begin{propriete}[Conséquences sur site (perte de productivité du sol)]
\begin{itemize}
    \item \textbf{Appauvrissement chimique :} Perte sélective des fines (argile, limon) et de la MO $\rightarrow$ baisse CEC, perte N, P, K, bases. Acidification. \textbf{Baisse des rendements} (10--50\% selon sévérité).
    \item \textbf{Dégradation physique :} Perte de l'horizon A (labourable) $\rightarrow$ affleurement horizon B (dur, compact, toxique Al/Mn). Baisse profondeur utile, réserve en eau $\downarrow$, infiltration $\downarrow$, ruissellement $\uparrow$ (cercle vicieux).
    \item \textbf{Destruction du capital foncier :} Ravins = perte de surface cultivable, morcellement parcelles, dépréciation valeur foncière. Enfouissement cultures/récoltes par dépôts grossiers (sable/gravier) en bas de pente.
    \item \textbf{Destruction d'infrastructures :} Routes (N3 Wouri, pistes rurales), ponts, canalisations, habitations, pipelines, barrages (envasement).
\end{itemize}
\end{propriete}

\begin{propriete}[Conséquences hors site (pollution, risques, coûts collectifs)]
\begin{itemize}
    \item \textbf{Envaseement des cours d'eau et barrages :} Diminution capacité de stockage (barrages de Lagdo, Bamendjing, Mapé, Mékin). Coûts de curage énormes. Montée des lits $\rightarrow$ inondations accrues (Logone, Bénoué, Wouri).
    \item \textbf{Pollution de l'eau :} Turbidité $\uparrow$ (coût potabilisation $\uparrow$), eutrophisation (N, P adsorbés sur sédiments) $\rightarrow$ pullulations algues/jacinthes (Lac Ossa, Lagdo), contamination pesticides.
    \item \textbf{Poussières atmosphériques (Éolienne) :} Santé publique (infections respiratoires, méningite, allergies), transport aérien perturbé, dépôt sur cultures/villes (Harmattan).
    \item \textbf{Perte de biodiversité :} Destruction habitats (ravins), colmatage frayères (poissons), perte ripisylves.
\end{itemize}
\end{propriete}

\begin{aretenir}[Coût économique de l'érosion au Cameroun]
La Banque Mondiale et le MINEPAT estiment le coût de la dégradation des terres (dont érosion majeure) à \textbf{4--8\% du PIB annuel} (pertes agricoles, infrastructures, santé, barrages). La restauration (reboisement, terrasses) coûte 10 à 50 fois moins cher que la réparation des dégâts (curage barrages, reconstruction routes, soins santé).
\end{aretenir}

\coursec{Les techniques de lutte contre l'érosion}

La lutte efficace repose sur une \textbf{approche intégrée} : \textbf{Protéger le sol} (couvert), \textbf{Ralentir l'agent} (ruissellement/vent), \textbf{Gérer l'eau/le vent} (aménagements). Le choix dépend du type d'érosion, de la pente, du sol et du contexte socio-économique.

\begin{propriete}[Principes fondamentaux de la conservation des sols et de l'eau (CSE)]
\begin{enumerate}
    \item \textbf{Maintenir un couvert permanent :} C'est la mesure n°1, universelle (hydrique + éolienne). Végétation vivante ou résidus morts (paillis). Objectif : $C \rightarrow 0$.
    \item \textbf{Réduire la longueur et la vitesse du ruissellement / du vent :} Par des obstacles perpendiculaires à l'écoulement/vent (courbes de niveau, terrasses, bandes enherbées, brise-vent, fascines).
    \item \textbf{Améliorer l'infiltration et la structure :} Labour minimal (SCV - Semis Direct sous Couvert Végétal), apport MO, agroforesterie.
    \item \textbf{Gérer les eaux de ruissellement concentré :} Fossés de dérivation, chutes à énergie dissipée, barrages de retenue.
    \item \textbf{Adapter l'usage du sol à sa capacité :} Mettre en défens/forêt les pentes $>$ 30--40\% ou sols très fragiles (Classement en "Terres de protection").
\end{enumerate}
\end{propriete}

\courssub{Techniques agronomiques et culturales (Première ligne de défense)}

\begin{center}
\begin{tabularx}{\linewidth}{|l|X|X|l|}
\hline
\rowcolor{popblueL}
\textbf{Technique} & \textbf{Principe / Mise en œuvre} & \textbf{Efficacité / Contexte Cameroun} & \textbf{Exemples locaux} \\ \hline
\cle{Cultures en courbes de niveau} & Labour, semis, billons \textbf{perpendiculaires à la pente} (suivent les isolignes). Crée des micro-barrages successifs. & Réduit vitesse ruissellement $\times$ 2--5. Efficace pente $<$ 10--15\%. Faible coût. & Maïs/arachide dans l'Ouest (Bamileke), Adamaoua. \\ \hline
\cle{Bandes enherbées / Bandes filtrantes} & Bandes d'herbes pérennes (Vetiver, Brachiaria, Pennisetum) larges de 1--3 m, espacées de 10--20 m selon pente. Piègent sédiments, infiltrent eau. & Très efficace érosion en nappe/rigoles. Vetiver : racines profondes (3--4 m), résistant feu/sécheresse. & \textbf{Vetiver} (Chrysopogon zizanioides) : standard national (MINADER). \\ \hline
\cle{Cultures de couverture / Engrais verts} & Plantes à croissance rapide, fort biomasses, semées en intercalaire ou relais (Mucuna, Crotalaria, Dolichos, Canavalia, Niébé). Couvrent sol inter-récolte. & Protègent \emph{splash}, nourrissent sol (MO, N si légumineuses), étouffent adventices. & \textbf{Mucuna pruriens} (Sud, maïs) ; \textbf{Crotalaria} (Nord, coton/maïs). \\ \hline
\cle{Paillage / Résidus de culture} & Couverture du sol par paille, tiges, feuilles (5--10 t/ms ha). Ne pas brûler ! & Amortit splash, réduit évaporation, nourrit vie du sol. Essentiel SCV. & Résidus maïs/sorgho/mil ; feuilles palmier/bananier. \\ \hline
\cle{Agroforesterie / Parc à Faidherbia} & Association arbres + cultures. \textbf{Faidherbia albida} (Cad, Gaïla) : feuillaison \emph{inversée} (feuilles saison sèche, nue saison pluies). & Ombrage modéré, apport N (fixation), MO, racines profondes, brise-vent, fourrage. & \textbf{Parcs à Faidherbia} (Nord, coton/sorgho) : modèle durable prouvé. \\ \hline
\cle{Semis Direct sous Couvert Végétal (SCV)} & Pas de labour. Semis direct dans résidus/couverture vivante. Herbicide ou rouleau-faucheur pour gérer couverture. & Zéro perturbation sol, MO $\uparrow\uparrow$, érosion $\approx$ 0. Transition technique exigeante. & Projets PDRT, PADFA, ONG (GIC) dans l'Ouest et Adamaoua. \\ \hline
\end{tabularx}
\end{center}

\courssub{Aménagements mécaniques / structures (Pour pentes fortes ou ruissellement concentré)}

\begin{center}
\begin{tabularx}{\linewidth}{|l|X|X|l|}
\hline
\rowcolor{poporangeL}
\textbf{Aménagement} & \textbf{Description / Dimensionnement} & \textbf{Rôle / Entretien} & \textbf{Contexte Cameroun} \\ \hline
\cle{Terrasses en gradins (Bench terraces)} & Transformation pente en marches planes (plateforme + talus). Travaux lourds (main d'œuvre/machine). Largeur plateforme 2--5 m. Talus enherbé (Vetiver) ou maçonné. & Supprime pente effective. Permet irrigation/riziculture. Entretien talus crucial (érosion talus = rupture). & \textbf{Région de l'Ouest} (Bamileke : "terrasses traditionnelles" riz/maïs), Monts Mandara. Investissement fort, main d'œuvre familiale. \\ \hline
\cle{Terrasses à talus enherbé (Grassed waterways / Fanya juu)} & Fossé + banquette en amont (terre jetée vers amont). Stabilisé par herbes (Vetiver). Espacement $E = 0,7 \times \sqrt{Pente\%}$ (m) ou formule Wischmeier. & Intercepte ruissellement, l'infiltre, piège sédiments. Moins coûteux que terrasses pleines. & Vulgarisé par MINADER/ONG (Projet PADFA, PNDP). Adapté pentes 10--30\%. \\ \hline
\cle{Fossés de dérivation / Couronnement} & Fossé en haut de versant ou bordure champ, pente faible (0,2--0,5\%), évacue eaux extérieures vers exutoire stable (ravin naturel, cours d'eau). & Protège le champ des eaux venues d'amont (ruissellement externe). & Indispensable en zone de piémont (Mont Cameroun, Ouest). \\ \hline
\cle{Barrages de retenue / Seuils (Check dams)} & Petits barrages (pierres, gabions, ciment, terre) dans thalwegs/rigoles actives. Hauteur $<$ 1--2 m. Série en cascade. & Réduisent pente effective du thalweg, favorisent dépôt, rechargent nappe, créent points d'eau. & \textbf{Barrages en pierres sèches} (Nord, Mayo Rey), \textbf{Gabions} (Wouri, routes). Gestion participative (COBAS). \\ \hline
\cle{Brise-vent (Haies brise-vent)} & Rangées d'arbres/arbustes pérennes, perpendiculaires vent dominant (NE = Harmattan). Porosité 40--50\% (pas mur étanche). Hauteur $H$, protection jusqu'à $15-20 H$ au vent. & Réduisent vitesse vent $>$ seuil déflation. Piègent poussières/sable. Microclimat (évapotranspiration $\downarrow$). & \textbf{Nord/Extrême-Nord} : Neem (Azadirachta), Acacia nilotica, Balanites, Prosopis. Double usage : bois, fourrage, pharmacopée. \\ \hline
\end{tabularx}
\end{center}

\begin{popfigure}
\begin{tikzpicture}[scale=0.85, every node/.style={font=\small}]
    % Coupe transversale pente avec terrasses et courbes de niveau
    \draw[popdark, very thick] (0,0) -- (14,0); % Base
    % Profil terrain naturel (pentu)
    \draw[poporange, thick, dashed] (0,0) -- (14,3.5) node[right, poporange] {Pente naturelle (25\%)};
    % Terrasses
    \foreach \i/\y in {0/0.5, 1/1.2, 2/1.9, 3/2.6, 4/3.3} {
        \draw[popgreen, very thick] (\i*3, \y) -- (\i*3+2.5, \y); % Plateforme
        \draw[popgreen, thick] (\i*3+2.5, \y) -- (\i*3+3, \y+0.5); % Talus
        \node[popgreen, above] at (\i*3+1.25, \y+0.2) {Plateforme cultivée};
        \node[popgreen!70!black, rotate=60] at (\i*3+2.7, \y+0.25) {Talus enherbé (Vetiver)};
    }
    % Courbes de niveau sur plateforme (petits traits)
    \foreach \i/\y in {0/0.5, 1/1.2, 2/1.9, 3/2.6, 4/3.3} {
        \foreach \x in {0.2,0.6,...,2.3} {
            \draw[popblue, thick] (\i*3+\x, \y-0.1) -- (\i*3+\x, \y+0.1);
        }
    }
    % Flèches ruissellement ralenti
    \draw[popblue, thick, -{Stealth[length=2mm]}] (0.5,0.8) -- (2.3,0.8);
    \draw[popblue, thick, -{Stealth[length=2mm]}] (3.5,1.5) -- (5.3,1.5);
    \node[popblue] at (1.5,1.1) {Ruissellement ralenti};
    % Infiltration
    \draw[popblue, dashed, -{Stealth[length=1.5mm]}] (1.5,0.5) -- (1.5,-0.5);
    \node[popblue, below] at (1.5,-0.6) {Infiltration $\uparrow$};
    % Fossé de dérivation en haut
    \draw[poppurple, thick, -{Stealth[length=2mm]}] (0,3.8) -- (14,3.8);
    \node[poppurple, above] at (7,4.0) {Fossé de dérivation (eaux extérieures)};
    % Légende
    \node[draw, fill=popcream, rounded corners, anchor=south west, font=\small, align=left] at (0.2,4.2) {
        \textbf{Coupe type : Terrasses à plateformes + Courbes de niveau + Fossé de couronnement}\\
        Pente initiale 25\% $\rightarrow$ Pente effective plateforme $\approx$ 0\%\\
        Talus stabilisés par \textbf{Vetiver} (racines 3-4 m)
    };
\end{tikzpicture}
\legende{Schémas de principe : Aménagement en terrasses à plateformes (culture plane) avec talus enherbés (Vetiver), labour en courbes de niveau sur chaque plateforme, et fossé de dérivation en amont pour détourner les eaux de ruissellement externes.}
\end{popfigure}

\courssub{Stratégies spécifiques selon le type d'érosion dominant}

\begin{itemize}
    \item \textbf{Zone Hydrique (Sud, Ouest, Littoral) :} Priorité \textbf{infiltration} et \textbf{ralentissement ruissellement}. \textbf{1.} Courbes de niveau + Bandes de Vetiver (pentes $<$ 20\%). \textbf{2.} Terrasses (pentes $>$ 20\%). \textbf{3.} Agroforesterie (arbres fertilité + fruitiers). \textbf{4.} Reboisement pentes $>$ 40\% / sources / berges (Ripisylves). \textbf{5.} Gestion bassins versants (Comités de bassin).
    \item \textbf{Zone Éolienne (Nord, Extrême-Nord) :} Priorité \textbf{couvert permanent} et \textbf{brise-vent}. \textbf{1.} \textbf{Paillage} obligatoire (résidus mil/sorgho/arachide) + \textbf{Zaï / Demi-lunes} (recharge eau + piège vent). \textbf{2.} \textbf{Brise-vent} (Haies vives Neem, Acacia, Balanites) tous les 100--200 m. \textbf{3.} \textbf{Parcs à Faidherbia} (agroforesterie sahélienne). \textbf{4.} Reboisement dunes fixes/mobiles (Acacia senegal, Balanites, Eucalyptus - avec nappe). \textbf{5.} Lutte contre surpâturage (parcours tournants, aires de pâturage aménagées).
    \item \textbf{Zone de Transition (Adamaoua) :} \textbf{Stratégie couplée} (Voir Exoresolu). Couvert permanent \textbf{annuel} (Mucuna/Crotalaria saison pluies + Paillis/Résidus saison sèche) + Brise-vent + Courbes de niveau.
\end{itemize}

\begin{savaistu}[Le Vetiver (Chrysopogon zizanioides) : "L'herbe miracle" de la lutte anti-érosive au Cameroun]
Introduit et vulgarisé massivement depuis les années 90 (Projet PADFA, MINADER, ONG), le Vetiver est devenu l'outil n°1 de la CSE paysanne.
\begin{itemize}
    \item \textbf{Racines :} Fasciculées, profondes (3--4 m en 1 an), très résistantes à la traction (renforcent talus comme armature béton).
    \item \textbf{Rusticité :} Supporte pH 3--11, salinité, submersion, sécheresse 6--8 mois, feu (repousse souche), métaux lourds.
    \item \textbf{Non invasif :} Stérile (ne produit pas de graines viables), ne concurrence pas cultures (racines verticales).
    \item \textbf{Usages multiples :} Huile essentielle (racines, parfum), fourrage (jeunes feuilles), artisanat (feuilles), phytoépuration (eaux usées).
\end{itemize}
\textbf{Plantation :} Boutures (éclats de touffe) tous les 10--15 cm sur la ligne de courbe de niveau. Formation haie dense en 6--12 mois.
\end{savaistu}

\coursec{Interprétation de documents et étude de cas}

\begin{methode}[Distinguer érosion hydrique et érosion éolienne sur un document (photo, schéma, terrain)]
\begin{enumerate}
    \item \textbf{Observer les formes du relief :} Présence de \textbf{ravins, rigoles, cônes de déjection, éboulis, dépôts alluviaux} $\rightarrow$ \textbf{Érosion hydrique}. Présence de \textbf{dunes, barkhanes, pavé éolien (graviers résiduels alignés), surface pelée lisse, creux de déflation} $\rightarrow$ \textbf{Érosion éolienne}.
    \item \textbf{Observer l'état de surface :} \textbf{Croûte de battance} lisse, brillante, fissurée en polygones (retrait argile) $\rightarrow$ Hydrique (battance). \textbf{Surface granuleuse, "peau d'orange", micro-dunes, dépôts de poussières fins} $\rightarrow$ Éolienne.
    \item \textbf{Observer la végétation :} Arbres déracinés, racines apparentes (mise à nu par lessivage/ruissellement) $\rightarrow$ Hydrique. Arbres aux troncs érodés/poncés/creusés au niveau du sol (abrasion par saltation), alignés perpendiculairement au vent dominant, "drapeau" (branches d'un côté) $\rightarrow$ Éolienne.
    \item \textbf{Contextualiser (climat, saison, région) :} Saison des pluies, zone humide/forestière, relief accidenté $\rightarrow$ Hydrique. Saison sèche, zone sahélienne/soudanienne, vent fort (Harmattan), sol sableux nu $\rightarrow$ Éolienne.
\end{enumerate}
\end{methode}

\begin{exemplebox}[Application de la méthode : Analyse de deux photos types Probatoire]
\textbf{Photo A :} Pente du Mont Cameroun (Buea, 800 m). Sol volcanique rougeâtre (andosol jeune), entaillé par des sillons profonds (1--2 m) convergents vers un talweg. Eau trouble brune au fond. Racines d'eucalyptus apparentes sur 50 cm. Pente forte ($>$ 30\%).
\textbf{Diagnostic :} \textbf{Érosion hydrique en ravins active.} Mécanisme : Pluies extrêmes (4000 mm/an) + Pente forte + Sol friable (cendres/ponces) + Déforestation (culture/ananas/eucalyptus) $\rightarrow$ Ruissellement concentré massif.
\textbf{Photo B :} Périphérie de Kousseri (Extrême-Nord, Logone), mars (saison sèche). Sol sableux clair (dunes fixées anciennes), aspect "pelé", pas de rigoles. Petites dunes mobiles (barkhanes naissantes) au pied des rares acacias (Acacia nilotica). Troncs d'acacias polis/creusés à 20--30 cm du sol (abrasion saltation). Ciel jaunâtre (poussières en suspension). Vent NE fort.
\textbf{Diagnostic :} \textbf{Érosion éolienne active (Déflation + Transport).} Mécanisme : Harmattan + Sol nu (récolte mil/sorgho, pâturage) + Texture sableuse + Absence brise-vent $\rightarrow$ Arrachage fines (suspension) + Abrasion cultures/arbres (saltation).
\end{exemplebox}

\begin{exoresolu}[Analyse d'une situation érosive complexe : Le cas de l'Adamaoua (Zone de transition)]
\textbf{Énoncé :} Un agriculteur de la région de l'Adamaoua (alt. 1100 m, pluviométrie 1500 mm/an, saison sèche marquée 5 mois, vent d'Harmattan modéré) observe sur son champ de maïs en pente (12\%) : (1) une croûte dure en surface après les premiers orages d'avril ; (2) des rigoles de 10 cm de large après les grosses pluies de juillet ; (3) en janvier, de la poussière soulevée par le vent du Nord-Est qui décape les jeunes plants de la culture de contre-saison (tomates irrigoutées). Identifier les types d'érosion en jeu et proposer une explication unifiée et des solutions adaptées.
\textbf{Solution détaillée :}
\begin{enumerate}
    \item \textbf{Identification des types d'érosion :}
        \begin{itemize}
            \item \textbf{Observations (1) \& (2) :} Croûte de battance $\rightarrow$ \emph{splash} (impact gouttes) ; Rigoles $\rightarrow$ \textbf{Érosion hydrique (stades nappe puis rigoles)}. Conditions : Pente 12\% + Pluies intenses (orages avril-juillet) + Sol nu (labour avril, pas de résidus) = Érosivité forte + Érodibilité accrue par labour.
            \item \textbf{Observation (3) :} Poussières en suspension, défoliation/décapitation jeunes plants (abrasion), vent NE (Harmattan), saison sèche (sol sec), sol nu (récolte maïs oct/nov, résidus brûlés ou exportés fourrage) = \textbf{Érosion éolienne (Suspension + Saltation abrasive)}.
        \end{itemize}
    \item \textbf{Explication unifiée (Couple érosif séquentiel) :} Le point commun structurel est la \textbf{nudité du sol pendant les périodes critiques} (avril-mai et nov-mars).
        \begin{itemize}
            \item \textbf{Saison des pluies (Avril--Oct) :} Labour + absence couverture $\rightarrow$ énergie pluie $\rightarrow$ battance + ruissellement $\rightarrow$ perte fines (argile, limon, MO) $\rightarrow$ dégradation structure.
            \item \textbf{Saison sèche (Nov--Mars) :} Résidus brûlés/exportés + absence culture de couverture $\rightarrow$ vent + sol sec/pulvérulent $\rightarrow$ déflation (perte limons/argiles résiduels + MO) + abrasion plants.
        \end{itemize}
        Le sol (probablement un \textbf{Nitisol} ou \textbf{Ferralsol} lessivé, rouge, profond mais sensible à la battance) perd sa structure (MO $\downarrow$), sa porosité (battance) et sa texture fine (argile $\downarrow$ par hydrique, limon/argile $\downarrow$ par éolien). C'est un \textbf{couple érosif hydrique/éolien séquentiel} typique des hautes terres soudano-guinéennes.
    \item \textbf{Leviers d'action prioritaires (Intégrés) :}
        \begin{enumerate}
            \item \textbf{Couvert permanent toute l'année (Clé de voûte) :}
                \begin{itemize}
                    \item Saison pluies : \textbf{Intercalaire maïs + Mucuna pruriens} (ou Crotalaria juncea) semé à la 2ème sarclée. Mucuna couvre sol après récolte maïs, fixe N, étouffe adventices.
                    \item Saison sèche : \textbf{Paillis de résidus maïs + Mucuna} (non brûlé, non exporté) $\rightarrow$ protège du vent, nourrit sol, conserve eau pour tomates.
                \end{itemize}
            \item \textbf{Aménagement pente :} \textbf{Courbes de niveau} rigoureuses + \textbf{Bandes de Vetiver} tous les 10--15 m (ralentit ruissellement + piège sédiments + brise vent rasant).
            \item \textbf{Agroforesterie :} Plantation lignes \textbf{Faidherbia albida} (10x10 m) : ombrage léger tomates saison sèche, apport N, brise-vent, feuilles saison sèche = paillis.
            \item \textbf{Gestion eau :} \textbf{Zaï / Demi-lunes} pour tomates (concentration eau + MO), irrigation goutte-à-goutte (économie eau, pas de ruissellement).
        \end{enumerate}
\end{enumerate}
\end{exoresolu}

\begin{exercice}[Entraînement Probatoire : Diagnostic et proposition de lutte]
\textbf{Document 1 :} Carte pédologique simplifiée d'un bassin versant de l'Ouest (Bamileke). Pentes moyennes 25\%. Sols : Ferralsols rouges profonds. Occupation : Cultures vivrières (maïs, haricot, taro) en labour itinérant (brûlis), pâturages extensifs, lambeaux de forêt galerie. Densité population : 200 hab/km².
\textbf{Document 2 :} Photo d'un champ en pente (15\%) après un orage d'avril. Sol nu, labouré dans le sens de la pente. Surface lisse, brillante, fissurée. Eau trouble s'écoulant en rigoles naissantes vers le bas.
\textbf{Document 3 :} Tableau de pertes en sol mesurées (t/ha/an) selon pratique culturale (essai station IRAD Dschang, pente 15\%, Ferralsol).
\begin{center}
\begin{tabular}{|l|c|}\hline
\textbf{Pratique} & \textbf{Perte sol (t/ha/an)} \\ \hline
Sol nu, labour sens pente & 42 \\ \hline
Maïs seul, labour sens pente & 28 \\ \hline
Maïs + Mucuna, labour courbes niveau & 3,5 \\ \hline
Bande Vetiver (10m) + Maïs courbes niveau & 1,2 \\ \hline
Forêt naturelle (Témoin) & 0,1 \\ \hline
\end{tabular}
\end{center}
\textbf{Questions :}
\begin{enumerate}
    \item À partir des Documents 1 et 2, identifier le type d'érosion dominant, ses formes visibles et les deux facteurs anthropiques principaux responsables.
    \item Interpréter les résultats du Document 3. Calculer le taux de réduction des pertes apporté par la combinaison "Maïs+Mucuna+Courbes de niveau" par rapport au "Sol nu". Quelle technique apporte le gain marginal le plus fort ?
    \item Proposer un plan de lutte intégré (3 mesures techniques + 1 mesure organisationnelle) pour ce bassin versant, en justifiant chaque mesure par un mécanisme physique précis.
\end{enumerate}
\end{exercice}

\begin{corrige}[Exercice n°1 : Diagnostic et proposition de lutte]
\begin{enumerate}
    \item \textbf{Type :} \textbf{Érosion hydrique} (zone Ouest, pluviométrie forte, relief accidenté, sols ferrallitiques).
    \textbf{Formes visibles (Doc 2) :} \textbf{Croûte de battance} (surface lisse, brillante, fissurée = scellement par \emph{splash}) + \textbf{Rigoles naissantes} (ruissellement concentré). Absence de ravins profonds = stade intermédiaire.
    \textbf{Facteurs anthropiques principaux :}
    \begin{itemize}
        \item \textbf{Labour dans le sens de la pente} (création de sillons préférentiels pour le ruissellement, augmentation vitesse $\times$ 3--5 vs courbes de niveau).
        \item \textbf{Mise à nu du sol (Brûlis + absence résidus/couverture)} au moment des premiers orages violents (avril) $\rightarrow$ énergie cinétique gouttes directe sur agrégats nus.
    \end{itemize}
    \item \textbf{Interprétation Doc 3 :}
    \begin{itemize}
        \item Perte sol nu = 42 t/ha/an (Érosion accélérée massive, $>$ tolérance renouvellement $\approx$ 1--2 t/ha/an).
        \item Maïs seul = 28 t/ha/an (Couverture partielle, mais labour pente annule bénéfice).
        \item Maïs+Mucuna+Courbes niveau = 3,5 t/ha/an \textbf{(Réduction 91,7\% vs sol nu)}. Calcul : $(42-3,5)/42 \times 100 = 91,7\%$.
        \item + Bandes Vetiver = 1,2 t/ha/an (Proche tolérance géologique).
        \item \textbf{Gain marginal le plus fort :} Passage "Maïs seul pente" (28) $\rightarrow$ "Maïs+Mucuna Courbes niveau" (3,5) = Gain de 24,5 t/ha/an. L'association \textbf{Couvert végétal (Mucuna) + Ralentissement hydraulique (Courbes niveau)} agit sur les deux leviers (Détachement + Transport).
    \end{itemize}
    \item \textbf{Plan de lutte intégré pour le bassin versant :}
    \begin{enumerate}
        \item \textbf{Mesure 1 (Agronomique - Source) :} \textbf{Généralisation SCV Maïs/Mucuna (ou Crotalaria) en courbes de niveau} sur toutes les parcelles cultivées.
            \textit{Justification :} Amortit \emph{splash} (couvert), maintient porosité (racines Mucuna), fixe N (légumineuse), fournit paillis saison sèche. Cible : Facteur $C$ (USLE) $\downarrow$ de 0,4 à $<$ 0,05.
        \item \textbf{Mesure 2 (Mécanique - Pente) :} \textbf{Implantation Bandes de Vetiver} tous les 10--15 m (suivant courbes de niveau) + \textbf{Terrasses à talus enherbé (Fanya juu)} sur pentes $>$ 20\%.
            \textit{Justification :} Réduit longueur de pente effective ($L \downarrow$), crée barrages physiques piégeant sédiments et infiltrant ruissellement. Cible : Facteur $LS \downarrow$ drastiquement.
        \item \textbf{Mesure 3 (Forestière - Amont/Berges) :} \textbf{Reboisement / Régénération naturelle assistée} des versants $>$ 35\% (non cultivables), des têtes de bassin (recharge nappe) et \textbf{Ripisylves} (10--20 m) le long cours d'eau.
            \textit{Justification :} Interception pluie (canopée), rugosité sol (litière), racines profondes (structure/infiltration), piégeage sédiments finaux (filtre tampon). Cible : Facteur $C \approx 0,001$ + protection hors site (eau).
        \item \textbf{Mesure 4 (Organisationnelle) :} \textbf{Création/Redynamisation Comités de Gestion de Bassin Versant (CGBV)} associant chefs traditionnels, GIC agriculteurs, éleveurs, MINEPAT/MINADER/COMMUNES.
            \textit{Justification :} L'érosion ne connaît pas les limites parcellaires. Nécessite concertation (pâturage vs cultures, entretien fossés/vetiver collectifs, règles brûlis, financement entretien terrasses). Garantit pérennité (gouvernance locale).
    \end{enumerate}
\end{enumerate}
\end{corrige}

\begin{aretenir}[Synthèse comparative : Érosion hydrique vs Érosion éolienne]
\begin{center}
\begin{tabularx}{\linewidth}{|l|X|X|}
\hline
\textbf{Critère} & \textbf{Érosion hydrique} & \textbf{Érosion éolienne} \\ \hline
\textbf{Agent moteur} & Eau (Gouttes \emph{splash} + Ruissellement) & Vent (Cisaillement + Turbulence) \\ \hline
\textbf{Condition déclenchante} & Pluie intense ($>$ 25 mm/h) + Pente + Sol nu/sec & Vent $>$ Seuil (5--7 m/s à 30 cm) + Sol sec + Nu \\ \hline
\textbf{Formes majeures (Érosion)} & Nappe $\rightarrow$ Rigoles $\rightarrow$ Ravins & Déflation (creux, cuvettes) \\ \hline
\textbf{Formes majeures (Dépôt)} & Cônes de déjection, Alluvions, Deltas & Dunes (Barkhanes, Seif), Loess \\ \hline
\textbf{Modes de transport} & Suspension (fines), Saltation, Traction (gros) & Suspension (poussières $<$ 0,1 mm), Saltation (sable 0,1--0,5 mm), Reptation ($>$ 0,5 mm) \\ \hline
\textbf{Sélectivité granulométrique} & Fines (argile, limon, MO) parties 1ères & \textbf{Extrêmement sélective} : Fines (argile, limon, MO) $\uparrow\uparrow$ parties \\ \hline
\textbf{Indicateur visuel clé (Sol)} & Croûte de battance, Rigoles, Ravins, Eaux troubles & Surface pelée, Pavé éolien (Lag gravel), Dunes, Troncs polis \\ \hline
\textbf{Zone Cameroun à risque} & Sud, Ouest, Littoral, Sud-Ouest, Adamaoua Sud & Extrême-Nord, Nord (Bénoué), Adamaoua Nord \\ \hline
\textbf{Mesure n°1 (Universelle)} & \textbf{Couvert végétal permanent} (Vivant/Mort) & \textbf{Couvert végétal permanent} (Vivant/Mort) \\ \hline
\textbf{Mesures spécifiques} & Courbes de niveau, Terrasses, Bandes Vetiver, Fossés, Reboisement & Brise-vent (Haies vives), Paillis, Zaï/Demi-lunes, Fixation dunes \\ \hline
\end{tabularx}
\end{center}
\end{aretenir}

\begin{aretenir}[Mots-clés du Probatoire à maîtriser pour cette leçon]
\cle{Érosion}, \cle{Splash}, \cle{Battance}, \cle{Ruissellement}, \cle{Rigole}, \cle{Ravin}, \cle{Déflation}, \cle{Saltation}, \cle{Suspension}, \cle{Reptation}, \cle{Pavé éolien}, \cle{Harmattan}, \cle{Érosivité}, \cle{Érodibilité}, \cle{Courbes de niveau}, \cle{Terrasses}, \cle{Bandes enherbées}, \cle{Vetiver}, \cle{Culture de couverture}, \cle{Paillis}, \cle{Agroforesterie}, \cle{Faidherbia}, \cle{SCV}, \cle{Brise-vent}, \cle{Zaï}, \cle{Envaseement}, \cle{Coût économique}.
\end{aretenir}

\coursec{Les facteurs favorisants l'érosion}

Après avoir distingué les deux grands types d'érosion — hydrique et éolienne — dans la section précédente, nous devons maintenant comprendre \textbf{ pourquoi } un même type d'érosion peut être négligeable dans un milieu et catastrophique dans un autre. Pourquoi une pluie d'avril sur les flancs du Mont Cameroun arrache-t-elle des tonnes de terre, alors qu'une pluie de même intensité sur la plaine de la Sanaga ne laisse presque aucune trace ? Pourquoi les vents d'harmattan soulèvent-ils des dunes de sable dans le Grand Nord, mais pas dans les forêts du Sud ? La réponse réside dans la combinaison de \textbf{facteurs naturels} (pente, climat, nature du sol) et de \textbf{facteurs anthropiques} (actions humaines) qui modifient la résistance du sol aux agents érosifs. C'est ce que nous allons analyser rigoureusement.

\courssub{Facteurs naturels : la pente, le climat et la nature du sol}

\textbf{1. La pente : le moteur gravitaire de l'érosion hydrique.}\par
Imaginez deux parcelles identiques, même sol, même pluviométrie, même couvert végétal : l'une est plate (marais de la Sanaga), l'autre est en forte pente (versant du Mont Cameroun). Après un orage violent, la parcelle en pente présente des ravines profondes, tandis que la plate n'a que de l'eau stagnante.

\begin{methode}[Démarche d'investigation : influence de la pente sur l'érosivité du ruissellement]
\textbf{Observation :} Sur un terrain en pente, l'eau s'écoule vite et creuse des sillons ; sur terrain plat, elle stagne ou s'infiltre.
\textbf{Hypothèse :} La vitesse du ruissellement (et donc son énergie cinétique érosive) augmente avec l'inclinaison de la pente.
\textbf{Expérience (parcelles expérimentales type Wischmeier) :} On installe des parcelles de même surface (ex. \SI{22}{m} de long), même sol, même couvert (sol nu), mais de pentes différentes (\SI{5}{\%}, \SI{10}{\%}, \SI{20}{\%}, \SI{30}{\%}). On simule une pluie de même intensité (ex. \SI{60}{mm/h}) et on mesure le débit de ruissellement et la quantité de terre perdue (en \SI{}{g/m^2} ou \SI{}{t/ha}).
\textbf{Résultats :} Le graphique ci-dessous synthétise les tendances classiques observées (équation universelle des pertes en terre - USLE, facteur topographique LS).
\textbf{Interprétation :} L'énergie cinétique du ruissellement est proportionnelle au carré de la vitesse ($E_c = \frac{1}{2}mv^2$). La vitesse augmente avec la racine carrée de la pente (loi de Manning-Strickler simplifiée). Donc l'érosivité augmente plus que proportionnellement à la pente.
\textbf{Conclusion :} \textbf{Plus la pente est forte, plus le ruissellement est rapide, plus son pouvoir de détachement et de transport des particules est grand.} La pente est le facteur naturel majeur de l'érosion hydrique.
\end{methode}

\begin{popfigure}
\begin{tikzpicture}
\begin{axis}[
    width=\linewidth,
    height=0.55\linewidth,
    axis lines=middle,
    xlabel={Pente (\% ou degrés)},
    ylabel={Vitesse du ruissellement / Perte de terre (unité relative)},
    xmin=0, xmax=40,
    ymin=0, ymax=100,
    xtick={0,10,20,30,40},
    ytick={0,25,50,75,100},
    domain=0:40,
    samples=100,
    grid=major,
    grid style={popline},
    tick label style={font=\small},
    label style={font=\small},
    clip=false,
    legend style={at={(0.02,0.98)}, anchor=north west, font=\small, fill=white, fill opacity=0.8, draw=popline}
]
\addplot[popcurve, thick, domain=0:40] {0.06*x^2}; % Approximation quadratique pour montrer l'emballement
\addlegendentry{Érosivité (pertes en terre $\propto$ pente$^{1.3\text{ à }1.5}$)}
\addplot[poporange, dashed, thick, domain=0:40] {1.5*x};
\addlegendentry{Vitesse ruissellement ($\propto \sqrt{\text{pente}}$)}
\node[popdark, anchor=south west, font=\small] at (axis cs:5,80) {Zone de culture raisonnable (< 15\%)};
\node[poppink, anchor=south west, font=\small] at (axis cs:25,80) {Zone à haut risque (> 25\%)};
\draw[poppink, thick, ->] (axis cs:25,75) -- (axis cs:30,60);
\draw[popgreen, thick, ->] (axis cs:5,75) -- (axis cs:10,60);
\end{axis}
\end{tikzpicture}
\legende{Évolution de la vitesse du ruissellement (linéaire avec $\sqrt{\text{pente}}$) et des pertes en terre (loi de puissance, exposant $\approx 1,4$) en fonction de la pente. Au-delà de 15-20 \%, l'érosion devient exponentielle. Source : adaptation modèle USLE/RUSLE.}
\end{popfigure}

\begin{propriete}[Loi de la pente (facteur topographique LS de l'USLE) -- Admise, justifiée par l'observation expérimentale ; non exigible en démonstration formelle au Probatoire]
L'érosion hydrique augmente selon une loi de puissance de la pente : les pertes en terre sont proportionnelles à $(\text{pente})^{1,3 \text{ à } 1,5}$. Une pente de 20 \% érode environ \textbf{4 à 6 fois plus} qu'une pente de 5 \%, toutes choses égales par ailleurs. Au Probatoire, on retient que \textbf{l'érosion augmente très vite avec la pente} ; on sait tracer et interpréter une courbe comme ci-dessus.
\end{propriete}

\textbf{2. L'agressivité du climat : pluies violentes et vents forts.}\par
Le climat fournit l'énergie motrice : l'énergie cinétique des gouttes de pluie (érosion par \emph{splash} ou battance) et l'énergie du vent.

\begin{itemize}
    \item \textbf{Pluviométrie intense et concentrée :} Au Cameroun, la zone équatoriale (Sud, Littoral) reçoit \SIrange{2000}{3000}{mm/an}, mais c'est l'\emph{intensité} horaire (mm/h) qui compte. Une goutte de \SI{5}{mm} tombant à \SI{9}{m/s} a une énergie cinétique 25 fois supérieure à une goutte de \SI{1}{mm} à \SI{4}{m/s}. Les orages tropicaux (avril-mai, septembre-octobre) déstructurent la surface du sol (battance), bouchent les pores (croûte de battance) et lancent les particules en l'air (érosion par splash), premier stade de l'érosion hydrique.
    \item \textbf{Vents forts en saison sèche (Harmattan) :} Dans le Grand Nord (Extrême-Nord, Nord, Adamaoua), l'harmattan souffle de novembre à mars avec des pointes $> \SI{50}{km/h}$. Sur sols nus, secs, sableux, le vent soulève les particules (déflation) : c'est l'érosion éolienne. La vitesse seuil de déflation pour du sable fin est $\approx \SIrange{15}{20}{km/h}$ à \SI{1}{m} du sol.
\end{itemize}

\begin{attention}[Piège fréquent : confondre pluviométrie totale et érosivité des pluies]
Une région à \SI{1500}{mm/an} bien répartis (bruine douce) érode moins qu'une région à \SI{1000}{mm/an} concentrée en 3 mois d'orages violents. \textbf{L'indice d'érosivité $R$ (facteur climatique de l'USLE) dépend de l'énergie cinétique $\times$ l'intensité maximale sur 30 min ($I_{30}$), pas du total annuel.} Ne pas écrire : « Il pleut beaucoup donc ça érode ». Écrire : « Les pluies sont intenses (orages tropicaux), donc forte énergie cinétique des gouttes, donc forte érosivité ».
\end{attention}

\textbf{3. La nature du sol : texture, structure et matière organique.}\par
Tous les sols ne résistent pas également à l'érosion. C'est la \textbf{sensibilité intrinsèque} (facteur $K$ de l'USLE).
\begin{itemize}
    \item \textbf{Sols sableux / limoneux peu structurés (ex. sols ferrugineux tropicaux lessivés du Nord, sols alluvionnaires bruts) :} Faible cohésion entre particules, faible teneur en matière organique, structure massive ou en plaques. Ils se déstructurent dès le premier choc de goutte (battance) $\rightarrow$ croûte de surface $\rightarrow$ ruissellement rapide $\rightarrow$ érosion massive. \cle{Sensibilité forte ($K$ élevé)}.
    \item \textbf{Sols argileux bien structurés, riches en humus (ex. sols ferrallitiques des forêts du Sud, sols volcaniques andosols du Mont Cameroun) :} Agrégats stables (complexe argilo-humique), porosité élevée, infiltration rapide. Ruissellement faible ou nul. \cle{Sensibilité faible ($K$ faible)}.
    \item \textbf{Sols gonflants (vertisols, plaine du Logone) :} Fissures en saison sèche $\rightarrow$ infiltration rapide au début de la pluie, mais gonflement rapide $\rightarrow$ ruissellement tardif mais puissant sur surface lisse.
\end{itemize}

\begin{exemplebox}[Exemple concret : Le bassin cotonnier du Nord-Cameroun]
Les sols sont souvent des \textbf{sols ferrugineux tropicaux lessivés} (texture sableuse à limoneuse, faible MO $< \SI{1}{\%}$, structure massive). Couplés à une pluviométrie de \SIrange{900}{1100}{mm} concentrée en 4 mois (orages violents) et à une topographie de piémonts (pentes \SIrange{5}{15}{\%}), ils présentent une \textbf{sensibilité à l'érosion très forte ($K > 0,3$)}. C'est pourquoi l'érosion y est un problème majeur pour la productivité cotonnière.
\end{exemplebox}

\courssub{Facteurs anthropiques : quand l'homme accélère l'érosion}

Les facteurs naturels définissent le \emph{potentiel} d'érosion. L'homme, par ses activités, réalise ce potentiel ou l'aggrave. On parle alors d'\textbf{érosion accélérée}.

\textbf{1. La déforestation et les feux de brousse : destruction du bouclier végétal.}\par
C'est la cause anthropique n°1 au Cameroun.
\begin{itemize}
    \item \textbf{Déforestation pour l'agriculture itinérante sur brûlis (Sud, Est), l'extension des cultures de rente (cacao, palmier à huile, hévéa), l'exploitation forestière (pistes de débardage), le bois de feu/charbon (périurbain Yaoundé/Douala).}
    \item \textbf{Feux de brousse (Nord, Adamaoua) :} Pratiqués pour le pâturage (herbe jeune), la chasse, ou par négligence. Ils détruisent la strate herbacée et la litière, laissant le sol nu face aux premiers orages violents (mai-juin) et aux vents d'harmattan.
\end{itemize}
Conséquence immédiate : disparition de l'interception pluviométrique, de la litière protectrice et du réseau racinaire. Le sol passe de « protégé » à « nu » en quelques heures/jours.

\textbf{2. Les pratiques agricoles inadaptées.}\par
\begin{itemize}
    \item \textbf{Culture dans le sens de la pente (labour haut-bas) :} Les sillons de labour deviennent des chenaux préférentiels pour le ruissellement. L'eau y prend de la vitesse et creuse des ravinements (ex. cultures de maïs/arachide sur piémonts de l'Adamaoua).
    \item \textbf{Labour excessif et travail du sol à l'envers :} Détruit la structure grumeleuse, enterre la matière organique de surface, crée un horizon de labour compact (semelle de labour) qui freine l'infiltration. Le sol devient poudreux $\rightarrow$ érodable au vent et à l'eau.
    \item \textbf{Culture sur brûlis (itinérante) :} Le feu minéralise la MO rapidement (pic de fertilité) mais détruit l'humus stable. Après 2-3 ans, le sol est épuisé, nu, érodé. Le paysan abandonne la parcelle (jachère longue) $\rightarrow$ dégradation du paysage (savane dégradée, \emph{Chromolaena odorata}).
    \item \textbf{Surpâturage (Nord, Adamaoua) :} Taux de chargement $>$ capacité de charge. Piétinement $\rightarrow$ compaction (densité apparente $> \SI{1,6}{g/cm^3}$) $\rightarrow$ infiltration nulle $\rightarrow$ ruissellement diffus $\rightarrow$ érosion en nappe puis ravinement. Détruit aussi la strate herbacée protectrice.
\end{itemize}

\begin{savaistu}[Le piétinement bovin : un compactage invisible mais meurtrier]
Au Nord-Cameroun, un troupeau de 50 zébus sur \SI{1}{ha} pendant la saison des pluies exerce une pression au sol équivalente à \SIrange{200}{300}{kPa} sous les sabots. En quelques semaines, la porosité macroporeuse (drainage) chute de \SI{15}{\%} à $< \SI{5}{\%}$. L'infiltration passe de \SI{50}{mm/h} à $< \SI{5}{mm/h}$. Le ruissellement augmente de \SI{300}{\%}, l'érosion de \SI{500}{\%}. Les pistes à bétail deviennent de véritables ravines (ex. autour de Maroua, Ngaoundéré).
\end{savaistu}

\courssub{Le rôle protecteur du couvert végétal : un écosystème « anti-érosion »}

Pourquoi la forêt ou une jachère dense protège-t-elle si bien ? Ce n'est pas seulement « les racines tiennent la terre ». C'est un \textbf{système multicouche} qui agit à chaque étape du processus érosif.

\begin{popfigure}
\begin{tikzpicture}[scale=0.9, every node/.style={font=\small}]
% Styles
\tikzset{
    soil/.style={fill=poporange!30, pattern=north east lines, pattern color=poporange!60},
    rock/.style={fill=popmuted!40, pattern=crosshatch, pattern color=popmuted},
    root/.style={popgreen, thick},
    water/.style={popblue, thick, ->},
    drop/.style={popblue, fill=popblue!20},
    canopy/.style={popgreen!60!black, fill=popgreen!20},
    litter/.style={popgold!60!black, fill=popgold!20},
    arrow/.style={-Stealth, thick, popdark},
    label/.style={popdark, anchor=west}
}

% ---- VERSANT BOISÉ (GAUCHE) ----
\begin{scope}[xshift=-0.5cm]
% Sol et roche
\draw[thick] (-6,0) -- (-6,-3) -- (0,-3) -- (0,0) -- cycle;
\fill[soil] (-6,0) rectangle (0,-1.5);
\fill[rock] (-6,-1.5) rectangle (0,-3);
% Litière
\draw[thick, popgold!60!black] (-6,0.1) -- (0,0.1);
\fill[litter] (-6,0) rectangle (0,0.15);
% Racines
\foreach \x in {-5.5,-4.5,-3.5,-2.5,-1.5,-0.5} {
    \draw[root] (\x,0) .. controls (\x,-0.5) and (\x-0.2,-1.5) .. (\x-0.3,-2.5);
    \draw[root] (\x,0) .. controls (\x+0.1,-0.8) and (\x+0.3,-2) .. (\x+0.2,-2.8);
}
% Arbre (tronc + canopée)
\draw[popgreen!60!black, line width=3pt] (-3,0) -- (-3,2.5);
\draw[canopy] (-4.5,2.5) ellipse (1.5 and 1);
\draw[canopy] (-3,3.2) ellipse (1.8 and 1.2);
\draw[canopy] (-1.5,2.5) ellipse (1.5 and 1);
% Gouttes de pluie interceptées
\foreach \i in {1,...,6} {
    \pgfmathsetmacro{\rx}{-5+rand*4}
    \pgfmathsetmacro{\ry}{3.5+rand*1}
    \fill[drop] (\rx,\ry) circle (0.08);
    \draw[popblue, thin] (\rx,\ry) -- (\rx,\ry-0.3);
}
% Ruissellement lent (flèches courbes)
\draw[water] (-5.5,0.2) .. controls (-5.5,0.5) and (-4,0.5) .. (-4,0.8) node[above, popdark, font=\small] {$v$ faible};
\draw[water] (-2,0.2) .. controls (-2,0.5) and (-1,0.5) .. (-1,0.8);
% Légende gauche
\node[popdark, align=center, font=\bfseries\small] at (-3,3.8) {VERSANT BOISÉ};
\node[label] at (0.3,2.5) {Canopée : interception gouttes};
\node[label] at (0.3,1.5) {Litière : éponge, infiltration};
\node[label] at (0.3,0.5) {Racines : fixation mécanique};
\node[label] at (0.3,-1) {Sol structuré, poreux};
\node[label] at (0.3,-2.5) {Roche mère};
\end{scope}

% ---- VERSANT DÉBOISÉ (DROITE) ----
\begin{scope}[xshift=7.5cm]
% Sol et roche
\draw[thick] (0,0) -- (0,-3) -- (6,-3) -- (6,0) -- cycle;
\fill[soil] (0,0) rectangle (6,-1.5);
\fill[rock] (0,-1.5) rectangle (6,-3);
% Croûte de battance
\draw[thick, poppink] (0,0.1) -- (6,0.1);
\fill[poppink!30] (0,0) rectangle (6,0.1);
% Racines résiduelles (peu)
\foreach \x in {1,3,5} {
    \draw[popgreen!50!black, thin, dashed] (\x,0) .. controls (\x,-0.3) and (\x,-1) .. (\x,-1.5);
}
% Gouttes de pluie impact direct (SPLASH)
\foreach \i in {1,...,8} {
    \pgfmathsetmacro{\rx}{rand*5}
    \pgfmathsetmacro{\ry}{3+rand*1}
    \fill[drop] (\rx,\ry) circle (0.08);
    \draw[popblue, thin] (\rx,\ry) -- (\rx,\ry-0.4);
    % Splash particles
    \foreach \j in {1,...,3} {
        \pgfmathsetmacro{\ax}{\rx+rand*0.5}
        \pgfmathsetmacro{\ay}{0.1+rand*0.3}
        \draw[poporange, very thin, ->] (\rx,0.1) -- (\ax,\ay);
    }
}
% Ruissellement rapide, concentré, ravinement
\draw[water, poppink] (1,0.2) -- (1.5,-0.5) -- (2,-1.2) node[right, poppink, font=\small] {Ravin (gully)};
\draw[water, poppink] (4,0.2) -- (4.5,-0.5) -- (5,-1.2);
\draw[water, poppink] (2.5,0.2) -- (3,-0.8);
% Flèches vitesse
\draw[arrow, poppink] (0.5,0.5) -- (1.5,0.5) node[above, poppink, font=\small] {$v$ forte};
\draw[arrow, poppink] (3.5,0.5) -- (4.5,0.5);
% Légende droite
\node[popdark, align=center, font=\bfseries\small] at (3,3.8) {VERSANT DÉBOISÉ};
\node[label, poppink] at (6.3,2.5) {Pas de canopée : splash direct};
\node[label, poppink] at (6.3,1.5) {Croûte de battance : imperméable};
\node[label, poppink] at (6.3,0.5) {Racines rares : sol non fixé};
\node[label, poppink] at (6.3,-1) {Ruissellement concentré};
\node[label, poppink] at (6.3,-2.5) {Ravinement, perte sol};
\end{scope}

% Flèche centrale comparatif
\draw[popdark, <->, thick, dashed] (-0.2,-3.5) -- (6.7,-3.5) node[midway, below, font=\small\bfseries] {Même pente, même pluie, même sol initial};
\end{tikzpicture}
\legende{Coupe schématique comparant un versant boisé (gauche) et un versant déboisé (droite) soumis à la même averse orageuse. À gauche : la canopée intercepte les gouttes, la litière absorbe l'eau, les racines fixent le sol $\rightarrow$ ruissellement lent, infiltré, eau claire. À droite : impact direct des gouttes (splash), formation d'une croûte de battance imperméable, ruissellement rapide et concentré creusant des ravines, eau chargée en sédiments (turbide).}
\end{popfigure}

\begin{definition}[Mécanismes de protection du couvert végétal]
Le couvert végétal (forêt, jachère, culture de couverture, paillis) agit par trois mécanismes complémentaires :
\begin{enumerate}
    \item \textbf{Interception pluviométrique (Canopée) :} Les feuilles et branches interceptent 15 à 30 \% des précipitations (jusqu'à 50 \% en forêt dense). L'énergie cinétique des gouttes est dissipée ; l'eau ruisselle le long des troncs (écoulement de tronc) ou tombe en gouttes grosses mais lentes (goutte-à-goutte), sans pouvoir de détachement.
    \item \textbf{Ralentissement du ruissellement (Litière + Tiges) :} La litière (feuilles mortes, débris) agit comme une éponge : elle absorbe l'eau, augmente la rugosité de surface (coefficient de Manning $n$ élevé), réduit la vitesse de surface ($v \propto 1/n$) et favorise l'infiltration (conductivité hydraulique $K_s$ maintenue élevée par la faune du sol).
    \item \textbf{Fixation mécanique du sol (Système racinaire) :} Les racines (fines et pivotantes) créent un \textbf{filet de renforcement} (effet « géotextile naturel »). Elles augmentent la cohésion apparente du sol ($c'$ en mécanique des sols) et créent des macropores (drainage). Une haie vetiver ou un arbre mature peut retenir des centaines de tonnes de terre sur un versant.
\end{enumerate}
\end{definition}

\begin{exemplebox}[Exemple chiffré : ordres de grandeur des pertes en terre (Cameroun, données IRAD/IRD)]
Sur un sol ferrugineux tropical en pente à 10 \%, sous pluie simulée de \SI{60}{mm/h} :
\begin{center}
\begin{tabularx}{\linewidth}{|l|X|X|}\hline
\textbf{Couvert végétal} & \textbf{Perte en terre (t/ha/an)} & \textbf{Coefficient de couverture $C$ (USLE)} \\ \hline
Forêt dense / Jachère ancienne ($>15$ ans) & $< \num{0,5}$ (souvent $< \num{0,1}$) & \num{0,001} -- \num{0,01} \\
Plantation pérenne (cacao, café, palmier) bien entretenue & \num{1} à \num{5} & \num{0,01} -- \num{0,1} \\
Culture annuelle (maïs, coton) \textbf{sans labour en courbes de niveau} & \num{20} à \num{60} & \num{0,2} -- \num{0,4} \\
Sol nu, labouré, en pente (jachère courte $<2$ ans) & \num{50} à \num{150} (jusqu'à \num{300}) & \num{0,5} -- \num{1,0} \\ \hline
\end{tabularx}
\end{center}
\textbf{Conclusion :} Passer d'un sol nu à un couvert forestier divise l'érosion par \textbf{100 à 1000}. C'est le levier le plus puissant de la lutte anti-érosive.
\end{exemplebox}

\begin{savaistu}[Les racines, un filet naturel haute performance]
Des mesures au pénétromètre et au cisaillement direct sur le versant du Mont Cameroun montrent que la cohésion du sol ($c'$) passe de \textbf{\SI{5}{kPa} (sol nu)} à \textbf{\SIrange{25}{40}{kPa}} sous forêt dense grâce aux racines fines ($< \SI{2}{mm}$) qui densifient l'horizon 0-30 cm. La résistance au cisaillement augmente de \SIrange{80}{150}{\%}. C'est pourquoi le \textbf{vetiver (\emph{Chrysopogon zizanioides})}, dont les racines descendent à \SIrange{3}{4}{m} en 1 an, est utilisé en haies anti-érosives sur les pentes de la Menoua et du Dschang : il réduit le ruissellement de \textbf{plus de \SI{80}{\%}} et les pertes en terre de \textbf{\SI{90}{\%}} par rapport au témoin nu.
\end{savaistu}

\courssub{Synthèse : l'érosion accélérée, fruit d'une interaction Nature-Homme}

L'érosion géologique (naturelle, lente, $\approx \SIrange{0,1}{1}{t/ha/an}$) est un processus lent de modelage des paysages. L'\textbf{érosion accélérée} que nous observons aujourd'hui (ravines de Bafia, ensablement du barrage de Lagdo, pertes de rendement dans le bassin cotonnier) résulte presque toujours de la \textbf{conjonction} :

\begin{center}
\textbf{Facteur naturel favorable (Pente + Pluie intense + Sol fragile)} \quad $\times$ \quad \textbf{Pression anthropique (Déforestation + Mauvaises pratiques + Surpâturage)}
\end{center}

\begin{propriete}[Règle générale : l'érosion accélérée = Facteur naturel $\times$ Facteur humain]
\begin{itemize}
    \item Sur pente forte ($> \SI{15}{\%}$) \textbf{ET} sol nu $\rightarrow$ Érosion catastrophique (ravinement, glissements).
    \item Sur pente forte \textbf{MAIS} couvert forestier dense $\rightarrow$ Érosion négligeable (ex. flancs du Mont Cameroun boisés).
    \item Sur pente faible \textbf{MAIS} sol nu, labouré, pluies violentes $\rightarrow$ Érosion diffuse en nappe (moins spectaculaire mais perte de fertilité majeure, ex. plaine du Logone, bassin cotonnier).
    \item \textbf{Le facteur humain est le levier d'action :} On ne peut pas changer la pente ni le climat, mais on \textbf{peut} restaurer le couvert végétal et adapter les pratiques culturales.
\end{itemize}
\end{propriete}

\begin{methode}[Identifier les causes de l'érosion dans une situation donnée (ex. photo, schéma, étude de cas)]
Suivez cette check-list systématique (méthode « PPPP » : \textbf{P}ente, \textbf{P}luie, \textbf{P}lante, \textbf{P}ratique) :
\begin{enumerate}
    \item \textbf{Analyser le relief (Pente) :} Forte ($> \SI{15}{\%}$) ? Modérée (\SIrange{5}{15}{\%}) ? Faible ($< \SI{5}{\%}$) ? Orientation des versants ? Présence de ravines existantes ?
    \item \textbf{Analyser le climat (Pluie/Vent) :} Zone à orages violents (Sud, Ouest) ? Zone à vents secs forts (Nord, Harmattan) ? Saisonnalité marquée ?
    \item \textbf{Analyser le sol (Nature) :} Texture (sableux = fragile) ? Structure (grumeleuse = stable / massive = fragile) ? Matière organique (couleur sombre = riche) ? Profondeur ?
    \item \textbf{Analyser le couvert végétal (Plante) :} Pourcentage de couverture au sol ($< \SI{30}{\%}$ = risque élevé) ? Type (forêt, jachère, culture, sol nu) ? Présence de litière ? Haies, bandes enherbées ?
    \item \textbf{Analyser les activités humaines (Pratique) :} Type de culture (itinérante, pérenne, intensive) ? Sens du labour (pente vs courbes de niveau) ? Pâturage (densité, parcours) ? Déforestation récente ? Pistes, routes non drainées ?
    \item \textbf{Synthétiser :} Croiser les facteurs. \emph{Exemple : « Érosion par ravinement due à la combinaison d'une pente de 20 \% (naturel), de sols sableux peu structurés (naturel), de la suppression de la forêt galerie pour le maïs (humain) et du labour dans le sens de la pente (humain). »}
\end{enumerate}
\end{methode}

\begin{attention}[Erreur fréquente au Probatoire : « La déforestation est la cause unique de l'érosion »]
\textbf{Faux.} La déforestation est un \emph{facteur déclenchant/aggravant} majeur, mais pas la seule cause.
\begin{itemize}
    \item Une pluie exceptionnelle (cyclone, orage stationnaire) sur une pente naturelle forte (falaises, volcans) érode même une forêt primaire (glissements de terrain, lahars). C'est de l'érosion \emph{naturelle} (ex. éboulements Mont Cameroun 1999/2000).
    \item L'érosion éolienne dans le Sahel est pilotée par la sécheresse climatique (vent + sol sec) avant même la surcharge pastorale.
    \item \textbf{Formulation correcte :} « L'érosion observée résulte de la conjugaison d'une forte érosivité climatique (pluies intenses / vents forts) et d'une forte érodibilité du sol (texture, pente), \emph{fortement amplifiée} par la mise à nu du sol par la déforestation et les cultures inadaptées. »
\end{itemize}
\end{attention}

\begin{exoresolu}[Analyse d'une situation érosive : Le bassin versant de la Menoua (Ouest Cameroun)]
\textbf{Énoncé :} Le bassin de la Menoua (affluent du Nkam) présente un réseau hydrographique dense, des pentes moyennes de \SIrange{15}{25}{\%}, des sols ferrallitiques lessivés (argileux mais peu profonds sur roche altérée), une pluviométrie de \SI{2000}{mm/an} bimodale. La population dense pratique l'agriculture intensive (maïs, haricot, pomme de terre) sur brûlis et labour moto-bineuse dans le sens de la pente. Les haies de vetiver ont été arrachées pour agrandir les champs. On observe de nombreux ravines (lavakas) et l'eau de la rivière est rouge (chargée en latérite) en saison des pluies. Identifiez les 4 facteurs principaux (2 naturels, 2 anthropiques) et proposez 3 techniques de lutte adaptées.

\textbf{Solution détaillée :}
\begin{enumerate}
    \item \textbf{Facteurs naturels :}
        \begin{itemize}
            \item \textbf{Pente forte (\SIrange{15}{25}{\%}) :} Facteur topographique majeur (LS élevé). Vitesse ruissellement élevée, transport de gros matériaux.
            \item \textbf{Agressivité climatique :} Pluviométrie élevée (\SI{2000}{mm}) + bimodalité = deux saisons d'érosion par an. Orages tropicaux intenses = forte érosivité $R$.
        \end{itemize}
    \item \textbf{Facteurs anthropiques :}
        \begin{itemize}
            \item \textbf{Labour dans le sens de la pente + absence de couverture (brûlis, arrachage vetiver) :} Sol nu pendant la mise en place des cultures (mars-avril, août-septembre) coïncidant avec les pics de pluies. Sillons = chenaux d'érosion. $C \approx 0,3-0,4$.
            \item \textbf{Pression foncière / intensification sans conservation :} Pas de jachère, pas de haies, pas de terrasses. Réduction de la matière organique $\rightarrow$ déstructuration.
        \end{itemize}
    \item \textbf{Techniques de lutte adaptées (3 minimum) :}
        \begin{enumerate}
            \item \textbf{Cultures en courbes de niveau (billons/ridelles perpendiculaires à la pente) :} Réduit la vitesse du ruissellement, favorise l'infiltration, piège les sédiments. \emph{Action sur le facteur pente (réduction effective de la longueur de pente $L$).}
            \item \textbf{Réimplantation de haies de vetiver (ou \emph{Tephrosia}, \emph{Calliandra}) en courbes de niveau tous les \SIrange{10}{15}{m} de dénivelé :} Barrière physique racinaire + ralentissement ruissellement + apport de paillis. \emph{Action sur le facteur couvert végétal (C $\downarrow$) et pente (L $\downarrow$).}
            \item \textbf{Agroforesterie / Jachère améliorée (ex. \emph{Mucuna}, \emph{Crotalaria} en intercalaire) :} Couvre le sol entre les rangs de maïs, fixe l'azote, racines pivotantes. \emph{Action sur le facteur couvert (C $\downarrow$) et fertilité.}
            \item \textit{(Bonus) Aménagement de terrasses en gradins (bench terraces) sur les pentes $> \SI{25}{\%}$ :} Transformation de la pente en marches plates. Coûteux en main-d'œuvre mais durable. \emph{Action directe sur le facteur pente (LS $\downarrow$ drastique).}
        \end{enumerate}
\end{enumerate}
\end{exoresolu}

\begin{aretenir}[Bilan : Les facteurs favorisants l'érosion]
\begin{itemize}
    \item \textbf{Naturels (Potentiel) :} Pente forte (énergie gravitaire), Pluies intenses/orages (énergie cinétique gouttes), Vents forts/secs (énergie éolienne), Sols sableux/peu structurés/pauvres en MO (faible cohésion).
    \item \textbf{Anthropiques (Déclencheur/Aggravant) :} Déforestation/Feux de brousse (suppression bouclier), Labour sens pente (chenaux), Labour excessif (destruction structure), Surpâturage (compaction + dénudation), Culture sur brûlis itinérante (cycle court, sol nu).
    \item \textbf{Bouclier naturel :} Couvert végétal (Canopée $\rightarrow$ Interception ; Litière $\rightarrow$ Infiltration/Rugosité ; Racines $\rightarrow$ Cohésion mécanique). Perte de terre : Forêt $< \SI{1}{t/ha/an}$ vs Sol nu $> \SI{50}{t/ha/an}$.
    \item \textbf{Érosion accélérée = Nature $\times$ Homme.} L'action porte sur le facteur humain (restauration couvert, pratiques culturales).
\end{itemize}
\end{aretenir}

\coursec{Les conséquences de l'érosion}

\courssub{Transition : des causes aux effets}

Dans la section précédente, nous avons identifié les \cle{facteurs favorisants} de l'érosion : la déforestation, la pente, les pratiques culturales inadaptées, l'intensité des pluies et l'action du vent. Nous comprenons maintenant \emph{comment} et \emph{pourquoi} le sol se détache et se transporte. Mais une question cruciale se pose pour les populations, les agriculteurs et les aménageurs du territoire : \textbf{quelles sont les répercussions concrètes de ce déplacement de particules ?} L'érosion n'est pas seulement un phénomène géologique abstrait ; elle menace directement la sécurité alimentaire, l'économie nationale et la santé des écosystèmes aquatiques. Observons d'abord une situation vécue dans nos campagnes.

\begin{experience}[Observation de terrain : le champ de M. Njoya à Dschang]
\textbf{Matériel :} Photos aériennes successives (2010, 2020) d'une parcelle en pente cultivée en maïs sans aménagement ; mesures de rendement à l'hectare sur 10 ans ; prélèvements d'eau dans le ruisseau en aval après une forte pluie.

\textbf{Observations :}
\begin{itemize}
    \item Sur les photos : apparition de sillons profonds (ravines) traversant la parcelle ; la couleur du sol passe du brun foncé (humifère) au jaune/rouge (horizon sous-jacent).
    \item Rendements : \SI{3,5}{t/ha} en 2010 $\rightarrow$ \SI{1,2}{t/ha} en 2020 (baisse de \SI{66}{\%}).
    \item Eau du ruisseau : trouble, couleur ocre, dépôt rapide de boue au fond du verre.
\end{itemize}

\textbf{Interprétation :} L'érosion hydrique a emporté la terre arable (riche en matière organique et éléments nutritifs). Le sol restant est moins fertile, d'où la chute des rendements. Les particules emportées chargent le cours d'eau : c'est l'\cle{ensablement}.

\textbf{Conclusion :} L'érosion dégrade le capital-sol \emph{sur place} (perte de fertilité) et dégrade l'environnement \emph{en aval} (pollution, colmatage). C'est cette double peine que nous allons analyser.
\end{experience}

\courssub{Conséquences sur le sol : appauvrissement et perte du capital fertilité}

\begin{definition}[Horizon fertile et capital-sol]
L'\cle{horizon fertile} (souvent l'horizon A ou humifère) est la couche superficielle du sol (\SI{20}{--30}{cm} en zone tropicale humide), riche en \cle{matière organique} (humus), en \cle{éléments nutritifs} (azote, phosphore, potassium, cations échangeables) et en \cle{biocénose du sol} (bactéries, champignons, vers de terre). C'est le \textbf{capital productif} de l'agriculteur. L'érosion agit comme un \emph{prélèvement forcé} sur ce capital.
\end{definition}

\begin{popfigure}
\begin{tikzpicture}[scale=0.9, every node/.style={font=\small}]
    % Coupe de sol initial (gauche)
    \draw[thick, fill=popcream] (0,0) rectangle (6,-5);
    % Horizon O (litière)
    \draw[fill=brown!60!black] (0,0) rectangle (6,-0.3);
    \node[align=center] at (3, -0.15) {\textbf{Horizon O} : Litière};
    % Horizon A (humifère)
    \draw[fill=brown!40!black] (0,-0.3) rectangle (6,-1.3);
    \node[align=center] at (3, -0.8) {\textbf{Horizon A} : Humifère, fertile, structuré};
    % Horizon B
    \draw[fill=orange!60!yellow] (0,-1.3) rectangle (6,-3.0);
    \node[align=center] at (3, -2.15) {\textbf{Horizon B} : Argileux, moins fertile};
    % Horizon C / Roche mère
    \draw[fill=gray!50] (0,-3.0) rectangle (6,-5.0);
    \node[align=center] at (3, -4.0) {\textbf{Horizon C / Roche mère}};
    
    \node[align=center, font=\bfseries] at (3, 0.5) {Sol sain (T0)};
    
    % Flèche temps
    \draw[->, thick, popblue] (7, -2.5) -- (9, -2.5) node[midway, above] {Érosion sur 10-20 ans};
    
    % Coupe de sol érodé (droite)
    \draw[thick, fill=popcream] (10,0) rectangle (16,-5);
    % Horizon O : disparu ou très fin
    \draw[fill=brown!60!black] (10,0) rectangle (16,-0.1);
    \node[align=center] at (13, -0.05) {Litière fine};
    % Horizon A : TRÈS AMINCI
    \draw[fill=brown!40!black] (10,-0.1) rectangle (16,-0.6);
    \node[align=center, text width=3cm] at (13, -0.35) {\textbf{Horizon A résiduel} : \textcolor{popdark}{Très aminci}, perte humus};
    % Horizon B : AFFLEURE
    \draw[fill=orange!60!yellow] (10,-0.6) rectangle (16,-3.0);
    \node[align=center] at (13, -1.8) {\textbf{Horizon B} : Affleure, compact, pauvre};
    % Horizon C
    \draw[fill=gray!50] (10,-3.0) rectangle (16,-5.0);
    \node[align=center] at (13, -4.0) {\textbf{Roche mère}};
    
    \node[align=center, font=\bfseries] at (13, 0.5) {Sol érodé (T1)};
    
    % Flèches d'érosion (particules partant)
    \foreach \y in {0.2, 0.5, 0.8} {
        \draw[->, poporange, thick, decorate, decoration={snake, amplitude=1pt, segment length=4pt}] (6, -\y) -- (7, -\y-0.2);
    }
    \node[poporange, align=center] at (6.5, -1.5) {Particules\\ espacées\\ humifères};
\end{tikzpicture}
\legende{Évolution d'un profil de sol sous l'effet de l'érosion hydrique de surface. L'horizon A (brun foncé), réservoir de fertilité, est sélectivement emporté, laissant affleurer l'horizon B (jaune/orange), plus argileux, compact et stérile. La régénération naturelle de l'horizon A prend des siècles.}
\end{popfigure}

\begin{propriete}[Sélectivité de l'érosion]
L'érosion n'emporte pas le sol de manière homogène. Elle est \cle{sélective} : elle privilégie les particules fines (argiles, limons, matière organique) qui sont les plus légères et les plus riches en nutriments. Les sables grossiers et graviers restent sur place. \textbf{Conséquence :} Le sol résiduel s'appauvrit en argile et en humus $\rightarrow$ perte de la structure grumeleuse $\rightarrow$ baisse de la capacité de rétention en eau (CRA) et de la capacité d'échange cationique (CEC).
\end{propriete}

\begin{aretenir}[Bilan sur le sol]
L'érosion entraîne :
\begin{enumerate}
    \item La diminution de l'épaisseur du sol arable (perte de \SI{1}{--50}{t/ha/an} selon la sévérité).
    \item La baisse drastique du taux de matière organique (souvent divisé par 2 à 3 en 10 ans).
    \item La dégradation des propriétés physiques : compaction, croûtage de surface, imperméabilisation $\rightarrow$ amplification du ruissellement (rétroaction positive).
\end{enumerate}
\end{aretenir}

\courssub{Conséquences sur l'agriculture et la sécurité alimentaire}

\begin{exemplebox}[Cas concret : Bassin arachidier du Nord-Cameroun (Région du Nord, Extrême-Nord)]
Dans le département du Mayo-Rey, la culture du coton et de l'arachide sur sols sableux (ferrugineux tropicaux lessivés) sans jachère ni haies vives a entraîné une érosion éolienne et hydrique sévère.
\begin{itemize}
    \item \textbf{Rendements :} Passés de \SI{1,8}{t/ha} (coton graine) dans les années 80 à \SI{0,6}{t/ha} actuellement sans engrais minéraux massifs.
    \item \textbf{Réponse des paysans :} Extension des surfaces cultivées (déforestation des dernières zones de parcours) pour maintenir la production totale $\rightarrow$ aggravation de l'érosion.
    \item \textbf{Insécurité alimentaire :} Les greniers se vident plus vite ; dépendance accrue aux importations de riz et de blé.
\end{itemize}
\end{exemplebox}

\begin{methode}[Lier érosion et rendement (Démarche d'investigation)]
Pour quantifier l'impact agronomique :
\begin{enumerate}
    \item \textbf{Observer :} Comparer deux parcelles voisines (même sol, même pente) : l'une protégée (haies, courbes de niveau), l'autre non.
    \item \textbf{Mesurer :} Épaisseur de l'horizon A (au pénétromètre ou à la bêche), taux de MO (perte au feu ou colorimétrie), rendement à la récolte.
    \item \textbf{Corréler :} Tracer le graphique Rendement = f(Épaisseur horizon A) ou Rendement = f(Taux MO).
    \item \textbf{Conclure :} La relation est souvent linéaire positive jusqu'à un seuil. En dessous de \SI{10}{--15}{cm} d'horizon A, le rendement s'effondre.
\end{enumerate}
\end{methode}

\begin{attention}[Piège fréquent : Confondre érosion et stérilité naturelle]
Un sol peut être naturellement pauvre (ex: sols ferrallitiques très lessivés du Sud-Cameroun, pauvres en bases) \textbf{sans être érodé}. L'érosion \emph{aggrave} une contrainte existante. Au Probatoire, si un document montre un sol pauvre, ne concluez pas systématiquement « c'est l'érosion » sans indice de perte de matériel (ravines, racines apparentes, dépôts en bas de pente).
\end{attention}

\courssub{Conséquences sur le relief : morphogenèse destructrice}

L'érosion sculpte le paysage. Lorsque le ruissellement se concentre, il creuse des formes caractéristiques.

\begin{definition}[Ravin, Badlands, Pédiment]
\begin{itemize}
    \item \cle{Ravin} : Sillon profond (profondeur $>$ \SI{30}{cm}, largeur $>$ \SI{50}{cm}) creusé par l'eau concentrée, impossible à combler par le labour ordinaire. Il morcelle la parcelle.
    \item \cle{Badlands} (Terres mauvaises) : Paysage profondément découpé par un réseau dense de ravines imbriquées, impropre à l'agriculture (ex: versants du Mont Cameroun, piémonts de l'Adamaoua, Mandara).
    \item \cle{Pédiment} : Surface d'érosion plane, rocheuse ou graveleuse, mise à nu après arrachage total du manteau d'altération.
\end{itemize}
\end{definition}

\begin{popfigure}
\begin{tikzpicture}[scale=0.85, every node/.style={font=\small}]
    % Profil initial : pente douce, sol continu
    \draw[thick, fill=popgreenL!50] (0,0) -- (10,0) -- (10,-0.5) -- (0,-0.5) -- cycle; % Sol
    \draw[thick, fill=brown!40!black] (0,0) -- (10,0) -- (10,-1) -- (0,-1) -- cycle; % Horizon A épais
    \draw[thick, fill=orange!50] (0,-1) -- (10,-1) -- (10,-3) -- (0,-3) -- cycle; % Roche
    \node at (5, -0.5) {Sol profond, continu};
    \node at (5, 0.5) {\textbf{Stade 1 : Pente convexe stable}};
    
    % Flèche temps
    \draw[->, thick, popblue] (11, -1.5) -- (13, -1.5);
    
    % Stade 2 : Apparition ravines
    \draw[thick, fill=popgreenL!50] (14,0) -- (24,0) -- (24,-0.5) -- (14,-0.5) -- cycle;
    \draw[thick, fill=brown!40!black] (14,0) -- (24,0) -- (24,-1) -- (14,-1) -- cycle;
    \draw[thick, fill=orange!50] (14,-1) -- (24,-1) -- (24,-3) -- (14,-3) -- cycle;
    % Ravines
    \fill[orange!70!black] (16,0) to[out=-90, in=180] (16.5,-1.5) to[out=0, in=-90] (17,0) -- cycle;
    \fill[orange!70!black] (20,0) to[out=-90, in=180] (20.8,-2) to[out=0, in=-90] (21.6,0) -- cycle;
    \node at (19, -0.5) {Ravines en V};
    \node at (19, 0.5) {\textbf{Stade 2 : Dissection (Ravinement)}};
    
    \draw[->, thick, popblue] (25, -1.5) -- (27, -1.5);
    
    % Stade 3 : Badlands
    \draw[thick, fill=orange!50] (28,0) -- (38,0) -- (38,-3) -- (28,-3) -- cycle;
    % Réseau dendritique de ravines
    \foreach \x/\y/\w in {30/0/2, 32/0/3, 34/0/1.5, 36/0/2.5, 29/0/1, 37/0/1} {
        \fill[orange!70!black] (\x,\y) to[out=-90, in=180] ({ \x + \w/2 }, -2.5) to[out=0, in=-90] ({ \x + \w }, \y) -- cycle;
    }
    \node at (33, -1.5) {Sol quasi absent};
    \node at (33, 0.5) {\textbf{Stade 3 : Badlands (Irréversible à court terme)}};
\end{tikzpicture}
\legende{Évolution morphologique d'un versant sous érosion hydrique concentrée. Le morcellement des terres cultivables (stade 2) précède la stérilisation totale (stade 3). La réhabilitation des badlands coûte extrêmement cher (génie civil + revégétalisation).}
\end{popfigure}

\courssub{Conséquences sur l'eau : ensablement, pollution et colmatage}

\begin{definition}[Ensablement et Sédimentation]
L'\cle{ensablement} est l'accumulation de sédiments (sables, limons, argiles) transportés par l'érosion, dans le lit des cours d'eau, les lacs, les retenues de barrages ou les zones humides. La \cle{sédimentation} est le processus de dépôt de ces particules en suspension lorsque la vitesse du courant diminue. C'est le \textbf{devenir « aval »} de l'érosion « amont ».
\end{definition}

\begin{exemplebox}[L'ensablement du bassin de la Bénoué et du barrage de Lagdo]
Le fleuve Bénoué, principal affluent du Niger au Cameroun, draine un vaste bassin versant incluant les hauts plateaux de l'Adamaoua (pentes fortes, sols ferrugineux, déforestation pour l'élevage et l'agriculture).
\begin{itemize}
    \item \textbf{Données :} Le débit solide (charriage + suspension) à Garoua est estimé à \SI{20}{--30} millions de tonnes/an.
    \item \textbf{Conséquence directe :} Le barrage hydroélectrique de \textbf{Lagdo} (mis en service 1982) a perdu plus de \SI{30}{\%} de sa capacité de stockage initiale (passé de \SI{7,7}{km^3} à $\approx$ \SI{5}{km^3}) à cause du colmatage par les boues. Cela réduit la production d'électricité (EDC) et la régulation des crues.
    \item \textbf{Navigation :} Le chenal navigable (Garoua - Niamey) s'ensable, freinant le transport fluvial.
    \item \textbf{Qualité de l'eau :} Turbidité extrême en saison des pluies (jusqu'à \SI{2000}{NTU}), coût de potabilisation élevé pour la SODECOTON/SNEC, perturbation de la faune ichtyologique (asphyxie des œufs, colmatage des frayères).
\end{itemize}
\end{exemplebox}

\begin{attention}[Portée spatiale : L'érosion n'a pas de frontières administratives]
Les sédiments produits sur les versants de l'Adamaoua (Région Adamaoua) polluent l'eau et colmatent le barrage à Lagdo (Région Nord) et perturbent la navigation jusqu'au Niger. \textbf{Au Probatoire :} Si on vous donne une carte du bassin versant, identifiez les zones \emph{sources} (amont, pentes fortes, déforestation) et les zones \emph{puits} (aval, plaines, barrages, deltas). La lutte doit se faire \emph{en amont} (protection des versants) pour protéger l'aval.
\end{attention}

\courssub{Conséquences sur les infrastructures et l'habitat}

L'érosion n'épargne pas le bâti. Les ravines progressent par érosion régressive (recul de la tête de ravin) et menacent les axes de communication.

\begin{savaistu}[Coût économique au Cameroun]
Selon des études de la Banque Mondiale et du MINEPAT, le coût annuel de la dégradation des sols (érosion + perte fertilité) est estimé entre \textbf{\SI{4}{--8}{\%} du PIB agricole}. La réhabilitation d'un km de route nationale emporté par une ravine (ex: RN1 au niveau de Nkongsamba, RN6 vers l'Est) coûte \SI{50}{--200} millions FCFA/km. Les glissements de terrain (coulées de boue) sur les pentes du Mont Cameroun (Buea, Limbe) ou de Yaoundé (quartiers Mvog-Ada, Nkolbisson) détruisent des habitations et font des victimes chaque saison des pluies.
\end{savaistu}

\begin{popfigure}
\begin{tikzpicture}[scale=0.9, every node/.style={font=\small}]
    % Coupe transversale vallée avec route et ravin
    \draw[thick, fill=popcream] (0,0) -- (12,0) -- (12,-4) -- (0,-4) -- cycle;
    % Sol
    \draw[fill=brown!30] (0,0) -- (12,0) -- (12,-1.5) -- (0,-1.5) -- cycle;
    \draw[fill=orange!40] (0,-1.5) -- (12,-1.5) -- (12,-4) -- (0,-4) -- cycle;
    
    % Route en haut de pente
    \draw[fill=gray!30, thick] (1,0) -- (5,0) -- (5,0.4) -- (1,0.4) -- cycle;
    \node at (3, 0.2) {Route Nationale};
    \draw[<->, popdark] (1,0.5) -- (5,0.5) node[midway, above] {Largeur 8m};
    
    % Ravin qui menace la route (érosion régressive)
    \fill[orange!60!black] (5,0) to[out=-90, in=30] (6.5,-2) to[out=150, in=-60] (5.5,-0.5) -- (5,0) -- cycle;
    \draw[thick, popdark] (5,0) to[out=-90, in=30] (6.5,-2) to[out=150, in=-60] (5.5,-0.5);
    \node[popdark, align=center] at (6.5, -0.5) {Tête de ravin\\ (progression $\rightarrow$)};
    
    % Maison en bas de pente (risque coulée)
    \draw[fill=popblueL, thick] (8, -3.5) rectangle (10.5, -2.5);
    \draw[fill=popblue] (9.25, -2.5) -- (10.5, -2.0) -- (9.25, -1.5) -- cycle; % Toit
    \node at (9.25, -3.0) {Habitat};
    \draw[->, thick, poporange] (7, -1.5) -- (9, -2.8) node[midway, above right] {Coulée de boue};
    
    % Rivière en bas
    \draw[fill=popblue!30] (0,-3.8) -- (12,-3.8) -- (12,-4) -- (0,-4) -- cycle;
    \node at (6, -3.9) {Rivière ensablée};
\end{tikzpicture}
\legende{Menaces de l'érosion sur les infrastructures : progression d'une ravinement par érosion régressive vers une route (haut de pente) et risque de coulée de boue sur l'habitat en bas de versant. La protection coûte moins cher que la réparation.}
\end{popfigure}

\courssub{Conséquences économiques et sociales : le cercle vicieux de la pauvreté}

C'est ici que la dimension humaine prend tout son sens. L'érosion n'est pas qu'une affaire de techniciens ; elle structure la société rurale.

\begin{popfigure}
\begin{tikzpicture}[
    scale=0.85,
    every node/.style={font=\small, align=center},
    box/.style={draw, thick, rounded corners, fill=popcream, minimum width=2.8cm, minimum height=1.2cm},
    arrow/.style={->, thick, popdark, >=Stealth}
]
    % Nœuds
    \node[box, fill=poporangeL, align=center] (E) at (0,0){\textbf{ÉROSION}\\ (Perte sol fertile)};
    \node[box, fill=poporangeL, align=center] (F) at (5,0){\textbf{BAISSE FERTILITÉ}\\ \& RENDEMENTS};
    \node[box, fill=poppinkL, align=center] (P) at (10,0){\textbf{PAUVRETÉ}\\ REVENUS AGRICOLES};
    \node[box, fill=poppinkL, align=center] (D) at (5,-3){\textbf{DÉFORESTATION}\\ NOUVELLES TERRES};
    \node[box, fill=popblueL, align=center] (I) at (0,-3){\textbf{INSÉCURITÉ}\\ ALIMENTAIRE};
    \node[box, fill=popblueL, align=center] (M) at (10,-3){\textbf{EXODE RURAL}\\ CONFLITS FONCIERS};
    
    % Flèches du cycle principal
    \draw[arrow] (E) -- (F) node[midway, above] {entraîne};
    \draw[arrow] (F) -- (P) node[midway, above] {provoque};
    \draw[arrow] (P) -- (D) node[midway, above] {pousse à};
    \draw[arrow] (D) -- (E) node[midway, left] {aggrave};
    
    % Conséquences sociales
    \draw[arrow] (P) -- (I) node[midway, left] {génère};
    \draw[arrow] (P) -- (M) node[midway, right] {provoque};
    \draw[arrow] (I) -- (M) node[midway, below] {alimente};
    
    % Flèche de rétroaction forte (cercle vicieux)
    \draw[arrow, poporange, very thick, bend left=40] (D) to node[midway, left, poporange] {\textbf{CERCLE VICIEUX}} (E);
\end{tikzpicture}
\legende{Le cercle vicieux de l'érosion et de la pauvreté rurale. La perte de fertilité (cause biophysique) entraîne la pauvreté (cause économique) qui force l'agriculteur à déforester de nouvelles terres marginales (réponse adaptative), relançant l'érosion. Briser ce cycle nécessite une intervention externe (aménagement, appui technique, sécurité foncière).}
\end{popfigure}

\begin{aretenir}[Synthèse des conséquences systémiques]
\begin{center}
\begin{tabularx}{\linewidth}{|l|X|X|}\hline
\textbf{Compartiment} & \textbf{Conséquences directes} & \textbf{Conséquences induites (systémiques)} \\ \hline
\textbf{Sol} & Perte horizon A, baisse MO, CEC, CRA & Stérilisation définitive (badlands), coût amendment élevé \\ \hline
\textbf{Agriculture} & Chute rendements, abandon parcelles & Insécurité alimentaire, importations, dépendance engrais \\ \hline
\textbf{Eau / Milieu aquatique} & Turbidité, ensablement lits/barrages, eutrophisation (N, P) & Perte production hydroélectrique (Lagdo), coût potabilisation, perte biodiversité \\ \hline
\textbf{Infrastructures} & Coupure routes, ponts, destruction habitats & Isolement zones rurales, coûts réhabilitation budgétaires lourds \\ \hline
\textbf{Social / Économique} & Perte revenus, exode rural, conflits fonciers & Urbanisation anarchique, pression sur forêts résiduelles, instabilité \\ \hline
\end{tabularx}
\end{center}
\end{aretenir}

\courssub{Exercice résolu : Analyse intégrée d'une situation d'érosion}

\begin{exoresolu}[Cas du bassin versant de la Menoua (Ouest-Cameroun)]
\textbf{Énoncé :} La Menoua prend sa source sur les hauts plateaux de l'Ouest (Bamendjou, Bafang) à \SI{1500}{m} d'altitude. La densité de population y dépasse \SI{200}{hab/km^2}. L'agriculture intensive (café, maïs, haricot) se pratique sur de fortes pentes ($> \SI{30}{\%}$) sans terrasses. Le fleuve charrie des quantités massives de sédiments vers le lac de retenue de \textbf{Mékinac} (barrage hydroélectrique projeté) et la Sanaga.
\begin{enumerate}
    \item Identifier, dans ce texte, 3 facteurs favorisants l'érosion.
    \item Citer 4 conséquences de cette érosion observables à trois échelles spatiales différentes (champ, fleuve, barrage).
    \item Expliquer pourquoi la construction de terrasses en courbes de niveau permet de casser le cercle vicieux identifié plus haut.
\end{enumerate}

\textbf{Solution détaillée :}
\begin{enumerate}
    \item \textbf{Facteurs favorisants (rappel section 3) :}
        \begin{itemize}
            \item \cle{Pente forte} ($> \SI{30}{\%}$) : énergie cinétique du ruissellement élevée.
            \item \cle{Pratiques culturales} : labour dans le sens de la pente, absence de haies/terrasses, sol nu après récolte.
            \item \cle{Pression démographique forte} : raccourcissement des jachères, culture de terres marginales (versants raides).
        \end{itemize}
    \item \textbf{Conséquences multi-échelles :}
        \begin{center}
        \begin{tabular}{|l|l|l|}\hline
        \textbf{Échelle} & \textbf{Lieu} & \textbf{Conséquence observable} \\ \hline
        Parcellaire (Champ) & Hauts plateaux & Ravines dans les champs, racines des caféiers apparentes, rendements maïs $< \SI{1}{t/ha}$. \\ \hline
        Bassin versant (Fleuve) & Vallée Menoua & Eau très trouble (turbidité $> \SI{1000}{NTU}$) en saison des pluies, lit surélevé par dépôts sableux, inondations fréquentes. \\ \hline
        Infrastructures (Barrage) & Retenue de Mékinac & Colmatage rapide de la retenue (perte volume mort), usure prématurée des turbines (abrasion par sables), coûts de dragage. \\ \hline
        \end{tabular}
        \end{center}
    \item \textbf{Rôle des terrasses (cassage du cercle vicieux) :}
        \begin{itemize}
            \item \textbf{Physique :} Réduction de la longueur de pente effective $\rightarrow$ baisse vitesse ruissellement $\rightarrow$ dépôt sédiments \emph{dans} la terrasse (piégeage) au lieu du fleuve.
            \item \textbf{Agronomique :} Maintien de l'horizon A sur place $\rightarrow$ fertilité conservée $\rightarrow$ rendements stables/augmentés $\rightarrow$ revenus maintenus.
            \item \textbf{Social :} L'agriculteur n'a plus besoin de défricher la forêt galerie ou les pentes supérieures pour compenser la perte de rendement. La pression foncière diminue.
        \end{itemize}
\end{enumerate}
\tcblower
\textbf{Conclusion :} La terrasse agit sur la \emph{cause biophysique} (vitesse eau) et brise le \emph{retour socio-économique} (pauvreté $\rightarrow$ déforestation). C'est une mesure \cle{curative} (répare) et \cle{préventive} (protège l'aval).
\end{exoresolu}

\begin{methode}[Interpréter un document (photo/schéma) montrant l'érosion : grille de lecture Probatoire]
Face à une photo aérienne, une coupe de sol ou un paysage de badlands :
\begin{enumerate}
    \item \textbf{Décrire objectivement :} Couleurs (sol nu rouge/jaune vs vert forêt), formes (ravines en V, dendritiques, terrasses), infrastructures (routes coupées, maisons menacées), eau (trouble, claire).
    \item \textbf{Identifier le type d'érosion :} Hydrique (ravines, gullies, érosion diffuse) ou Éolienne (dunes, dépôts sableux sur cultures, racines aériennes).
    \item \textbf{Repérer les facteurs visibles :} Pente, état de la couverture végétale (nu, culture, forêt), traces de labour (sens pente vs courbes niveau), présence/absence d'aménagements.
    \item \textbf{Lister les conséquences visibles :} Perte sol (épaisseur), ensablement (dépôts aval), destruction (maison, route), pollution eau.
    \item \textbf{Proposer des mesures adaptées :} Lier chaque facteur/conséquence à une technique (reboisement, terrasses, fascines, haies vives, enrochement, barrage de retenue).
\end{enumerate}
\end{methode}

\begin{exercice}[Application : Diagnostic d'une photo de ravinement]
\textit{On vous présente une photo prise dans la région de Ngaoundéré (Adamaoua) montrant un large ravin (largeur \SI{15}{m}, profondeur \SI{8}{m}) coupant une pâture à \textit{Imperata cylindrica}. Les bords du ravin sont instables, des éboulis sont visibles au fond. En arrière-plan, on distingue des champs de maïs labourés dans le sens de la pente jusqu'au bord du ravin.}
\begin{enumerate}
    \item Qualifier le stade d'érosion observé.
    \item Citer deux facteurs anthropiques visibles sur le document qui ont déclenché/aggravé ce ravin.
    \item Proposer deux techniques de lutte adaptées pour stabiliser ce ravin et protéger les champs en amont.
    \item Quelle conséquence majeure ce ravin a-t-il sur le fleuve Bénoué situé \SI{50}{km} en aval ?
\end{enumerate}
\end{exercice}

\begin{corrige}[Exercice : Diagnostic d'une photo de ravinement]
\begin{enumerate}
    \item \textbf{Stade :} \cle{Ravinement avancé} (ravin majeur, dimensions $>$ \SI{2}{m} profondeur/largeur, bords instables). C'est une forme irréversible par simple labour ; nécessite du génie végétal/génie civil.
    \item \textbf{Facteurs anthropiques :}
        \begin{itemize}
            \item Labour \emph{dans le sens de la pente} (création de sillons préférentiels pour le ruissellement concentré).
            \item Absence de \cle{bande enherbée} ou de \cle{haie vive} en bordure de ravin (cultivation jusqu'au bord = pas de zone tampon).
        \end{itemize}
    \item \textbf{Techniques adaptées :}
        \begin{enumerate}
            \item \textbf{Aménagement du ravin (curatif) :} Construction de \cle{seuils de ravine} (en pierres sèches, gabions ou béton) pour casser la pente longitudinale, stopper l'érosion régressive et piéger les sédiments. Revégétalisation des berges (bouturage de \textit{Gliricidia}, \textit{Leucaena}, vetiver).
            \item \textbf{Protection de l'amont (préventif) :} Mise en place de \cle{cultures en courbes de niveau} + \cle{haies vives} (ex: \textit{Calliandra}, \textit{Flemingia}) sur les champs en amont pour infiltrer l'eau avant qu'elle n'atteigne le ravin.
        \end{enumerate}
    \item \textbf{Conséquence sur la Bénoué :} Apport massif de \cle{sédiments fins} (argiles, limons) et \cle{d'éléments nutritifs} (azote, phosphore des engrais) $\rightarrow$ \cle{ensablement} du lit mineur, \cle{eutrophisation} risque (prolifération jacinthes d'eau), colmatage frayères, augmentation coûts traitement eau pour populations aval.
\end{enumerate}
\end{corrige}

\coursec{Les techniques de lutte contre l'érosion}

\paragraph{Introduction : un impératif pour la durabilité des agroécosystèmes camerounais.}
Après avoir analysé les conséquences dramatiques de l'érosion sur la fertilité des sols (perte des horizons superficiels riches en humus et en éléments nutritifs), la qualité de l'eau (turbidité, eutrophisation des barrages de Songloulou, Lagdo, Memve'ele) et les infrastructures (ensablement des routes, affouillements des ponts), une question s'impose : \textbf{comment protéger durablement le sol ?} La réponse ne réside pas dans une unique solution miracle, mais dans l'association raisonnée de techniques adaptées au milieu physique (pente, climat, type de sol, sensibilité à la battance) et aux contraintes socio-économiques des populations (tenure foncière, main-d'œuvre, capital). Le principe directeur unit toutes ces méthodes : \textbf{maintenir ou rétablir un couvert végétal permanent} et \textbf{ralentir la vitesse du ruissellement} pour favoriser l'infiltration et recharger la réserve utile.

\courssub{Principes généraux de la conservation des sols (CES)}

L'érosion hydrique et éolienne trouvent leur force motrice dans l'énergie cinétique de l'eau qui ruisselle ou du vent qui balaie la surface nue. Toute technique de lutte vise donc à agir sur les facteurs de l'équation universelle des pertes en sol (type USLE/RUSLE) : \emph{R} (érosivité pluies/vent), \emph{K} (érodibilité sol), \emph{LS} (longueur et pente du versant), \emph{C} (couvert/cultures), \emph{P} (pratiques de soutien).

\begin{enumerate}
    \item \textbf{Dissiper l'énergie cinétique des agents érosifs :} en interceptant les gouttes de pluie avant qu'elles n'atteignent le sol (couvert végétal dense, paillage, résidus de culture) ou en brisant la vitesse du vent à hauteur du sol (haies brise-vent, rideaux forestiers).
    \item \textbf{Augmenter la rugosité de surface :} pour créer une résistance hydraulique au ruissellement diffus (labour en courbes de niveau, résidus de culture en surface, cordons de pierres, billons).
    \item \textbf{Améliorer la structure et la porosité du sol (réduire K) :} par l'apport régulier de matière organique (fumure, compost, engrais verts, paillage) et la limitation du travail du sol (labour minimal, semis direct sous couvert végétal - SCV) pour favoriser l'agrégation stable, l'activité macrofaunique (vers de terre) et l'infiltration plutôt que le ruissellement.
    \item \textbf{Raccourcir la longueur effective du versant (réduire L) :} par des ouvrages de ralentissement (terrasses, seuils, haies vives, bandes enherbées) qui découpent la pente en segments courts où le ruissellement n'atteint pas sa vitesse critique.
    \item \textbf{Stabiliser mécaniquement les versants :} par le système racinaire pivotant et fasciculé des plantes ligneuses (reboisement, agroforesterie, haies vives) ou par des ouvrages de génie civil (terrasses à mur, seuils en gabions, barrages de ravines).
\end{enumerate}

\begin{definition}[Conservation des sols (CES) / Conservation de l'eau et des sols (CES/DRS)]
Ensemble de pratiques agronomiques (culturales), forestières (agroforestières) et mécaniques (aménagements) visant à prévenir ou réduire la dégradation des sols par l'érosion, à maintenir ou restaurer leur fertilité (physique, chimique, biologique) et à gérer durablement les ressources en eau (infiltration, réserve utile) et en biomasse.
\end{definition}

\begin{attention}[Piège fréquent : confondre « aménagement » et « gestion durable »]
Construire des terrasses ou planter des arbres ne suffit pas. Une terrasse non entretenue (brèche sur le merlon, colmatage du fossé de drainage, absence de couvert sur la plateforme) devient un piège : l'eau s'accumule en amont du merlon, la pression hydrostatique augmente, le merlon cède et l'érosion devient catastrophique (érosion en entonnoir, rupture de terrasse en cascade). \textbf{La durabilité exige l'entretien régulier (curage fossés, réfection merlons) et l'association systématique avec un couvert végétal permanent (herbe, culture, paillage).}
\end{attention}

\courssub{Le reboisement : la racine comme ancre et le feuillage comme parapluie}

\paragraph{Observation.} Sur les flancs du Mont Cameroun (zone volcanique, pluviométrie > 4000 mm/an) ou dans les monts Mandara (zone soudano-sahélienne, sols ferrugineux tropicaux lessivés), les versants dénudés présentent des ravines profondes et des glissements de terrain, tandis que les zones boisées voisines, sous la même pluviométrie ou le même vent, montrent des sols stables, profonds et une hydrologie régulée.

\paragraph{Hypothèse.} La végétation (arbres, arbustes, strate herbacée) protège le sol par une double action : interception de l'énergie cinétique des pluies par le houppier (réduction de la battance) et fixation mécanique du sol par le système racinaire (cohésion, porosité).

\begin{experience}[Mise en évidence du rôle fixateur des racines et de l'infiltration]
\textbf{Matériel :} Deux bacs identiques (1 m $\times$ 0,5 m $\times$ 0,3 m) remplis du même sol meuble (latéritique, horizon B), l'un planté de jeunes arbres fixateurs d'azote (ex. \emph{Acacia mangium}, \emph{Leucaena leucocephala}, \emph{Calliandra calothyrsus}) depuis 6 à 8 mois (racines développées), l'autre laissé nu. Pulvérisateur d'eau calibré (simulation pluie d'intensité 30 à 50 mm/h), récupérateurs d'eau de ruissellement et d'infiltration. Pente réglée à 15\%.
\textbf{Protocole :} Arroser les deux bacs pendant 20 min à intensité constante. Recueillir et mesurer : volume d'eau ruisselée, volume d'eau infiltrée, turbidité (quantité de sédiments en suspension par filtration/séchage/poids).
\textbf{Observations :} Le bac nu produit un ruissellement rapide, abondant (80-90\% de l'eau), trouble (chargé en sédiments : 50-200 g/L). Le bac planté produit peu de ruissellement (10-20\%), clair (< 5 g/L sédiments), avec un fort taux d'infiltration (80-90\%).
\textbf{Interprétation :} Les racines créent une porosité macroporeuse (galeries racinaires, chenilles de vers de terre attirés par les exsudats) qui draine l'eau vers la profondeur (infiltration). Le réseau racinaire forme un « filet » mécanique retenant les agrégats du sol (cohésion). Le feuillage et la litière interceptent l'énergie cinétique des gouttes (zéro battance, pas de croûte de battance).
\textbf{Conclusion :} Le reboisement est la technique \emph{par excellence} pour la restauration des versants fortement dégradés (sols minces, pentes fortes) et la protection des bassins versants d'alimentation en eau potable (captages de Yaoundé - Mfoundi, Douala - Kwa Kwa, Ngaoundéré - Mbé).
\end{experience}

\begin{definition}[Reboisement et restauration forestière]
Action de reconstituer un peuplement forestier sur une zone déboisée ou dégradée, par plantation d'essences adaptées au stationnement (climat, sol, altitude) : \textbf{essences indigènes de préférence} pour la biodiversité et la résilience (\emph{Prunus africana}, \emph{Podocarpus latifolius}, \emph{Terminalia superba}, \emph{Entandrophragma} spp. en zone humide ; \emph{Khaya senegalensis}, \emph{Pterocarpus erinaceus} en zone soudanienne) ; ou \textbf{essences exotiques à croissance rapide} pour la production de bois/énergie et la fixation rapide en zone sèche (\emph{Eucalyptus} spp., \emph{Pinus} spp., \emph{Acacia} spp., \emph{Prosopis juliflora} -- attention au caractère invasif de \emph{Prosopis}). Objectifs : fixer les sols, réguler l'hydrologie (effet éponge), séquestrer le carbone, reconstituer la biodiversité et fournir des produits forestiers (bois, fruits, pharmacopée).
\end{definition}

\begin{savaistu}[La Grande Muraille Verte et l'opération « Sahel Vert » au Cameroun]
Lancée en 2007 par l'Union Africaine, la \textbf{Grande Muraille Verte pour le Sahara et le Sahel} vise à restaurer 100 millions d'hectares de terres dégradées d'ici 2030, en plantant une mosaïque d'arbres, d'arbustes et d'herbes sur 8 000 km de Dakar à Djibouti. Au Cameroun, le programme s'inscrit dans l'opération historique « \textbf{Sahel Vert} » (années 70-80, relancée) dans les régions du Nord et de l'Extrême-Nord : fixation des dunes mobiles par \emph{Prosopis juliflora} et \emph{Balanites aegyptiaca}, création de parcs à karité (\emph{Vitellaria paradoxa}) et de néré (\emph{Parkia biglobosa}), établissement de ceintures vertes autour des villages (Mora, Kousseri, Waza) et des axes routiers. C'est la plus grande initiative de reboisement participatif au monde, alliant lutte anti-érosive, sécurité alimentaire (produits forestiers non ligneux : gomme arabique, fruits, fourrage) et adaptation au changement climatique.
\end{savaistu}

\begin{popfigure}
\begin{tikzpicture}[scale=0.9, every node/.style={font=\small}]
    % Sol nu - érosion éolienne
    \begin{scope}[xshift=0cm]
        \draw[popdark, thick] (0,0) -- (5,0); % Surface du sol
        \draw[poporange, decorate, decoration={snake, amplitude=1mm, segment length=3mm}, -{Stealth[length=2mm]}] (0.5,1.5) -- (4.5,1.5) node[midway, above] {Vent dominant};
        % Particules en saltation
        \foreach \x in {0.5,1.2,2.0,2.8,3.5,4.2} {
            \draw[poporange, -{Stealth[length=1.5mm]}, thick] (\x,1.0) -- (\x+0.3,0.2);
            \draw[poporange, -{Stealth[length=1.5mm]}, thick] (\x+0.2,0.8) -- (\x+0.5,0.1);
        }
        \node[popdark, align=center] at (2.5,-1.0) {Sans haie : érosion\\ éolienne massive (défletion)};
    \end{scope}
    
    % Avec haie brise-vent
    \begin{scope}[xshift=7.5cm]
        \draw[popdark, thick] (0,0) -- (5,0); % Surface du sol
        % Haie (trois rangées décalées)
        \foreach \x in {2.0, 2.3, 2.6} {
            \draw[popgreen!60!popdark, thick] (\x,0) -- (\x,2.5);
            \foreach \y in {0.3,0.8,1.3,1.8,2.3} {
                \draw[popgreen!60!popdark, fill=popgreen!30] (\x-0.15,\y) ellipse (0.15 and 0.08);
                \draw[popgreen!60!popdark, fill=popgreen!30] (\x+0.15,\y) ellipse (0.15 and 0.08);
            }
        }
        % Vent arrivant (amont)
        \draw[poporange, decorate, decoration={snake, amplitude=1.5mm, segment length=4mm}, -{Stealth[length=2mm]}] (0.5,2.8) -- (1.7,2.8) node[midway, above] {Vent fort};
        % Tourbillon au-dessus haie (zone de compression)
        \draw[poporange, dashed, thick] (2.3,2.5) .. controls (2.3,3.2) and (3.0,3.2) .. (3.0,2.5);
        \node[poporange, font=\tiny] at (2.65,3.0) {Compression};
        % Vent ralenti au sol (aval) - zone protégée
        \draw[poporange!50, -{Stealth[length=2mm]}] (3.0,0.3) -- (4.5,0.3) node[midway, above] {Vent ralenti ($< V_{seuil}$)};
        \draw[poporange!50, -{Stealth[length=2mm]}] (3.0,0.8) -- (4.5,0.8);
        % Flèches sédimentation
        \draw[poporange, -{Stealth[length=1.5mm]}] (2.8,0.5) -- (2.8,0.1);
        \node[popdark, align=center] at (2.5,-1.0) {Avec haie brise-vent :\\ protection sur 10-20 H (H = hauteur haie)};
        % Cote H
        \draw[<->, popdark] (2.6,0) -- (2.6,2.5) node[midway, right] {H};
    \end{scope}
\end{tikzpicture}
\legende{Schéma de principe d'une haie brise-vent efficace (3 rangées d'arbres à feuillage persistant, ex. \emph{Casuarina equisetifolia}, \emph{Azadirachta indica} (neem), \emph{Eucalyptus camaldulensis}). La haie oblige le vent à s'élever (compression), créant une zone de dépression au vent (amont) et une zone de calme relatif sous le vent (aval) sur une distance de 10 à 20 fois la hauteur de la haie (H). La vitesse du vent au sol chute sous le seuil d'entrainement des particules de sol (défletion/saltation stoppées).}
\end{popfigure}

\courssub{Les terrasses : dompter la pente par l'aménagement horizontal}

Sur les fortes pentes (supérieures à 15-20\%), le ruissellement acquiert une vitesse érosive telle que les seules pratiques culturales (courbes de niveau, paillage) ne suffisent plus à garantir la sécurité du sol. Il faut \textbf{modifier la topographie} pour réduire la longueur effective du versant (facteur L) et la pente effective (facteur S).

\begin{definition}[Terrasse de culture (banc culturel / terrasse en gradins)]
Aménagement d'une pente forte en une succession de marches horizontales (plateformes de culture) retenues par un merlon (bourrelet de terre compacté, stabilisé par végétation) ou un mur de soutènement (pierres sèches, béton, gabions), séparées par un fossé de collecte/drainage ou un merlon végétalisé. Elle transforme une pente continue en une série de plats horizontaux (pente effective $\approx$ 0-1\% sur la plateforme).
\end{definition}

\begin{popfigure}
\begin{tikzpicture}[scale=0.85, every node/.style={font=\small}]
    % Paramètres
    \def\H{0.9} % Hauteur de marche (risberme)
    \def\L{2.8} % Largeur de plateforme (marche)
    \def\N{5}   % Nombre de marches
    
    % Rocher mère / base
    \draw[popdark, thick] (0,0) -- (\N*\L,0);
    
    % Boucle pour dessiner les terrasses (coupe transversale)
    \foreach \i in {0,...,\N} {
        \pgfmathsetmacro{\prevx}{\i*\L}
        \pgfmathsetmacro{\prevy}{\i*\H}
        \pgfmathsetmacro{\nextx}{(\i+1)*\L}
        \pgfmathsetmacro{\nexty}{(\i+1)*\H}
        
        % Plateforme (sol cultivé, épaisseur ~0.6H)
        \draw[popdark, thick, fill=popgreen!10] (\prevx,\prevy) rectangle (\nextx,\prevy+0.6*\H);
        
        % Culture sur plateforme (représentation schématique)
        \foreach \k in {1,2,3} {
            \pgfmathsetmacro{\cx}{\prevx + 0.25*\L + \k*0.2*\L}
            \draw[popgreen!60!popdark, thick] (\cx,\prevy+0.6*\H) -- (\cx,\prevy+0.9*\H);
        }
        
        % Flèche infiltration (bleue)
        \pgfmathsetmacro{\ix}{\prevx + 0.5*\L}
        \pgfmathsetmacro{\iy}{\prevy + 0.75*\H}
        \draw[popblue, -{Stealth[length=2mm]}, thick] (\ix,\iy) -- (\ix,\prevy+0.3*\H) node[midway, right] {Infiltration};
        
        % Merlon / Mur de soutènement (sauf pour la dernière marche tout en haut)
        \ifnum\i<\N
            % Merlon en terre végétalisé (recommandé)
            \draw[poporange!80!popdark, thick, fill=poporange!20] (\nextx,\prevy) -- (\nextx,\nexty) -- (\nextx+0.15,\nexty) -- (\nextx+0.15,\prevy) -- cycle;
            % Herbe sur merlon (Vetiver)
            \foreach \g in {0.1,0.3,0.5,0.7,0.9} {
                \draw[popgreen!60!popdark, thick] (\nextx+0.075,\prevy+\g*\H) -- (\nextx+0.2,\prevy+\g*\H+0.1);
            }
            \node[popdark, rotate=90, anchor=south, font=\tiny] at (\nextx+0.2,\prevy+0.4*\H) {Merlon végétalisé (Vetiver)};
        \fi
    }
    
    % Ruissellement ralenti sur la dernière plateforme (haut du versant)
    \draw[popblue, decorate, decoration={snake, amplitude=1mm, segment length=4mm}, -{Stealth[length=3mm]}, thick] (0.2,\N*\H+0.4) -- (\N*\L-0.2,\N*\H+0.4) node[midway, above] {Ruissellement ralenti, non érosif};
    
    % Pente initiale (fantôme)
    \draw[popmuted, dashed, thick] (0,0) -- (\N*\L,\N*\H) node[midway, above right, sloped, font=\small] {Pente initiale non aménagée};
    
    % Légende position
    \node[popdark, align=center] at (\N*\L/2, -1.0) {Coupe transversale d'un versant aménagé en terrasses à merlons végétalisés (type « bancs culturels »)};
\end{tikzpicture}
\legende{Principe de la terrasse : la pente forte est découpée en plateformes horizontales (largeur L = 3-5 m, hauteur H = 0,8-1,5 m). L'eau de pluie s'infiltre sur chaque plateforme au lieu de ruisseler violemment. Le merlon (ou mur) retient la terre et l'eau. \textbf{Exemples camerounais :} terrasses traditionnelles des monts Mandara (peuls Mafa/Kapsiki) pour le mil/sorgho ; terrasses caféicoles et maraîchères des hauts plateaux de l'Ouest (Bamiléké, Bamoun) ; terrasses de riziculture en zone de montagne (Nord-Ouest).}
\end{popfigure}

\begin{methode}[Choisir la technique anti-érosive adaptée au terrain (Démarche diagnostique)]
\textbf{Étape 1 : Diagnostiquer la pente (facteur topographique majeur).}
\begin{itemize}
    \item \textbf{Pente faible ($< 5\%$) :} Couvert végétal permanent, paillage, labour minimal / semis direct (SCV), bandes enherbées tous les 20-30 m.
    \item \textbf{Pente modérée ($5\% - 15/20\%$) :} \textbf{Cultures en courbes de niveau} (obligatoire) + haies vives sur courbes de niveau + bandes enherbées (vetiver) + paillage.
    \item \textbf{Pente forte ($> 15/20\%$) :} \textbf{Terrasses} (merlon végétalisé ou mur) + reboisement des risbermes (talus) + cultures en courbes de niveau \emph{sur} les plateformes.
\end{itemize}

\textbf{Étape 2 : Analyser le climat et le type d'érosion dominant.}
\begin{itemize}
    \item \textbf{Zone soudano-sahélienne (Nord, Extrême-Nord) :} Érosion éolienne + hydrique (orages violents). $\rightarrow$ Haies brise-vent, cordons de pierres (diguettes), zaï/tassa, reboisement \emph{Acacia/Prosopis/Balanites}, agroforesterie parc à karité/néré.
    \item \textbf{Zone humide/subhumide (Sud, Littoral, Ouest, Adamaoua) :} Érosion hydrique diffuse et concentrée (ruissellement, ravines). $\rightarrow$ Terrasses, courbes de niveau, agroforesterie dense, cultures de couverture (mucuna, pueraria, crotalaire), haies vives fixatrices d'azote.
\end{itemize}

\textbf{Étape 3 : Évaluer les moyens et la sécurité foncière (facteur humain).}
\begin{itemize}
    \item \textbf{Faibles moyens / tenure précaire (location, prêt) :} Techniques agronomiques réversibles, peu coûteuses : courbes de niveau, paillage, zaï, haies vives (boutures), cordons de pierres (pierres du champ).
    \item \textbf{Moyens importants / tenure sécurisée (titre foncier, héritage stable) :} Ouvrages mécaniques lourds : terrasses à mur (pierres, béton), seuils en gabions, barrages de ravines, drainage souterrain.
\end{itemize}

\textbf{Étape 4 : Associer systématiquement (Approche intégrée CES/DRS).} \emph{Une terrasse sans herbe sur le merlon s'effondre. Une courbe de niveau sans paillage sur sol nu batte. Le reboisement sans protection feu pâturage échoue.}
\end{methode}

\courssub{Cultures en courbes de niveau : suivre le relief pour retenir l'eau}

\paragraph{Observation.} Dans les champs de maïs des hauts plateaux de l'Ouest (Bafoussam, Dschang, Bangangté), on observe deux types de sillons : certains descendent droit dans la pente (érosion visible : ravines, dépôts alluvionnaires en bas de versant), d'autres suivent le contour de la colline (sol stable, humidité persistante, rendements supérieurs).

\begin{popfigure}
\begin{tikzpicture}[scale=0.9, every node/.style={font=\small}]
    % Vue de dessus : Culture DANS LA PENTE (Érosion)
    \begin{scope}[xshift=0cm, yshift=0cm]
        \draw[popdark, thick] (0,0) rectangle (6,4);
        \node[popdark, align=center] at (3,4.4) {\textbf{Culture dans le sens de la plus grande pente}};
        \node[popred, align=center] at (3,4.15) {(Érosion hydrique maximale)};
        % Flèches pente (ruissellement concentré dans sillons)
        \foreach \y in {0.5,1.5,2.5,3.5} {
            \draw[popred, -{Stealth[length=2mm]}, thick] (3,\y) -- (3,\y-0.4) node[midway, right, font=\tiny] {Ruissellement concentré};
        }
        % Sillons parallèles à la pente (labour)
        \foreach \x in {0.8,1.6,2.4,3.2,4.0,4.8,5.6} {
            \draw[poporange!80!popdark, thick] (\x,0) -- (\x,4);
        }
        % Ravines formées dans les sillons
        \draw[popred, thick, decorate, decoration={zigzag, amplitude=1.5pt}] (1.5,3.8) -- (1.2,0.2);
        \draw[popred, thick, decorate, decoration={zigzag, amplitude=1.5pt}] (4.5,3.8) -- (4.8,0.2);
        \node[popred, font=\tiny] at (1.5,2) {Ravine};
        \node[popred, font=\tiny] at (4.5,2) {Ravine};
    \end{scope}
    
    % Vue de dessus : Culture EN COURBES DE NIVEAU (Protection)
    \begin{scope}[xshift=8cm, yshift=0cm]
        \draw[popdark, thick] (0,0) rectangle (6,4);
        \node[popdark, align=center] at (3,4.4) {\textbf{Culture en courbes de niveau}};
        \node[popgreen!60!popdark, align=center] at (3,4.15) {(Protection optimale du sol)};
        % Courbes de niveau (lignes de niveau)
        \foreach \y/\c in {0.5/1, 1.2/2, 2.0/3, 2.9/4, 3.6/5} {
            \draw[popgreen!60!popdark, thick] (0,\y) .. controls (2,\y+0.1) and (4,\y-0.1) .. (6,\y);
            \node[popgreen!60!popdark, font=\small] at (6.3,\y) {CN \c};
        }
        % Flèches infiltration (bleues) entre les courbes
        \foreach \y in {0.85,1.6,2.45,3.25} {
            \draw[popblue, -{Stealth[length=2mm]}, thick] (3,\y+0.1) -- (3,\y-0.1) node[midway, right, font=\tiny] {Infiltration};
        }
        % Haie vive sur une courbe de niveau (ex CN 2)
        \draw[popgreen!40!popdark, thick, fill=popgreen!10] (0,1.2) .. controls (2,1.3) and (4,1.1) .. (6,1.2) .. controls (4,1.35) and (2,1.05) .. (0,1.35);
        \node[popgreen!40!popdark, font=\tiny, align=center] at (3,1.27) {Haie vive fixatrice d'azote\\(\emph{Leucaena}, \emph{Calliandra})};
    \end{scope}
\end{tikzpicture}
\legende{Vue de dessus comparant deux itinéraires culturaux sur une même pente. \textbf{Gauche :} Sillons parallèles à la plus grande pente $\rightarrow$ l'eau s'engouffre dans les sillons, accélère par gravité, creuse des ravines (érosion linéaire concentrée). \textbf{Droite :} Sillons perpendiculaires à la pente (courbes de niveau) $\rightarrow$ chaque rangée de culture agit comme un micro-barrage, l'eau stagne temporairement, s'infiltre, l'érosion est stoppée. L'association avec une haie vive sur la courbe de niveau renforce l'effet barrière physique, apporte de la biomasse (paillage, fourrage) et fixe l'azote.}
\end{popfigure}

\begin{definition}[Culture en courbes de niveau]
Pratique culturale consistant à effectuer \textbf{tous} les travaux du sol (labour, semis, binage, buttage, récolte) le long de lignes de niveau (courbes de niveau), c'est-à-dire perpendiculairement à la plus grande pente. Chaque rangée de culture forme un petit barrage transverse qui ralentit le ruissellement diffus, favorise l'infiltration et piège les sédiments. Le traçage se fait à l'aide d'un \textbf{niveau à eau} (tuyau transparent de 10-20 m rempli d'eau, principe des vases communicants) ou d'un \textbf{niveau en A} (artisanal, bois et fil à plomb, précision $\pm$ 2-3 cm).
\end{definition}

\begin{exemplebox}[Les courbes de niveau sur les Hauts Plateaux de l'Ouest (Cameroun) : un savoir-faire paysan scientifiquement validé]
Dans la région de l'Ouest (Bamiléké, Bamoun), l'agriculture intensive sur pentes volcaniques (andosols profonds, très fertiles mais très sensibles à l'érosion par dispersion d'argile) a généralisé la culture en courbes de niveau pour le maïs, le haricot, la pomme de terre, le chou et le café arabica. Les agriculteurs tracent les courbes à l'aide du niveau à eau ou du niveau en A. Les billons suivent le contour. Souvent, une haie vive de \emph{Leucaena leucocephala} ou \emph{Calliandra calothyrsus} est plantée sur la courbe de niveau (espacement 5-10 m entre haies selon la pente) : elle fixe l'azote atmosphérique (engrais vert), fournit du fourrage protéiné pour l'élevage en stabulation (vaches laitières, porcs) et du bois de feu, tout en consolidant la \textbf{terrasse vivante} (ou \textit{fanya juu} progressif) qui se forme par accumulation naturelle de sédiments en amont de la haie. C'est un modèle d'agroforesterie paysanne durable.
\end{exemplebox}

\courssub{Techniques complémentaires et agroécologiques (la « boîte à outils » du paysan)}

Au-delà des structures lourdes (terrasses) et du reboisement, un panel de techniques « douces », peu coûteuses et applicables partout, constitue la base de l'Agriculture de Conservation (AC) :

\begin{itemize}
    \item \textbf{Bandes enherbées (ou bandes filtrantes / tampons) :} Bandes d'herbacées pérennes à racines profondes et fasciculées (vetiver \emph{Chrysopogon zizanioides}, \emph{Pennisetum purpureum} herbe éléphant, \emph{Brachiaria} spp., \emph{Andropogon gayanus} gamba) laissées en place ou plantées tous les 10-20 m le long des courbes de niveau. Elles filtrent l'eau ruisselante, retiennent les sédiments fins (argile, limon, matière organique, pesticides) et freinent la vitesse du ruissellement. Le \textbf{vetiver} est la référence : racines pivotantes très profondes (3-4 m), touffe dense non invasive (stérile), résistant à la sécheresse, au feu, aux inondations et aux pics de toxicité (métaux lourds).
    \item \textbf{Haies vives (haies agroforestières) :} Rangées d'arbustes/lianes (ex. \emph{Gliricidia sepium}, \emph{Flemingia macrophylla}, \emph{Acacia angustissima}, \emph{Leucaena}, \emph{Calliandra}, \emph{Senna siamea}) plantées serrées (0,25-0,50 m sur le rang) sur les courbes de niveau. Rôle multiple : barrière physique anti-érosive, apport de biomasse par taille (paillage, engrais vert), fixation biologique d'azote (légumineuses), fourrage (protéines), bois de feu, piquets. Elles transforment la pente en succession de micro-terrasses (formation de terrasses vivantes / \textit{fanya juu}).
    \item \textbf{Paillage (mulch) :} Couvrir le sol nu entre les rangs de culture avec des résidus organiques (paille de maïs/riz/sorgho, feuilles de bananier, tontes d'herbe, résidus de mucuna, sciure). Effets majeurs : absorption totale de l'énergie cinétique des gouttes (zéro battance, pas de croûte de surface), maintien de l'humidité du sol (réduction évaporation), régulation thermique, nourrit la vie du sol (vers de terre $\rightarrow$ macroporosité, bioturbation). \emph{Indispensable en Agriculture de Conservation (Semis Direct sous Couvert Végétal - SCV).}
    \item \textbf{Labour minimal / Semis direct (SCV) :} Ne pas retourner le sol (pas de charrue, pas de herse). Ouvrir seulement une fente de semis (semoir direct à disques ou manuel type « jab planter » / « mata-biche »). Préserve la structure agrégée, la macrofaune, le réseau de pores continus et le couvert de résidus en surface. Réduit l'érosion de 80 à 95\% par rapport au labour conventionnel à nu.
    \item \textbf{Cordons de pierres (diguettes / seuils de pierres) :} En zone sahélienne (Nord Cameroun : Diamaré, Mayo-Danay, Bénoué), alignements de pierres de 20-40 cm de haut, posées sur les courbes de niveau (espacement 10-20 m selon pente). Freinent le ruissellement diffus, favorisent la sédimentation en amont (création de micro-terrasses naturelles), augmentent l'infiltration (stockage eau pour la culture). Peu coûteux (main-d'œuvre familiale, pierres du champ/glacis).
    \item \textbf{Zaï / Tassa (trous de plantation) :} Trous de 20-30 cm de diamètre, 15-20 cm de profondeur, creusés en quinconce (espacement 0,8 $\times$ 0,8 m) sur sol dégradé, croûté, imperméable (glacis, zipellé). Amendés avec fumure/organique (compost, fumier). Concentrent l'eau de ruissellement et la fertilité au pied de la plante (mil, sorgho, maïs, niébé). Technique paysanne ancestrale (Yatenga, Burkina Faso) adoptée avec succès au Nord-Cameroun (projet PADL, ONG).
\end{itemize}

\begin{exoresolu}[Choix d'une stratégie anti-érosive intégrée pour une exploitation à Dschang (Hauts Plateaux de l'Ouest)]
\textbf{Énoncé.} Un jeune agriculteur s'installe sur une pente de 18\% à Dschang (alt. 1400 m, pluviométrie 2000 mm/an, sols volcaniques andosols profonds mais très sensibles à l'érosion par dispersion). Il veut cultiver maïs, haricot, pomme de terre et installer un petit élevage porcin (10 truies). Proposez un plan de conservation intégré (CES/DRS) argumenté.
\textbf{Solution détaillée.}
\begin{enumerate}
    \item \textbf{Diagnostic physique :} Pente forte ($>15\%$), pluviométrie élevée (érosivité R forte), sol friable (andosol, érodibilité K forte, dispersion argile), risque érosif \textbf{très fort}. Labour conventionnel à nu = perte de sol garantie ($> 50$ t/ha/an).
    \item \textbf{Aménagement structurel (pente forte) :} Construction de \textbf{terrasses à merlon végétalisé (type banc culturel / fanya chini)}. Creuser le fossé de drainage en amont (courbe de niveau), rejeter la terre en aval pour former le merlon (largeur base 1 m, hauteur 0,8-1 m). Largeur plateforme utile 3-4 m. \textbf{Stabilisation immédiate du merlon} par plantation en double rang de \textbf{Vetiver} (racines 3-4 m, ancre le merlon) ou \emph{Pennisetum purpureum} (herbe éléphant, fourrage pour porcs).
    \item \textbf{Pratiques culturales sur plateformes :} \textbf{Cultures en courbes de niveau} (suivant le merlon). Association culturale : Maïs + Haricot grimpant (couvert rapide, azote). Rotation culturale : Pomme de terre (récolte) $\rightarrow$ \textbf{Mucuna pruriens} (engrais vert rampant, couvre-sol total 3-4 mois, étouffe adventices, fixe azote, biomasse énorme) $\rightarrow$ Maïs/Haricot.
    \item \textbf{Agroforesterie / Haies vives sur merlons :} Plantation de \emph{Calliandra calothyrsus} ou \emph{Leucaena leucocephala} en double rang sur le merlon (espacement 0,5 m). Taille régulière (3-4 fois/an) $\rightarrow$ fourrage protéiné (20-25\% MAT) pour porcs + paillage au pied des cultures (restitution NPK) + bois de feu.
    \item \textbf{Gestion de l'eau hors champ :} Fossé de dérivation (courbe de niveau) en haut du champ pour capter le ruissellement venant de la piste/route/voisin nu et l'évacuer en sécurité vers un exutoire végétalisé (fossé enherbé, zone humide).
    \item \textbf{Règle « Sol nu zéro » :} Paillage systématique (résidus maïs, mucuna, herbe éléphant). Semis direct si possible (semoir manuel type « jab planter » pour maïs/haricot dans le mulch).
    \item \textbf{Intégration Élevage :} Lisier/fumier de porcherie $\rightarrow$ compostage $\rightarrow$ amendement trous de plantation (pomme de terre) et zaï si zones dégradées. Boucle nutriments fermée.
\end{enumerate}
\textbf{Justification globale :} Cette combinaison (Terrasse + Courbes de niveau + Haie vive/Vetiver + Paillage + SCV + Agroforesterie fourragère + Gestion eau hors champ) traite la cause physique (vitesse eau), protège la ressource capitale (sol/eau), produit de la biomasse (fourrage, fertilité, énergie) et s'intègre à l'activité génératrice de revenus (élevage). C'est le modèle \textbf{« Agriculture de Conservation sur Terrasse »} promu par l'IRAD, le MINADER et les ONG (ex. GIC Agro-Pastoraux, APRODEP) dans l'Ouest.
\end{exoresolu}

\courssub{Rôle de l'État, des collectivités et de l'éducation : une responsabilité partagée}

La lutte contre l'érosion dépasse le champ individuel. C'est un \textbf{bien public} : qualité de l'eau des barrages hydroélectriques (Songloulou, Lagdo, Memve'ele, Nachtigal), navigabilité du Wouri (port de Douala), prévention des glissements de terrain meurtriers (Bafoussam, Limbe, Kribi), préservation de la biodiversité (parcs nationaux).

\begin{itemize}
    \item \textbf{Cadre légal et réglementaire :} Loi n°96/12 du 5 août 1996 relative à la gestion de l'environnement (Évaluation d'Impact Environnemental et Social - EIES obligatoire pour grands aménagements). Code forestier (Loi n°94/01) : reboisement obligatoire post-exploitation forestière, protection des forêts de protection (bassins versants, pentes $> 60\%$). Loi sur l'orientation agricole.
    \item \textbf{Programmes nationaux et projets :} Opération « \textbf{Sahel Vert} » (Nord/Extrême-Nord, fixation dunes, parcs agroforestiers). Programme National de Reboisement (objectif 1 million d'arbres/an, essence locales). Restauration des berges du Lac Tchad et des bassins versants (Projet PADL - Projet d'Appui au Développement Local, Projet PIRVA - Projet d'Intégration et de Résilience des Vallées). Initiative « \textbf{Paysages Forestiers Restaurés} » (Bonn Challenge, engagement Cameroun 12 millions ha).
    \item \textbf{Appui technique et vulgarisation :} IRAD (recherche variétés améliorées, techniques CES adaptées), MINADER (vulgarisation, conseillers agricoles, centres de formation), MINEPDED (suivi environnemental, police de l'eau/forêt), ONG internationales (WWF, CIFOR, ICRAF pour l'agroforesterie, GIZ, FAO).
    \item \textbf{Éducation environnementale et changement de comportement :} Intégration dans les programmes SVT (cette leçon !), clubs « Environnement » et « Jeunes Reboiseurs » dans les lycées et collèges, journées nationales de l'arbre (1er août), opérations de reboisement scolaire et citoyen. \textbf{Objectif pédagogique :} changer les mentalités : le sol n'est pas un support inerte inépuisable mais un \textbf{patrimoine vivant non renouvelable à l'échelle humaine} (formation 1 cm = 100-1000 ans).
\end{itemize}

\begin{aretenir}[Synthèse opérationnelle : Adapter la technique au milieu (Tableau de décision)]
\begin{center}
\begin{tabularx}{\linewidth}{|l|X|X|X|}\hline
\textbf{Milieu / Contrainte principale} & \textbf{Technique structurelle principale} & \textbf{Techniques associées indispensables (Agronomiques/Végétales)} & \textbf{Exemple camerounais type} \\ \hline
Forte pente ($>20\%$), sol profond, zone humide (Ouest, SW, NO) & Terrasses (merlon végétalisé ou mur pierre/béton) & Vetiver/Pennisetum sur merlon, Courbes de niveau sur plateforme, Haies vives (Calliandra/Leucaena), Paillage, SCV & Monts Mandara (mil), Hauts Plateaux Ouest (café, pomme de terre, maïs), Monts Bakossi \\ \hline
Pente modérée ($5-15\%$), zone humide/subhumide (Adamaoua, Est, Sud, Littoral) & Courbes de niveau (traçage niveau à eau/niveau en A) & Haies vives fixatrices d'azote, Bandes enherbées (Vetiver), Paillage, SCV, Agroforesterie (cacao/arbres d'ombrage) & Région de l'Ouest (maïs/haricot), Adamaoua (patate douce/maïs), Sud (cacao/bananier) \\ \hline
Pente faible ($<5\%$), zone soudano-sahélienne (sec, venteux, sols sableux/latéritiques) & Haies brise-vent (rideaux), Cordons de pierres (diguettes), Zaï/Tassa & Reboisement \emph{Acacia/Prosopis/Balanites}, Paillage, Agroforesterie parc à karité/néré, Rotation céréales/légumineuses & Extrême-Nord (Diamaré, Mayo-Danay, Logone), Nord (Bénoué, Faro) \\ \hline
Versants très dégradés, sols minces, protection bassin versant / captage eau & Reboisement (essences indigènes prioritaires) & Enrichissement (sous-bois), Protection anti-feu (pare-feu), Coupes contrôlées, Gestion participative (forêts communautaires) & Bassin du Sanaga (barrages), Mont Cameroun (parc national), Monts Bakossi, Bassin du Lac Tchad \\ \hline
\end{tabularx}
\end{center}
\textbf{Règle d'or universelle :} \emph{Aucune technique, aussi sophistiquée soit-elle, ne fonctionne durablement sans couvert végétal permanent (vivant ou mort) et sans entretien régulier. La durabilité naît de l'association \textbf{Sol + Eau + Végétal + Humain}.}
\end{aretenir}

\begin{exercice}[Analyse d'un aménagement anti-érosif traditionnel au Sahel]
On observe sur une photo aérienne d'un versant des monts Mandara (Extrême-Nord, zone soudano-sahélienne, pluviométrie 700-900 mm/an) un réseau régulier de lignes sombres parallèles aux courbes de niveau, séparées de 10 à 15 m. Au sol, ces lignes correspondent à des alignements de pierres sèches de 30-40 cm de haut, derrière lesquels le sol est plus haut de 20 cm qu'en aval (formation de micro-terrasses). Les champs sont cultivés en mil (\emph{Pennisetum glaucum}) et sorgho (\emph{Sorghum bicolor}).
\begin{enumerate}
    \item Identifiez la technique principale mise en œuvre. Justifiez votre réponse par deux éléments précis de l'observation (photo + terrain).
    \item Expliquez le mécanisme physique par lequel cette technique réduit l'érosion hydrique (ruissellement diffus) et améliore le bilan hydrique de la culture.
    \item Pourquoi cette technique est-elle particulièrement adaptée au contexte soudano-sahélien (climat, ressources disponibles, système de production) ?
    \item Proposez deux améliorations agroécologiques concrètes pour augmenter l'efficacité anti-érosive et la productivité agricole de ce système.
\end{enumerate}
\end{exercice}

\begin{corrige}[Exercice n°1 : Cordons de pierres (Diguettes) dans les Monts Mandara]
\begin{enumerate}
    \item \textbf{Technique principale :} \textbf{Cordons de pierres} (appelés localement \textbf{diguettes} ou seuils de pierres).
    \textbf{Justifications :}
    \begin{enumerate}
        \item Alignements de pierres sèches disposés \textbf{suivant les courbes de niveau} (perpendiculaires à la plus grande pente), espacés régulièrement (10-15 m).
        \item \textbf{Accumulation de sédiments en amont} de chaque cordon : le sol est rehaussé de 20 cm, prouvant le piégeage efficace des particules charriées par le ruissellement diffus (formation de micro-terrasses progressives / \textit{fanya chini}).
    \end{enumerate}
    \item \textbf{Mécanisme physique :}
    \begin{itemize}
        \item \textbf{Rugosité et barrage :} Les pierres augmentent fortement la rugosité de surface (facteur $n$ de Manning) et créent une succession de micro-barrages transversaux au flux.
        \item \textbf{Dissipation d'énergie :} L'eau ruisselant en nappe (ruissellement diffus) heurte l'obstacle, perd brutalement de l'énergie cinétique (vitesse $\downarrow$, nombre de Froude $\downarrow$), passe en régime de fluage.
        \item \textbf{Sédimentation et infiltration :} La capacité de transport du flux chute $\rightarrow$ dépôt des sédiments (sable, limon, matière organique) en amont du cordon (colmatage progressif). L'eau stagne temporairement $\rightarrow$ temps d'infiltration accru $\rightarrow$ recharge de la réserve utile du sol (stockage « eau verte » pour le mil/sorgho).
    \end{itemize}
    \item \textbf{Adaptation au contexte soudano-sahélien :}
    \begin{itemize}
        \item \textbf{Matériaux locaux disponibles :} Pierres abondantes sur les glacis et piémonts des monts Mandara (grès, quartzites). Main-d'œuvre familiale disponible (organisation collective \emph{tontine} / \emph{gawal}). Faible coût financier (pas d'intrants achetés).
        \item \textbf{Adapté au régime pluviométrique :} Pluies intenses, courtes, orageuses (fortes intensités horaires $\rightarrow$ ruissellement violent). Les cordons captent ce ruissellement « gratuit » pour l'infiltrer (supplément hydrique crucial pour le cycle cultural court du mil/sorgho).
        \item \textbf{Résistance au feu :} Pas de partie végétale aérienne sèche inflammable en saison sèche (contrairement aux haies vives/herbes), résistant aux feux de brousse fréquents.
        \item \textbf{Tenure foncière :} Adapté aux terres collectives/gérées par le chef de terre, ne nécessite pas de titre foncier individuel.
    \end{itemize}
    \item \textbf{Améliorations agroécologiques (Intensification écologique) :}
    \begin{enumerate}
        \item \textbf{Végétalisation du cordon (Haie vive sur pierres) :} Planter \emph{Andropogon gayanus} (gamba) ou \emph{Vetiver} \emph{sur} le cordon (entre les pierres) et/ou \emph{Acacia senegal} (gommier) / \emph{Ziziphus mauritiana} (jujubier) \emph{en aval} immédiat. \textbf{Avantages :} Les racines ancrent définitivement les pierres (stabilité mécanique), la partie aérienne filtre mieux les fins sédiments (argile, MO), fournit fourrage (gamba), gomme arabique (acacia), fruits/jujubes (jujubier), bois. Transforme l'ouvrage minéral en ouvrage vivant auto-entretenu.
        \item \textbf{Zaï associé en amont du cordon :} Creuser des trous de plantation (Zaï : 20-30 cm $\varnothing$, 15-20 cm prof.) en quinconce dans la zone de dépôt sédimentaire en amont du cordon, y déposer du compost/fumier de l'élevage (disponible localement). \textbf{Avantages :} Concentre l'eau (ruissellement capté par cordon) + la fertilité (compost) au pied de la plante $\rightarrow$ levée assurée, enracinement profond, rendements doublés/triplés (500-800 kg/ha $\rightarrow$ 1500-2500 kg/ha mil/sorgho).
        \item \textbf{(Bonus) Paillage inter-rangs :} Couvrir la surface entre les cordons avec les résidus de mil/sorgho (tiges, feuilles) après récolte. Zéro battance, économie d'eau, nutrition sol.
    \end{enumerate}
\end{enumerate}
\end{corrige}

\begin{attention}[Erreurs classiques à ne pas commettre au Probatoire (SVT 1ère C/TI)]
\begin{itemize}
    \item \textbf{Confondre courbe de niveau et ligne de plus grande pente.} La courbe de niveau est \textbf{perpendiculaire} à la pente. Labourer/semer \emph{dans le sens de la pente} (plus grande pente) crée des chenaux préférentiels qui \textbf{aggravent} l'érosion (facteur P $> 1$). C'est l'erreur n°1.
    \item \textbf{Oublier l'entretien dans la définition d'une technique.} « Faire des terrasses » $\neq$ « Lutter contre l'érosion ». Une terrasse sans merlon végétalisé, sans fossé de drainage propre, sans culture en courbes de niveau \emph{sur} la plateforme, est un danger (rupture en cascade). \textbf{Entretien = partie intégrante de la technique.}
    \item \textbf{Croire que « Reboisement » = Monoculture d'Eucalyptus ou de Pin.} L'eucalyptus (ex. \emph{E. camaldulensis}, \emph{E. globulus}) a des allélopathies, assèche fortement la nappe (transpiration énorme), n'apporte peu de litière décomposable (C/N élevé), acidifie le sol, brûle très bien (risque incendie). \textbf{Privilégier les essences indigènes ou agroforestières multipropos (fixatrices d'azote, fourragères, mellifères).}
    \item \textbf{Négliger l'amont (approche bassin versant).} Protéger son champ en aval sans gérer le ruissellement venant de la piste, du voisin nu, de la route, du pâturage en amont = échec assuré. Il faut une approche \textbf{bassin versant / paysage} (concertation villageoise, plans d'aménagement).
    \item \textbf{Confondre « Zaï » et « Trou de plantation classique ».} Le Zaï est creusé \textbf{en saison sèche} sur sol dur/croûté, amendé \textbf{avant} la pluie, pour capter le ruissellement. Le trou de plantation classique est fait au moment du semis/planting en saison humide sur sol meuble.
\end{itemize}
\end{attention}

\coursec{Interprétation de documents et étude de cas}

Après avoir étudié les types d'érosion, leurs facteurs favorisants, leurs conséquences et les principales techniques de lutte (reboisement, terrasses, cultures en courbes de niveau, brise-vent, fascines), il est indispensable de savoir mobiliser ces connaissances face à des documents réels : photographies de terrain, schémas de versant, cartes d'occupation du sol ou données chiffrées. Cette section vous apprend une démarche rigoureuse, étape par étape, pour analyser tout document d'érosion, formuler un diagnostic précis et proposer des solutions adaptées au contexte camerounais, conformément aux attendus du Probatoire.

\courssub{Méthode d'analyse d'une photographie montrant un phénomène érosif}

Une photographie fige un instant du paysage. Pour en extraire l'information scientifique, il ne suffit pas de « regarder » : il faut observer systématiquement, décrire avec un vocabulaire technique précis, expliquer les processus en jeu et enfin proposer des mesures correctives. La démarche suivante, éprouvée aux épreuves du Probatoire, structure votre raisonnement.

\begin{methode}[Méthode en 4 étapes pour interpréter un document sur l'érosion]
\textbf{1. Observer} : Balayer l'image du premier plan à l'arrière-plan. Repérer les couleurs (sol nu, végétation, eau), les formes (ravins, dunes, dépôts), les infrastructures (routes, habitations). Noter l'échelle si elle est indiquée (ombre d'un arbre, taille d'une maison, échelle graphique).\\
\textbf{2. Décrire} : Rédiger une description objective en utilisant les termes techniques : \cle{ravinement}, \cle{gullies} (synonyme de ravins), \cle{dépôts alluviaux}, \cle{colluvions}, \cle{surface dénudée}, \cle{pente}, \cle{ligne de crête}, \cle{ruissellement diffus}, \cle{ruissellement concentré}. Préciser la localisation relative (amont/aval, pied de versant, haut de versant).\\
\textbf{3. Expliquer} : Relier les formes observées aux processus érosifs. Identifier le \cle{type d'érosion} (hydrique : ruissellement diffus $\rightarrow$ érosion en nappe ; ruissellement concentré $\rightarrow$ érosion en sillons \emph{rills} puis ravinement \emph{gullies} ; éolienne : \cle{déflation}, \cle{saltation}, formation de \cle{dunes mobiles}). Citer les \cle{facteurs favorisants} visibles : déforestation, pente forte, labour dans le sens de la pente, absence de haies, surpâturage, climat (saison des pluies intenses, harmattan).\\
\textbf{4. Proposer} : Choisir une ou plusieurs techniques de lutte adaptées aux caractéristiques du milieu (pente, nature du sol, climat, pression foncière, main-d'œuvre disponible). Justifier chaque choix en liant la technique aux contraintes du terrain : « terrasses en pierres sèches car pente > 15\% et sol argileux profond » ; « reboisement en essences locales à croissance rapide et racines pivotantes pour stabiliser les berges des ravins » ; « brise-vent perpendiculaires au vent dominant (harmattan NE-SO) pour réduire la vitesse du vent au sol sous le seuil de transport éolien ».
\end{methode}

\begin{exemplebox}[Exemple de réponse rédigée type Probatoire à partir d'une photo de ravin]
\textbf{Énoncé :} La photographie montre un versant des collines de Yaoundé après une forte pluie. On observe de larges ravins en V, des dépôts de sable et d'argile au pied du versant, une zone déboisée en amont et quelques maisons en aval.\\
\textbf{Réponse :} 
\begin{enumerate}
\item \textbf{Type d'érosion :} Érosion hydrique par ruissellement concentré (formation de ravins / \emph{gullies}).
\item \textbf{Causes visibles :} Déforestation de la partie supérieure du versant (sol nu), pente marquée ($\approx 20\,\%$), absence de travaux de conservation (ni terrasses, ni courbes de niveau), forte intensité pluviométrique (saison des pluies, gouttes de pluie érodantes).
\item \textbf{Conséquences :} Perte de la couche arable (horizon A), enfoncement des ravins qui menacent les habitations en aval par érosion régressive, colmatage des cours d'eau par les sédiments (envasement), dégradation de la qualité de l'eau.
\item \textbf{Mesures proposées :} 
\begin{itemize}
\item Reboisement immédiat de la zone déboisée avec des essences à racines profondes et pivotantes (\emph{Acacia mangium}, \emph{Acacia auriculiformis}, \emph{Eucalyptus} spp.) pour réduire le ruissellement et fixer le sol.
\item Construction de terrasses en pierres sèches ou en terre battue (selon disponibilité matériaux) sur les parties cultivées pour réduire la longueur effective de pente.
\item Aménagement de fascines (branchages entrelacés) ou de seuils en pierres sèches/gabions en travers des ravins existants pour freiner l'érosion régressive (remontée vers l'amont).
\item Sensibilisation des riverains à l'agroforesterie (cultures associées à des arbres fixateurs d'azote comme \emph{Gliricidia sepium}, \emph{Leucaena leucocephala}) pour concilier production agricole et protection durable.
\end{itemize}
\end{enumerate}
\end{exemplebox}

\courssub{Méthode d'analyse d'un schéma de versant}

Un schéma de versant est une représentation simplifiée en coupe transversale (souvent avec exagération verticale). Il met en évidence la géométrie de la pente, la répartition de la végétation, le sens du ruissellement (flèches bleues) et les formes d'érosion (ravins, cônes de déjection). L'analyse suit la même logique en 4 étapes, avec une attention particulière au \cle{profil longitudinal} et aux \cle{unités morphologiques} (crête, pente, pied de versant, glacis).

\begin{methode}[Analyse d'un schéma de versant]
\textbf{1. Observer} : Identifier l'orientation du schéma (amont/aval, nord/sud si indiqué). Repérer la ligne de crête, la pente (angle ou pourcentage si échelle), les zones boisées / nues, les cours d'eau, les ouvrages existants.\\
\textbf{2. Décrire} : Noter la présence de ravins (incision en V), de dépôts (cônes de déjection au pied de versant), d'éventuels ouvrages (terrasses, seuils, bandes enherbées). Mesurer si possible la pente grâce à l'échelle graphique.\\
\textbf{3. Expliquer} : Relier la pente et le couvert végétal à l'énergie du ruissellement. Une pente $> 15\,\%$ sans végétation génère des vitesses érosives élevées (énergie cinétique $\propto v^2$). Les ravins indiquent une érosion régressive (remontée vers l'amont par recul des têtes de ravins).\\
\textbf{4. Proposer} : Adapter les techniques : courbes de niveau pour pentes $5-15\,\%$, terrasses pour $> 15\,\%$ ou sols peu profonds, reboisement sur les versants très dégradés (pente $> 30\,\%$ ou sols squelettiques), bandes enherbées en pied de versant pour piéger les sédiments avant le cours d'eau.
\end{methode}

\begin{popfigure}
\begin{tikzpicture}[scale=0.9, every node/.style={font=\small}]
% Sol et substrat
\draw[thick, popdark] (0,0) -- (12,0) -- (12,-6) -- (0,-6) -- cycle;
% Surface du sol (profil topographique)
\draw[very thick, popdark] (0,3) -- (12,0);
% Ravins (incisions)
\draw[poporange, thick] (2,2.5) .. controls (3,1.5) and (4,1) .. (5,0);
\draw[poporange, thick] (6,2) .. controls (7,1) and (8,0.5) .. (9,0);
% Zone déboisée (sol nu)
\fill[poporange!30] (3,2) rectangle (7,2.8);
% Végétation (arbres stylisés) - Zone boisée amont
\foreach \x in {0.5,1,1.5} {
  \fill[popgreen!70] (\x,2.9) circle (0.12);
  \draw[popgreen!70] (\x,2.9) -- (\x,3.4);
}
% Végétation - Zone aval résiduelle
\foreach \x in {8,8.5,9,10,10.5} {
  \fill[popgreen!70] (\x,2.9) circle (0.12);
  \draw[popgreen!70] (\x,2.9) -- (\x,3.4);
}
% Habitations en aval
\draw[popblue, thick] (11.2,-0.5) rectangle (11.8,0.2);
\draw[popblue, thick] (11.2,0.2) -- (11.5,0.6) -- (11.8,0.2);
% Flèches ruissellement
\draw[->, popblue, thick] (4,2.3) -- (4,1);
\draw[->, popblue, thick] (7,1.8) -- (7,0.5);
\draw[->, popblue, thick] (10,1) -- (10,0.2);
% Légendes directes
\node[above, align=center] at (1,3.5) {Zone boisée\\ (crête / haut de versant)};
\node[above, align=center] at (5,3.0) {Zone déboisée\\ (pente forte $\approx 22\,\%$)};
\node[below, align=center] at (5,-0.8) {Ravins (gullies)\\ érosion régressive};
\node[right, align=left] at (12,0) {Pied de versant :\\ cône de déjection,\\ dépôts alluviaux,\\ habitations menacées};
\node[above, align=center] at (10,1.3) {Ruissellement\\ concentré};
% Échelle (indicative)
\draw[|-|, thick] (0.5,-5) -- (2.5,-5) node[midway, below] {200 m};
\end{tikzpicture}
\legende{Schéma-type d'un versant érodé (coupe transversale) : pente, ravins (gullies), zones déboisées, habitations en aval. Noter l'exagération verticale typique des coupes.}
\end{popfigure}

\courssub{Étude de cas 1 : Versant déboisé des collines de Yaoundé — Diagnostic et propositions}

\begin{experience}[Diagnostic du versant de Nkolbisson (Yaoundé, Centre)]
\textbf{Contexte :} Les collines autour de Yaoundé (alt. 700–900 m, climat équatorial de transition, pluviométrie $\approx 1\,600$ mm/an) subissent une pression urbaine et agricole forte. Le document (photo aérienne ou schéma ci-dessus) montre un versant déboisé sur 15 ha, pente moyenne $22\,\%$, sol ferralitique (latéritique) argilo-sableux, faiblement structuré. En saison des pluies (mars–novembre), des ravins de 2–3 m de profondeur se forment et menacent des habitations précaires en contrebas.\\
\textbf{Observations :} 
\begin{itemize}
\item Sol nu sur 70\% de la surface (cultures de maïs en lignes de pente, jachères courtes $< 2$ ans).
\item Ravins en V multiples, convergents vers un collecteur principal (érosion régressive active).
\item Dépôts sédimentaires (colluvions) au pied du versant, colmatage d'un petit ruisseau affluent de la Mfoundi.
\item Absence de tout aménagement anti-érosif (ni courbes de niveau, ni haies, ni terrasses).
\end{itemize}
\textbf{Interprétation :} L'érosion hydrique est de type \cle{ravinement} (gullies) par ruissellement concentré. Les facteurs aggravants sont : la déforestation (perte de l'interception pluviométrique et du rôle fixateur des racines), la pente $> 20\,\%$ (vitesse de ruissellement élevée), la texture du sol (faible cohésion quand saturé, sensibilité au dispersion), et l'intensité pluviométrique (pointes $> 50$ mm/h en saison des pluies).\\
\textbf{Propositions techniques (hiérarchisées selon urgence/efficacité/coût) :}
\begin{enumerate}
\item \textbf{Urgence (stopper l'érosion régressive) :} Stabilisation des têtes de ravins par des seuils en pierres sèches ou en gabions (grillages métalliques remplis de pierres) pour créer un seuil de base local et arrêter la remontée des ravins.
\item \textbf{Court terme (ralentir le ruissellement dans les ravins) :} Mise en place de fascines (branchages entrelacés piquetés) en travers des ravins actifs, associées à des plantations de vétiver (\emph{Chrysopogon zizanioides}) sur les berges (racines profondes $> 2$ m, touffes denses).
\item \textbf{Moyen terme (restaurer le couvert végétal) :} Reboisement de la zone déboisée avec des essences locales ou naturalisées à croissance rapide, racines pivotantes et pouvoir fixateur d'azote (\emph{Acacia auriculiformis}, \emph{Leucaena leucocephala}, \emph{Gliricidia sepium}) en bandes de 10 m le long des lignes de plus grande pente.
\item \textbf{Long terme (aménagement durable des parties cultivées) :} Aménagement de terrasses en pierres sèches (si pierres disponibles) ou en terre battue (renforcées par vétiver sur le talus) sur les parties cultivées (pente $> 15\,\%$) et promotion de l'agroforesterie (maïs + \emph{Gliricidia sepium} en haies vives sur courbes de niveau) pour maintenir un couvert permanent.
\item \textbf{Accompagnement socio-institutionnel :} Sensibilisation des occupants via les comités de développement villageois (CDV), appui technique du MINADER (agriculture) et du MINEPDED (environnement), recherche de financement (FEICOM, projets PNUD/Banque Mondiale).
\end{enumerate}
\end{experience}

\courssub{Étude de cas 2 : Champ en pente à Bafoussam — Choix entre terrasses et courbes de niveau}

\begin{exoresolu}[Choix de la technique anti-érosive à Bafoussam (Ouest, 1\,400 m alt.)]
\textbf{Énoncé :} Un exploitant cultive du maïs et des haricots sur une parcelle de 2 ha, pente moyenne $12\,\%$, sol volcanique (andosol) profond, bien drainé, pluviométrie $1\,800$ mm/an. Il hésite entre terrasses en terre battue et cultures en courbes de niveau.\\
\textbf{Solution détaillée :}
\begin{enumerate}
\item \textbf{Analyse du milieu :} Pente modérée ($12\,\%$), sol perméable mais sensible au lessivage et à l'érosion par splash (impact gouttes), forte pluviométrie, culture annuelle labourée (sol nu partiel).
\item \textbf{Critères de choix :}
\begin{itemize}
\item \textbf{Terrasses en terre battue :} Coût élevé en main-d'œuvre (création et entretien annuel des murets), perte de surface cultivable (emprise des murets $\approx 10-15\,\%$), adaptées aux pentes $> 15\,\%$ ou sols peu profonds. Ici, la pente de $12\,\%$ ne justifie pas l'investissement initial lourd.
\item \textbf{Cultures en courbes de niveau :} Faible coût (traçage au niveau à bulle, niveau à eau ou GPS), maintiennent la surface cultivable, réduisent la vitesse du ruissellement de $50-70\,\%$, favorisent l'infiltration, compatibles avec le labour animal ou motorisé (pratique des \emph{ridger} / billons traditionnels dans l'Ouest).
\end{itemize}
\item \textbf{Recommandation :} Adopter les \textbf{cultures en courbes de niveau} associées à des \textbf{bandes enherbées} (largeur 1 m tous les 10–15 m le long des courbes de niveau) et à un \textbf{paillage} résidus de maïs/haricots. Compléter par une haie vive de \emph{Calliandra calothyrsus} ou \emph{Leucaena leucocephala} en bordure aval pour filtrer les sédiments résiduels.
\item \textbf{Justification :} La pente de $12\,\%$ se situe dans le domaine optimal des courbes de niveau ($5-15\,\%$). Le sol volcanique profond permet une bonne infiltration si le ruissellement est ralenti. La technique est socialement acceptée dans la région (pratique traditionnelle des billons de culture) et peu coûteuse.
\end{enumerate}
\end{exoresolu}

\courssub{Étude de cas 3 : Zone sahélienne de l'Extrême-Nord — Lutte contre l'érosion éolienne}

Dans le département du Logone-et-Chari (Extrême-Nord, climat sahélien, pluviométrie $< 600$ mm/an), l'harmattan (vent sec et violent de NE $\rightarrow$ SO, novembre à mars, vitesses $> 30$ km/h) favorise l'érosion éolienne. Les sols sableux (arénosols), nus après la récolte du mil et du sorgho (saison sèche longue de 7-8 mois), subissent la \cle{déflation} (enlèvement des particules fines $< 0,1$ mm : argiles, limons, matière organique) et la formation de \cle{dunes mobiles} (barchanes) qui envahissent les champs, les pistes et les villages.

\begin{experience}[Dispositif anti-éolien intégré à Kousseri (Logone-et-Chari)]
\textbf{Observation :} Photos satellites et terrain montrent des couloirs de déflation alignés selon le vent dominant (NE–SO), des dunes barchanes de 2–3 m de haut progressant vers les villages et la RN1.\\
\textbf{Diagnostic :} Érosion éolienne par déflation (érosion de surface) et transport saltatoire (grains de sable $0,1-0,5$ mm). Facteurs aggravants : vent fort persistant, sol nu $> 6$ mois/an, absence de haies brise-vent, surpâturage (destruction du couvert herbacé, piétinement croûtes de surface).\\
\textbf{Solutions intégrées (mécanique + biologique + gestion) :}
\begin{enumerate}
\item \textbf{Fixation mécanique (effet immédiat) :} Pose de palissades en branchages (technique des \emph{brise-vent} traditionnels) perpendiculaires au vent dominant, espacées de 10–15 m (10 à 15 fois la hauteur), pour réduire la vitesse du vent au sol sous le seuil de transport ($< 5$ m/s).
\item \textbf{Fixation biologique (effet durable) :} Plantation de haies vives en double rang (espacement 2 m sur le rang, 1 m entre rangs) d'essences xérophiles : \emph{Prosopis juliflora} (croissance très rapide, racines profondes, attention caractère envahissant à surveiller), \emph{Balanites aegyptiaca} (espèce locale, fourragère, fruit comestible), \emph{Acacia senegal} (gommier). Semis direct de \emph{Aristida pungens} (drinn, graminée pérenne locale) pour stabiliser les interdunes.
\item \textbf{Gestion des pâturages (cause racine) :} Rotation des parcours, clôture temporaire des zones en régénération (repos pâturant 2-3 ans), promotion de l'élevage en parcs (zéro pâturage / stabulation) avec apport de fourrage (cultures fourragères : \emph{Pennisetum purpureum}, résidus de récolte) pour réduire le piétinement et la déflation.
\item \textbf{Aménagement hydraulique d'appoint :} Petits barrages de retenue (mares de \emph{bourgou} - \emph{Echinochloa stagnina}) pour maintenir l'humidité du sol en saison sèche et favoriser la végétation spontanée / cultures de décrue.
\end{enumerate}
\textbf{Résultat attendu :} Réduction de $80\,\%$ du transport éolien en 3 ans, récupération de 50 ha de terres cultivables, protection des infrastructures (route RN1, école, centre de santé).
\end{experience}

\courssub{Rédaction d'une conclusion argumentée : Relier causes, conséquences et solutions}

Toute réponse à une question d'interprétation (photo, schéma, carte) doit se terminer par une conclusion synthétique qui articule le chaînage causal complet : \textbf{Facteurs} $\rightarrow$ \textbf{Processus érosif} $\rightarrow$ \textbf{Impacts} $\rightarrow$ \textbf{Mesures adaptées}. Utilisez la structure suivante pour garantir les points de « rédaction scientifique » et « argumentation » au Probatoire.

\begin{methode}[Structure de la conclusion argumentée]
\textbf{1. Rappel du diagnostic :} « Le document montre une érosion hydrique de type ravinement sur un versant déboisé de $22\,\%$ de pente, sol ferralitique argilo-sableux. »\\
\textbf{2. Chaînage causal :} « La déforestation a supprimé l'interception pluviométrique et le rôle fixateur des racines ; la pente forte a accéléré le ruissellement diffus en ruissellement concentré, créant des ravins qui régressent vers l'amont par érosion régressive. »\\
\textbf{3. Conséquences majeures chiffrées si possible :} « Perte estimée de 15 cm de sol arable en 5 ans (perte fertilité), menace directe sur 12 habitations en aval, colmatage du cours d'eau (envasement, eutrophisation). »\\
\textbf{4. Solutions priorisées et justifiées par le milieu :} « En urgence, seuils en gabions aux têtes de ravins (stopper la régression). À moyen terme, reboisement en \emph{Acacia auriculiformis} + terrasses en pierres sèches sur les parties cultivées (pente $> 15\,\%$). À long terme, agroforesterie (maïs + \emph{Gliricidia}) pour pérenniser le couvert. »\\
\textbf{5. Conditions de réussite (gouvernance) :} « La mise en œuvre nécessite l'implication effective des populations (CDV), l'appui technique continu du MINADER/MINEPDED et un financement pérenne (FEICOM, budget communal, partenaires). »
\end{methode}

\begin{attention}[À l'examen, toujours justifier la technique proposée par les caractéristiques du milieu (pente, climat, sol, pression foncière)]
Ne proposez jamais une technique « par cœur » sans l'ancrer dans le document. Exemple d'erreur : « Je propose le reboisement » sans dire pourquoi. Le correcteur attend : « Compte tenu de la pente de $25\,\%$ (forte) et de l'absence de végétation, le reboisement en essences à racines pivotantes est la technique la plus adaptée pour stabiliser le sol et réduire le ruissellement. » De même, ne confondez pas \textbf{courbes de niveau} (pente $5-15\,\%$, sol profond, culture annuelle) et \textbf{terrasses} (pente $> 15\,\%$ ou sol peu profond, culture pérenne ou forte valeur ajoutée). La précision du vocabulaire (\cle{ravinement}, \cle{déflation}, \cle{saltation}, \cle{cône de déjection}, \cle{érosion régressive}, \cle{érosion en nappe}) est notée et valorisée.
\end{attention}

\begin{savaistu}[Le savais-tu ?]
Les questions d'interprétation de photos et de schémas représentent environ 30\% des points au Probatoire SVT (séries C et TI). Les sujets portent souvent sur : (1) une photo de ravin dans les collines de l'Ouest (Bafoussam, Dschang) ou du Centre (Yaoundé), (2) un schéma de versant avec flèches de ruissellement et exagération verticale, (3) une carte d'occupation du sol montrant l'avancée des dunes dans l'Extrême-Nord (Logone-et-Chari, Mayo-Danay). Maîtriser la méthode en 4 étapes (Observer – Décrire – Expliquer – Proposer) et connaître les seuils de pente pour choisir entre courbes de niveau ($5-15\,\%$) et terrasses ($> 15\,\%$) sont des atouts décisifs pour obtenir la note maximale.
\end{savaistu}

\begin{popfigure}
\begin{center}
\begin{tabularx}{\linewidth}{|l|X|X|X|}\hline
\rowcolor{popblueL}
\textbf{Type d'érosion} & \textbf{Causes principales (Facteurs)} & \textbf{Conséquences majeures} & \textbf{Techniques de lutte adaptées} \\\hline
Hydrique : \newline
- Ruissellement diffus (érosion en nappe) \newline
- Ruissellement concentré (érosion en sillons \emph{rills}, ravinement \emph{gullies}) & Déforestation, pente forte ($>15\,\%$), labour dans le sens de la pente, sols nus (jachères courtes), fortes pluies (intensité $> 25$ mm/h), sols sensibles (ferralitiques, latéritiques) & Perte couche arable (horizon A), ravins (gullies), sédimentation cours d'eau (envasement, eutrophisation), glissements de terrain, perte fertilité, menace habitats & Reboisement (essences locales, racines pivotantes), terrasses (pente$>15\,\%$), cultures en courbes de niveau ($5-15\,\%$), fascines, seuils/gabions (têtes ravins), bandes enherbées, paillage, agroforesterie \\\hline
Éolienne : \newline
- Déflation (perte fines) \newline
- Saltation (transport sable) \newline
- Dunes mobiles (barchanes) & Vent fort (harmattan $>30$ km/h), sol sableux nu ($>6$ mois/an), surpâturage, absence de haies, sécheresse, cultures itinérantes sur brûlis & Perte fertilité (argiles, MO), dunes mobiles envahissant champs/pistes/villages, tempêtes de sable (santé), désertification, perte biodiversité & Brise-vent mécaniques (palissades branchages), haies vives (\emph{Prosopis}, \emph{Balanites}, \emph{Acacia senegal}), fixation dunes (\emph{Aristida pungens}), Régénération Naturelle Assistée (RNA), gestion pâturages (rotation, parcs), mares de bourgou \\\hline
\end{tabularx}
\end{center}
\legende{Tableau de synthèse : types d'érosion, causes, conséquences, techniques de lutte adaptées au contexte camerounais.}
\end{popfigure}

\begin{exercice}[Entraînement : Analyse d'un document photographique]
La photographie ci-contre (imaginée) montre un champ de coton en pente ($18\,\%$) dans la région de Garoua (Nord, climat soudano-sahélien), après un orage violent en août. On voit des sillons d'érosion parallèles au labour, des dépôts de sable au bas du champ, et une haie d'euphorbes (\emph{Euphorbia balsamifera}) clairsemée en bordure aval.\\
\textbf{Consignes :} Appliquez la méthode en 4 étapes (Observer, Décrire, Expliquer, Proposer) pour rédiger une réponse complète (diagnostic, causes, conséquences, propositions). Justifiez précisément le choix entre terrasses et courbes de niveau au vu de la pente.
\end{exercice}

\begin{corrige}[Exercice n°1]
\textbf{1. Observer :} Pente marquée ($\approx 18\,\%$), labour en lignes de pente (sillons parallèles à la plus grande pente), sol nu entre les plants de coton (stade jeune), dépôts sableux (colluvions) en pied de champ, haie d'euphorbes (peu dense, hauteur $\approx 1$ m) en bordure aval.\\
\textbf{2. Décrire :} Érosion hydrique par ruissellement diffus concentré dans les sillons de labour (érosion en sillons / \emph{rill erosion}). Dépôts de colluvions au bas du versant (pied de pente). Haie vive résiduelle peu efficace.\\
\textbf{3. Expliquer :} Le labour dans le sens de la pente crée des voies préférentielles pour l'écoulement de l'eau (ruissellement concentré artificiel). La pente de $18\,\%$ donne une énergie cinétique suffisante pour détacher et transporter les particules (pouvoir érosif $\propto v^2$). L'absence de couvert végétal entre les rangs expose le sol à l'impact direct des gouttes (érosion par \emph{splash}, dispersion des agrégats). La haie d'euphorbes freine partiellement les sédiments grossiers mais est insuffisante (peu dense, racines superficielles).\\
\textbf{4. Proposer :} 
\begin{itemize}
\item \textbf{Terrasses en terre battue} (renforcées par vétiver sur le talus) \textbf{car pente $> 15\,\%$} : elles réduisent la longueur effective de pente, favorisent l'infiltration et retiennent les sédiments. C'est la technique de référence pour cette pente.
\item \textbf{Courbes de niveau} comme alternative transitoire si la main-d'œuvre/financement manquent pour les terrasses : tracer les lignes de niveau au niveau à bulle/eau, labourer et semer perpendiculairement à la pente. Efficacité moindre sur pente $> 15\,\%$ (risque de rupture des billons).
\item \textbf{Bandes enherbées} (largeur 1 m, \emph{Pennisetum purpureum} ou \emph{Brachiaria}) tous les 10 m le long des courbes de niveau/terrasses pour filtrer le ruissellement.
\item \textbf{Renforcement de la haie aval} : planter une double rangée de \emph{Gliricidia sepium} ou \emph{Leucaena leucocephala} (racines profondes, fixatrices d'azote) en aval pour piéger les sédiments résiduels et fournir du fourrage/bois.
\end{itemize}
\textbf{Justification du choix terrasses vs courbes de niveau :} La pente de $18\,\%$ place la parcelle dans le domaine des terrasses (seuil pédagogique et technique $15\,\%$). Cependant, le coton est une culture annuelle ; les terrasses en terre battue sont réversibles et moins coûteuses que les terrasses en pierres. Si l'exploitant ne peut pas investir dans la création des terrasses (main-d'œuvre importante), les courbes de niveau + bandes enherbées + paillage constituent une solution transitoire efficace mais nécessitant un entretien rigoureux chaque campagne.
\end{corrige}

\newpage

\coursec{Activité d'intégration}

\coursec{Lutte contre l'érosion}

\courssub{Objectifs de l'activité}
Cette activité d'intégration vise à mobiliser l'ensemble des connaissances et savoir-faire acquis au cours de la leçon sur la \cle{lutte contre l'érosion}. Elle permet à l'élève de :
\begin{itemize}
    \item Identifier et analyser les facteurs favorisant l'\cle{érosion} sur un versant réel.
    \item Prévoir les conséquences environnementales et socio-économiques d'une érosion non maîtrisée.
    \item Concevoir un plan d'aménagement cohérent en sélectionnant et en localisant judicieusement les techniques de lutte adaptées (reboisement, terrasses, cultures en \cle{courbes de niveau}, bandes enherbées).
    \item Justifier scientifiquement chaque choix technique en lien avec les processus érosifs (ruissellement, \cle{ravinement}, érosion diffuse).
    \item Communiquer ses propositions sous forme de schémas légendés et d'arguments structurés.
\end{itemize}

\courssub{Situation : Le versant de Nkolbisson menacé}
Situé au sud-ouest de Yaoundé, le quartier de Nkolbisson s'étend sur un versant fortement pentu du bassin versant de la Mfoundi. L'urbanisation rapide et l'extension des cultures vivrières ont profondément modifié le couvert végétal originel. Une étude topographique récente réalisée par la communauté urbaine met en évidence la configuration suivante (voir le schéma simplifié ci-dessous) :

\begin{popfigure}
\begin{tikzpicture}[scale=0.9, every node/.style={font=\small}]
    % Échelle et repères
    \draw[->, thick] (-0.5,0) -- (13,0) node[right] {Distance horizontale (m)};
    \draw[->, thick] (0,-0.5) -- (0,7) node[above] {Altitude (m)};
    \foreach \x in {0,2,...,12} \draw (\x,0.1) -- (\x,-0.1) node[below] {\x0};
    \foreach \y in {0,1,...,6} \draw (0.1,\y) -- (-0.1,\y) node[left] {\y0};
    
    % Profil du versant (courbe)
    \draw[thick, popdark] plot[smooth, tension=0.7] coordinates {
        (0,6) (2,5.5) (4,4.5) (6,3.2) (8,2.2) (10,1.3) (12,0.5)
    };
    
    % Zones
    % Sommet déboisé (0-3)
    \fill[poporange!30, opacity=0.5] (0,6) -- plot[smooth, tension=0.7] coordinates {(0,6) (2,5.5) (4,4.5)} -- (4,0) -- (0,0) -- cycle;
    \node[align=center, poporange!80!black] at (2,5.8) {\textbf{Zone A : Sommet}\\Déboisement total\\Pente $\approx$ 25\%};
    
    % Champ manioc sens pente (3-7)
    \fill[popgreen!30, opacity=0.5] (4,4.5) -- plot[smooth, tension=0.7] coordinates {(4,4.5) (6,3.2) (8,2.2)} -- (8,0) -- (4,0) -- cycle;
    \draw[pattern=north east lines, pattern color=popgreen!80!black] (4,4.5) -- plot[smooth, tension=0.7] coordinates {(4,4.5) (6,3.2) (8,2.2)} -- (8,0) -- (4,0) -- cycle;
    \node[align=center, popgreen!70!black] at (6,3.5) {\textbf{Zone B : Mi-versant}\\Champ de manioc\\Culture \textbf{dans le sens de la pente}\\Pente $\approx$ 15-20\%};
    % Flèches sens pente
    \foreach \x in {4.5,5.5,6.5,7.5} {
        \draw[-{Stealth[length=2mm]}, thick, popgreen!70!black] (\x, {3.2 - 0.55*(\x-4)}) -- ++(0.8, -0.44);
    }
    
    % Ravin naissant (vers x=9)
    \draw[thick, poppink, decorate, decoration={zigzag, segment length=4, amplitude=1.5}] (9,2.5) -- (9.5,1.5);
    \node[poppink, align=center] at (9.5,2.8) {\textbf{Ravin naissant}\\Incision active};
    
    % Habitations (bas versant)
    \fill[popblue!30, opacity=0.5] (10,1.3) -- (12,0.5) -- (12,0) -- (10,0) -- cycle;
    \draw[popblue] (10.2,0.2) rectangle (11.8,1.0);
    \draw[popblue] (10.2,0.2) -- (11,0.7) -- (11.8,0.2);
    \node[align=center, popblue!80!black] at (11,0.2) {\textbf{Zone C : Bas-versant}\\Habitations menacées\\Pente $\approx$ 10\%};
    
    % Flèche ruissellement général
    \draw[-{Stealth[length=3mm]}, very thick, popblue] (1,5.8) -- (11,0.8) node[midway, above, sloped, align=center, font=\tiny] {Ruissellement concentré $\rightarrow$};
\end{tikzpicture}
\legende{Profil topographique simplifié du versant de Nkolbisson (Yaoundé) montrant les trois zones critiques : A (sommet déboisé), B (champ de manioc cultivé dans le sens de la pente), C (habitations en contrebas menacées par un ravin naissant).}
\end{popfigure}

\textbf{Données chiffrées clés :}
\begin{itemize}
    \item \textbf{Zone A (Sommet, 0--300 m horizontal) :} Pente moyenne \textbf{25\%}. Ancienne forêt galerie rasée il y a 3 ans pour du bois de chauffe. Sol nu, battu par les pluies violentes (Yaoundé : \SI{1600}{mm/an}, pointe en septembre-octobre). Absence de litière.
    \item \textbf{Zone B (Mi-versant, 300--800 m) :} Pente moyenne \textbf{17\%}. Champ de manioc (\emph{Manihot esculenta}) en culture pure. Labour et billons réalisés \textbf{parallèles à la ligne de plus grande pente} (sens de l'écoulement de l'eau). Sol sableux-argileux, sensible à la \cle{battance}.
    \item \textbf{Zone C (Bas-versant, 800--1200 m) :} Pente moyenne \textbf{10\%}. Habitat dense non planifié. Un ravin en forme de V s'est initié à la limite Zone B/Zone C (à \SI{800}{m} environ), s'élargissant de \SI{0,5}{m/an} selon les riverains. Les fondations des maisons en bas de pente sont menacées.
    \item \textbf{Climat :} Saison des pluies longue (mars-juin) et courte (sept-nov). Intensités pluvieuses fréquentes $> \SI{30}{mm/h}$.
\end{itemize}

\courssub{Consignes}
Travailler par groupes de 4 à 5 élèves. Chaque groupe rend un rapport écrit (2 pages max) et un schéma du versant aménagé (format A3). Les questions guident votre démarche :

\begin{enumerate}
    \item \textbf{Diagnostic des causes (Identification) :} Sur le schéma fourni, repérez et entourez les \textbf{facteurs favorisant l'érosion} (naturels et anthropiques) pour chaque zone (A, B, C). Classez-les en facteurs \emph{climatiques}, \emph{topographiques}, \emph{édaphiques} (liés au sol) et \emph{anthropiques}. Expliquez le mécanisme par lequel chaque facteur agit (ex : « Pente 25\% $\rightarrow$ vitesse du ruissellement $\rightarrow$ pouvoir érosif »).
    \item \textbf{Prévision des conséquences (Anticipation) :} Si aucune action n'est entreprise, décrivez l'évolution probable du versant à 5 ans et 10 ans. Citez au moins \textbf{4 conséquences distinctes} : sur le sol (perte de fertilité), sur l'hydrologie (crues, pollution), sur les infrastructures (maisons, routes) et sur la biodiversité.
    \item \textbf{Proposition d'aménagement (Résolution) :} Proposez un \textbf{zonage des techniques de lutte} en justifiant \textbf{chaque choix} par le processus érosif qu'il combat :
    \begin{itemize}
        \item Zone A : Technique(s) prioritaire(s) ? Pourquoi au sommet ?
        \item Zone B : Technique(s) adaptée(s) à l'agriculture ? Comment les mettre en œuvre concrètement (espacement, orientation) ?
        \item Zone C : Technique(s) de protection des biens et personnes ? Liaison avec le ravin.
    \end{itemize}
    \item \textbf{Représentation graphique (Communication) :} Dessinez le \textbf{versant aménagé} (profil ou vue de dessus). Légendez obligatoirement : les nouvelles structures (terrasses, bandes enherbées, reboisement), le sens des cultures, le tracé du ravin stabilisé, les écoulements d'eau maîtrisés. Indiquez les largeurs et espacements estimés.
    \item \textbf{Argumentaire de soutenance :} Préparez une présentation orale de 3 minutes pour défendre votre plan face à un « comité de quartier » (le reste de la classe). Insistez sur le rapport coût/efficacité et l'acceptabilité sociale (maintien de la production de manioc).
\end{enumerate}

\courssub{Ressources à mobiliser}
Pour réaliser cette activité, vous devez réactiver les savoirs suivants (rappel synthétique) :

\begin{aretenir}[Rappel des techniques de lutte et conditions d'application]
\textbf{1. Reboisement / Boisement (Zone de départ / Sommet)}
\begin{itemize}
    \item \textbf{Rôle :} Interception des gouttes de pluie (réduction énergie cinétique), augmentation de l'infiltration par les racines/litière, fixation mécanique du sol.
    \item \textbf{Essences adaptées au contexte Yaoundé :} \emph{Acacia mangium}, \emph{Leucaena leucocephala} (croissance rapide, fixation azote), \emph{Terminalia superba} (Fraké, essence locale), \emph{Gmelina arborea}. Association avec des herbacées (Pueraria, Brachiaria) pour couverture rapide.
    \item \textbf{Condition :} Nécessaire sur pente $> 20\%$ et sols nus.
\end{itemize}

\textbf{2. Cultures en courbes de niveau (Zone culturale / Mi-versant)}
\begin{itemize}
    \item \textbf{Principe :} Labour, semis, billons \textbf{perpendiculaires à la pente} (suivent les courbes de niveau). Crée des micro-barrages qui ralentissent le ruissellement, favorisent l'infiltration et retiennent les particules.
    \item \textbf{Mise en œuvre :} Traçage au niveau (niveau à bulle, niveau à eau, cadre en A). Espacement entre lignes : $E = \frac{100}{\text{pente en \%}} \times k$ (k $\approx$ 0,5 à 1 pour cultures annuelles). Pour 17\% : $E \approx \SI{3}{m}$ à \SI{6}{m}.
    \item \textbf{Association :} Bande enherbée (vetiver, citronnelle) tous les 10-20 m ou en bordure de parcelle.
\end{itemize}

\textbf{3. Terrasses (Pentes fortes / Zones instables)}
\begin{itemize}
    \item \textbf{Type :} Terrasses à gradins (bench terraces) pour pentes $> 15-20\%$ si culture intensive. Terrasses en talus (risbermes) pour stabilisation physique.
    \item \textbf{Coût :} Élevé (travail du sol, main d'œuvre). Réservé aux zones à fort enjeu ou pentes extrêmes.
\end{itemize}

\textbf{4. Bandes enherbées / Haies vives / Fascines (Bas de versant / Ravin)}
\begin{itemize}
    \item \textbf{Bande enherbée (Vetiver \emph{Chrysopogon zizanioides}) :} Racines pivotantes profondes ($> \SI{3}{m}$), tiges rigides freinent l'eau, résiste à la sécheresse et au feu. Plantation en lignes serrées (10-15 cm) sur courbe de niveau.
    \item \textbf{Haie vive :} \emph{Gliricidia sepium}, \emph{Calliandra calothyrsus} (fourrage + piquets + azote).
    \item \textbf{Fascines / Seuils :} Branches entrelacées piquetées dans le fond du ravin pour casser la pente longitudinale et piéger les sédiments (lutte contre l'incision verticale).
\end{itemize}

\textbf{5. Règle générale d'aménagement intégré (Bassin versant) :}
\textit{« Protéger le haut pour sauver le bas. »} Agir en priorité sur la zone de production du ruissellement (sommet) et la zone de transit (mi-versant) avant le traitement curatif du ravin (bas).
\end{aretenir}

\courssub{Démarche et corrigé commenté}
\begin{exoresolu}[Corrigé type et guide de l'enseignant]
\noindent\textbf{1. Diagnostic des causes (Identification des facteurs)}

\begin{center}
\begin{tabularx}{\linewidth}{|l|X|X|}\hline
\textbf{Zone} & \textbf{Facteurs identifiés} & \textbf{Mécanisme érosif associé} \\ \hline
\textbf{A (Sommet)} & \textbf{Naturels :} Pente forte (25\%), Pluviométrie intense (énergie cinétique gouttes). \newline \textbf{Anthropiques :} Déforestation totale (sol nu), Absence de litière, Compactage (piétinement/bois de chauffe). & \textbf{Battance} : Gouttes détruisent la structure superficielle $\rightarrow$ croûte de battance $\rightarrow$ infiltration nulle. \newline \textbf{Ruissellement de surface (Hortonien) :} Déclenché dès le début de la pluie. Vitesse élevée (pente) $\rightarrow$ pouvoir de détachement et transport maximal. \newline \textbf{Érosion diffuse (interrill) $\rightarrow$ début de ravinement.} \\ \hline
\textbf{B (Mi-versant)} & \textbf{Naturels :} Pente moyenne (17\%), Sol sableux-argileux sensible. \newline \textbf{Anthropiques :} \textbf{Culture dans le sens de la pente} (billons parallèles à la pente), Monoculture manioc (couvert lent, sol nu longtemps), Labour répété (semelle de labour). & \textbf{Effet « gouttière » :} Les billons dans le sens de la pente canalisent l'eau $\rightarrow$ concentration du ruissellement $\rightarrow$ érosion linéique (rigoles) rapide. \newline \textbf{Apport sédimentaire massif} depuis Zone A + production locale $\rightarrow$ surcharge solide. \\ \hline
\textbf{C (Bas-versant)} & \textbf{Naturels :} Rupture de pente (convergence écoulements), Ravin existant (tête de ravin active). \newline \textbf{Anthropiques :} Habitat en zone à risque, Imperméabilisation (toits, cours), Absence de drainage. & \textbf{Ravinement (Gully erosion) :} Incision verticale (régression amont) + élargissement latéral. \newline \textbf{Dépôts alluvionnaires} (colmatage, pollution). \newline \textbf{Menace directe} : Affouillement fondations, glissements de terrain. \\ \hline
\end{tabularx}
\end{center}

\noindent\textbf{2. Prévision des conséquences (Scénario « Business as usual » à 10 ans)}

\begin{enumerate}
    \item \textbf{Dégradation irréversible du sol (Zone A \& B) :} Perte de l'horizon A (humus, argiles, éléments nutritifs). Sol réduit à l'horizon B (latéritique, stérile, durci). Perte de productivité agricole $> 50\%$. Abandon des champs.
    \item \textbf{Aggravation du ravin (Zone C) :} Régression amont de la tête de ravin (avancée de \SI{2}{m} à \SI{5}{m/an} en sols meubles). Atteinte des premières habitations d'ici 3-5 ans. Élargissement du ravin ($> \SI{10}{m}$ large) $\rightarrow$ coupure de voies d'accès.
    \item \textbf{Risques hydrologiques et sanitaires :} Crues éclaires plus violentes et fréquentes en aval (Mfoundi). Pollution des eaux par turbidité, pesticides, déchets domestiques (choléra, typhoïde). Colmatage des ouvrages hydrauliques (barrage d'Akoaman, prise d'eau).
    \item \textbf{Perte de biodiversité et de services écosystémiques :} Disparition de la faune du sol (vers de terre, termites), perte de la régulation climatique locale (îlot de chaleur), perte de ressources (bois, plantes médicinales).
\end{enumerate}

\noindent\textbf{3. Proposition d'aménagement zoné (Plan de sauvetage)}

\begin{center}
\begin{tabularx}{\linewidth}{|l|X|l|}\hline
\textbf{Zone} & \textbf{Techniques proposées \& Mise en œuvre} & \textbf{Justification scientifique} \\ \hline
\textbf{A : Sommet (Priorité 1)} & \textbf{Reboisement d'urgence + Zaï / Demi-lunes.} \newline - Plantation \emph{Acacia mangium} (3m x 3m) + \emph{Vetiver} en lignes de pied de pente. \newline - Travail du sol : Zaï (trous \SI{30}{cm} diam, \SI{20}{cm} prof, fumure) pour casser la croûte et infiltrer. \newline - Interdiction de coupe (convention locale). & \textbf{Casser l'énergie à la source.} \newline Arbres : Interception (30\% pluie), Racines (porosité, structure), Litière (protection). \newline Zaï : Infiltration localisée, arrêt ruissellement diffus. \newline \textit{Objectif : Réduire le débit de pointe sortant de Zone A de 60-80\%.} \\ \hline
\textbf{B : Mi-versant (Priorité 2)} & \textbf{Courbes de niveau + Bandes enherbées (Vetiver) + Agroforesterie.} \newline - Traçage courbes de niveau (cadre en A) tous les \SI{4}{m} (pente 17\%). \newline - Billons \textbf{perpendiculaires à la pente}. \newline - Bande de Vetiver (\SI{1}{m} large) tous les \SI{20}{m} ou à chaque rupture de pente. \newline - Introduction \emph{Gliricidia} en haies vives (alignement courbes de niveau) : tuteur manioc + paillis + azote. & \textbf{Maîtriser le transit.} \newline Courbes de niveau : Ralentissement ($V \downarrow$), Infiltration $\uparrow$, Dépôt sédiments $\uparrow$. \newline Vetiver : Barrage végétal dense, racines profondes (ancrage), non invasif. \newline Agroforesterie : Maintien production manioc + amélioration sol (C/N, structure). \textit{Acceptabilité paysanne forte.} \\ \hline
\textbf{C : Bas-versant / Ravin (Priorité 3 - Curatif)} & \textbf{Stabilisation ravin + Bandes tampons + Drainage.} \newline - \textbf{Fond de ravin :} Seuils en fascines (bois local/bambou) ou gabions (si fonds dispo) tous les \SI{5}{m} dénivelé. \newline - \textbf{Bords :} Plantation Vetiver (3 lignes serrées) + \emph{Bambou} (racines fasciculées) en haut des berges. \newline - \textbf{Habitations :} Bande enherbée \SI{5}{m} large au-dessus des maisons. Drainage périphérique (fossés en courbes de niveau) vers exutoire stable. & \textbf{Protéger les vies et biens / Arrêter la régression.} \newline Seuils : Casser la pente longitudinale du ravin $\rightarrow$ Vitesse $\downarrow$ $\rightarrow$ Dépôt $\uparrow$ $\rightarrow$ Colmatage naturel. \newline Végétalisation berges : Cohésion mécanique (racines), buttoir physique. \newline Bande tampon : Filtration ultime avant habitat. Drainage : Éviter saturation sol (glissement). \\ \hline
\end{tabularx}
\end{center}

\noindent\textbf{4. Schéma du versant aménagé (Description pour le dessin)}
Voir la figure ci-après représentant le profil aménagé.

\noindent\textbf{5. Éléments de l'argumentaire de soutenance (Points clés à faire valoir)}
\begin{itemize}
    \item \textbf{Efficacité prouvée :} Courbes de niveau + Vetiver réduisent pertes en sol de 70-90\% (réf. IRAD, essais Dschang/Yaoundé).
    \item \textbf{Coût maîtrisé :} Reboisement (pépinière communautaire, main d'œuvre locale), Courbes de niveau (main d'œuvre familiale, cadre en A coût négligeable), Vetiver (boutures gratuites/peu chères). Pas de génie civil lourd (gabions) sauf si ravin $> \SI{3}{m}$ profond.
    \item \textbf{Maintien production :} Haies \emph{Gliricidia} = tuteurs manioc + engrais vert + fourrage chèvres. Pas de perte de surface cultivable (bandes étroites).
    \item \textbf{Durabilité / Appropriation :} Convention locale (chef de quartier, chef de terre) pour protection Zone A. Formation « École Champs Paysans » (ECP) pour traçage courbes de niveau. Comité de suivi ravin (riverains).
    \item \textbf{Co-bénéfices :} Bois d'œuvre/énergie (Zone A à 5-7 ans), Séquestration C, Embellissement cadre de vie, Réduction risques inondations centre-ville Yaoundé.
\end{itemize}
\tcblower
\textbf{Commentaires pédagogiques pour l'enseignant :}
\begin{itemize}
    \item \textbf{Erreur fréquente Q1 :} Confondre « facteur » et « conséquence ». Ex : « Le ravin est un facteur ». \textbf{Correction :} Le ravin est une conséquence/forme d'érosion. Les facteurs sont ce qui le cause (pente, sol nu, concentration eau).
    \item \textbf{Erreur fréquente Q3 :} Proposer des terrasses en pierre (bench terraces) en Zone B. \textbf{Correction :} Coût/travail disproportionné pour du manioc. Les courbes de niveau (billons) suffisent pour pente 17\% et culture annuelle. Réserver terrasses en dur aux pentes $> 30\%$ ou cultures pérennes (thé, café).
    \item \textbf{Erreur fréquente Q3 :} Oublier le lien Zone A $\rightarrow$ Zone B. \textbf{Correction :} Si Zone A n'est pas traitée, le débit arrivant en Zone B submerge les courbes de niveau. \textit{Approche bassin versant obligatoire.}
    \item \textbf{Erreur fréquente Q4 :} Dessiner les billons encore dans le sens de la pente. \textbf{Vérification :} Vérifier la perpendicularité aux courbes de niveau (donc horizontaux sur le profil).
    \item \textbf{Niveau d'exigence Probatoire :} L'élève doit utiliser le vocabulaire précis : \emph{battance, ruissellement hortonien, érosion diffuse/linéique/ravinement, courbes de niveau, bande enherbée, tête de ravin, régression amont, infiltration, pouvoir érosif}. Le schéma doit être légendé (échelle, nord, symboles).
\end{itemize}
\end{exoresolu}

% Figure déplacée hors de la boîte exoresolu pour respecter la non-imbrication
\begin{popfigure}
\begin{tikzpicture}[scale=0.85, every node/.style={font=\small}]
    % Axes
    \draw[->, thick] (-0.5,0) -- (13,0) node[right] {Distance (m)};
    \draw[->, thick] (0,-0.5) -- (0,7) node[above] {Altitude (m)};
    \foreach \x in {0,2,...,12} \draw (\x,0.1) -- (\x,-0.1) node[below] {\x0};
    \foreach \y in {0,1,...,6} \draw (0.1,\y) -- (-0.1,\y) node[left] {\y0};
    
    % Profil initial (pointillés)
    \draw[thick, dashed, popmuted] plot[smooth, tension=0.7] coordinates {
        (0,6) (2,5.5) (4,4.5) (6,3.2) (8,2.2) (10,1.3) (12,0.5)
    };
    \node[popmuted, align=left, font=\tiny] at (1,6.2) {Profil initial};
    
    % NOUVEAU PROFIL AMÉNAGÉ
    % Zone A : Reboisement (surface rugueuse)
    \draw[thick, popgreen!70!black] plot[smooth, tension=0.7] coordinates {
        (0,6.1) (2,5.6) (4,4.6)
    };
    \fill[popgreen!20, opacity=0.6] (0,6.1) -- plot[smooth, tension=0.7] coordinates {(0,6.1) (2,5.6) (4,4.6)} -- (4,0) -- (0,0) -- cycle;
    \node[align=center, popgreen!70!black] at (2,5.9) {\textbf{ZA : Reboisement}\\Acacia + Vetiver pied pente\\Zaï (trous fumés)};
    % Arbres symboliques
    \foreach \x in {0.5, 1.5, 2.5, 3.5} {
        \draw[popgreen!70!black] (\x, {5.9 - 0.25*\x}) -- ++(0,0.3);
        \draw[popgreen!70!black, fill=popgreen!30] (\x, {6.2 - 0.25*\x}) circle (0.15);
    }
    
    % Zone B : Courbes de niveau (terrasses légères / billons)
    \draw[thick, popblue!70!black] plot[smooth, tension=0.7] coordinates {
        (4,4.6) (6,3.3) (8,2.3)
    };
    \fill[popblue!15, opacity=0.6] (4,4.6) -- plot[smooth, tension=0.7] coordinates {(4,4.6) (6,3.3) (8,2.3)} -- (8,0) -- (4,0) -- cycle;
    \node[align=center, popblue!70!black] at (6,3.6) {\textbf{ZB : Courbes de niveau}\\Billons $\perp$ pente (esp. 4m)\\Haies Gliricidia + Bandes Vetiver 20m};
    % Symboles billons perpendiculaires (traits courts horizontaux)
    \foreach \y in {4.4, 4.0, 3.6, 3.2, 2.8, 2.4} {
        \draw[popblue!70!black, thick] (4.2, \y) -- (7.8, \y);
    }
    % Haies Gliricidia (points)
    \foreach \y in {4.2, 3.4, 2.6} {
        \foreach \x in {4.5, 5.5, 6.5, 7.5} {
            \fill[poppurple] (\x, \y) circle (0.05);
        }
    }
    
    % Zone C : Ravin stabilisé + Habitations protégées
    \draw[thick, poporange!70!black] plot[smooth, tension=0.7] coordinates {
        (8,2.3) (10,1.4) (12,0.6)
    };
    \fill[poporange!15, opacity=0.6] (8,2.3) -- plot[smooth, tension=0.7] coordinates {(8,2.3) (10,1.4) (12,0.6)} -- (12,0) -- (8,0) -- cycle;
    
    % Ravin stabilisé (forme U large, végétalisé)
    \draw[thick, popgreen!70!black] (9,2.2) .. controls (9.3,1.8) and (9.5,1.2) .. (9.5,0.8);
    \draw[thick, popgreen!70!black] (9.5,2.2) .. controls (9.2,1.8) and (9,1.2) .. (9,0.8);
    \fill[popgreen!30] (9,2.2) .. controls (9.3,1.8) and (9.5,1.2) .. (9.5,0.8) -- (9.5,0) -- (9,0) -- cycle;
    \node[align=center, popgreen!70!black, font=\tiny] at (9.5,2.5) {Ravin stabilisé\\\textbf{Seuils fascines} + \textbf{Vetiver 3 lignes}\\Bambou berges};
    
    % Habitations
    \draw[popblue] (10.2,0.2) rectangle (11.8,1.0);
    \draw[popblue] (10.2,0.2) -- (11,0.7) -- (11.8,0.2);
    \node[align=center, popblue!80!black] at (11,0.2) {\textbf{ZC : Habitat protégé}\\Bande tampon Vetiver 5m\\Drainage périphérique};
    % Bande tampon
    \fill[popgreen!40!yellow, opacity=0.5] (10,1.4) -- (10,0.8) -- (12,0) -- (12,0.6) -- cycle;
    \draw[pattern=north west lines, pattern color=popgreen!70!black] (10,1.4) -- (10,0.8) -- (12,0) -- (12,0.6) -- cycle;
    
    % Flèches ruissellement maîtrisé
    \draw[-{Stealth[length=2mm]}, thick, popblue] (1,5.9) -- (3.5,4.8) node[midway, above, sloped, font=\tiny] {Infiltration $\uparrow$};
    \draw[-{Stealth[length=2mm]}, thick, popblue] (5,4.0) -- (7.5,2.8) node[midway, above, sloped, font=\tiny] {Ruissellement ralenti};
    \draw[-{Stealth[length=2mm]}, thick, popblue] (9.5,1.5) -- (11,0.5) node[midway, above, sloped, font=\tiny] {Écoulement filtré};
\end{tikzpicture}
\legende{Profil du versant de Nkolbisson APRÈS aménagement : Zone A reboisée (rugosité, infiltration), Zone B cultivée en courbes de niveau (billons perpendiculaires, haies vives, bandes de Vetiver), Zone C avec ravin stabilisé (seuils, végétalisation) et bande tampon protectrice pour l'habitat.}
\end{popfigure}

\courssub{Bilan de l'activité}
Cette activité d'intégration a permis de mettre en évidence que la lutte contre l'érosion ne se résume pas à une technique unique, mais relève d'une \textbf{stratégie intégrée à l'échelle du bassin versant (ou du versant)}. Les points clés acquis sont :

\begin{enumerate}
    \item \textbf{La hiérarchie des interventions :} La priorité absolue est la \textbf{protection de la zone de production du ruissellement (sommet, Zone A)} par le reboisement et la restauration de la capacité d'infiltration (Zaï). Traiter le ravin (Zone C) sans traiter le haut est voué à l'échec (coût infini, régression continue).
    \item \textbf{L'adaptation des techniques au contexte biophysique et socio-économique :} En Zone B (agricole), les \textbf{courbes de niveau} (zéro coût matériel, main d'œuvre locale) associées aux \textbf{bandes de Vetiver} et à l'\textbf{agroforesterie (Gliricidia)} constituent le meilleur compromis entre conservation des sols et maintien de la production vivrière (manioc). Les terrasses en dur sont exclues pour des raisons de rentabilité.
    \item \textbf{La complémentarité des fonctions :} Chaque technique remplit une fonction précise : \emph{Interception/Infiltration} (Arbres Zone A), \emph{Ralentissement/Rétention} (Courbes de niveau + Vetiver Zone B), \emph{Stabilisation mécanique/Filtration} (Seuils + Vetiver + Bambou Zone C).
    \item \textbf{La dimension humaine :} La réussite technique suppose une \textbf{gouvernance locale} (conventions, comités de suivi, Écoles Champs Paysans). La technique la plus performante échoue si elle n'est pas appropriée par les usagers (paysans, riverains).
    \item \textbf{Compétences mobilisées :} Observation/Analyse de documents (schéma, données chiffrées) $\rightarrow$ Raisonnement causal (Facteurs $\rightarrow$ Processus $\rightarrow$ Conséquences) $\rightarrow$ Conception de solutions (Choix techniques argumentés) $\rightarrow$ Communication graphique et orale.
\end{enumerate}

\begin{attention}[Pièges à éviter pour le Probatoire]
\begin{itemize}
    \item Ne pas confondre \textbf{terrasses (génie civil, modification topographie)} et \textbf{cultures en courbes de niveau (pratique culturale, suit topographie)}.
    \item Ne pas citer « plantation d'arbres » sans préciser \textbf{l'essence, l'espacement, la strate (arbre/herbacée) et la localisation} (sommet, bordure, haie).
    \item Oublier que le \textbf{Vetiver} est une herbe stérile (non invasive) aux racines pivotantes, idéale pour les bandes enherbées, contrairement aux bambous (traçants, invasifs) qu'on réserve aux berges de ravin.
    \item Négliger le \textbf{drainage} en bas de versant : un sol saturé en eau (nappe montante) perd sa cohésion $\rightarrow$ glissement de terrain. Le drainage périphérique de l'habitat est vital.
\end{itemize}
\end{attention}

\begin{savaistu}[Le Vetiver : « L'herbe miracle » contre l'érosion au Cameroun]
Originaire d'Inde, \emph{Chrysopogon zizanioides} (Vetiver) est introduit au Cameroun depuis les années 90 (projet PNVRA, IRAD). Ses atouts uniques :
\begin{itemize}
    \item Système racinaire \textbf{fasciculé, vertical, profond (3-4 m)} : ancre le sol, ne concurrence pas les cultures voisines (racines superficielles).
    \item Tiges \textbf{rigides, érigées, denses} : résistent au courant fort, filtrent l'eau, piègent les sédiments (formation de terrasses naturelles au fil du temps).
    \item \textbf{Stérile} (ne produit pas de graines viables) : ne devient pas une mauvaise herbe invasive.
    \item Tolérance extrême : pH 3-11, salinité, sécheresse (\SI{6}{mois} sans pluie), feu, métaux lourds, inondation temporaire.
    \item Usages multiples : Huile essentielle (parfumerie - racines), fourrage (jeunes feuilles), paillis, artisanat (chapeaux, nattes), phytoremédiation.
\end{itemize}
Au Cameroun, le « Réseau Vetiver Cameroun » (RVC) et l'IRAD diffusent des plants certifiés (cultivar \emph{Monto} ou \emph{Sunshine}). C'est la technologie de référence \textbf{faible coût / haute efficacité} pour la stabilisation des ravins, protection des berges, bandes enherbées en agriculture et protection des infrastructures routières (coupures, remblais).
\end{savaistu}

\newpage

\coursec{Méthodes et exercices résolus}

\courssub{Méthodes et exercices résolus}

\begin{methode}[Identifier les causes de l'érosion dans une situation donnée]
    \textbf{Objectif :} À partir d'un texte, d'un tableau de données ou d'une photo, repérer les facteurs naturels et anthropiques qui déclenchent ou aggravent l'érosion.
    \begin{enumerate}
        \item \textbf{Analyser le support :} Lire attentivement le texte ou observer l'image (pente, état de la végétation, type de sol, activité humaine, climat).
        \item \textbf{Classer les facteurs :} Distinguer les \textbf{facteurs naturels} (pente forte, pluviométrie intense/violente, nature du sol : sableux/argileux, vent fort) des \textbf{facteurs anthropiques} (déforestation, surpâturage, labour dans le sens de la pente, cultures itinérantes sur brûlis, urbanisation non maîtrisée).
        \item \textbf{Établir les liens de causalité :} Relier chaque facteur observé au type d'érosion concerné (hydrique : ruissellement, ravines ; éolienne : déflation, dunes).
        \item \textbf{Rédiger la réponse :} Citer explicitement les facteurs identifiés en les classant, et expliquer le mécanisme (ex: « La déforestation expose le sol nu à l'impact direct des gouttes de pluie qui détruisent la structure du sol \emph{(battance)}, favorisant le ruissellement et l'érosion hydrique »).
    \end{enumerate}
    \textbf{Réflexe Probatoire :} Ne jamais dire « c'est à cause de la pluie ». Dire : « L'intensité des pluies (facteur climatique) combinée à l'absence de couverture végétale (facteur anthropique) favorise le détachement des particules par battance. »
\end{methode}

\begin{exoresolu}[Analyse des causes d'une érosion hydrique dans les Hauts-Plateaux de l'Ouest]
    \textbf{Énoncé :} Dans la région des Hauts-Plateaux de l'Ouest-Cameroun (Bamenda, Dschang), on observe une multiplication de ravines et une turbidité accrue des cours d'eau (Mifi, Noun) en saison des pluies. Les pentes y sont fortes (15 à 30\%), les sols sont des \textbf{sols ferrallitiques lessivés}, et la densité de population élevée a entraîné une déforestation massive pour l'agriculture (maïs, haricot) et l'habitat. Les paysans labourent traditionnellement dans le sens de la pente maximale pour « faciliter l'écoulement de l'eau ».
    \textit{Identifiez et classez les facteurs favorisant l'érosion hydrique dans cette situation. Expliquez le mécanisme de formation des ravines.}
    \tcblower
    \textbf{Solution détaillée :}
    \begin{enumerate}
        \item \textbf{Facteurs naturels (physiques) :}
        \begin{itemize}
            \item \textbf{Relief :} Pentes fortes (15–30\%). La gravité accélère le ruissellement de surface, augmentant son énergie cinétique et son pouvoir érosif (transport de particules grosses).
            \item \textbf{Climat :} Pluviométrie élevée (1500–2500 mm/an) avec des orages violents en saison des pluies. L'intensité pluviométrique ($I$ en mm/h) dépasse souvent l'infiltration, générant un ruissellement hortonien.
            \item \textbf{Nature du sol :} \textbf{Sols ferrallitiques lessivés} (argileux, riches en oxydes de fer). Bien qu'agrégés, ils sont sensibles à la \emph{battance} (éclatement des agrégats sous l'impact des gouttes) et au colmatage de surface, réduisant l'infiltration.
        \end{itemize}
        \item \textbf{Facteurs anthropiques (humains) :}
        \begin{itemize}
            \item \textbf{Déforestation massive :} Destruction de la forêt naturelle pour l'agriculture et l'habitat. Perte de la litière et du système racinaire qui protégeaient le sol de l'impact direct des gouttes et maintenaient la porosité.
            \item \textbf{Pratiques culturales inadaptées :} Labour dans le sens de la pente (sillons parallèles à la pente). Cela crée des « autoroutes » pour l'eau, concentrant le ruissellement au fond des sillons au lieu de le freiner.
            \item \textbf{Surpâturage / Feux de brousse :} Maintien du sol nu sur de longues périodes.
        \end{itemize}
        \item \textbf{Mécanisme de formation des ravines (érosion linéaire) :}
        \begin{enumerate}
            \item \textbf{Battance :} Les gouttes de pluie frappent le sol nu $\rightarrow$ éclatement des agrégats $\rightarrow$ formation d'une croûte de battance (colmatage).
            \item \textbf{Ruissellement diffus :} L'eau ne s'infiltre plus, elle ruisselle en nappe mince sur la croûte, emportant les fines particules (érosion diffuse / interrill).
            \item \textbf{Concentration :} Le labour dans le sens de la pente ou les micro-dépressions naturelles concentrent l'eau en filets.
            \item \textbf{Incision :} La vitesse et le débit augmentent dans les filets $\rightarrow$ pouvoir érosif maximal $\rightarrow$ creusement de sillons (rill erosion) qui s'élargissent et s'approfondissent pour devenir des \textbf{ravines} (gully erosion) irréversibles par le labour simple.
        \end{enumerate}
    \end{enumerate}
    \textbf{Conclusion :} L'érosion observée résulte de la conjugaison de facteurs naturels défavorables (pente, pluies intenses, sols sensibles) et surtout de pratiques humaines destructrices (déforestation, labour en pente) qui ont supprimé les freins naturels (végétation, rugosité du sol).
\end{exoresolu}

\begin{methode}[Proposer des techniques adaptées de lutte contre l'érosion]
    \textbf{Objectif :} Choisir et justifier des techniques de conservation des sols (CES) adaptées au type d'érosion, au contexte biophysique et socio-économique.
    \begin{enumerate}
        \item \textbf{Diagnostiquer le type d'érosion :} Hydrique (diffuse, linéaire/ravines, massifs/glissements) ou Éolienne (déflation, ensablement).
        \item \textbf{Analyser le contexte :} Pente (faible < 5\%, moyenne 5–15\%, forte > 15\%), type de sol, pluviométrie, disponibilité en main-d'œuvre/outils, acceptabilité paysanne.
        \item \textbf{Sélectionner les techniques par catégorie (approche intégrée) :}
        \begin{itemize}
            \item \textbf{Agronomiques / Végétales (base) :} Couverture permanente (cultures de couverture, paillage), agroforesterie (haies vives, arbres fertilitaires), rotation culturale, cultures associées. \emph{But :} Protéger le sol de l'impact des gouttes, augmenter l'infiltration, fixer le sol par les racines.
            \item \textbf{Mécaniques / Structurelles (appui sur forte pente) :} \textbf{Courbes de niveau} (labour, semis, billons perpendiculaires à la pente), \textbf{Fascines} (branches enterrées), \textbf{Bancs / Terrasses} (bench terraces, fanya juu/fanya chini), \textbf{Digues de pierre} (cordons pierreux), \textbf{Bassins de rétention}. \emph{But :} Ralentir l'eau, la faire infiltrer, réduire la longueur de pente effective.
            \item \textbf{Forestières :} Reboisement (espèces locales à croissance rapide, racines pivotantes), enrésinement, protection des sources et berges (végétation ripicole).
        \end{itemize}
        \item \textbf{Justifier le choix :} Lier la technique au facteur combattu (ex: « Courbes de niveau pour réduire la vitesse du ruissellement sur pente moyenne » ; « Haies vives de \emph{Calliandra} ou \emph{Leucaena} pour fixer les berges et fournir du fourrage »).
        \item \textbf{Préciser l'entretien :} Une technique non entretenue (terrasse effondrée, haie non taillée) aggrave l'érosion (effet barrage puis rupture brutale).
    \end{enumerate}
    \textbf{Réflexe Probatoire :} Ne pas lister toutes les techniques. En choisir 2 à 3 \textbf{pertinentes} pour le cas précis et les \textbf{justifier scientifiquement}. Citer des espèces locales (ex: \emph{Acacia}, \emph{Leucaena}, \emph{Vetiver}, \emph{Brachiaria}).
\end{methode}

\begin{exoresolu}[Proposition de techniques pour un bassin versant dégradé au Nord-Cameroun]
    \textbf{Énoncé :} Dans le département du Mayo-Rey (Nord-Cameroun), un bassin versant de 500 ha présente une érosion hydrique sévère (ravines profondes, sols dénudés, turbidité de la Bénoué) et une érosion éolienne (vents d'harmattan violents, dunes mobiles menaçant les champs). Le sol est sableux (\textbf{sol ferrugineux tropical lessivé}), la pente moyenne (5–10\%), pluviométrie 900 mm/an (saison unique courte). La population pratique l'agriculture extensive (mil, sorgho, coton) avec brûlis et labour au tracteur dans le sens de la pente. Proposez un plan de lutte intégré (3 techniques minimum) en justifiant chaque choix.
    \tcblower
    \textbf{Solution détaillée :}
    \textbf{Diagnostic :} Érosion hydrique (ravines) due au labour en pente + sol sableux sensible + pluies intenses courtes. Érosion éolienne (déflation) due au sol nu en saison sèche + vents forts + texture sableuse. Contexte : main-d'œuvre familiale, faible capital, besoin de fourrage/bois.
    \vspace{0,2cm}
    \textbf{Technique 1 : Culture en courbes de niveau + Billons attachés (Fanya Juu / Fanya Chini) — \emph{Lutte hydrique principale}.}
    \begin{itemize}
        \item \textbf{Principe :} Labour, semis et formation de billons (buttes de terre) \textbf{perpendiculairement à la pente} (suivant les courbes de niveau). Tous les 10–20 m, on creuse une tranchée (fanya chini) et on rejette la terre en amont pour former une butte (fanya juu) plantée d'herbe.
        \item \textbf{Justification :} Sur pente moyenne (5–10\%), cela réduit la \textbf{longueur de pente effective} et la \textbf{vitesse du ruissellement} ($v \propto \sqrt{\text{pente}}$). L'eau stagne dans les tranchées $\rightarrow$ \textbf{infiltration maximale} (recharge nappe) $\rightarrow$ arrêt de l'érosion diffuse et prévention de l'initiation de nouvelles ravines. Adapté au labour tracteur (marquage au niveau à bulle ou cadre en A).
    \end{itemize}
    \textbf{Technique 2 : Haies vives pérennes en courbes de niveau (Agroforesterie) — \emph{Lutte hydrique + éolienne + fertilité}.}
    \begin{itemize}
        \item \textbf{Espèces :} \emph{Leucaena leucocephala}, \emph{Calliandra calothyrsus}, \emph{Gliricidia sepium} (légumineuses fixatrices d'azote, croissance rapide, taillables) + \emph{Vetiveria zizanioides} (Vetiver) en bas de pente/berges (racines pivotantes > 3 m, touffe dense).
        \item \textbf{Justification :} Double rôle. (1) \textbf{Barrière physique :} Ralentit ruissellement (hydrique) et vent (éolien) au niveau du sol. (2) \textbf{Amélioration sol :} Apport de matière organique (feuilles), fixation azote $\rightarrow$ fertilité $\rightarrow$ meilleure couverture culturale future. (3) \textbf{Produits :} Fourrage (saison sèche), bois de feu, piquets. \emph{Vetiver} indispensable pour colmater les têtes de ravines existantes.
    \end{itemize}
    \textbf{Technique 3 : Paillage / Résidus de culture + Zéro labour (Semis direct sous couverture végétale - SDCV) — \emph{Lutte éolienne + hydrique diffuse + fertilité}.}
    \begin{itemize}
        \item \textbf{Principe :} Ne pas brûler les résidus (tiges mil/sorgho, coton). Les laisser au sol ou importer paille/herbe. Semer directement dans le mulch (semoir direct ou pochettes).
        \item \textbf{Justification :} (1) \textbf{Couverture permanente :} Élimine la \emph{battance} (protection impact gouttes) et le \textbf{vent au ras du sol} (vitesse nulle à la surface du mulch $\rightarrow$ arrêt déflation). (2) \textbf{Économie d'eau :} Réduit évaporation (crucial 900 mm/an). (3) \textbf{Vie du sol :} Nourrit vers/termites $\rightarrow$ porosité $\rightarrow$ infiltration. (4) \textbf{Coût faible :} Pas de terrassement lourd, utilise résidus locaux.
    \end{itemize}
    \textbf{Complément indispensable :} \textbf{Protection des berges de la Bénoué} par reboisement en \emph{Acacia nilotica}, \emph{Mitragyna inermis} (espèces ripicoles locales) sur 20–50 m de large pour filtrer les apports érosifs et stabiliser la rive.
    \vspace{0,2cm}
    \textbf{Conclusion :} L'association \textbf{Courbes de niveau (structure) + Haies vives/Vetiver (végétal pérenne) + Paillage/SDCV (agronomique)} forme un système résilient, peu coûteux, productif (fourrage, bois, rendements) et adapté au contexte sahélien.
\end{exoresolu}

\begin{methode}[Interpréter un schéma ou une photo montrant un phénomène érosif]
    \textbf{Objectif :} Décrire rigoureusement un document iconographique (photo terrain, coupe schématique, série de photos temporelles) pour en extraire le processus érosif et son stade.
    \begin{enumerate}
        \item \textbf{Identifier le document :} Nature (photo aérienne, photo au sol, coupe schématique, carte), échelle approximative, localisation si connue.
        \item \textbf{Décrire objectivement (ce qu'on VOIT) :}
        \begin{itemize}
            \item \textbf{Relief :} Pente (forte/douce), formes (ravines en V, en U, dunes, plaques de déflation, cônes de déjection).
            \item \textbf{Sol/Surface :} Nu ? Végétalisé ? Croûte de battance visible (brillante, lisse) ? Pierres en surface (pavement éolien/hydrique) ? Racines apparentes (dénudation) ?
            \item \textbf{Eau/Vent :} Ruissellement visible (filets, couleur trouble/chargée) ? Traces de vent (dunes, ripples, dépôts au pied d'obstacles) ?
            \item \textbf{Végétation :} Type (forêt, culture, jachère, nue), densité, alignement (haies, courbes de niveau ?).
            \item \textbf{Aménagements :} Terrasses, digues, fascines, barrages ? État (fonctionnels, rompus, colmatés) ?
        \end{itemize}
        \item \textbf{Interpréter (ce qu'on COMPREND) :}
        \begin{itemize}
            \item \textbf{Type d'érosion :} Hydrique (diffuse : plaques lisses ; linéaire : ravines/gouilles ; massifs : glissements) ou Éolienne (déflation : surface déflée, racines nues ; accumulation : dunes, nébkhas).
            \item \textbf{Stade :} Initial (battance, croûte), Avancé (ravines actives, têtes de ravines en recul), Terminale (sol squelette, rocher affleurant, badlands) ou Stabilisée (végétation recolonise, terrasses entretenues).
            \item \textbf{Dynamique :} Active (eau trouble, parois verticales nues, éboulis frais) ou Fossile/Stabilisée (parois végétalisées, fond plat colmaté).
        \end{itemize}
        \item \textbf{Conclure / Proposer :} Résumer le processus dominant et suggérer une action prioritaire visible ou manquante sur le document.
    \end{enumerate}
    \textbf{Vocabulaire précis à utiliser :} \emph{Battance, croûte de battance, ruissellement diffus/linéaire, ravinement, régression rétrograde (tête de ravine), déflation, pavement, nébkha, cône de déjection, colmatage.}
\end{methode}

\begin{exoresolu}[Interprétation d'une coupe schématique d'une ravine en évolution]
    \textbf{Énoncé :} Le document ci-dessous représente une coupe transversale d'une ravine active dans une zone de sols ferrugineux tropicaux lessivés (pente 12\%).
    \begin{popfigure}
    \begin{tikzpicture}[scale=0.8]
        % Sol initial
        \fill[poporange!30] (-5,0) -- (5,0) -- (5,-0.5) -- (-5,-0.5) -- cycle;
        \draw[thick] (-5,0) -- (5,0);
        % Ravine
        \fill[white] (-1,0) .. controls (-1,-1.5) and (1,-1.5) .. (1,0) -- cycle;
        \draw[thick, popdark] (-1,0) .. controls (-1,-1.5) and (1,-1.5) .. (1,0);
        % Eau au fond
        \fill[popblue!50] (-0.5,-1.3) .. controls (-0.5,-1.4) and (0.5,-1.4) .. (0.5,-1.3) -- cycle;
        \draw[popblue] (-0.5,-1.3) .. controls (-0.5,-1.4) and (0.5,-1.4) .. (0.5,-1.3);
        % Flèches ruissellement
        \draw[->, thick, popblue] (-4,0.2) -- (-4,-0.5) node[right] {Ruissellement diffus};
        \draw[->, thick, popblue] (4,0.2) -- (4,-0.5) node[left] {Ruissellement diffus};
        \draw[->, thick, popblue] (0,-0.2) -- (0,-1.2) node[right] {Concentration};
        % Racines apparentes
        \draw[thin, popgreen!50!popdark] (-1.2, -0.1) -- (-1.3, -0.6);
        \draw[thin, popgreen!50!popdark] (1.2, -0.1) -- (1.3, -0.6);
        % Légendes internes
        \node at (-3, -0.2) {Sol nu, croûte de battance};
        \node at (3, -0.2) {Sol nu, croûte de battance};
        \node at (0, -1.6) {Fond ravin: eau trouble, graviers};
        \node at (0, 0.5) {Surface culturale};
    \end{tikzpicture}
    \legende{Coupe schématique d'une ravine active en forme de V dans un sol ferrugineux tropical lessivé (pente 12\%). On note la convergence du ruissellement diffus vers le thalweg, la croûte de battance en surface, les racines apparentes sur les bords et l'eau trouble chargée en sédiments au fond.}
    \end{popfigure}
    \textit{1. Décrivez les éléments visibles sur la coupe. 2. Identifiez le type et le stade d'érosion. 3. Expliquez le rôle de la croûte de battance et de la pente dans la genèse de cette forme. 4. Proposez une technique pour stabiliser cette ravine.}
    \tcblower
    \textbf{Solution détaillée :}
    \begin{enumerate}
        \item \textbf{Description objective :}
        \begin{itemize}
            \item \textbf{Surface :} Horizon supérieur (0 à -0,5 m) uniforme, colorié en orange (sol ferrugineux), sans végétation visible (« Sol nu »), mention « croûte de battance ».
            \item \textbf{Forme centrale :} Une incision symétrique en « V » large (environ 2 m au sommet) et profonde (1,5 m), aux parois nues, verticales à sub-verticales.
            \item \textbf{Fond :} Présence d'eau bleue trouble (chargée en sédiments) et de graviers/galets.
            \item   \textbf{Flèches bleues :} Convergence du ruissellement depuis les bords vers le centre (concentration).
            \item \textbf{Détails :} Racines d'arbres/herbes sectionnées et apparentes sur les bords supérieurs de la ravine (témoins du niveau initial du sol).
        \end{itemize}
        \item \textbf{Identification :}
        \begin{itemize}
            \item \textbf{Type :} \textbf{Érosion hydrique linéaire (ravinement / \emph{gully erosion})}.
            \item \textbf{Stade :} \textbf{Actif / En expansion.} Preuves : parois nues et verticales (instables), eau trouble (transport solide actif), racines apparentes (recul récent des bords), absence de végétation colonisatrice au fond.
        \end{itemize}
        \item \textbf{Explication génétique (Rôle battance + pente) :}
        \begin{enumerate}
            \item \textbf{Battance :} Sol nu + gouttes de pluie violentes $\rightarrow$ éclatement agrégats $\rightarrow$ formation \textbf{croûte de battance} (horizon de surface colmaté, porosité quasi nulle).
            \item \textbf{Imperméabilisation :} Croûte $\rightarrow$ \textbf{infiltration quasi nulle} $\rightarrow$ \textbf{ruissellement hortonien} immédiat et massif sur toute la surface (flèches latérales).
            \item \textbf{Pente (12\%) + Concentration :} La pente confère de l'énergie cinétique. Le ruissellement diffus converge vers la ligne de plus grande pente (thalweg naturel ou sillon de labour) $\rightarrow$ \textbf{concentration hydraulique}.
            \item \textbf{Incision :} Le filet d'eau concentré a un pouvoir érosif (force de traction) bien supérieur au ruissellement diffus. Il entame le sol sous la croûte (horizon moins structuré) $\rightarrow$ creusement rapide en « V » par érosion régressive (tête de ravine qui remonte) et élargissement par éboulements des parois.
        \end{enumerate}
        \item \textbf{Technique de stabilisation prioritaire :}
        \textbf{Barrages de ralentissement (en pierres sèches ou gabions) + Revégétalisation par Vetiver.}
        \begin{itemize}
            \item \textbf{Barrages :} Tous les 5–10 m de dénivelé dans le fond. But : \textbf{casser la pente longitudinale} du fond de ravine $\rightarrow$ réduire la vitesse de l'eau $\rightarrow$ provoquer la \textbf{dépôt des alluvions} (colmatage progressif du fond) $\rightarrow$ relever le niveau de base local.
            \item \textbf{Vetiver (Chrysopogon zizanioides) :} Planté en lignes denses sur les barrages et sur les parois (boutures). Racines pivotantes > 3 m $\rightarrow$ \textbf{ancrage mécanique} des parois + \textbf{filtration} de l'eau qui s'infiltre par les parois (évite la pression interstitielle cause d'éboulements).
            \item \textbf{Amont :} Obligatoire : Courbes de niveau + Haies vives sur le bassin versant amont pour \textbf{réduire l'apport liquide et solide} à la source. Sans cela, les barrages se colmatent trop vite ou sont contournés.
        \end{itemize}
    \end{enumerate}
    \textbf{Conclusion :} Le document illustre le passage de l'érosion diffuse (battance) à l'érosion linéaire concentrée (ravine) sous l'effet combiné de la dégradation de la surface du sol (croûte) et du relief. La lutte curative (barrages + Vetiver) doit impérativement s'accompagner de la lutte préventive en amont (CES agronomiques).
\end{exoresolu}

\begin{methode}[Analyser l'efficacité d'une technique de lutte (Courbes de niveau / Terrasses) à partir de données expérimentales]
    \textbf{Objectif :} Interpréter un tableau de résultats (pertes en sol, pertes en eau, rendement) comparant une parcelle témoin (labour pente) et une parcelle traitée (courbes de niveau, terrasses, haies).
    \begin{enumerate}
        \item \textbf{Lire le tableau :} Identifier les modalités (Témoin vs Traité), les variables (Perte en sol t/ha/an, Perte en eau mm/an, Rendement kg/ha), la durée de l'essai (années), le contexte (pente, sol, culture).
        \item \textbf{Calculer les indicateurs d'efficacité :}
        \begin{itemize}
            \item \textbf{Taux de réduction des pertes en sol (TRS) :} $TRS(\%) = \frac{Perte_{Témoin} - Perte_{Traité}}{Perte_{Témoin}} \times 100$.
            \item \textbf{Taux de réduction du ruissellement (TRR) :} $TRR(\%) = \frac{Ruiss_{Témoin} - Ruiss_{Traité}}{Ruiss_{Témoin}} \times 100$.
            \item \textbf{Gain de rendement :} $\Delta R = Rend_{Traité} - Rend_{Témoin}$ (ou $\%$).
        \end{itemize}
        \item \textbf{Interpréter physiquement :} Lier la réduction du ruissellement (infiltration accrue) à la réduction de l'érosion (moins d'énergie de transport) et au gain de rendement (plus d'eau + nutriments retenus).
        \item   \textbf{Critiquer la méthode :} Durée de l'essai (effet nouveauté vs durabilité), entretien des ouvrages, charge de travail, coût/bénéfice.
    \end{enumerate}
    \textbf{Formules clés :} $Ruissellement (\%) = \frac{Hauteur\ d'eau\ ruissellée}{Hauteur\ de\ pluie\ totale} \times 100$. $Perte\ en\ sol = Concentration\ (g/L) \times Volume\ ruissellement\ (L/ha)$.
\end{methode}

\begin{exoresolu}[Évaluation de l'efficacité de billons en courbes de niveau à Dschang]
    \textbf{Énoncé :} Un essai de 3 ans (2018–2020) a été conduit à Dschang (Ouest, pente 18\%, sol ferrallitique, maïs) comparant deux modalités : Témoin (Labour motorisé dans le sens de la pente) et Traité (Billons en courbes de niveau + paillage paille de maïs 3 t/ha). Résultats moyens annuels :
    \begin{center}
    \begin{tabular}{|l|c|c|c|}\hline
    \textbf{Paramètre} & \textbf{Témoin (Pente)} & \textbf{Traité (Courbes de niveau + Paillis)} & \textbf{Unité} \\ \hline
    Pluie totale & 1850 & 1850 & mm/an \\ \hline
    Ruissellement & 420 & 95 & mm/an \\ \hline
    Perte en sol & 42 & 3,8 & t/ha/an \\ \hline
    Rendement grain & 2,1 & 3,8 & t/ha \\ \hline
    \end{tabular}
    \end{center}
    \textit{1. Calculez le pourcentage de ruissellement par rapport à la pluie pour chaque modalité. 2. Calculez les taux de réduction du ruissellement (TRR) et des pertes en sol (TRS). 3. Interprétez l'impact sur le rendement. 4. Ce dispositif est-il durable sans entretien ? Justifiez.}
    \tcblower
    \textbf{Solution détaillée :}
    \begin{enumerate}
        \item \textbf{Pourcentage de ruissellement (Coefficient de ruissellement $C_r$) :}
        \begin{align*}
            C_{r,Témoin} &= \frac{420}{1850} \times 100 = \mathbf{22,7\%} \\
            C_{r,Traité} &= \frac{95}{1850} \times 100 = \mathbf{5,1\%}
        \end{align*}
        \item \textbf{Taux de réduction :}
        \begin{align*}
            TRR &= \frac{420 - 95}{420} \times 100 = \frac{325}{420} \times 100 = \mathbf{77,4\%} \\
            TRS &= \frac{42 - 3,8}{42} \times 100 = \frac{38,2}{42} \times 100 = \mathbf{91,0\%}
        \end{align*}
        \item \textbf{Interprétation :}
        \begin{itemize}
            \item \textbf{Hydrologie :} Le ruissellement passe de 22,7\% à 5,1\% de la pluie. Les billons + paillis créent une \textbf{rugosité de surface} et un \textbf{stockage temporaire} (micro-bassins) qui favorisent l'infiltration (ruissellement ralenti, temps de contact sol/eau augmenté). L'eau stockée est disponible pour la culture.
            \item \textbf{Pédologie :} L'érosion chute de 91\% (42 $\rightarrow$ 3,8 t/ha/an). La valeur 3,8 t/ha/an reste au-dessus du seuil de tolérance géologique souvent cité (~1–2 t/ha/an pour sols tropicaux profonds), mais représente une amélioration majeure. Le paillis est clé : il absorbe l'énergie cinétique des gouttes (zéro battance) et freine le flux laminaire.
            \item \textbf{Agronomie :} Rendement +81\% (2,1 $\rightarrow$ 3,8 t/ha). Gain dû à : (1) \textbf{Eau} : +325 mm infiltrés (disponibles pour la plante, critique en fin de cycle). (2) \textbf{Nutriments} : Matière organique et éléments minéraux (N, P, K) retenus dans le champ au lieu d'être exportés (42 t de sol perdues = énorme perte de fertilité). (3) \textbf{Structure} : Paillis $\rightarrow$ activité biologique $\rightarrow$ porosité.
        \end{itemize}
        \item \textbf{Durabilité et entretien :} \textbf{NON, pas durable sans entretien rigoureux.}
        \begin{itemize}
            \item \textbf{Billons :} S'affaissent par tassement (pluie, passage homme/machine), érosion latérale, bioturbation. Doivent être \textbf{reconstitués chaque campagne} (ou chaque saison si double culture).
            \item \textbf{Paillis :} Se décompose (termites, champignons) ou est emporté/brûlé. Doit être \textbf{renouvelé chaque cycle cultural} (3 t/ha/an). Si pénurie de résidus $\rightarrow$ échec.
            \item \textbf{Risque majeur :} Si un billon casse (surverse) sur une longue distance, l'eau se concentre dans la brèche $\rightarrow$ \textbf{initiation d'une ravine catastrophique} (effet « barrage qui cède »). L'entretien préventif (rebourrage des brèches) est vital.
        \end{itemize}
    \end{enumerate}
    \textbf{Conclusion :} La technique est \textbf{très efficace} (TRS 91\%, TRR 77\%, rendement quasi doublé) mais \textbf{exigeante en main-d'œuvre et biomasse}. Son adoption paysanne dépend de la disponibilité de paille (concurrence fourrage/paillis) et de la maîtrise du tracé des courbes de niveau (cadre en A, niveau à bulle).
\end{exoresolu}

\begin{methode}[Différencier érosion hydrique et éolienne et proposer la lutte adaptée au Sahel]
    \textbf{Objectif :} Distinguer les processus, les formes et les techniques spécifiques à l'érosion éolienne (déflation/ensablement) vs hydrique, dans le contexte sahélien/soudano-sahélien (Nord Cameroun).
    \begin{enumerate}
        \item \textbf{Identifier l'agent moteur :} Eau (pluie, ruissellement) vs Vent (harmattan, orages secs).
        \item \textbf{Reconnaître les formes types :}
        \begin{itemize}
            \item \textbf{Hydrique :} Battance (croûte), érosion diffuse (plaques lisses), linéaire (ravines, gouilles), massifs (glissements), dépôts (cônes de déjection, alluvions banc de sable dans lit mineur).
            \item \textbf{Éolienne :} \textbf{Déflation} (creusement : cuvettes de déflation, racines nues, pavement éolien = graviers résiduels), \textbf{Transport} (saltation, suspension), \textbf{Accumulation} (dunes mobiles, barkhanes, nébkhas = dunes fixées sur végétation, ensablement champs/maisons/pistes).
        \end{itemize}
        \item \textbf{Facteurs déclenchants spécifiques :} Sol nu + texture sableuse/limoneuse + vent > seuil de fluidité (~5–6 m/s au ras du sol) + absence de barrières.
        \item \textbf{Techniques spécifiques anti-éoliennes (en + des CES hydriques) :}
        \begin{itemize}
            \item \textbf{Brise-vent / Haies brise-vent :} Périmétriques (pourtour champ) ou intra-parcellaires (espacement 10–20 H, H = hauteur haie). Espèces : \emph{Azadirachta indica} (Neem), \emph{Acacia nilotica}, \emph{Balanites aegyptiaca}, \emph{Prosopis juliflora} (attention invasive). Réduisent vitesse vent au sol $\rightarrow$ dépôt particules.
            \item \textbf{Fixation des dunes mobiles :} \textbf{Palissage} (branches, paillis en quadrillage 1x1m ou losanges) pour briser vitesse vent au sol + \textbf{Plantation fixatrice} (\emph{Leptadenia pyrotechnica}, \emph{Calotropis procera}, \emph{Euphorbia balsamifera}, \emph{Vetiver} en bas de dune).
            \item \textbf{Couverture permanente (SDCV/Paillage) :} Seule technique qui lutte \textbf{simultanément} contre battance (hydrique) et déflation (éolienne) en maintenant une rugosité et une couverture au ras du sol.
        \end{itemize}
    \end{enumerate}
    \textbf{Piège :} Ne pas confondre « pavement éolien » (graviers résiduels après déflation des fines) et « croûte de battance » (colmatage par les fines). Le pavement *protège* le sol sous-jacent ; la croûte *aggrave* le ruissellement.
\end{methode}

\begin{exoresolu}[Cas complexe : Érosion éolienne et ensablement dans le Logone-et-Chari]
    \textbf{Énoncé :} Dans le département du Logone-et-Chari (Extrême-Nord, zone sahélienne, pluviométrie ~500 mm/an, vents d'harmattan violents Nov–Mars), des dunes mobiles (barkhanes) progressent vers le Sud-Est à ~15 m/an, ensablant des champs de mil, obstruant la piste Kousseri-Waza et menaçant le village de Blangoua. Le sol est sableux (dunes fixées anciennes remobilisées). La jachère est quasi inexistante (pression foncière), le pâturage intensif. Proposez une stratégie de lutte intégrée en 3 phases (Urgence, Stabilisation, Durabilité) en justifiant les espèces et techniques.
    \tcblower
    \textbf{Solution détaillée :}
    \textbf{Diagnostic :} Érosion éolienne dominante (déflation source dunes + transport + accumulation/ensablement). Facteurs : Vent fort (Harmattan), Sol sableux nu (pas de végétation, surpâturage, culture extensive), Pente faible mais fetch (longueur de parcours vent) immense. Absence de brise-vent.
    \vspace{0,2cm}
    \textbf{Phase 1 : URGENCE — Protection immédiate des biens/village (Mois 0–6).}
    \begin{itemize}
        \item \textbf{Barrières mécaniques provisoires :} Pose de \textbf{palissades en branches (palissage)} ou filets géotextiles/bâches perpendiculaires au vent dominant (N-E $\rightarrow$ S-O) en amont immédiat du village et de la piste. Espacement : 5–10 m. But : Créer une rugosité artificielle $\rightarrow$ dépôt sable *avant* le village/piste (créer un « piège à sable » contrôlé).
        \item \textbf{Dégagement mécanique :} Enlèvement régulier du sable sur la piste (chargeuses) + dépotage en zone de sacrifice en aval.
    \end{itemize}
    \textbf{Phase 2 : STABILISATION — Fixation des dunes mobiles sources (Année 1–3).}
    \begin{itemize}
        \item \textbf{Quadrillage (Palissage définitif) :} Sur les dunes mobiles actives (source), pose de branches (palmiers, acacias morts) ou paillis en \textbf{losanges ou carrés de 1 à 2 m de côté}, enfoncés de 30 cm. Hauteur 30–40 cm. But : Réduire vitesse vent au sol à < seuil de transport (~0,2 m/s) sur toute la surface de la dune.
        \item \textbf{Revégétalisation pionnière (Fixation biologique) :} Semis/Plantation \textbf{dans les alvéoles du quadrillage} d'espèces psammophiles (aimant le sable), rustiques, non appétées (protégées du pâturage) :
        \begin{itemize}
            \item \emph{Leptadenia pyrotechnica} (Nguiguiss) : Racines adventives sur tiges couchées, très résistant sécheresse/pâturage.
            \item \emph{Calotropis procera} (Pommier de Sodome) : Croissance rapide, graines anémochores, latex toxique (anti-pâturage).
            \item \emph{Euphorbia balsamifera} : Arbuste succulent, excellente fixation.
            \item \emph{Vetiveria zizanioides} (Vetiver) : En bas de pente des dunes / pieds de dune (racines profondes, touffe dense barrière physique).
        \end{itemize}
        \item \textbf{Interdiction stricte du pâturage / coupe} sur les dunes en fixation (clôture sociale ou physique).
    \end{itemize}
    \textbf{Phase 3 : DURABILITÉ — Aménagement du territoire \& Agroforesterie (Année 3+).}
    \begin{itemize}
        \item \textbf{Réseau de haies brise-vent pérennes :} Implantation de \textbf{haies vives multi-étagées} (Arbres + Arbustes + Herbes) perpendiculaires au vent dominant (NE-SO), espacées de \textbf{10 à 15 fois la hauteur adulte} (ex: Haie 5 m $\rightarrow$ espacement 50–75 m).
        \item \textbf{Espèces haies :} \emph{Azadirachta indica} (Neem, croissance rapide, ombrage, bois, médicinal), \emph{Acacia nilotica} (Gommier, gomme arabique, fourrage gousses, azote), \emph{Balanites aegyptiaca} (Dattier du désert, fruit comestible, huile, très rustique), \emph{Ziziphus mauritiana} (Jujubier, fruit, fourrage).
        \item \textbf{Intégration culturale (Systèmes Agroforestiers) :} \textbf{Cultures en couloirs (Alley cropping)} entre les haies. Mil / Sorgho / Niébé en courbes de niveau (si micro-pente) ou billons attachés. \textbf{Paillage obligatoire} résidus mil/sorgho (ne pas brûler !).
        \item \textbf{Gestion pastorale :} Définition de couloirs de transhumance, points d'eau éloignés des zones de fixation, pâturage tournants, valorisation fourrage haies (tailles).
        \item \textbf{Gouvernance :} Convention locale (COGER) pour gestion commune ressources (bois, pâture, eau), sanctions feux de brousse / coupe abusive.
    \end{itemize}
    \textbf{Conclusion :} La lutte anti-éolienne au Sahel ne peut être que \textbf{systémique} : fixation physique d'urgence $\rightarrow$ fixation biologique des sources (dunes) $\rightarrow$ aménagement durable de l'espace par haies brise-vent + cultures couvertes + gouvernance foncière/pastorale. Le \textbf{Vetiver} et le \textbf{Neem} sont les « piliers » techniques ; la \textbf{convention villageoise} est le pilier social.
\end{exoresolu}

\begin{methode}[Réaliser un diagnostic d'érosion et proposer un plan d'action (Étude de cas complète)]
    \textbf{Objectif :} Synthétiser l'ensemble des compétences : décrire un site, identifier types/causes, proposer techniques priorisées, planifier.
    \begin{enumerate}
        \item \textbf{Description du site (Fiche de terrain) :} Localisation (région, bassin versant), Climat (pluie, vent), Relief (pente \%), Sol (type, profondeur, texture), Végétation (couverture, type), Occupation du sol (cultures, jachère, habitat), Hydrographie.
        \item \textbf{Constat d'érosion (Preuves) :} Types observés (hydrique/éolienne), Formes (ravines, dunes, croûtes, racines nues), Étendue (% bassin touché), Activité (actif/fossile), Dégâts (champs perdus, infrastructures menacées, eau trouble).
        \item \textbf{Analyse causale (Arbre des causes) :} Causes profondes (démographie, tenure foncière, pauvreté, politiques) $\rightarrow$ Causes directes (déforestation, labour pente, surpâturage, feux) $\rightarrow$ Facteurs aggravants naturels (pente, pluies intenses, sols fragiles).
        \item \textbf{Plan d'action priorisé (Logique bassin versant) :}
        \begin{enumerate}
            \item \textbf{Protection / Conservation (Haut de bassin / Sources) :} Reboisement zones sources, têtes de ravines, pentes > 30\% (mise en défens). \emph{Techniques :} Reboisement, fascines, terrasses.
            \item \textbf{Aménagement / Production (Moyen de bassin / Terres arables) :} CES sur champs : Courbes de niveau, Haies vives (agroforesterie), SDCV/Paillis, Rotation/Associations. \emph{But :} « Produire plus, produire mieux, dégrader moins ».
            \item \textbf{Correction / Récupération (Bas de bassin / Ravines actives) :} Barrages (pierres, gabions, terre), Vetiver, Reboisement berges. \emph{Coûteux, dernier recours.}
            \item \textbf{Lutte éolienne (si zone concernée) :} Brise-vent, fixation dunes, palissage.
        \end{enumerate}
        \item \textbf{Indicateurs de suivi-évaluation :} Taux de couverture végétale, Perte en sol (t/ha/an), Niveau nappe / débit baseflow, Rendements, Revenu paysan, Taux d'adoption techniques.
    \end{enumerate}
    \textbf{Conseil rédaction :} Structurez la réponse en parties clairement titrées. Utilisez le vocabulaire technique précis. Justifiez chaque technique par le facteur qu'elle combat.
\end{methode}

\begin{exoresolu}[Plan de lutte intégré pour le bassin de la Mifi (Bafoussam)]
    \textbf{Énoncé :} Le bassin versant de la Mifi (affluent du Noun, Ouest Cameroun) alimente en eau la ville de Bafoussam. Il présente une dégradation avancée : ravines profondes en amont (zone de pâturage/feux), érosion diffuse sévère sur cultures (maïs, pomme de terre, labour pente), turbidité extrême de l'eau en saison des pluies (coût traitement eau élevé), colmatage du barrage de retenue. Pentes 10–25\%. Sols ferrallitiques lessivés. Population dense. Élaborer un plan d'action structuré en 3 zones (Haut, Moyen, Bas) avec techniques, espèces, acteurs.
    \tcblower
    \textbf{Solution détaillée :}
    \textbf{1. DIAGNOSTIC RAPIDE}
    \begin{itemize}
        \item \textbf{Types érosion :} Hydrique diffuse (battance, croûtes sur champs), Linéaire active (ravines profondes en zone pâturage/feux, régression rétrograde rapide), Massifs localisés (glissements sur pentes > 25\% argiles gonflantes).
        \item \textbf{Causes :} Déforestation (bois énergie, extension cultures), Feux de brousse annuels (pâturage), Labour motorisé systématique dans la pente (pomme de terre), Surpâturage (zones communes), Absence de CES.
        \item \textbf{Enjeux :} Ressource en eau potable (Bafoussam), Perte fertilité (sécurité alimentaire), Infrastructures (routes, barrage), Biodiversité.
    \end{itemize}
    \textbf{2. PLAN D'ACTION PAR ZONES (Logique « Haut $\rightarrow$ Bas »)}
    \begin{center}
    \begin{tabularx}{\linewidth}{|l|X|X|X|}\hline
    \rowcolor{popblueL}
    \textbf{ZONE} & \textbf{HAUT DE BASSIN (Sources, pentes > 20\%, pâturages, forêts résiduelles)} & \textbf{MOYEN DE BASSIN (Terres arables, pentes 10–20\%, forte densité pop.)} & \textbf{BAS DE BASSIN (Vallée, berges Mifi, barrage, zone urbaine périurbaine)} \\ \hline
    \textbf{Objectif Principal} & \textbf{PROTÉGER / INFILTRER} (Château d'eau) & \textbf{PRODUIRE DURABLEMENT / CONSERVER} (Cœur productif) & \textbf{CORRIGER / ÉPURER / PROTÉGER INFRA} \\ \hline
    \textbf{Techniques Prioritaires} &
    \begin{itemize}
        \item \textbf{Mise en défens} (Arrêt feux, pâturage contrôlé) zones sources \& têtes ravines.
        \item \textbf{Reboisement} espèces locales pionnières (\emph{Prunus africana}, \emph{Podocarpus}, \emph{Alnus}, \emph{Acacia} spp) + \emph{Vetiver} têtes ravines.
        \item \textbf{Fascines / Banquettes} en courbes de niveau pour stabiliser pentes fortes avant reboisement.
        \item \textbf{Barrages de ralentissement} (pierres sèches) dans ravines actives.
    \end{itemize} &
    \begin{itemize}
        \item \textbf{Courbes de niveau obligatoires} (traçage niveau à bulle/cadre A) pour tout labour.
        \item \textbf{Billons attachés (Fanya Juu/Chini)} + \textbf{Paillis} (résidus patate/maïs, herbes).
        \item \textbf{Haies vives agroforestières} en courbes de niveau (esp. 10–15 m) : \emph{Calliandra}, \emph{Leucaena}, \emph{Tephrosia}, \emph{Flemingia} (fourrage, azote, piquets, paillis).
        \item \textbf{Rotation / Association} : Maïs + Niébé / Soja ; Pomme de terre $\rightarrow$ Céréale + Légumineuse.
        \item \textbf{Agroforesterie fruitière} : Avocatier, Prunier africain (\emph{Dacryodes}), Safoutier (revenu + couverture pérenne).
    \end{itemize} &
    \begin{itemize}
        \item \textbf{Revégétalisation berges (50 m min)} : \emph{Alchornea cordifolia}, \emph{Raphia}, \emph{Alnus}, \emph{Vetiver} pied de berge. Filtration nitrates/pesticides/sédiments.
        \item \textbf{Barrages de fond / Seuils} (gabions, béton cyclopéen) dans grands collecteurs ravines pour protéger ponts/routes/prise eau.
        \item \textbf{Bassins de décantation / Zones humides tampons} en amont prise d'eau / barrage.
        \item \textbf{Assainissement urbain} : Gestion déchets, eaux usées (pollution point source).
    \end{itemize} \\ \hline
    \textbf{Acteurs Clés} & MINEPDED, MINFOF, Communes, Chefferies, ONG (WWF, CARE), Groupements pastoraux & MINADER, IRAD, Communes, OP (Coopératives, GIC), Projets (PADFA, ACEFA), Lycées agricoles & MINEE (Camwater), MINEPDED, Commune Bafoussam, Riverains, Entreprises BTP \\ \hline
    \end{tabularx}
    \end{center}
    \textbf{3. MESURES TRANSVERSALES (Conditions de réussite)}
    \begin{enumerate}
        \item \textbf{Gouvernance foncière / Convention locale :} Sécurisation tenure (certificats) pour inciter investissement long terme (haies, arbres). Règlementation feux / pâturage / coupes.
        \item \textbf{Appui-conseil / Formation :} Vulgarisation CES (écoles champs paysannes), formation traçage courbes de niveau (cadre en A), pépinières villageoises (Vetiver, arbres).
        \item \textbf{Incitations économiques :} Paiement pour Services Environnementaux (PSE) — Ville de Bafoussam / Camwater paie amont pour eau qualité/quantité. Crédit intrants (semences haies, outils) conditionné à adoption CES.
        \item \textbf{Suivi-Évaluation :} Postes de jaugeage (débit, turbidité), Parcelles témoins (pertes sol/eau), Imagerie drone/satellite (couverture végétale, extension ravines), Enquêtes socio-éco (revenus, adoption).
    \end{enumerate}
    \textbf{Conclusion :} Ce plan applique le principe de \textbf{gestion intégrée des ressources en eau et des sols (GIRES) à l'échelle du bassin versant}. La priorité absolue est la \textbf{protection des zones sources (Haut)} pour réduire l'apport liquide/solide à la source, suivie de l'\textbf{intensification écologique (Moyen)} pour concilier production et conservation, et enfin de la \textbf{protection curative ciblée (Bas)}. Sans convention sociale et incitations économiques, les techniques techniques seules échoueront.
\end{exoresolu}

\begin{attention}[Les 5 erreurs qui coûtent des points au Probatoire]
    \begin{enumerate}
        \item \textbf{Confondre « Érosion » et « Dégradation des sols » :} L'érosion est le \textbf{transport} de particules (par eau ou vent). La dégradation inclut aussi la perte de fertilité chimique (lessivage, minéralisation), la compaction, la salinisation, l'acidification \emph{sans} transport massif. \textit{Au Probatoire :} Si la question porte sur l'érosion, ne parlez pas seulement de « perte de fertilité », parlez de \textbf{perte physique du sol} (tonnes/ha/an).
        \item \textbf{Dire « Planter des arbres » sans précision :} « Reboisement » n'est pas une réponse suffisante. Il faut préciser : \textbf{Où ?} (haut de bassin, berges, courbes de niveau), \textbf{Quelles espèces ?} (locales, adaptées : \emph{Vetiver} pour ravines/berges, \emph{Calliandra/Leucaena} pour haies fourragères, \emph{Neem/Acacia} pour brise-vent Sahel, \emph{Prunus/Podocarpus} pour montagnards), \textbf{Comment ?} (densité, protection feu/pâturage, entretien). \textit{Erreur fatale :} Proposer \emph{Eucalyptus} ou \emph{Pinus} en zone de source d'eau (pompage nappe) ou en haie vive dans les champs (allélopathie, concurrence eau).
        \item \textbf{Ignorer l'entretien des ouvrages mécaniques :} Une terrasse, un billon, un barrage de pierres \textbf{sans entretien} devient un danger (rupture $\rightarrow$ ravine catastrophique). \textit{Toujours mentionner :} « À condition d'un entretien régulier (rebourrage, décolmatage, taille haies) ». C'est un critère d'évaluation de la durabilité.
        \item \textbf{Mélanger les techniques anti-hydriques et anti-éoliennes sans discernement :} En zone sahélienne (Nord), l'érosion éolienne domine en saison sèche. Proposer \textbf{seulement} des courbes de niveau (efficaces pour l'eau) sans \textbf{brise-vent, paillage permanent, fixation dunes} montre une incompréhension du milieu. \textit{Le paillage / SDCV est la SEULE technique efficace sur les DEUX fronts simultanément.}
        \item \textbf{Calculer le ruissellement ou la perte en sol sans unités ou sans formule :} Si un tableau donne Pluie (mm) et Ruissellement (mm), le coefficient de ruissellement $C_r = \frac{Ruiss}{Pluie} \times 100$ (\%). Si on vous donne Volume ruissellement (m$^3$) et Concentration (g/L), Perte sol (kg) = $Vol \times Conc$. \textit{Sanction :} Résultat brut sans unité (ex: « 0,22 » au lieu de « 22\% » ou « 22 \% ») ou formule non écrite = 0 point sur le calcul.
    \end{enumerate}
\end{attention}

\newpage

\coursec{Exercices}

\courssub{Je vérifie mes connaissances}

\begin{exercice}[QCM : Types et facteurs de l'érosion]
Parmi les affirmations suivantes, une seule est exacte. La recopier en la complétant.
\begin{enumerate}
    \item L'érosion éolienne est le principal type d'érosion observé dans la région du Littoral au Cameroun.
    \item La déforestation protège le sol contre l'érosion hydrique en laissant les racines en place.
    \item \textbf{L'érosion hydrique est favorisée par la pente, l'absence de couverture végétale et l'intensité des pluies.}
    \item Les terrasses de culture augmentent la vitesse de ruissellement sur les versants.
\end{enumerate}
\end{exercice}

\begin{exercice}[Vrai ou Faux ? Justifiez]
Dire si les affirmations suivantes sont vraies ou fausses. Justifier brièvement chaque réponse.
\begin{enumerate}
    \item Le labour dans le sens de la pente maximale réduit les risques d'érosion hydrique.
    \item Le reboisement avec des essences à croissance rapide (eucalyptus, acacia) est une technique efficace de lutte contre l'érosion.
    \item L'érosion n'a d'impact que sur la couche arable ; le sous-sol n'est jamais atteint.
    \item Les cultures en courbes de niveau consistent à labourer perpendiculairement à la pente pour créer des mini-barrages naturels.
\end{enumerate}
\end{exercice}

\begin{exercice}[Définitions et associations]
\begin{enumerate}
    \item Définir précisément les termes suivants : \textbf{érosion}, \textbf{ruissellement}, \textbf{ravine}, \textbf{déforestation}.
    \item Relier chaque type d'érosion (colonne A) à son agent principal (colonne B) et à un facteur favorisant majeur (colonne C).
    \begin{center}
    \begin{tabularx}{\linewidth}{|X|X|X|}\hline
    \textbf{A : Type d'érosion} & \textbf{B : Agent érosif} & \textbf{C : Facteur favorisant} \\ \hline
    Érosion hydrique & Vent & Pente forte + sol nu \\ \hline
    Érosion éolienne & Eau de ruissellement & Aridité + absence de haies brise-vent \\ \hline
    \end{tabularx}
    \end{center}
\end{enumerate}
\end{exercice}

\begin{exercice}[Calcul direct : Taux d'érosion]
Une parcelle agricole de \SI{2}{\hectare} située sur une pente de 12\% à Dschang a perdu une couche de terre moyenne de \SI{5}{\milli\metre} d'épaisseur après une saison des pluies particulièrement intense. La masse volumique apparente du sol est de \SI{1,3}{\tonne\per\cubic\metre}.
\begin{enumerate}
    \item Calculer le volume de terre perdu (en \si{\cubic\metre}).
    \item Calculer la masse de terre perdue (en tonnes).
    \item Exprimer le taux d'érosion en \si{t.ha^{-1}.an^{-1}}. Ce taux est-il tolérable selon les normes (seuil souvent admis < \SI{10}{t.ha^{-1}.an^{-1}}) ?
\end{enumerate}
\end{exercice}

\courssub{Je m'entraîne}

\begin{exercice}[Analyse de photographies : Ravines urbaines à Douala]
On observe deux photographies prises dans un quartier en pente de Douala (quartier \emph{Makepe} ou \emph{Ndogbong}) après une forte averse.
\begin{itemize}
    \item \textbf{Photo 1 :} Un ravin profond de 3 à 4 mètres coupe la route en terre battue. Les berges sont verticales et instables. Des déchets plastiques sont visibles au fond.
    \item \textbf{Photo 2 :} En bas du ravin, un dépôt de boue et de sable bouche un caniveau en béton. L'eau stagne autour.
\end{itemize}
\begin{enumerate}
    \item Décrire le phénomène érosif visible sur la Photo 1. Préciser le type d'érosion.
    \item Identifier \textbf{trois} causes probables (facteurs favorisants) de la formation de ce ravin en milieu urbain.
    \item Expliquer le lien entre le ravin (Photo 1) et le bouchage du caniveau (Photo 2).
    \item Proposer \textbf{deux} mesures d'urgence et \textbf{deux} mesures de prévention durable pour ce quartier.
\end{enumerate}
\end{exercice}

\begin{exercice}[Choix technique : Cas d'un agriculteur à Foumban]
Monsieur Njoya cultive du maïs et des haricots sur une parcelle de \SI{1,5}{\hectare} située sur un versant de \SI{15}{\%} de pente (environ 8,5°) dans la région de l'Ouest (sol ferrallitique, saison des pluies longue et intense). Il laboure actuellement dans le sens de la pente. Il hésite entre deux techniques d'aménagement : \textbf{les terrasses en gradins} (terrasses étroites soutenues par des murets en pierre) et \textbf{les cultures en courbes de niveau} (avec bandes enherbées).
\begin{enumerate}
    \item Rappeler le principe de chacune de ces deux techniques.
    \item Comparer leurs avantages et inconvénients dans le contexte de M. Njoya (coût de mise en œuvre, main-d'œuvre, efficacité sur sol ferrallitique, entretien).
    \item Formuler un conseil argumenté : quelle technique choisir ? Justifier en fonction de la pente, du type de sol, des cultures et des moyens du paysan.
\end{enumerate}
\end{exercice}

\begin{exercice}[Déforestation et inondations : Cas du Mont Cameroun]
Les pentes du Mont Cameroun subissent une déforestation rapide pour l'expansion des plantations (palmier à huile, bananier) et l'habitat. En aval, les villes de Limbé, Bota et Idenau subissent des inondations et coulées de boue récurrentes (ex. catastrophe de 2001 à Limbé).
\begin{enumerate}
    \item À l'aide d'un schéma annoté (coupe transversale du volcan : sommet $\rightarrow$ pied), expliquer le mécanisme par lequel la déforestation en amont aggrave les crues et coulées de boue en aval.
    \item Citer \textbf{trois} services écosystémiques rendus par la forêt naturelle du Mont Cameroun qui sont perdus lors de sa destruction.
    \item Proposer une stratégie de gestion intégrée du bassin versant (amont/aval) associant autorités locales, planteurs et populations riveraines.
\end{enumerate}
\end{exercice}

\begin{exercice}[Tableau de synthèse : Types, facteurs, conséquences, lutte]
Compléter le tableau ci-dessous en utilisant les mots-clés de la leçon. Une case peut contenir plusieurs éléments.
\begin{center}
\begin{tabularx}{\linewidth}{|>{\bfseries}l|X|X|X|}\hline
\textbf{Type d'érosion} & \textbf{Facteurs favorisants principaux} & \textbf{Conséquences majeures (sol / environnement)} & \textbf{Techniques de lutte adaptées} \\ \hline
Érosion hydrique en feuille (diffuse) & & & \\ \hline
Érosion hydrique concentrée (rigoles, ravines) & & & \\ \hline
Érosion éolienne (deflation) & & & \\ \hline
\end{tabularx}
\end{center}
\end{exercice}

\courssub{Je me prépare au Probatoire}

\begin{exercice}[Sujet de synthèse : Lutte contre l'érosion et développement durable de l'agriculture camerounaise]
\textbf{Contexte :} Le Cameroun perd annuellement des millions de tonnes de terre arable par érosion (hydrique au Sud et à l'Ouest, éolienne au Nord). Le Programme National de Lutte contre la Désertification (PAN/LCD) et les projets d'agriculture climato-intelligente (ACI) intègrent la conservation des sols.

\textbf{Consigne :} \textbf{Montrer que la lutte contre l'érosion est une condition indispensable du développement durable de l'agriculture camerounaise.}

Votre réponse structurée en trois parties distinctes comportera :
\begin{enumerate}
    \item \textbf{Les causes :} Analyser les causes naturelles et surtout anthropiques de l'érosion des sols agricoles au Cameroun (citer des exemples régionaux précis : Nord, Ouest, Littoral, Adamaoua).
    \item \textbf{Les conséquences :} Détailler les impacts sur la productivité agricole (rendements, sécurité alimentaire), sur l'environnement (colmatage barrages, pollution, biodiversité) et sur l'économie rurale (revenus, migration).
    \item \textbf{Les solutions :} Présenter une stratégie globale combinant techniques agronomiques (cultures en courbes de niveau, agroforesterie, couverts végétaux), techniques mécaniques (terrasses, fascines, barrages filtrants), et mesures institutionnelles (politiques foncières, vulgarisation, éducation). Illustrer par des exemples de réussite (ex. projet PADFA, agriculture de conservation dans le Nord).
\end{enumerate}
\end{exercice}

\begin{exercice}[Interprétation de schéma : Versant aménagé]
Le schéma ci-dessous représente une coupe transversale d'un versant aménagé dans la région de l'Ouest (type \emph{bamiléké} moderne). L'élève doit placer les légendes numérotées (1 à 7) sur le schéma et expliquer le rôle de chaque aménagement.

\begin{popfigure}
\begin{tikzpicture}[scale=0.9, every node/.style={font=\small}]
% Sol de base
\fill[poporange!30] (0,0) -- (14,0) -- (14,-0.5) -- (0,-0.5) -- cycle;
% Pente
\fill[poporange!40] (0,0) -- (14,2.5) -- (14,2) -- (0,-0.5) -- cycle;

% Terrasses (3 niveaux)
\draw[thick, popdark] (0,0) -- (14,0); % Base
\draw[thick, popdark] (0.5,0.8) -- (13.5,0.8); % Niveau 1
\draw[thick, popdark] (1,1.6) -- (13,1.6); % Niveau 2
\draw[thick, popdark] (1.5,2.4) -- (12.5,2.4); % Niveau 3 (sommet)

% Murets de soutènement (pierres)
\foreach \x in {0.5, 1, 1.5} {
    \fill[popmuted] (\x,0) rectangle (\x+0.1,0.8);
    \fill[popmuted] (\x+0.1,0.8) rectangle (\x+0.2,1.6);
    \fill[popmuted] (\x+0.2,1.6) rectangle (\x+0.3,2.4);
}

% Cultures (lignes horizontales sur terrasses)
\foreach \y in {0.4, 1.2, 2.0} {
    \foreach \x in {1, 2, 3, 4, 5, 6, 7, 8, 9, 10, 11, 12} {
        \draw[popgreen!70!popdark] (\x,\y-0.1) -- (\x,\y+0.1);
    }
}

% Haies vives (verticales sur les bords des terrasses)
\foreach \x in {1, 4, 7, 10, 13} {
    \draw[popgreen!50!popdark, thick] (\x,0) -- (\x,0.8);
    \draw[popgreen!50!popdark, thick] (\x+0.1,0.8) -- (\x+0.1,1.6);
    \draw[popgreen!50!popdark, thick] (\x+0.2,1.6) -- (\x+0.2,2.4);
}

% Zone reboisée (sommet)
\fill[popgreen!30] (1,2.4) -- (12.5,2.4) -- (13,3.5) -- (0.5,3.5) -- cycle;
% Arbres
\foreach \x in {2, 4, 6, 8, 10, 12} {
    \node[circle, fill=popgreen!60!popdark, minimum size=6pt, inner sep=0] at (\x, 3.0) {};
    \node[circle, fill=popgreen!40!popdark, minimum size=4pt, inner sep=0] at (\x+0.3, 2.8) {};
}

% Ruissellement (flèches bleues)
\draw[->, popblue, thick] (13.5, 2.5) -- (13.5, 2.0) node[midway, right] {Ruissellement ralenti};
\draw[->, popblue, thick] (13, 1.7) -- (13, 1.2);
\draw[->, popblue, thick] (12.5, 0.9) -- (12.5, 0.4);

% Points de repère pour légendes (cercles numérotés invisibles, l'élève les place)
% 1: Muret de soutènement (pied terrasse)
% 2: Haie vive (bord terrasse)
% 3: Culture en courbes de niveau (sur terrasse)
% 4: Fascine / Boudin végétal (bas de pente - non dessiné, à imaginer ou ajouter)
% 5: Zone reboisée (sommet)
% 6: Drain de pied de versant (bas)
% 7: Allée de circulation / chemin de ronde

% Marqueurs pour l'exercice (lettres A-G)
\node[circle, draw=popink, thick, fill=popink!20, minimum size=18pt, font=\bfseries] at (0.3, 0.4) {A};
\node[circle, draw=popink, thick, fill=popink!20, minimum size=18pt, font=\bfseries] at (13.2, 0.4) {B};
\node[circle, draw=popink, thick, fill=popink!20, minimum size=18pt, font=\bfseries] at (13.2, 1.2) {C};
\node[circle, draw=popink, thick, fill=popink!20, minimum size=18pt, font=\bfseries] at (0.7, 3.2) {D};
\node[circle, draw=popink, thick, fill=popink!20, minimum size=18pt, font=\bfseries] at (7, 0.4) {E};
\node[circle, draw=popink, thick, fill=popink!20, minimum size=18pt, font=\bfseries] at (7, 1.2) {F};
\node[circle, draw=popink, thick, fill=popink!20, minimum size=18pt, font=\bfseries] at (7, 2.0) {G};

\end{tikzpicture}
\legende{Coupe schématique d'un versant aménagé en terrasses (région de l'Ouest). Les lettres A à G indiquent des éléments à identifier et commenter.}
\end{popfigure}

\begin{enumerate}
    \item Associer chaque lettre (A à G) à l'aménagement correspondant parmi la liste :
    \begin{multicols}{2}
    \begin{enumerate}
        \item Muret de soutènement en pierres sèches
        \item Haie vive (ex. \textit{Leucaena}, \textit{Calliandra}, \textit{Euphorbia})
        \item Culture en courbes de niveau (labour perpendicular à la pente)
        \item Zone de reboisement / boisement de protection (sommet de versant)
        \item Fascine (boudin de branchages) ou bande enherbée de pied de pente
        \item Drain d'évacuation des eaux de ruissellement (fossé enherbé)
        \item Chemin de ronde / allée de circulation pour l'entretien
    \end{enumerate}
    \end{multicols}
    \item Pour \textbf{chacun} des 7 aménagements identifiés, expliquer \textbf{précisément} son rôle mécanique ou biologique dans la lutte contre l'érosion (ex. : ralentir la vitesse du ruissellement, augmenter l'infiltration, fixer le sol par les racines, piéger les sédiments, drainer l'excès d'eau sans érosion).
    \item Pourquoi l'association de ces techniques (approche intégrée) est-elle plus efficace qu'une technique unique ?
\end{enumerate}
\end{exercice}

\begin{exercice}[Étude de cas : Barrage de Lagdo et colmatation]
Le barrage hydroélectrique de Lagdo (Bénoué, Nord-Cameroun) voit sa capacité de stockage réduire à cause du colmatage par les sédiments charriés par la Bénoué et ses affluents (Mayo Rey, Faro, etc.). L'érosion des bassins versants amont (déforestation, agriculture sur brûlis, surpâturage, orpaillage) est la cause principale.
\begin{enumerate}
    \item Calculer le volume annuel de sédiments piégés si le barrage a perdu 20\% de sa capacité initiale (8 milliards de m³) en 40 ans. Exprimer en \si{m^3.an^{-1}}.
    \item Expliquer le mécanisme physique : pourquoi l'érosion en amont remplit le barrage en aval ? (Notions : transport solide, énergie de la pente, rupture de pente au lac).
    \item Citer \textbf{quatre} conséquences négatives du colmatation du barrage de Lagdo sur : la production d'électricité, l'irrigation (périmètres de la SEMRY), la navigation fluviale, les écosystèmes aquatiques.
    \item Proposer un plan d'action concerté (État, EDC, SEMRY, populations riveraines, ONG) pour réduire l'apport sédimentaire à la source (bassins versants). Distinguer mesures curatives (dragage) et préventives (CES/DRS, reboisement, agroforesterie, gestion des pâturages).
\end{enumerate}
\end{exercice}

\newpage

\coursec{Corrigés des exercices}

\courssub{Je vérifie mes connaissances}

\begin{corrige}[Exercice 1 — QCM : Types et facteurs de l'érosion]
L'affirmation exacte est la \textbf{numéro 3} : \emph{« L'érosion hydrique est favorisée par la pente, l'absence de couverture végétale et l'intensité des pluies. »}

\medskip
\textbf{Justifications des distracteurs :}
\begin{itemize}
    \item \textbf{1. Faux.} La région du Littoral (Douala, Kribi) est soumise à un climat équatorial très humide avec de fortes pluies. L'érosion \textbf{hydrique} (ruissellement, ravines) y domine largement. L'érosion éolienne est caractéristique des zones arides et semi-arides (Grand Nord : Extrême-Nord, Nord).
    \item \textbf{2. Faux.} La \textbf{déforestation} \emph{aggrave} l'érosion hydrique. La suppression de la couverture végétale expose le sol nu à l'impact direct des gouttes de pluie (érosion par \emph{splash}) et supprime l'effet freinant des racines et de la litière sur le ruissellement. Les racines mortes se décomposent et ne maintiennent plus la structure du sol.
    \item \textbf{4. Faux.} Les \textbf{terrasses de culture} (ou terrasses en gradins) transforment une pente continue en une succession de plates-formes horizontales séparées par des murets. Elles \textbf{réduisent fortement la vitesse du ruissellement} (raccourcissement de la longueur de pente, rupture de pente), favorisent l'infiltration et piègent les sédiments.
\end{itemize}
\end{corrige}

\begin{corrige}[Exercice 2 — Vrai ou Faux ? Justifiez]
\begin{enumerate}
    \item \textbf{Faux.} Le labour \textbf{dans le sens de la pente maximale} crée des sillons qui agissent comme des chenaux préférentiels pour l'écoulement de l'eau. Cela \textbf{accélère le ruissellement}, augmente son pouvoir érosif (énergie cinétique) et favorise la naissance de rigoles puis de ravines. La bonne pratique est le labour \textbf{perpendiculaire à la pente} (cultures en courbes de niveau).
    \item \textbf{Vrai.} Le \textbf{reboisement} avec des essences à croissance rapide (eucalyptus, acacia, \textit{Leucaena}, \textit{Calliandra}) reconstitue rapidement une couverture végétale. Le feuillage intercepte la pluie (réduction de l'érosion par \emph{splash}), le système racinaire dense structure le sol (porosité, agrégation) et fixe les particules, la litière absorbe l'eau. C'est une technique de \textbf{protection biologique} majeure, surtout en amont des bassins versants.
    \item \textbf{Faux.} L'érosion commence par la couche arable (horizon A, riche en humus et nutriments), mais si elle n'est pas stoppée, elle attaque l'horizon B (sous-sol) et peut creuser jusqu'au rocher mère (horizon C). Les \textbf{ravines} et \textbf{coulées de boue} témoignent de cette atteinte aux horizons profonds, avec perte irréversible du capital sol.
    \item \textbf{Vrai.} C'est la définition même des \textbf{cultures en courbes de niveau}. Le labour et les lignes de semis suivent les lignes d'égale altitude (courbes de niveau). Les sillons ainsi formés agissent comme une multitude de \textbf{mini-barrages} qui retiennent l'eau de pluie, augmentent le temps d'infiltration et réduisent la vitesse du ruissellement.
\end{enumerate}
\end{corrige}

\begin{corrige}[Exercice 3 — Définitions et associations]
\begin{enumerate}
    \item \textbf{Définitions précises :}
    \begin{definition}[Érosion]
    Ensemble des processus naturels ou anthropiques de \textbf{détachement}, de \textbf{transport} et de \textbf{dépôt} des particules du sol (ou de la roche) sous l'action d'agents physiques : l'eau (ruissellement, vagues), le vent (deflation), la glace (gel-dégel, glaciers) ou la gravité (éboulements).
    \end{definition}
    \begin{definition}[Ruissellement]
    Écoulement superficiel de l'eau de pluie (ou de fonte des neiges) à la surface du sol, qui se produit lorsque l'intensité pluviométrique dépasse la \textbf{capacité d'infiltration} du sol (sol saturé, croûte de battance, sol gelé, imperméabilisation). Il peut être \emph{diffus} (nappe mince) ou \emph{concentré} (rigoles, ravines, cours d'eau).
    \end{definition}
    \begin{definition}[Ravine]
    Forme d'érosion hydrique \textbf{concentrée} et avancée, creusée par le ruissellement dans un thalweg ou sur un versant. Creux profond (\textbf{> 30--50 cm}, infranchissable par le matériel agricole), aux parois souvent verticales et instables, qui évolue par \textbf{régression amont} (recul de la tête de ravine) et élargissement latéral.
    \end{definition}
    \begin{definition}[Déforestation]
    Réduction durable, voire destruction, de la \textbf{couverture forestière} naturelle par l'action humaine (agriculture sur brûlis, exploitation forestière non durable, collecte de bois de feu, extension urbaine, mines) sans régénération naturelle ni reboisement compensatoire suffisant. Elle expose le sol nu aux agents érosifs.
    \end{definition}

    \item \textbf{Associations correctes (tableau complété) :}
    \begin{center}
    \begin{tabularx}{\linewidth}{|>{\bfseries}X|X|X|}\hline
    \textbf{A : Type d'érosion} & \textbf{B : Agent érosif} & \textbf{C : Facteur favorisant majeur} \\ \hline
    Érosion hydrique & Eau de ruissellement (pluie, cours d'eau) & Pente forte + sol nu (déforestation, labour sens pente) \\ \hline
    Érosion éolienne & Vent (deflation, saltation, suspension) & Aridité (sol sec) + absence de haies brise-vent + sol nu \\ \hline
    \end{tabularx}
    \end{center}
\end{enumerate}
\end{corrige}

\begin{corrige}[Exercice 4 — Calcul direct : Taux d'érosion]
\textbf{Données :}
$S = \SI{2}{\hectare} = \SI{20\,000}{\square\metre}$ \quad $e = \SI{5}{\milli\metre} = \SI{0,005}{\metre}$ \quad $\rho_{\text{app}} = \SI{1,3}{\tonne\per\cubic\metre}$ \quad $\Delta t = \SI{1}{an}$ (une saison des pluies)

\begin{enumerate}
    \item \textbf{Volume de terre perdu ($V$) :}
    \[
    V = S \times e = \SI{20\,000}{\square\metre} \times \SI{0,005}{\metre} = \boxed{\SI{100}{\cubic\metre}}
    \]

    \item \textbf{Masse de terre perdue ($M$) :}
    \[
    M = V \times \rho_{\text{app}} = \SI{100}{\cubic\metre} \times \SI{1,3}{\tonne\per\cubic\metre} = \boxed{\SI{130}{\tonne}}
    \]

    \item \textbf{Taux d'érosion ($T$) :}
    \[
    T = \frac{M}{S_{\text{ha}} \times \Delta t} = \frac{\SI{130}{\tonne}}{\SI{2}{\hectare} \times \SI{1}{an}} = \boxed{\SI{65}{t.ha^{-1}.an^{-1}}}
    \]

    \textbf{Interprétation :} Ce taux (\SI{65}{t.ha^{-1}.an^{-1}}) est \textbf{très largement supérieur} au seuil de tolérance généralement admis (\textbf{< \SI{10}{t.ha^{-1}.an^{-1}}}). Il correspond à une \textbf{érosion sévère à très sévère}. La parcelle perd l'équivalent de \SI{5}{\milli\metre} de terre par an ; à ce rythme, la couche arable (20--30 cm) disparaîtrait en 40 à 60 ans. Une intervention urgente (courbes de niveau, haies, couverts) est indispensable.
\end{enumerate}
\end{corrige}

\courssub{Je m'entraîne}

\begin{corrige}[Exercice 5 — Analyse de photographies : Ravines urbaines à Douala]
\begin{enumerate}
    \item \textbf{Phénomène visible (Photo 1) :} Une \textbf{ravine} (gully erosion) profonde de 3--4 m coupe la route en terre battue. Les berges verticales et instables indiquent une érosion active par \textbf{régression amont} (la tête de la ravine recule vers le haut du versant) et éboulements latéraux. La présence de déchets plastiques au fond montre l'usage du ravin comme décharge.
    \textbf{Type d'érosion :} \textbf{Érosion hydrique concentrée} (stade ravine), conséquence d'un ruissellement de surface non maîtrisé.

    \item \textbf{Trois causes probables (facteurs favorisants) en milieu urbain (Douala, pente, climat équatorial) :}
    \begin{enumerate}
        \item \textbf{Imperméabilisation et concentration du ruissellement :} Urbanisation anarchique (toits, routes, cours bétonnées) sans réseau de drainage adapté. L'eau de pluie est concentrée vers les points bas (rues en pente) avec une énergie érosive décuplée.
        \item \textbf{Sol nu et instable :} Absence de végétation (arbres d'alignement, gazon, haies) sur les talus et emprises. Sol latéritique (ferralsol) sensible à la dispersion (argiles kaolinitiques, faible cohésion) dès qu'il est mis à nu.
        \item \textbf{Obstruction des exutoires naturels :} Décharges sauvages (plastiques, gravats) dans les talwegs et caniveaux existants. Lors des fortes averses, l'eau déborde, érode les berges et creuse de nouveaux chemins (ravines).
    \end{enumerate}

    \item \textbf{Lien Photo 1 $\rightarrow$ Photo 2 :} La ravine (Photo 1) est une \textbf{source massive de sédiments} (argiles, sables, graviers, déchets). Le ruissellement torrentiel charrie ce \textbf{transport solide} en aval. Au pied de la pente (rupture de pente), la vitesse chute brutalement dans le caniveau en béton (Photo 2) : \textbf{dépôt (colmatage)} de la charge solide (boue, sable). Le caniveau bouché ne peut plus évacuer l'eau $\rightarrow$ \textbf{stagnation, inondations locales, insalubrité}.

    \item \textbf{Mesures proposées :}
    \begin{itemize}
        \item \textbf{Urgence (curatif) :}
        \begin{enumerate}
            \item \textbf{Débouchage et curage} immédiat du caniveau (Photo 2) et du fond de ravine pour rétablir l'écoulement.
            \item \textbf{Stabilisation d'urgence} de la tête et des berges de la ravine : pose de \textbf{gabions} (cages grillagées remplies de pierres), \textbf{fascines} (boudins de branchages) ou \textbf{palplanches} pour stopper la régression amont et protéger la route.
        \end{enumerate}
        \item \textbf{Prévention durable (structurel) :}
        \begin{enumerate}
            \item \textbf{Aménagement hydraulique du bassin versant :} Création de \textbf{bassins de rétention/détention} en amont, dimensionnement et réalisation de \textbf{caniveaux en béton armés} (avec grilles anti-déchets) tout le long des voiries, exutoires vers le fleuve Wouri ou le marais.
            \item \textbf{Végétalisation et assainissement :} Reboisement des talus (vetiver, bambous, arbres à racines pivotantes), plantation de haies vives en courbes de niveau, mise en place d'une \textbf{collecte régulière des ordures} pour éviter l'obstruction des drains, plan d'urbanisme respectant les lignes d'écoulement naturel.
        \end{enumerate}
    \end{itemize}
\end{enumerate}
\end{corrige}

\begin{corrige}[Exercice 6 — Choix technique : Cas d'un agriculteur à Foumban]
\begin{enumerate}
    \item \textbf{Principes :}
    \begin{itemize}
        \item \textbf{Terrasses en gradins (terrasses étroites à murets) :} Aménagement mécanique lourd. Découpage du versant en \textbf{plates-formes horizontales} (berces) de largeur réduite (2--5 m), soutenues par des \textbf{murets en pierres sèches} (ou en terre végétalisée). Transforment une pente unique en une succession de marches. Raccourcissent la longueur de pente effective à la largeur de la berce, annulent la pente longitudinale $\rightarrow$ infiltration maximale, ruissellement quasi nul.
        \item \textbf{Cultures en courbes de niveau (avec bandes enherbées) :} Technique agronomique. Labour, semis, binage \textbf{perpendiculairement à la pente maximale} (suivant les courbes de niveau). Les sillons et les \textbf{bandes enherbées} (largeur 1--2 m, espèces : \textit{Brachiaria}, \textit{Panicum}, \textit{Vetiver}, légumineuses) placées à intervalles réguliers (tous les 10--20 m selon la pente) forment des \textbf{barrages filtrants} vivants qui ralentissent l'eau, piègent les sédiments et augmentent l'infiltration.
    \end{itemize}

    \item \textbf{Comparaison dans le contexte de M. Njoya (Pente 15\%, Sol ferrallitique, Pluies intenses, Paysan moyen) :}

    \begin{center}
    \begin{tabularx}{\linewidth}{|>{\bfseries}l|X|X|}\hline
    \textbf{Critère} & \textbf{Terrasses en gradins (murets pierres)} & \textbf{Courbes de niveau + Bandes enherbées} \\ \hline
    \textbf{Coût mise en œuvre} & \textbf{Très élevé.} Main-d'œuvre intense (terrassement, transport pierres, construction murets). Nécessite souvent main-d'œuvre salariée ou groupement. & \textbf{Faible.} Changement de pratique cultural (pas de terrassement). Coût = semences bandes enherbées + main-d'œuvre familiale pour tracer les courbes (niveau à cadre/bourdon). \\ \hline
    \textbf{Main-d'œuvre / Entretien} & Entretien \textbf{régulier et contraignant} : réparation murets (éboulements), curage fossés, gestion de l'eau sur berces. Risque rupture muret = catastrophe locale. & Entretien \textbf{simple} : fauche/roulage bandes enherbées (paillis), respect du sens labour. Auto-entretenu par la végétation. \\ \hline
    \textbf{Efficacité sol ferrallitique} & \textbf{Excellente} si murets stables. Sol ferrallitique profond $\rightarrow$ bonne réserve en eau sur berces horizontales. Risque : lessivage latéral si muret perméable. & \textbf{Bonne à très bonne.} Bandes enherbées structurent le sol (racines), apportent MO, fixent N (légumineuses). Efficacité prouvée sur pentes < 20\% en zone soudano-guinéenne. \\ \hline
    \textbf{Surface cultivable} & \textbf{Perte de surface} (emprise murets + fossés $\approx$ 15--20\% de la parcelle). & \textbf{Perte minime} (bandes enherbées $\approx$ 5--10\%, valorisables : fourrage, paillis, piquets). \\ \hline
    \textbf{Adaptation cultures (Maïs/Haricot)} & Adapté (riz, maraîchage, maïs). Difficile pour cultures hautes (maïs) sur berces étroites (ombrage muret, vent). & \textbf{Idéal.} Maïs/Haricot en rotation/association sur billons courbes. Bandes enherbées en bordure ou intercalaires. \\ \hline
    \end{tabularx}
    \end{center}

    \item \textbf{Conseil argumenté :}
    \textbf{Choix recommandé : \emph{Cultures en courbes de niveau associées à des bandes enherbées (vetiver/brachiaria) et haies vives (Calliandra/Leucaena) en bordure de parcelle.}}

    \textbf{Justification :}
    \begin{itemize}
        \item \textbf{Pente 15\% :} Seuil où les courbes de niveau + bandes enherbées restent très efficaces (recommandées jusqu'à 20--25\% si sol profond). Terrasses justifiées > 20--25\%.
        \item \textbf{Sol ferrallitique (Ouest) :} Profond, bien drainé, sensible à l'érosion par dispersion. Les bandes enherbées (racines fasciculées) restructurent la surface, augmentent la MO et l'activité biologique (vers de terre) $\rightarrow$ durabilité de la fertilité.
        \item \textbf{Moyens du paysan :} Faible investissement initial (pas de pierres à acheter/transporter), main-d'œuvre familiale suffisante, autonomie technique (apprentissage rapide du tracé au niveau à cadre/bourdon).
        \item \textbf{Intensification écologique :} Bandes enherbées = fourrage (élevage), paillis (maïs), bois de feu (haies), azote (légumineuses). Système \textbf{agroforestier} résilient (agriculture climato-intelligente).
    \end{itemize}
    \begin{attention}[Point de vigilance]
    Sur pente 15\% avec pluies intenses (Ouest), les courbes de niveau \emph{seules} peuvent être insuffisantes en cas d'orage exceptionnel. L'association \textbf{obligatoire} avec des \textbf{bandes enherbées pérennes} (vetiver tous les 10--15 m) et des \textbf{haies vives} en limite de parcelle est la condition de l'efficacité durable. Le tracé rigoureux des courbes de niveau (niveau à eau ou à cadre) est la clé de la réussite.
    \end{attention}
\end{enumerate}
\end{corrige}

\begin{corrige}[Exercice 7 — Déforestation et inondations : Cas du Mont Cameroun]
\begin{enumerate}
    \item \textbf{Schéma annoté et mécanisme (Coupe transversale Sommet $\rightarrow$ Pied) :}
    \begin{popfigure}
    \begin{tikzpicture}[scale=0.7, every node/.style={font=\small}]
    % Profil du volcan
    \fill[poporange!30] (0,0) -- (16,0) -- (14,5) -- (10,7) -- (6,7) -- (2,5) -- cycle;
    \fill[popgreen!20] (2,5) -- (6,7) -- (10,7) -- (14,5) -- (13,4) -- (11,3) -- (9,3.5) -- (7,3.5) -- (5,3) -- (3,2.5) -- cycle;
    
    % Forêt naturelle (amont - gauche)
    \foreach \x in {1, 1.5, 2.2, 2.8, 3.5} {
        \node[circle, fill=popgreen!60!popdark, minimum size=8pt, inner sep=0] at (\x, 2.5+0.5*rand) {};
        \node[circle, fill=popgreen!40!popdark, minimum size=5pt, inner sep=0] at (\x+0.2, 2.2+0.5*rand) {};
    }
    \node[align=center, popgreen!60!popdark] at (2, 4.5) {Forêt naturelle\\dense, humus épais};
    
    % Zone déforestée / plantations (milieu)
    \fill[poporange!40] (4,3) -- (12,5.5) -- (11,4.5) -- (5,2.5) -- cycle;
    \foreach \x in {5, 7, 9, 11} {
        \node[circle, fill=popgreen!50!poporange, minimum size=6pt, inner sep=0] at (\x, 3.5+0.2*rand) {};
    }
    \node[align=center, poporange!80!popdark] at (8, 5) {Plantations\\Palmier/Hévéa/Banane\\Sol nu entre rangs};
    
    % Zone urbaine / pied de pente (aval - droite)
    \fill[popmuted!30] (12,0) -- (16,0) -- (14,2) -- (12,2) -- cycle;
    \draw[popmuted, thick] (12.5,0) -- (12.5,1.5) node[midway, right] {Villes (Limbé, Bota)};
    \draw[popmuted, thick] (13.5,0) -- (13.5,1.2);
    \draw[popmuted, thick] (14.5,0) -- (14.5,1);
    
    % Pluie
    \foreach \x in {2, 4, 6, 8, 10, 12, 14} {
        \draw[popblue, thick, ->] (\x, 7.5) -- (\x, 6.5);
    }
    \node[popblue] at (8, 7.8) {Pluies intenses (4000--10000 mm/an)};
    
    % Infiltration / Ruissellement FORÊT
    \draw[popblue, ->, thick] (2, 4) -- (2, 3) node[midway, left, align=center]{Infiltration\\rapide};
    \draw[popblue, dashed, ->] (2, 4) -- (1.5, 3.5) node[midway, above, align=center]{Ruissellement\\faible};
    
    % Ruissellement / Érosion PLANTATIONS
    \draw[popblue, ->, very thick] (7, 4.5) -- (7, 3.5) node[midway, right, align=center]{Ruissellement\\CONCENTRÉ \& RAPIDE};
    \draw[poporange, ->, thick] (7, 3.5) -- (7, 2.5) node[midway, right] {Érosion en rigoles/ravines};
    \draw[poporange, ->, thick] (7, 2.5) -- (10, 1.5) node[midway, below, align=center]{Transport solide\\boues, blocs};
    
    % Dépôt / Inondation AVAL
    \draw[popblue, ->, thick] (13, 2) -- (13, 1) node[midway, right, align=center]{Dépôt brutal\\(cônes de déjection)};
    \draw[popblue, <->, thick] (12.5, 0.5) -- (14.5, 0.5) node[midway, below] {Inondations, coulées de boue};
    
    % Flèches gravité
    \draw[popdark, ->] (8, 6) -- (8, 5.5) node[midway, right] {Pente forte (énergie)};
    \end{tikzpicture}
    \legende{Mécanisme : Déforestation amont $\rightarrow$ Perte pouvoir tampon $\rightarrow$ Ruissellement concentré + Érosion massive $\rightarrow$ Transport solide $\rightarrow$ Rupture de pente au pied $\rightarrow$ Dépôt (colmatage) + Inondations aval.}
    \end{popfigure}

    \textbf{Explication du mécanisme :}
    \begin{enumerate}
        \item \textbf{Amont (Forêt naturelle) :} Couvert dense + litière épaisse + racines $\rightarrow$ \textbf{Infiltration maximale}, ruissellement diffus lent, érosion quasi nulle. Rôle « éponge ».
        \item \textbf{Milieu (Déforestation/Plantations) :} Sol nu entre rangs, routes d'exploitation, perte de MO $\rightarrow$ \textbf{Battance, croûte de surface} $\rightarrow$ Infiltration chuté $\rightarrow$ \textbf{Ruissellement de surface concentré, rapide, érosif}. Naissance de rigoles/ravines sur les pistes et lignes de pente. \textbf{Érosion hydrique massive} (détachement + transport).
        \item \textbf{Aval (Pied de pente / Villes) :} \textbf{Rupture de pente brutale} (volcan $\rightarrow$ plaine côtière). Perte soudaine de l'énergie cinétique du flux $\rightarrow$ \textbf{Dépôt brutal de la charge solide} (alluvions, boues, blocs) : cônes de déjection. Colmatage des lits mineurs $\rightarrow$ \textbf{Débordements, inondations, coulées de boue} dans les quartiers (Limbé 2001).
    \end{enumerate}

    \item \textbf{Trois services écosystémiques perdus :}
    \begin{enumerate}
        \item \textbf{Régulation hydrique (effet éponge) :} Stockage de l'eau de pluie dans l'humus et la porosité du sol forestier, restitution lente (maintien débit étiage, amortissement crues).
        \item \textbf{Protection des sols / Lutte anti-érosive :} Interception pluie (canopée), fixation sol (racines), structure stable (MO, faune du sol). Perte $\rightarrow$ érosion accélérée $\times$ 10 à 100.
        \item \textbf{Stockage de carbone / Atténuation changement climatique :} Biomasse aérienne et racinaire énorme (forêt tropicale dense). Déforestation $\rightarrow$ émission \ce{CO2} massive + perte puits de carbone.
        \item \textit{(Autres acceptables : Biodiversité habitat/refuge espèces endémiques ; Ressources NTFP/bois/médecine ; Régulation microclimat locale).}
    \end{enumerate}

    \item \textbf{Stratégie de gestion intégrée du bassin versant (Amont/Aval) :}
    \begin{center}
    \begin{tabularx}{\linewidth}{|>{\bfseries}l|X|X|}\hline
    \textbf{Zonage} & \textbf{Acteurs clés} & \textbf{Mesures concrètes} \\ \hline
    \textbf{AMONT} (Pentes fortes > 30\%, sources) & MINFOF, MINEPDED, ONG, Communautés riveraines & \textbf{Protection intégrale :} Arrêt déforestation, reboisement espèces natives, création aires protégées/forêts communautaires. \textbf{Agroforesterie} dans zones tampons (café/cacao sous couvert). \\ \hline
    \textbf{MILIEU} (Plantations agro-industrielles CDC/SOCAPALM, petits exploitants) & Planteurs (CDC, petits), MINADER, MINEPDED & \textbf{Agriculture de conservation :} Couverts végétaux permanents (Mucuna, Pueraria), haies vives en courbes de niveau, non-labour, gestion intégrée ravines (barrages filtrants, fascines). \textbf{Zéro brûlis.} \\ \hline
    \textbf{AVAL} (Villes Limbé/Bota/Idenau, zone côtière) & Mairies, MINHDU, MINTP, Protection Civile, Populations & \textbf{Aménagement hydraulique :} Caniveaux dimensionnés (crue centennale), bassins de rétention/détention, curage régulier, \textbf{rejet des déchets interdit}. \textbf{Plan d'occupation des sols (POS)} : zones non aedificandi en zones inondables. Système d'alerte précoce (SAP). \\ \hline
    \textbf{GOUVERNANCE} (Bassin versant entier) & Préfet (Comité de bassin), MINEPDED, MINADER, MINFOF, CTD, ONG, Privés & \textbf{Contrat de bassin versant / Schéma Directeur d'Aménagement et de Gestion de l'Eau (SDADE).} \textbf{PSE (Paiement pour Services Environnementaux)} : planteurs/riverains rémunérés pour pratiques vertueuses (haies, reboisement). Fonds dédié (Fonds Forestier, Fonds Environnement). \\ \hline
    \end{tabularx}
    \end{center}
\end{enumerate}
\end{corrige}

\begin{corrige}[Exercice 8 — Tableau de synthèse : Types, facteurs, conséquences, lutte]
\begin{center}
\begin{tabularx}{\linewidth}{|>{\bfseries}l|X|X|X|}\hline
\textbf{Type d'érosion} & \textbf{Facteurs favorisants principaux} & \textbf{Conséquences majeures (sol / environnement)} & \textbf{Techniques de lutte adaptées} \\ \hline
Érosion hydrique en feuille (diffuse) & Pluies intenses (énergie cinétique gouttes), \textbf{sol nu} (absence couvert végétal, résidus), structure instable (faible MO, battance), pente faible à modérée, labour fin & Perte \textbf{uniforme} de la couche arable (horizon A), appauvrissement progressif en \textbf{éléments nutritifs (N, P, K), MO}, baisse capacité de rétention d'eau, baisse rendements (insidieuse, peu visible) & \textbf{Protection biologique du sol :} Couverts végétaux permanents (engrais verts, résidus), \textbf{non-labour / semis direct}, paillage, rotation culturale, haies vives en courbes de niveau, agroforesterie \\ \hline
Érosion hydrique concentrée (rigoles, ravines) & \textbf{Concentration du ruissellement} (talwegs, pistes, labour sens pente), \textbf{pente forte}, sol nu, absence d'aménagements, crues soudaines & \textbf{Creusement profond} (rigoles > ravines), perte massive de sol (horizons A+B+C), \textbf{destruction infrastructures} (routes, champs, habitats), \textbf{colmatage aval} (barrages, cours d'eau), paysages dégradés (badlands) & \textbf{Aménagements mécaniques + biologiques :} Terrasses (gradins, benched), \textbf{fascines}, \textbf{barrages filtrants} (pierres, gabions) dans talwegs, reboisement têtes de ravines, \textbf{drains d'évacuation} enherbés, endiguement \\ \hline
Érosion éolienne (deflation) & \textbf{Aridité / Saison sèche} (sol sec, pulvérulent), \textbf{vent fort} (harmattan), \textbf{sol nu} (surpâturage, labour, récolte résidus), absence haies brise-vent, texture fine (sable, limon) & Perte sélective des \textbf{particules fines} (argiles, limons, MO, nutriments) $\rightarrow$ appauvrissement chimique/physique. \textbf{Ensablement} cultures, habitations, infrastructures. Tempêtes de sable (santé, transport). Dégradation nappes phréatiques. & \textbf{Fixation + Protection :} \textbf{Haies brise-vent} pérennes (Casuarina, Prosopis, Euphorbia, Leucaena), \textbf{couverts végétaux} résidus, \textbf{non-labour}, agriculture de conservation, \textbf{reboisement} dunes/bas-fonds, techniques zaï/demi-lunes (infiltration), pâturage tournant \\ \hline
\end{tabularx}
\end{center}
\end{corrige}

\courssub{Je me prépare au Probatoire}

\begin{corrige}[Exercice 9 — Sujet de synthèse : Lutte contre l'érosion et développement durable de l'agriculture camerounaise]
\textbf{Barème indicatif :} \emph{Total 8 points} (Causes : 2,5 pts ; Conséquences : 2,5 pts ; Solutions : 3 pts). \textbf{Points de vigilance :} Exemples régionaux précis obligatoires ; articulation causes $\rightarrow$ conséquences $\rightarrow$ solutions ; vocabulaire technique (CES/DRS, agroforesterie, colmatation, PAN/LCD, ACI) ; structure en 3 parties distinctes.

\medskip
\textbf{Introduction :} L'érosion des sols (hydrique au Sud/Centre/Est/Ouest, éolienne au Nord/Adamaoua) constitue la \textbf{menace n°1} pour la durabilité de l'agriculture camerounaise. Elle détruit le capital naturel (sol, eau, biodiversité) base de la production. La lutte anti-érosive n'est pas une option technique mais une \textbf{condition \emph{sine qua non}} du développement durable (sécurité alimentaire, revenus ruraux, résilience climatique).

\medskip
\textbf{1. Les causes : Analyse multi-factorielle (Naturelles $\ll$ Anthropiques)}

\begin{itemize}
    \item \textbf{Facteurs naturels (prédisposants) :}
    \begin{itemize}
        \item \textbf{Climat :} Pluies intenses, érosives (indice Fournier élevé) au Sud (bimodal), Ouest (montagnes), Littoral. Saison sèche longue + Harmattan violent (Nord, Adamaoua) $\rightarrow$ sols secs, pulvérulents.
        \item \textbf{Relief :} Pentes fortes (Monts Cameroun, Manengouba, Bamboutos, Mandara, Adamaoua) $\rightarrow$ énergie cinétique ruissellement / vent.
        \item \textbf{Sols :} Ferrallitiques (Sud/Ouest) : lessivés, faiblement structurés, sensibles à la dispersion. Ferrugineux tropicaux (Nord/Adamaoua) : croûte de battance, faible MO. Sols volcaniques jeunes (Mont Cameroun) : instables.
    \end{itemize}
    \item \textbf{Facteurs anthropiques (déclenchants/aggravants) — \emph{Responsabilité majeure} :}
    \begin{itemize}
        \item \textbf{Déforestation massive :} Bois de feu (90\% énergie ménages), agriculture sur brûlis (\emph{ankara}, \emph{bilong}), exploitation forestière illégale, extension plantations (palmier, hévéa, banane) $\rightarrow$ \textbf{Sud, Est, Littoral, Mont Cameroun}.
        \item \textbf{Pratiques culturales inadaptées :} Labour dans le sens de la pente (Ouest, Adamaoua), absence de couverts/résidus, monocultures (coton, maïs) sans rotation, labour excessif (pulvérisation structure).
        \item \textbf{Surpâturage / Feux de brousse :} Nord (Logone, Bénoué), Adamaoua. Destruction couvert herbacé, piétinement (croûte), feux tardifs (sol nu mois avant pluies).
        \item \textbf{Pression foncière / Urbain :} Raccourcissement jachères (Ouest, Nord-Ouest), agriculture sur pentes marginales, urbanisation anarchique sans drainage (Douala, Yaoundé $\rightarrow$ ravines urbaines).
        \item \textbf{Orpaillage artisanal :} Est (Bétaré-Oya), Nord (Mayo Rey) $\rightarrow$ destruction totale profil sol, mercure.
    \end{itemize}
\end{itemize}

\medskip
\textbf{2. Les conséquences : Impacts systémiques (Productivité -- Environnement -- Économie)}

\begin{itemize}
    \item \textbf{Sur la productivité agricole \& Sécurité alimentaire :}
    \begin{itemize}
        \item Perte horizon arable (réservoir nutriments, eau, biodiversité du sol) $\rightarrow$ \textbf{baisse rendements} (maïs -30\%, coton -40\% zones érodées Nord).
        \item Dégradation physique (compaction, perte structure) $\rightarrow$ travail du sol plus difficile/coûteux.
        \item \textbf{Insécurité alimentaire} structurelle (régions Nord, Adamaoua, Extrême-Nord) : dépendance aide alimentaire, importations (riz, blé).
    \end{itemize}
    \item \textbf{Sur l'environnement :}
    \begin{itemize}
        \item \textbf{Colmatage barrages :} Lagdo (Bénoué), Maga (Logone), Lom Pangar (Lom), Bamendjin (Noun) $\rightarrow$ perte capacité stockage (Lagdo -20\% en 40 ans), évaporation accrue, eutrophisation.
        \item \textbf{Pollution eaux :} Turbidité, pesticides, nitrates $\rightarrow$ coût potabilisation, mortalité piscicole.
        \item \textbf{Perte biodiversité :} Destruction habitats (forêts, zones humides), érosion génétique (variétés locales adaptées).
    \end{itemize}
    \item \textbf{Sur l'économie rurale \& Social :}
    \begin{itemize}
        \item Baisse revenus paysans (intrants plus chers pour compenser perte fertilité), endettement.
        \item \textbf{Exode rural / Migration :} Jeunes quittent terres dégradées $\rightarrow$ vieillissement population agricole, main-d'œuvre rare/coûteuse.
        \item Conflits fonciers (agriculteurs/éleveurs Nord), coûts curatifs énormes (dragage Lagdo, réfection routes).
    \end{itemize}
\end{itemize}

\medskip
\textbf{3. Les solutions : Stratégie globale intégrée (Technique + Institutionnelle)}

\begin{methode}[Stratégie nationale de conservation des sols (CES/DRS)]
\begin{enumerate}
    \item \textbf{Techniques agronomiques (Protection biologique, faible coût, paysannes) :}
    \begin{itemize}
        \item \textbf{Cultures en courbes de niveau} + \textbf{Bandes enherbées} (Vetiver, Brachiaria, Stylosanthes) : Ouest, Nord-Ouest, Adamaoua.
        \item \textbf{Agroforesterie / Parc à Faidherbia albida} (Nord, Extrême-Nord) : ombrage, azote, fourrage, bois, racines pivotantes.
        \item \textbf{Couverts végétaux permanents / Semis direct sous couvert (SCV) :} Mucuna, Canavalia, Crotalaria (cotton basin Nord, maïs Sud). Zéro labour.
        \item \textbf{Haies vives multi-étagées} (Calliandra, Leucaena, Gliricidia) : brise-vent, piquets, fourrage, fertilité.
    \end{itemize}
    \item \textbf{Techniques mécaniques (Investissement lourd, pentes fortes/concentrées) :}
    \begin{itemize}
        \item \textbf{Terrasses en gradins (benched terraces) + murets pierres :} Hauts plateaux de l'Ouest (Bamiléké), Monts Mandara.
        \item \textbf{Fascines / Boudins végétaux / Barrages filtrants :} Talwegs, têtes de ravines, bas de pente.
        \item \textbf{Demi-lunes / Zai / Cordons pierreux :} Nord (zaï Burkina importé), zones sahéliennes (récupération terres nues \emph{glacis}).
        \item \textbf{Drains d'évacuation enherbés / Bassins de rétention :} Zones urbaines, pieds de versant.
    \end{itemize}
    \item \textbf{Mesures institutionnelles \& politiques (Levier d'adoption) :}
    \begin{itemize}
        \item \textbf{Sécurisation foncière (Code foncier 2022, titres/contrats) :} Incite à l'investissement long terme (terrasses, arbres).
        \item \textbf{Vulgarisation de proximité :} Relance \textbf{CNV (Centres Nationaux de Vulgarisation)}, Ecoles Champs Paysans (ECP), conseillers agricoles formés CES/DRS.
        \item \textbf{Éducation à l'environnement :} Intégration programmes scolaires (SVT, Géographie), médias, chefs traditionnels.
        \item \textbf{Financement innovant :} \textbf{PADFA} (Projet d'Appui au Développement de la Filière Agricole -- Nord), \textbf{PRODESA}, \textbf{ACI} (Agriculture Climato-Intelligente -- Banque Mondiale/FAO), \textbf{PSE} (Paiement Services Environnementaux), micro-crédit vert.
        \item \textbf{Cadre légal :} Application stricte Loi Forestière (1994), Loi sur l'Environnement (1996), PAN/LCD (Plan d'Action National de Lutte Contre la Désertification).
    \end{itemize}
\end{enumerate}
\end{methode}

\medskip
\textbf{Exemples de réussite à citer :}
\begin{itemize}
    \item \textbf{Nord (Bénoué, Mayo-Rey) :} \textbf{PADFA / ACI Coton} : Adoption massive SCV (Mucuna, Brachiaria) + Zaï + Haies vives $\rightarrow$ Rendements coton/maïs doublés, sols régénérés, femmes soulagées (moins de binage).
    \item \textbf{Ouest (Bamiléké) :} \textbf{Terrasses traditionnelles modernisées} (pierres sèches, haies vives, fumure) + Agroforesterie café/bananier $\rightarrow$ Densités rurales records soutenues durablement.
    \item \textbf{Adamaoua / Nord :} \textbf{Parcs à Faidherbia} (régénération naturelle assistée - RNA) + Gestion pastorale (coupes-feu, points d'eau) $\rightarrow$ Réhabilitation sols ferrugineux, coexistence agriculture/élevage.
\end{itemize}

\medskip
\textbf{Conclusion :} La lutte contre l'érosion au Cameroun exige une \textbf{approche paysagère (bassin versant)}, combinant savoir-faire paysan (terrasses Bamiléké, Zaï, Faidherbia) et innovations scientifiques (SCV, ACI), portée par une \textbf{volonté politique forte} (foncier, vulgarisation, financement). C'est le socle de la \textbf{Stratégie Nationale de Développement (SND30)} et de l'\textbf{Agenda 2063} de l'UA : « L'Afrique que nous voulons » ne se construira pas sur des sols qui s'envolent ou s'écoulent.
\end{corrige}

\begin{corrige}[Exercice 10 — Interprétation de schéma : Versant aménagé (Type Bamiléké moderne)]
\begin{enumerate}
    \item \textbf{Association Lettres $\rightarrow$ Aménagements :}
    \begin{center}
    \begin{tabularx}{\linewidth}{|c|X|X|}\hline
    \textbf{Lettre} & \textbf{Aménagement identifié} & \textbf{Justification visuelle (schéma)} \\ \hline
    \textbf{A} & \textbf{1. Muret de soutènement en pierres sèches} & Pied de la première terrasse (côté gauche), soutènement vertical visible. \\ \hline
    \textbf{B} & \textbf{2. Haie vive} (\textit{Leucaena, Calliandra, Euphorbia}) & Bord droit de la terrasse inférieure, alignement vertical d'arbustes. \\ \hline
    \textbf{C} & \textbf{6. Drain d'évacuation des eaux (fossé enherbé)} & Bord droit de la terrasse moyenne, position exutoire latéral (flèches bleues). \\ \hline
    \textbf{D} & \textbf{4. Zone de reboisement / boisement de protection (sommet)} & Zone supérieure plate, couverte d'arbres (cercles verts). \\ \hline
    \textbf{E} & \textbf{3. Culture en courbes de niveau} (labour $\perp$ pente) & Surface plane de la terrasse inférieure, motifs de cultures (traits verts). \\ \hline
    \textbf{F} & \textbf{5. Fascine / Bande enherbée de pied de pente} & \emph{Position théorique} : bas de versant / pied de la première terrasse (marqueur centré bas). \\ \hline
    \textbf{G} & \textbf{7. Chemin de ronde / allée de circulation} & \emph{Position théorique} : sur la terrasse supérieure (largeur, accès entretien murets/haies). \\ \hline
    \end{tabularx}
    \end{center}
    \begin{attention}[Note sur le schéma]
    Le schéma fourni est une coupe schématique simplifiée. Les repères F et G ne correspondent pas à des objets dessinés explicitement (fascine, chemin) mais indiquent leur \textbf{emplacement standard} dans un système bamiléké complet : F au pied du versant (exutoire final), G sur la berce la plus haute (desserte).
    \end{attention}

    \item \textbf{Rôle précis de chaque aménagement dans la lutte anti-érosive :}
    \begin{enumerate}
        \item \textbf{Muret de soutènement (A) :} \textbf{Rôle mécanique structurel.} Supporte la terrasse, annule la pente longitudinale sur la berce, résiste à la poussée du sol. \textbf{Freine/stoppe le ruissellement} descendant, force l'infiltration sur la plate-forme. Piège les sédiments en amont du muret.
        \item \textbf{Haie vive (B) :} \textbf{Rôle biologique + mécanique.} Racines profondes/fasciculées $\rightarrow$ \textbf{fixation du sol} (cohésion), porosité (infiltration). Partie aérienne $\rightarrow$ \textbf{ralentit le ruissellement} (rugosité), \textbf{piège les sédiments} (filtration), brise-vent. Apport MO (feuilles, tailles). \textit{Leucaena/Calliandra} : fixatrices \ce{N2}.
        \item \textbf{Culture en courbes de niveau (E) :} \textbf{Rôle agronomique.} Sillons $\perp$ pente = \textbf{mini-barrages} successifs. Réduit \textbf{vitesse ruissellement} (facteur $V^2$ dans pouvoir érosif), augmente \textbf{temps d'infiltration}, répartit l'eau uniformément sur la berce.
        \item \textbf{Zone reboisée sommet (D) :} \textbf{Rôle régulateur (amont).} « Château d'eau » : interception pluie, infiltration maximale, alimentation nappes/sources. \textbf{Protection de la tête de bassin} : empêche naissance ravines, stabilise le sommet du versant. Biodiversité, bois, carbone.
        \item \textbf{Fascine / Bande enherbée pied de pente (F) :} \textbf{Rôle filtre / épuration finale.} Dernier barrage avant l'exutoire (cours d'eau, route). \textbf{Piège les derniers sédiments} (colmatage aval), épure l'eau (nitrates, pesticides), stabilise le pied de versant (érosion régressive).
        \item \textbf{Drain d'évacuation (C) :} \textbf{Rôle hydraulique sécuritaire.} Évacue l'\textbf{excès d'eau} (crue) \textbf{sans érosion} (fond enherbé, pente douce, vérins). Empêche la rupture des murets par surpression/saturation. Dirige l'eau vers exutoire stable (fossé collecteur, talweg végétalisé).
        \item \textbf{Chemin de ronde (G) :} \textbf{Rôle logistique / entretien.} Permet \textbf{accès piéton/ânes} pour surveillance, réparation murets, taille haies, récolte fourrage/bois, curage drains. \textbf{Condition de la durabilité} du système (entretien préventif).
    \end{enumerate}

    \item \textbf{Pourquoi l'approche intégrée (association) est plus efficace qu'une technique unique ?}
    \begin{aretenir}[Principes de la Défense et Restauration des Sols (DRS) intégrée]
    \begin{enumerate}
        \item \textbf{Action sur tous les maillons de la chaîne érosive :} \textbf{Amont (D)} : réduire production ruissellement/sédiments (infiltration). \textbf{Versant (A, B, E, G)} : ralentir, infiltrer, piéger sur place. \textbf{Aval (F, C)} : épurer, évacuer sans dégâts.
        \item \textbf{Redondance / Sécurité :} Si une technique faillit (ex. muret fissuré, haie malade), les autres assurent la fonction (cultures, drain, fascine). Résilience face aux événements extrêmes (changement climatique).
        \item \textbf{Synergies biologiques :} Haies + Bandes enherbées + Reboisement = \textbf{corridors écologiques}, habitat auxiliaires (prédateurs ravageurs), pollinisateurs, MO continue $\rightarrow$ fertilité durable.
        \item \textbf{Optimisation coûts/bénéfices :} Techniques légères (courbes, haies) sur grandes surfaces, techniques lourdes (murets, drains) seulement où nécessaire (pentes fortes, concentration). Adapté aux moyens paysans.
        \item \textbf{Multifonctionnalité (Développement Durable) :} Production (cultures, fourrage, bois, fruits) + Protection (sol, eau, biodiversité) + Social (entraide entretien, chemins) + Économique (revenus diversifiés, résilience).
    \end{enumerate}
    \end{aretenir}
    \textbf{En résumé :} Une technique seule (ex. terrasses sans haies ni drains) conduit à l'échec (éboulements, saturation, rupture muret en aval, ravine nouvelle). L'association \textbf{intégrée} (mécanique + biologique + hydraulique + foncière) est la \textbf{seule garantie de durabilité} du système agraire sur le long terme. C'est l'esprit du \emph{paysannat bamiléké} modernisé par la science (DRS/CES).
\end{enumerate}
\end{corrige}

\begin{corrige}[Exercice 11 — Étude de cas : Barrage de Lagdo et colmatation]
\begin{enumerate}
    \item \textbf{Calcul du volume annuel de sédiments piégés :}
    \textbf{Données :}
    $V_{\text{initial}} = \SI{8 \times 10^9}{\cubic\metre}$ \quad Perte = $20\%$ \quad $\Delta t = \SI{40}{an}$

    \textbf{Volume total perdu (colmaté) :}
    \[
    V_{\text{perdu}} = \num{0,20} \times \SI{8 \times 10^9}{\cubic\metre} = \SI{1,6 \times 10^9}{\cubic\metre}
    \]

    \textbf{Volume annuel moyen piégé :}
    \[
    V_{\text{an}} = \frac{V_{\text{perdu}}}{\Delta t} = \frac{\SI{1,6 \times 10^9}{\cubic\metre}}{\SI{40}{an}} = \boxed{\SI{4 \times 10^7}{m^3.an^{-1}}}
    \]

    \item \textbf{Mécanisme physique : Érosion amont $\rightarrow$ Colmatage aval}
    \begin{enumerate}
        \item \textbf{Production (Bassins versants amont : Mayo Rey, Faro, Bénoué amont) :} Déforestation, agriculture sur brûlis, surpâturage, orpaillage $\rightarrow$ \textbf{Sol nu} $\rightarrow$ \textbf{Érosion hydrique massive} (feuille + rigoles + ravines). Détachement particules (argiles, sables, graviers, blocs).
        \item \textbf{Transport (Réseau hydrographique) :} L'eau de ruissellement charrie cette \textbf{charge solide} (transport en suspension : argiles/silts ; charriage : sables/graviers ; roulement : galets/blocs). L'énergie de la \textbf{pente} (gradient longitudinal du fleuve) maintient la capacité de transport.
        \item \textbf{Rupture de pente / Dépôt (Lac de retenue / Barrage) :} Arrivée au barrage, la pente devient \textbf{nulle} (plan d'eau). La vitesse de l'écoulement chute brutalement ($\approx 0$). L'énergie cinétique nécessaire au transport disparaît $\rightarrow$ \textbf{Dépôt immédiat de toute la charge solide} (loi de Stokes / Hjulström). Les particules les plus grossières (sables) se déposent en amont du lac (cônes de déjection), les fines (argiles) vont jusqu'au barrage.
        \item \textbf{Colmatage progressif :} Accumulation annuelle ($\sim \SI{4 \times 10^7}{\cubic\metre}$) $\rightarrow$ \textbf{Réduction volume mort puis volume utile} $\rightarrow$ Perte capacité de régulation (crue/étiage), production hydroélectrique, irrigation.
    \end{enumerate}

    \item \textbf{Quatre conséquences négatives du colmatage du barrage de Lagdo :}
    \begin{enumerate}
        \item \textbf{Production d'électricité (EDC / AES-SONEL) :} Perte volume utile $\rightarrow$ \textbf{baisse hauteur de chute nette} $\rightarrow$ puissance turbine réduite. Risque d'\textbf{arrêt turbines} en saison sèche (volume mort atteint). Coût \textbf{énergie thermique de substitution} (fioul, gaz) très élevé. Pertes de revenus EDC/État.
        \item \textbf{Irrigation (SEMRY - Société d'Expansion et de Modernisation de la Riziculture de Yagoua) :} Réduction volume d'eau stocké $\rightarrow$ \textbf{insécurité hydrique} pour les périmètres irrigués (riz, maïs, sorgho) en saison sèche. Baisse rendements rizicoles, menace sécurité alimentaire région Nord/Extrême-Nord. Conflits d'usage eau (énergie vs irrigation vs pêche vs eau potable).
        \item \textbf{Navigation fluviale (Bénoué) :} Colmatage du lit mineur en aval du barrage (bras morts, îlots sableux) + régulation débit altérée $\rightarrow$ \textbf{difficultés navigation} (transport marchandises, personnes, pêche) sur fleuve Bénoué (Garoua - Lagdo - Tchad). Impact économique régional.
        \item \textbf{Écosystèmes aquatiques / Biodiversité :} \textbf{Turbidité chronique} (argiles en suspension) $\rightarrow$ mortalité alevins, asphyxie branchies, réduction photosynthèse (phytoplancton). \textbf{Modification substrat} (vases vs graviers) $\rightarrow$ perte frayères (poissons migrateurs : \textit{Alestes}, \textit{Lates}, \textit{Heterotis}). \textbf{Eutrophisation} (nutriments piégés dans sédiments) $\rightarrow$ pullulations algues, anoxie. Perte biodiversité zones humides (Ramsar : plaine d'inondation Bénoué).
    \end{enumerate}

    \item \textbf{Plan d'action concerté (État, EDC, SEMRY, Populations, ONG) :}
    \begin{center}
    \begin{tabularx}{\linewidth}{|>{\bfseries}l|X|X|}\hline
    \textbf{Volet} & \textbf{Mesures CURATIVES (Urgence / Symptôme)} & \textbf{Mesures PRÉVENTIVES (Source / Durable)} \\ \hline
    \textbf{1. Dragage / Désenvasement} & \textbf{Dragage mécanique/hydraulique} ciblé (volume mort, prise d'eau turbines, entrée canaux SEMRY). \textbf{Valorisation sédiments :} amendement sols acides (chaux), briqueterie, construction routes. \textbf{Financement :} Fonds commun État/EDC/SEMRY/BAD. & \textbf{Non durable seul.} Coût récurrent énorme ($\sim$ milliards FCFA/an). Doit accompagner mesures préventives. \\ \hline
    \textbf{2. Conservation Eaux/Sols (CES/DRS) Bassins Versants} & \textbf{Stabilisation urgente} têtes de ravines actives (gabions, fascines, reboisement d'urgence) dans sous-bassins prioritaires (Mayo Rey, Faro). & \textbf{Cœur du plan :} \textbf{Généralisation CES/DRS} à l'échelle des bassins versants (Mayo Rey, Faro, Bénoué amont). \\
    & & \begin{itemize}
        \item \textbf{Agriculture de conservation (AC/SCV) :} Non-labour, couverts permanents (Mucuna, Brachiaria, Stylosanthes), rotation Coton/Maïs/Légumineuses. Cible : Planteurs coton (SODECOTON), vivriers.
        \item \textbf{Agroforesterie / Parc à Faidherbia / Haies vives :} Intégration arbres (Faidherbia, Acacia, Leucaena) dans champs + haies courbes de niveau. Fixation N, MO, brise-vent, racines.
        \item \textbf{Aménagements mécaniques ciblés :} Cordons pierreux, demi-lunes, zaï (zones dégradées glacis), terrasses (piémonts Mandara/Adamaoua), barrages filtrants (talwegs).
        \item \textbf{Reboisement / Régénération Naturelle Assistée (RNA) :} Zones sources, berges (ripisylves), corridors. Essences locales (Khaya, Afzelia, Parkia, Anogeissus).
    \end{itemize} \\ \hline
    \textbf{3. Gestion Pâturages / Élevage} & \textbf{Coupes-feu précoces}, points d'eau stratégiques pour répartir troupeaux. & \textbf{Pâturage tournant / différé}, aménagement parcours (réserves fourragères, cultures fourragères), \textbf{sédentarisation progressive} (pâturages améliorés), lutte feux tardifs (brigades villageoises). \\ \hline
    \textbf{4. Orpaillage / Mines} & \textbf{Fermeture sites illégaux} en zones sensibles (têtes de bassin). & \textbf{Formalisation / Encadrement} (coopératives), \textbf{techniques sans mercure}, \textbf{réhabilitation sites} (remblai, reboisement) obligation cahier charges. \\ \hline
    \textbf{5. Gouvernance / Financement / Suivi} & \textbf{Comité de Gestion Bassin Versant Bénoué (CGBV)} opérationnel (Préfet, MINEPDED, MINADER, MINFOF, MINEE, EDC, SEMRY, CTD, ONG, Chefs traditionnels, Planteurs, Éleveurs). & \begin{itemize}
        \item \textbf{Schéma Directeur d'Aménagement et de Gestion de l'Eau (SDADE) Bénoué} mis à jour \& appliqué.
        \item \textbf{PSE (Paiement Services Environnementaux) :} Rémunération planteurs/éleveurs pour pratiques vertueuses (haies, SCV, RNA) financé par \textbf{redevances eau EDC/SEMRY} + Fonds Environnement + Partenaires (BM, AFD, UE, GIZ).
        \item \textbf{Observatoire Sédimentaire :} Suivi bathymétrique barrage (annuel), turbidité fleuve (continu), taux érosion bassins (modélisation RUSLE/USLE + terrain).
        \item \textbf{Éducation / Sensibilisation :} Radios locales, écoles, chefs coutumiers, leaders d'opinion. « Protégeons le barrage, c'est notre avenir ».
    \end{itemize} \\ \hline
    \end{tabularx}
    \end{center}

    \begin{aretenir}[Clé de la réussite : Intégration Amont-Aval]
    Le colmatage de Lagdo est un \textbf{problème de bassin versant}, pas un problème de barrage. Aucune mesure curative (dragage) ne sera rentable sans \textbf{réduction drastique de l'apport sédimentaire à la source}. Cela exige une \textbf{mobilisation foncière, technique, financière et politique} sur 20--30 ans. L'alignement des intérêts (EDC/SEMRY = eau ; Paysans = sol/fertilité ; État = électricité/riz/sécurité) via le \textbf{PSE} est le levier institutionnel majeur.
    \end{aretenir}
\end{enumerate}
\end{corrige}

\newpage

\coursec{Fiche bilan}

\begin{popfigure}
\centering
\begin{tikzpicture}[
    mindmap,
    grow cyclic,
    every node/.style={concept, execute at begin node=\small\bfseries},
    root concept/.append style={concept color=popblue, minimum size=3.5cm, font=\large\bfseries},
    level 1 concept/.append style={minimum size=2.8cm, level distance=4.5cm, sibling angle=360/6},
    level 2 concept/.append style={minimum size=2.0cm, level distance=3.5cm},
    concept color=popblue!70,
    text=white,
    font=\small
]
\node[root concept, align=center]{Lutte contre\\ l'érosion}
    child[concept color=poporange] { node[concept] {Types d'érosion}
        child { node[concept, concept color=poporange!60, align=center]{Hydrique\\ (ruissellement, ravinement)} }
        child { node[concept, concept color=poporange!60, align=center]{Éolienne\\ (défletion, dépôts)} }
    }
    child[concept color=popgreen] { node[concept] {Facteurs favorisants}
        child { node[concept, concept color=popgreen!60, align=center]{Naturels\\ (pente, climat, sol nu)} }
        child { node[concept, concept color=popgreen!60, align=center]{Anthropiques\\ (déforestation, surpâturage, labour)} }
    }
    child[concept color=poppurple] { node[concept] {Conséquences}
        child { node[concept, concept color=poppurple!60, align=center]{Sur le sol\\ (perte fertilité, minéralisation)} }
        child { node[concept, concept color=poppurple!60, align=center]{Hors site\\ (envasement, inondations)} }
        child { node[concept, concept color=poppurple!60, align=center]{Socio-éco\\ (rendements, migrations)} }
    }
    child[concept color=poppink] { node[concept] {Techniques de lutte}
        child { node[concept, concept color=poppink!60, align=center]{Agronomiques\\ (courbes niveau, agroforesterie)} }
        child { node[concept, concept color=poppink!60, align=center]{Mécaniques\\ (terrasses, fascines, banquettes)} }
        child { node[concept, concept color=poppink!60, align=center]{Forestières\\ (reboisement, brise-vent)} }
    }
    child[concept color=popgold] { node[concept] {Démarche scientifique}
        child { node[concept, concept color=popgold!60, align=center]{Observer\\ (photo, terrain)} }
        child { node[concept, concept color=popgold!60] {Analyser causes} }
        child { node[concept, concept color=popgold!60] {Proposer solutions} }
    }
    child[concept color=popdark] { node[concept] {Contexte Cameroun}
        child { node[concept, concept color=popdark!60, align=center]{Mont Cameroun\\ (pentes, pluies)} }
        child { node[concept, concept color=popdark!60, align=center]{Grand Nord\\ (éolienne, désertification)} }
        child { node[concept, concept color=popdark!60, align=center]{Barrages\\ (envasement Lagdo, Memve'ele)} }
    };
\end{tikzpicture}
\legende{Carte mentale récapitulative de la leçon NCT08 : types, facteurs, conséquences, techniques de lutte, démarche et exemples camerounais.}
\end{popfigure}

\begin{bilan}
\courssub{Définitions clés}
\begin{itemize}
    \item \textbf{Érosion} : Détachement, transport et dépôt des particules du sol sous l'action d'agents naturels (eau, vent) ou anthropiques. Perte irréversible de la couche arable.
    \item \textbf{Érosion hydrique} : Causée par l'eau de pluie (éclaboussure, ruissellement diffus, ruissellement concentré $\rightarrow$ ravinement, gullying). Dominante au Sud-Cameroun (Mont Cameroun, Ouest).
    \item \textbf{Érosion éolienne} : Causée par le vent (défletion, saltation, suspension, dépôt). Dominante au Grand Nord (régions de l'Extrême-Nord, Nord, Adamaoua).
    \item \textbf{Courbe de niveau} : Ligne joignant les points d'une même altitude. Base du labour en courbes de niveau et de la construction des terrasses.
    \item \textbf{Agroforesterie} : Association délibérée d'arbres (ou arbustes) et de cultures (ou élevage) sur une même parcelle pour limiter l'érosion et améliorer la fertilité.
\end{itemize}

\courssub{Facteurs favorisants (Tableau de synthèse)}
\begin{center}
\begin{tabularx}{\linewidth}{|l|X|X|}\hline
\rowcolor{popblueL}
\textbf{Catégorie} & \textbf{Facteurs naturels} & \textbf{Facteurs anthropiques} \\ \hline
\textbf{Climat} & Pluies intenses/orages (érosivité), vents violents (Sahélien) & -- \\ \hline
\textbf{Relief} & Pentes fortes ($> 10\%$), longues pentes (Mont Cameroun, Mandara) & Déforestation des versants, cultures sur pente sans aménagement \\ \hline
\textbf{Sol} & Texture légère (sableux), faible agrégation, faible teneur en MO & Labour excessif (destruction structure), brûlis, jachères raccourcies \\ \hline
\textbf{Végétation} & -- & Déforestation (bois de feu, agriculture), surpâturage, cultures sarclées (maïs, coton) sans couverture \\ \hline
\end{tabularx}
\end{center}

\courssub{Conséquences majeures}
\begin{itemize}
    \item \textbf{Sur site (champ) :} Perte horizons superficiels (A, E) $\rightarrow$ baisse fertilité (N, P, K, MO), perte capacité de rétention d'eau, affleurement roches/cailloux, formation ravines/gullies (irrécupérables).
    \item \textbf{Hors site (aval/vent) :} Envasement cours d'eau et barrages (Lagdo, Memve'ele, Mape $\rightarrow$ perte capacité stockage, coûts curage), pollution eaux (turbidité, pesticides), colmatage zones humides, tempêtes de sable (Nord), inondations par réhausse lits mineurs.
    \item \textbf{Socio-économiques :} Baisse rendements agricoles, insécurité alimentaire, exode rural, coûts réhabilitation élevés.
\end{itemize}

\courssub{Techniques de lutte (Classification fonctionnelle)}
\begin{center}
\begin{tabularx}{\linewidth}{|l|X|X|}\hline
\rowcolor{popgreenL}
\textbf{Type} & \textbf{Techniques} & \textbf{Principe / Exemple Cameroun} \\ \hline
\textbf{Agronomiques} & Labour en courbes de niveau ; Bandes enherbées ; Agroforesterie (haies vives, \emph{Leucaena}, \emph{Gliricidia}) ; Couverts végétaux (mucuna, crotalaire) ; Paillage & Ralentir ruissellement, augmenter infiltration, protéger sol impact gouttes. Adapté cultures vivrières (maïs, manioc) Ouest/Centre. \\ \hline
\textbf{Mécaniques} & Terrasses (en gradins, à revers) ; Fascines (branchages) ; Banquettes ; Digues filtrantes (pierres) ; Bassins de rétention & Obstacles physiques au ruissellement concentré. Terrasses : rizières montagnes Ouest, caféiers Mont Cameroun. Fascines/banquettes : ravines naissantes. \\ \hline
\textbf{Forestières / Fixation} & Reboisement (eucalyptus, pin, acacia) ; Brise-vent (rangées d'arbres perpendiculaires vent dominant) ; Protection berges (bambous, \emph{Raphia}) & Recouvrir durablement le sol, fixer particules, briser force vent. Brise-vent : périmètres irrigués Nord (SEMRY, SODERCAM), fixation dunes Moungo. \\ \hline
\end{tabularx}
\end{center}

\courssub{Méthodes express (Savoir-faire probatoire)}
\begin{methode}[Identifier les causes dans une situation donnée]
\begin{enumerate}
    \item Observer le document (photo, texte, carte) : type d'érosion visible (ravines, dunes, dépôts), pente, état végétation, activités humaines.
    \item Classer les facteurs : Naturels (climat, pente, sol) vs Anthropiques (déforestation, labour, pâturage).
    \item Établir le lien de causalité : \emph{« La déforestation a mis le sol nu face aux pluies intenses sur pente forte $\rightarrow$ ruissellement concentré $\rightarrow$ ravinement. »}
\end{enumerate}
\end{methode}

\begin{methode}[Proposer des techniques adaptées]
\begin{enumerate}
    \item Diagnostiquer : Type érosion (hydrique/éolienne) ? Degré (débutant/avancé) ? Pente ? Contexte socio-éco (moyens paysans) ?
    \item Choisir la combinaison optimale : \textbf{Protection végétale (base)} + \textbf{Aménagement physique (si pente forte/ravines)}.
    \item Exemples : Pente $< 10\%$ + cultures annuelles $\rightarrow$ Courbes de niveau + haies vives. Pente $> 20\%$ + ravines $\rightarrow$ Terrasses + reboisement versants + fascines fond ravin. Zone sahélienne ventée $\rightarrow$ Brise-vent + paillage + cultures associées.
\end{enumerate}
\end{methode}

\begin{methode}[Interpréter un schéma/photo érosif]
\begin{enumerate}
    \item Légender les formes observées : splash (éclaboussure), rill (rigole), gully (ravine), dune, pavement éolien.
    \item Indiquer le sens du transport (flèches eau/vent) et zones dépôt (cône déjection, bas-fond, barrage).
    \item Identifier les indices de gravité : profondeur ravines, couleur eau (chargée), arbres déracinés, affleurement horizon B/C.
    \item Conclure sur l'urgence et la technique prioritaire.
\end{enumerate}
\end{methode}
\end{bilan}

\begin{aretenir}[Checklist avant le Probatoire]
\begin{itemize}
    \item $\square$ Définir l'érosion et distinguer clairement érosion hydrique (mécanismes : splash, ruissellement, ravinement) et érosion éolienne (défletion, transport, dépôt).
    \item $\square$ Citer \textbf{3 facteurs naturels} (pente, érosivité pluies, texture sol) et \textbf{3 facteurs anthropiques} (déforestation, labour à contre-pente, surpâturage) en les reliant au contexte camerounais (Mont Cameroun, Grand Nord).
    \item $\square$ Expliquer la conséquence \textbf{sur site} (perte horizon argilo-humique, baisse rendement) et \textbf{hors site} (envasement barrages Lagdo/Memve'ele, colmatage, tempêtes sable).
    \item $\square$ Différencier les \textbf{3 catégories de techniques} : Agronomiques (courbes niveau, agroforesterie), Mécaniques (terrasses, fascines, banquettes), Forestières (reboisement, brise-vent) et donner un exemple d'application pour chacune.
    \item $\square$ Justifier le choix d'une technique en fonction de la \textbf{pente}, du \textbf{type d'érosion} et des \textbf{moyens du paysan} (coût, main-d'œuvre).
    \item $\square$ Analyser une photo de ravinement : identifier la tête, les parois, le fond, le cône de déjection ; proposer fascines/banquettes + reboisement.
    \item $\square$ Analyser une photo de dune/érosion éolienne : identifier le sens du vent, la zone de déflation, la zone d'accumulation ; proposer brise-vent + fixation dune (palmiers, euphorbes).
    \item $\square$ Expliquer le rôle de la \textbf{matière organique} et de la \textbf{structure grumeleuse} dans la résistance du sol à l'érosion (agrégats stables).
    \item $\square$ Connaître les grands programmes camerounais : Plan National de Lutte contre la Désertification (PAN/LCD), reboisement « Opération Green Sahel », protection bassins versants barrages.
    \item $\square$ Maîtriser le vocabulaire précis : \emph{érosivité, érodibilité, courbe de niveau, terrasses à revers, fascine, banquette, brise-vent, envasement, gullying}.
    \item $\square$ Savoir réaliser un croquis annoté d'un aménagement anti-érosif (terrasse, courbe de niveau, brise-vent) avec flèches indiquant sens eau/vent et infiltration.
\end{itemize}
\end{aretenir}

\findecours

\end{document}
